% 第15章：SLAM与稠密建图
\chapter{SLAM与稠密建图}
\label{ch:15}

\paragraph*{本章阅读锚点}
本章的锚点是在线建图。原始 3DGS 离线优化完整图像集合；SLAM 方法必须边估位姿边更新高斯地图，并控制漂移、实时性和闭环一致性。


同时定位与建图（Simultaneous Localization and Mapping, SLAM）是移动机器人自主导航的核心能力：在探索未知环境的同时估计自身位姿并构建环境地图。传统 SLAM 系统通常使用稀疏点云或网格表示地图，无法提供直观的场景可视化。3DGS 将地图从稀疏几何提升为可实时渲染的照片级三维场景。这使得 SLAM 系统不仅能定位和建图，还能直接输出可用于下游任务的丰富场景表示。

3DGS-SLAM 的核心优势在于\textbf{表示统一}：同一组高斯基元同时服务于跟踪（tracking）和建图（mapping）。跟踪利用可微渲染的梯度回传优化位姿，建图通过光度/深度损失联合优化高斯参数。与传统方法相比，3DGS-SLAM 的输出不仅是点云或网格，更是可实时渲染的高保真三维场景，可直接用于 AR/VR、数字孪生、机器人导航等下游应用。

本章围绕 3DGS 在 SLAM 与稠密建图中的应用展开。第~\ref{sec:15_1}~节讨论基于 RGB-D 传感器的 SLAM 系统，这是 3DGS-SLAM 研究最密集的方向。第~\ref{sec:15_2}~节讨论更具挑战性的单目与纯 RGB SLAM。第~\ref{sec:15_3}~节介绍融合惯性测量单元（IMU）、LiDAR、雷达、事件相机等多模态传感器的 SLAM 方法。第~\ref{sec:15_4}~节关注视觉定位与重定位——利用 3DGS 作为场景表示进行精确的~6~自由度（6-DoF）位姿估计。第~\ref{sec:15_5}~节讨论语义 SLAM 与建图，将开放词汇理解融入稠密重建。第~\ref{sec:15_6}~节处理动态场景 SLAM 的特殊挑战。第~\ref{sec:15_7}~节讨论系统级优化与工程问题。第~\ref{sec:15_8}~节覆盖大规模场景和多智能体协同 SLAM。

\begin{figure}[htbp]
\centering
\resizebox{\textwidth}{!}{\begin{tikzpicture}[
  route/.style={rectangle, rounded corners=5pt, draw=#1, fill=#1!10,
    text width=3cm, minimum height=0.9cm, align=center, font=\small},
  method/.style={rectangle, rounded corners=3pt, draw=#1, fill=#1!6,
    font=\scriptsize, text width=2.8cm, align=center, minimum height=0.6cm},
  arr/.style={-{Stealth[length=4pt]}, techgray!60},
  >=Stealth
]

  \node[route=techgray, text width=4cm, font=\small\bfseries] (root) at (33.2, 5.6) {SLAM与稠密建图方法分类};
  \node[route=techblue] (b0) at (6.9, 3.1) {RGB-D SLAM};
  \node[method=techblue] (b0_l0) at (1.5, 0.6) {基础系统\\{\tiny SplaTAM, Photo-SLAM, CG-SLAM}};
  \node[method=techblue] (b0_l1) at (5.1, 0.6) {鲁棒加速\\{\tiny S$^3$LAM, 2DGS-SLAM, G2S-ICP}};
  \node[method=techblue] (b0_l2) at (8.7, 0.6) {恶劣环境\\{\tiny MipSLAM, EGG-Fusion}};
  \node[route=techgreen] (b1) at (15.9, 3.1) {单目/RGB SLAM};
  \node[method=techgreen] (b1_l0) at (12.3, 0.6) {经典优化\\{\tiny MonoGS, SplaTAM-RGB}};
  \node[method=techgreen] (b1_l1) at (15.9, 0.6) {前馈范式\\{\tiny Depth-SLAM, FeedForward-GS}};
  \node[route=techorange] (b2) at (23.5, 3.1) {多模态SLAM};
  \node[method=techorange] (b2_l0) at (19.6, 0.6) {IMU融合\\{\tiny GS-IMU, Inertial-GS}};
  \node[method=techorange] (b2_l1) at (23.5, 0.6) {LiDAR融合\\{\tiny LiDAR-GS, LIV-GS}};
  \node[route=techpurple] (b3) at (31.3, 3.1) {视觉定位};
  \node[method=techpurple] (b3_l0) at (27.4, 0.6) {渲染驱动\\{\tiny Render-Pose, Pose-GS}};
  \node[method=techpurple] (b3_l1) at (31.3, 0.6) {特征场定位\\{\tiny FeatureField-GS, GeomGS}};
  \node[route=techblue] (b4) at (35.1, 3.1) {语义SLAM};
  \node[method=techblue] (b4_l0) at (35.1, 0.6) {开放词汇\\{\tiny Semantic-GS, LangSplat-SLAM}};
  \node[method=techblue] (b4_l1) at (38.8, 0.6) {实时建图\\{\tiny SplatBridge, Online-SemGS}};
  \node[route=techgreen] (b5) at (42.4, 3.1) {动态场景SLAM};
  \node[method=techgreen] (b5_l0) at (42.4, 0.6) {运动分割\\{\tiny Dynamic-GS, MotionSeg-SLAM}};
  \node[method=techgreen] (b5_l1) at (46.0, 0.6) {动态物体\\{\tiny DynSLAM-GS, Obj-GS-SLAM}};
  \node[route=techorange] (b6) at (49.6, 3.1) {系统优化};
  \node[method=techorange] (b6_l0) at (49.6, 0.6) {工程优化\\{\tiny GS-ICP-SLAM, Fast-GS-SLAM}};
  \node[method=techorange] (b6_l1) at (53.5, 0.6) {CUDA加速\\{\tiny CUDA-GS, RealTime-SLAM}};
  \node[route=techpurple] (b7) at (57.4, 3.1) {大规模建图};
  \node[method=techpurple] (b7_l0) at (57.4, 0.6) {多智能体\\{\tiny MultiAgent-GS, Swarm-GS}};
  \node[method=techpurple] (b7_l1) at (60.7, 0.6) {分布式\\{\tiny DistGS-SLAM, CloudGS}};
  \draw[arr] (root) -- (b0);
  \draw[arr] (root) -- (b1);
  \draw[arr] (root) -- (b2);
  \draw[arr] (root) -- (b3);
  \draw[arr] (root) -- (b4);
  \draw[arr] (root) -- (b5);
  \draw[arr] (root) -- (b6);
  \draw[arr] (root) -- (b7);
  \draw[arr] (b0) -- (b0_l0);
  \draw[arr] (b0) -- (b0_l1);
  \draw[arr] (b0) -- (b0_l2);
  \draw[arr] (b1) -- (b1_l0);
  \draw[arr] (b1) -- (b1_l1);
  \draw[arr] (b2) -- (b2_l0);
  \draw[arr] (b2) -- (b2_l1);
  \draw[arr] (b3) -- (b3_l0);
  \draw[arr] (b3) -- (b3_l1);
  \draw[arr] (b4) -- (b4_l0);
  \draw[arr] (b4) -- (b4_l1);
  \draw[arr] (b5) -- (b5_l0);
  \draw[arr] (b5) -- (b5_l1);
  \draw[arr] (b6) -- (b6_l0);
  \draw[arr] (b6) -- (b6_l1);
  \draw[arr] (b7) -- (b7_l0);
  \draw[arr] (b7) -- (b7_l1);
\end{tikzpicture}}
\caption{SLAM与稠密建图方法分类}
\label{fig:ch15_slam_taxonomy}
\end{figure}

% ======================================================================
\section{RGB-D SLAM}
\label{sec:15_1}

RGB-D 传感器同时提供彩色图像和深度图，是 3DGS-SLAM 研究最基础的传感器配置。深度信息为高斯基元的初始化和几何约束提供了直接监督，使 RGB-D SLAM 成为 3DGS-SLAM 中跟踪精度和重建质量最稳定的方向。本节从奠基性系统出发，按技术路线组织 26 篇论文。这些工作覆盖抗锯齿与不确定性建模、几何表示改进、实时加速和恶劣环境鲁棒性等核心进展。

\begin{figure}[htbp]
\centering
\resizebox{\textwidth}{!}{\begin{tikzpicture}[
  box/.style={rectangle, rounded corners=4pt, draw=#1, fill=#1!10,
    text width=3.2cm, align=center, font=\small, minimum height=1cm},
  sbox/.style={rectangle, rounded corners=3pt, draw=#1!60, fill=#1!5,
    font=\scriptsize, text width=2.4cm, align=center, minimum height=0.5cm},
  arr/.style={-{Stealth[length=6pt]}, thick, techgray},
  darr/.style={-{Stealth[length=6pt]}, thick, techgray, dashed},
  >=Stealth
]

  \node[box=techblue] (s1) at (4.1, 0.6) {传感器输入(RGB-D/单目/IMU)};
  \node[box=techgreen] (s2) at (12.2, 0.6) {前端跟踪(帧间配准)};
  \node[box=techorange] (s3) at (19.7, 0.6) {局部建图(高斯优化/ADC)};
  \node[box=techpurple] (s4) at (27.8, 0.6) {后端优化(位姿图/全局BA)};
  \node[box=techblue] (s5) at (35.3, 0.6) {闭环检测(场景匹配)};
  \node[box=techgreen] (s6) at (41.9, 0.6) {全局地图(高斯场)};
  \draw[arr] (s1) -- (s2);
  \draw[arr] (s2) -- (s3);
  \draw[arr] (s3) -- (s4);
  \draw[arr] (s4) -- (s5);
  \draw[arr] (s5) -- (s6);
  \draw[darr] (s5) to[bend right=30] (s4);
  \draw[darr] (s6) to[bend right=30] node[midway, font=\scriptsize, text=techgray] {地图$\rightarrow$跟踪} (s2);
\end{tikzpicture}}
\caption{3DGS-SLAM 系统架构：前端跟踪、局部建图、后端优化与闭环检测}
\label{fig:ch15_slam_system}
\end{figure}

\subsection{基础系统与几何表示}
\label{ss:15_1}

本节介绍两条技术路线：（1）奠基性系统；（2）抗锯齿与不确定性建模。

\subsubsection{奠基性系统}

本路线介绍 3DGS-SLAM 的两个奠基性系统——SplaTAM 和 Photo-SLAM。它们分别确立了基于轮廓感知的在线致密化范式和基于 ORB 特征的解耦建图范式。3DGS-SLAM 的起点可以追溯到两个几乎同时出现的奠基性工作。SplaTAM（2023.12）~\cite{2312_02126} 首次将 3D 高斯基元作为 RGB-D SLAM 的唯一三维表示，统一了跟踪、建图和渲染。具体而言，它修改了标准 3DGS 的场景表示与致密化机制：用各向同性球状高斯替代完整协方差椭球高斯，用基于渲染轮廓的在线致密化替代基于梯度统计的自适应致密化，而 $\alpha$-混合渲染管线保持不变。

\textbf{问题与动机。}在 SplaTAM 之前，NeRF-SLAM 方法（如 NICE-SLAM、Co-SLAM）使用隐式多层感知机（Multilayer Perceptron, MLP）或三平面编码表示场景，需要逐射线采样和体渲染，单帧渲染耗时数百毫秒；且无法显式判断某个空间位置是否已被充分观测。即使是标准 3DGS（59 维参数/高斯，含球谐系数和完整 3D 协方差矩阵），也并非为在线 SLAM 设计：其自适应致密化依赖离线全局视角统计，不适合帧序到达的在线设定；椭球高斯的旋转和缩放参数增加了优化维度，在稀疏视角下容易过拟合。因此，SplaTAM 的核心思路是：对 3DGS 表示做面向在线 SLAM 的最小化简化，同时引入轮廓感知机制解决"不知道哪里没建好"的根本问题。

\textbf{简化表示。}SplaTAM 将每个高斯简化为 8 个参数：位置 $\boldsymbol{\mu} \in \mathbb{R}^3$、RGB 颜色 $\mathbf{c} \in \mathbb{R}^3$（省略球谐函数，采用视角无关颜色）、半径 $r$（各向同性球状，省略完整协方差矩阵）和不透明度 $o \in [0,1]$。这一简化使参数量从 59 维降至 8 维，大幅减少在线优化的变量数目。各向同性假设的设计动机不仅是简化。论文的消融实验表明，在 SLAM 的在线帧序设定下，各向同性与各向异性高斯在跟踪精度和重建质量上差异微小（ScanNet++ 上 PSNR 差距小于 0.1 dB），但前者速度更快、内存效率更高。

\textbf{可微渲染与轮廓图像。}在渲染端，SplaTAM 不仅渲染 RGB 图像和深度图，还引入了\textbf{轮廓图像}（silhouette image）的概念。直觉上，轮廓图像是渲染过程的一个"副产品"——它不记录颜色或深度，只记录每个像素被高斯"覆盖了多少"。SplaTAM 同时输出三种图像：颜色图像 $C(\mathbf{p})$（与标准 3DGS 相同的 $\alpha$-混合公式）；深度图像 $D(\mathbf{p}) = \sum_{i=1}^{n} d_i f_i(\mathbf{p}) \prod_{j=1}^{i-1}(1 - f_j(\mathbf{p}))$（将颜色替换为各高斯在相机坐标系下的深度 $d_i$，几何含义为沿射线的加权平均深度）；以及轮廓图像：

\begin{equation*}
S(\mathbf{p}) = \sum_{i=1}^{n} f_i(\mathbf{p}) \prod_{j=1}^{i-1}(1 - f_j(\mathbf{p})),
\end{equation*}

其中 $f_i(\mathbf{p}) = o_i \exp\left(-\frac{\|\mathbf{p} - \boldsymbol{\mu}_i^{2D}\|^2}{2(r_i^{2D})^2}\right)$ 为第 $i$ 个高斯的 2D 投影贡献，$\boldsymbol{\mu}_i^{2D} = K E_t \boldsymbol{\mu}_i$ 和 $r_i^{2D} = r_i f / d_i$ 分别为高斯中心在像素空间的投影位置和投影半径（$K$ 为相机内参矩阵，$E_t$ 为帧 $t$ 的外参矩阵，$f$ 为焦距，$d_i = (E_t \boldsymbol{\mu}_i)_z$ 为高斯在相机坐标系下的深度）。轮廓值 $S(\mathbf{p})$ 接近 1 表示该像素已被充分覆盖，接近 0 则表示该区域尚无高斯或覆盖不足。简言之，轮廓值量化了"地图对每个像素的了解程度"。

\textbf{轮廓感知跟踪。}这一设计是 SplaTAM 的关键洞察：轮廓掩码为系统提供了\textbf{逐像素的认知不确定性}（epistemic uncertainty），使跟踪阶段能仅在高置信区域计算损失，避免未建图区域产生误导梯度。具体而言，系统先用恒速模型（constant velocity model）初始化位姿 $E_{t+1} = E_t + (E_t - E_{t-1})$——即假设相机运动速度在连续帧间保持不变，为梯度优化提供一个合理的起始点。然后通过可微渲染最小化光度-深度联合损失：

\begin{equation*}
\mathcal{L}_t = \sum_{\mathbf{p}} \mathbb{1}(S(\mathbf{p}) > 0.99) \cdot \left(\|\hat{D}(\mathbf{p}) - D(\mathbf{p})\|_1 + 0.5 \|\hat{C}(\mathbf{p}) - C(\mathbf{p})\|_1\right),
\end{equation*}

即仅对轮廓超过 0.99 的像素施加约束。论文中颜色损失权重 0.5 是经验选定的——观测到深度 $\hat{D}$ 的典型残差范围（0.002--0.006）远小于颜色 $\hat{C}$ 的残差范围（0.01--0.03），需要降权颜色项以避免其主导梯度。消融实验证实了轮廓掩码的必要性：不使用掩码时跟踪完全失败；阈值从 0.99 降至 0.5（仅保留"大致可见"的像素）时误差增大 5 倍。

\textbf{在线致密化。}标准 3DGS 使用基于梯度统计的自适应致密化（在累积梯度超过阈值的位置克隆或分裂高斯），但该机制依赖离线多轮训练的全局视角统计，不适合帧序到达的在线 SLAM。SplaTAM 转而使用基于轮廓和深度的在线致密化——直觉是：哪里地图"没覆盖到"或"深度不对"，就在哪里添加新高斯。具体通过致密化掩码判断哪些像素需要新增高斯：

\begin{equation*}
M(\mathbf{p}) = \mathbb{1}(S(\mathbf{p}) < 0.5) + \mathbb{1}\!\left(\hat{D}_{\mathrm{GT}}(\mathbf{p}) < D(\mathbf{p})\right) \cdot \mathbb{1}\!\left(\|D(\mathbf{p})\|_1 > 50 \cdot \mathrm{MDE}\right),
\end{equation*}

其中第一项 $S < 0.5$ 捕获"覆盖不足"的区域（轮廓值低意味着该像素前几乎没有高斯）；第二项捕获"深度不一致"的区域（真实深度比渲染深度更近，说明观测到了当前地图中不存在的新几何，且误差超过中值深度误差的 50 倍以排除噪声）。掩码为正的像素处新增高斯：初始不透明度设为 $0.5$，半径设为 $r = d_{\mathrm{GT}} / f$（即在图像平面上恰好占据约一个像素的大小），位置由像素深度反投影到 3D 空间确定。

\textbf{建图与关键帧管理。}建图阶段固定相机位姿，仅优化高斯参数。与跟踪阶段不同，建图损失不使用轮廓掩码（因为建图的目标是优化地图中所有区域的质量），并额外加入 SSIM 损失以提升结构相似性。关键帧选择策略基于当前帧点云与各候选关键帧视锥的重叠度：保留间隔为 $n$ 帧的关键帧集合，每轮选取 $k$ 帧（含当前帧、最近关键帧和 $k-2$ 个重叠度最高的历史关键帧）联合优化，以控制计算量。此外，系统周期性剔除无效高斯：不透明度接近零（对渲染几乎无贡献）或半径过大（膨胀到不合理尺寸）的高斯被移除。

\textbf{实验表现。}在 Replica 数据集上，SplaTAM 将轨迹误差从先前最佳 Point-SLAM 的 0.52 cm 降至 0.36 cm（减少 30\%），渲染 PSNR 比 NICE-SLAM 和 Vox-Fusion 高约 10 dB。在 ScanNet++ 数据集上，Point-SLAM 和 ORB-SLAM3 因帧间大位移和纹理缺失完全跟踪失败，而 SplaTAM 仍保持 1.2 cm 的轨迹误差，新视角合成达到 24.41 dB PSNR。在 TUM-RGBD 上（深度稀疏且运动模糊严重），SplaTAM 将轨迹误差从 Point-SLAM 的 8.92 cm 降至 5.48 cm（减少近 40\%）。这些结果确立了 3DGS-SLAM 的可行性基准。

与 SplaTAM 将跟踪与建图统一在同一可微渲染管线中的范式不同，Photo-SLAM（2023.11）~\cite{2311_16728} 采用了跟踪与建图解耦的架构。其设计动机在于：SplaTAM 的纯梯度下降位姿优化收敛域有限，在快速运动或纹理缺失时容易陷入局部最优；而传统特征点方法（如 ORB-SLAM3）在位姿估计上已非常成熟且鲁棒——在 TUM 数据集上，ORB-SLAM3 的定位精度甚至优于 DROID-SLAM、Go-SLAM 等学习方法。因此 Photo-SLAM 将定位任务完全交给 ORB 特征前端的因子图优化（用 Levenberg--Marquardt 算法最小化重投影误差），将照片级渲染交给后端 3DGS 优化，两者通过\textbf{超基元地图}（hyper primitives map）桥接。

\textbf{超基元地图与 SplaTAM 简化高斯的差异。}与 SplaTAM 将高斯参数压缩为 8 维各向同性球状高斯不同，Photo-SLAM 的超基元保留了完整 3DGS 表达能力。每个超基元存储位置 $\mathbf{P} \in \mathbb{R}^3$、旋转 $\mathbf{r} \in SO(3)$、缩放 $\mathbf{s} \in \mathbb{R}^3$、密度 $\alpha \in \mathbb{R}^1$、球谐系数 $\text{SH} \in \mathbb{R}^{16}$（对应 0--1 阶球谐，可建模视角相关颜色），以及 ORB 描述子 $\mathbf{O} \in \mathbb{R}^{256}$。换言之，Photo-SLAM 不需要简化高斯表示，因为定位任务已由前端 ORB 模块承担，高斯参数只需优化渲染质量。渲染公式与标准 3DGS 的 $\alpha$-混合一致：$C(R,t) = \sum_{i} c_i \alpha_i \prod_{j=1}^{i-1}(1-\alpha_j)$。光度损失采用 $\mathcal{L} = (1-\lambda)|I_r - I_{\text{gt}}|_1 + \lambda(1-\text{SSIM}(I_r, I_{\text{gt}}))$，用固定学习率 0.2 的随机梯度下降（Stochastic Gradient Descent, SGD）优化。

\textbf{几何致密化：利用 ORB 特征引导高斯增生。}SplaTAM 用渲染轮廓和深度不一致性驱动在线致密化，而 Photo-SLAM 利用前端已有的 ORB 特征实现一种互补的致密化策略。具体而言，帧中约 70\% 的 ORB 特征点在几何建图阶段未能与 3D 点建立对应关系（即"非活跃"特征点），但这些点恰好分布在纹理丰富、需要更多高斯建模的区域。Photo-SLAM 在每个关键帧创建时，主动将这些非活跃 ORB 特征点转化为临时超基元：在 RGB-D 模式下直接用深度值投影；在单目模式下用最近邻活跃特征点的深度估计；在双目模式下用立体匹配估计深度。消融实验表明，去掉几何致密化后 PSNR 下降 2.028 dB（单目）和 1.662 dB（RGB-D）。

\textbf{高斯金字塔训练。}Photo-SLAM 进一步引入\textbf{高斯金字塔}（Gaussian-Pyramid）渐进训练方法：默认使用 3 级金字塔（$n=2$），训练从最高层级（最粗分辨率）开始，随迭代增加逐级降至原始分辨率。形式化地，在训练时间段 $t_0$ 至 $t_n$ 中依次使用 $L(I^{n}_r, \sigma^n(I_{\text{gt}}))$, $L(I^{n-1}_r, \sigma^{n-1}(I_{\text{gt}}))$, $\ldots$, $L(I^{0}_r, \sigma^{0}(I_{\text{gt}}))$ 监督，其中 $\sigma^n$ 表示高斯金字塔第 $n$ 层的下采样图像。这一设计的核心洞察是：先在低分辨率上让高斯快速收敛到大致结构（此时每个像素的感受野覆盖更大区域，梯度更稳定），再逐级精化细节，避免高分辨率下高频噪声导致的优化不稳定。消融实验显示 3 级金字塔取得最佳效果，4 级反而因过拟合低层特征而恶化。值得注意的是，Photo-SLAM 从单张图像自生成多分辨率金字塔，与 BungeeNeRF 等需要多分辨率输入图像的方法不同，无需额外数据。

\textbf{回环闭合。}Photo-SLAM 还引入了回环闭合机制。该机制利用词袋模型（Bag of Words, BoW）的 DBoW2 库检测回环后，通过相似变换（similarity transformation）校正局部关键帧位姿和超基元位置，再由照片级建图线程进一步消除由里程计漂移引起的重影伪影。SplaTAM 不具备回环闭合能力，这是 Photo-SLAM 在长序列场景下的重要优势。

\textbf{实验表现。}在 Replica 数据集上，Photo-SLAM 的渲染速度约 300 FPS（RTX 4090），比 ESLAM 和 Go-SLAM 快约 300 倍，PSNR 在 RGB-D 和单目设定下均达到当时最优水平。在嵌入式平台 Jetson AGX Orin 上仍保持约 100 FPS 的渲染速度和实时跟踪（$>40$ FPS），跟踪速度比 ESLAM 快约 6 倍、比 Co-SLAM 快约 3 倍，且跟踪延迟的标准差最小。Photo-SLAM 还支持单目、双目和 RGB-D 三种相机模式。这是当时绝大多数 NeRF-SLAM 和 3DGS-SLAM 所不具备的灵活性，展示了 3DGS-SLAM 在机器人应用中的落地潜力。

\subsubsection{几何表示改进与表面重建}

本路线介绍通过改进 3D 高斯的几何表示来提升 SLAM 表面重建质量的方法，涵盖 2D 高斯、表面法向约束和 TSDF 融合等技术。

奠基性系统展示了 3DGS-SLAM 的可行性，但标准 3D 高斯的椭球体表示在表面重建精度上存在固有局限。这一小节聚焦于如何通过改进几何表示来提升表面重建质量。标准 3D 高斯是三维椭球体，在表示薄壁和平面表面时存在固有几何冗余。用一组 3D 椭球去拟合一张薄墙或桌面，需要大量高斯堆叠才能近似平面效果。

\textbf{从 3D 椭球到 2D Surfel 的动机。}下文的 CG-SLAM 通过各向同性正则化和深度不确定性建模约束 3D 高斯的几何行为，但本质上仍在椭球体表示的框架内工作。椭球在射线方向上的厚度始终存在，导致深度渲染是多个椭球沿射线的加权平均，而非对表面的精确采样。S$^3$LAM（2025.07）~\cite{2507_20854} 采取了更根本的路线：直接放弃 3D 椭球体，改用 2D Gaussian Surfel（表面元素，Surface Element）作为基元，从表示层面消除射线方向的几何冗余。这一设计选择的核心动机是：3D 椭球在射线方向上始终具有非零厚度，导致渲染深度是多个椭球的加权平均而非对表面的精确采样；而 2D Surfel 的厚度为零，深度直接由射线与平面的交点决定，从数学上消除了这种模糊性。

\textbf{2D Surfel 表示。}与 3D 椭球不同，每个 2D Surfel 是一个有朝向的二维平面高斯，由中心 $\mathbf{p}_k$、两个切向量 $\mathbf{t}_u, \mathbf{t}_v$ 及对应缩放因子 $s_u, s_v$ 参数化。法向量由叉积 $\mathbf{t}_w = \mathbf{t}_u \times \mathbf{t}_v$ 定义，旋转矩阵为 $R = [\mathbf{t}_u, \mathbf{t}_v, \mathbf{t}_w]$。Surfel 上的局部点参数化为 $P_k(u,v) = \mathbf{p}_k + s_u \mathbf{t}_u u + s_v \mathbf{t}_v v$，直觉上是一个以 $\mathbf{p}_k$ 为中心、沿 $\mathbf{t}_u, \mathbf{t}_v$ 方向展开的局部平面片。渲染时，像素射线与 Surfel 平面求交得到局部坐标 $(u,v)$，然后计算高斯透明度 $G_k(u) = \exp(-u^2/2 - v^2/2)$，颜色、深度和法向量通过 $\alpha$-混合 渲染：$C(\mathbf{x}) = \sum_k \tau_k c_k$，$D(\mathbf{x}) = \sum_k \tau_k z_k / \sum_k \tau_k$，$N(\mathbf{x}) = \sum_k \tau_k \mathbf{t}_w^k$。与 3D 高斯的关键区别在于：Surfel 的深度是射线与平面交点的精确深度（而非椭球中心的投影深度），且法向量天然编码了表面朝向信息。

\textbf{自适应渲染：深度失真驱动的确定性切换。}在线 SLAM 中高斯基元数量有限且多视角覆盖不充分，直接使用 $\alpha$-混合 加权平均可能在边缘和角落等尖锐几何处产生不完整的表面重建——$\alpha$-混合 会将不同深度的 Surfel 混合，模糊真实的表面位置。S$^3$LAM 引入\textbf{深度失真项}（depth distortion）作为逐像素不确定性的度量：

\begin{equation}
    D_d = \sum_{i,j} \tau_i \tau_j |z_i - z_j|
\end{equation}

其中 $\tau_i$ 为像素处的 $\alpha$-混合 权重，$z_i$ 为第 $i$ 个射线--Surfel 交点的深度。直觉上，$D_d$ 衡量的是：在该像素处，被渲染到的各个 Surfel 之间深度的加权离散程度——$D_d$ 大说明多个深度差异显著的 Surfel 竞争同一像素，几何不收敛。S$^3$LAM 以阈值 $D_d > \theta$（$\theta = 5 \times 10^{-6}$）识别不确定像素，对这些像素切换渲染策略：不再使用 $\alpha$-混合 加权平均，而是直接选取 blending 权重最大的 Surfel（dominant surfel），以其射线--平面交点确定深度 $D_c$ 和法向量 $N_c$。为避免切换策略在大平面上产生空洞，仅在 $D > D_c$（即 $\alpha$-混合 深度比 dominant surfel 深度更大）时才执行替换，确保平面完整性。这一策略可在前向和反向传播中无缝集成到 Surfel 光栅化管线中，计算开销极小。

\textbf{Surfel 级位姿雅可比与径向梯度。}CG-SLAM 对 EWA splatting 推导了相机位姿关于渲染损失的完整解析梯度，但其深度渲染的雅可比仅包含 $\alpha$-混合 权重对位姿的导数（因为 3D 高斯深度直接取椭球中心的投影，不涉及交点计算）。S$^3$LAM 的 Surfel 深度渲染是射线与定向平面交点的 $\alpha$-混合，其位姿雅可比包含一项额外的\textbf{径向梯度}分量：

\begin{equation*}
    \frac{\partial D(\mathbf{x})}{\partial T_{CW}} = \frac{\partial D}{\partial \{\tau_{i,\mathbf{x}}\}} \frac{\partial \{\tau_{i,\mathbf{x}}\}}{\partial T_{CW}} + \frac{\partial D}{\partial P(u(\mathbf{x}))} \left(\frac{\partial \mathbf{p}_i}{\partial T_{CW}} + \underbrace{\frac{\partial \mathbf{t}_r}{\partial T_{CW}}}_{\text{径向梯度}}\right)
\end{equation*}

其中 $P(u(\mathbf{x}))$ 为射线--Surfel 交点的相机坐标，$\mathbf{t}_r = R(u_s s_u \mathbf{t}_u + v_s s_v \mathbf{t}_v)$ 为从 Surfel 中心到射线--Surfel 交点的径向向量。径向梯度在统计意义上引导相机旋转使射线"靠近"或"远离"Surfel 中心，从而将重建表面对齐到真实观测。论文的收敛域分析实验证实了这一点：在仅使用 $3 \times 3$ 训练视角的受限设置下，S$^3$LAM 在初始位姿偏离目标 1.6 m 时仍能成功收敛，而 3DGS-based 位姿优化在此距离上几乎完全失败。这解释了 S$^3$LAM 在 ScanNet++ 等帧间视点剧变序列上的鲁棒跟踪表现——例如在 S1 序列的第 363--364 帧之间发生大幅旋转导致其他方法跟踪失败时，S$^3$LAM 仍能正确收敛。

\textbf{建图损失与 Surfel 管理。}建图优化同时渲染颜色、深度和法向量，使用三路损失的加权组合：

\begin{equation*}
    \mathcal{L}_\text{map} = \lambda_C \|C_t - \hat{C}_t\|_1 + \lambda_D \|D_t - \hat{D}_t\|_1 + \lambda_N (1 - N_t \cdot \hat{N}_t)
\end{equation*}

其中 $\lambda_C = 1.0, \lambda_D = 1.0, \lambda_N = 0.1$，法向量损失采用余弦相似度。Surfel 新增基于三条准则——透射率超过阈值 $\tau_T = 0.5$（覆盖不足）、深度误差超过 $\tau_d = 0.1$、颜色误差超过 $\tau_c = 0.1$——满足任一准则的像素均匀采样后投影到 3D 空间生成新 Surfel。错误 Surfel 通过累积平均误差超过双倍阈值来淘汰。关键帧选择基于相对旋转和相对平移阈值。在 RTX 3090 上，S$^3$LAM 在 Replica 数据集的 3D 重建精度（Accuracy/Completion/F1）上达到或超越所有 3DGS-based 方法，与最佳 NeRF-based 方法 Point-SLAM 相当，但计算效率和内存消耗显著更优。

S$^3$LAM 证明了 2D Surfel 在在线 SLAM 中的潜力，但其缺乏回环检测与全局校正机制，长序列中跟踪漂移会不断累积。2DGS-SLAM（2025.06）~\cite{2506_00970} 保留了 S$^3$LAM 的 2D Surfel 表示和基于射线的深度渲染框架，但在三个关键维度上做了系统性扩展。

\textbf{动机。}S$^3$LAM 虽已验证了 2D Surfel 在在线 SLAM 中的可行性，却缺少回环检测与全局校正，长序列中漂移不断累积，最终导致场景表示与真实环境严重偏离。

\textbf{基于李代数的 SE(3) 位姿雅可比。}S$^3$LAM 的位姿雅可比包含径向梯度分量以改善收敛域；2DGS-SLAM 则采用不同的推导路径——在 Lie 代数 $se(3)$ 上推导完整的解析雅可比，将位姿梯度分解为中心点梯度 $\partial L / \partial \boldsymbol{\mu}_c$ 和旋转矩阵梯度 $\partial L / \partial R_c$ 两路，通过 Lie 群上的 $\oplus$ 运算投影到同一切空间后求和，保证优化始终在 $SE(3)$ 流形上进行，避免使用四元数--旋转矩阵转换时的流形外偏差。跟踪损失采用法向量遮罩 $M_n(\mathbf{n}_r^\top \mathbf{r} > 0)$ 过滤背向面片，并辅以不透明度遮罩 $M_o(O > 0.95)$ 排除未充分重建区域。

\textbf{回环检测与全局地图变形。}这是 S$^3$LAM 完全不具备的后端能力。2DGS-SLAM 利用 MASt3R 基础模型提取视觉 Transformer 编码器特征，通过 ASMK（Aggregated Selective Match Kernel）量化为稀疏二值表示进行高效图像检索。回环验证流程依次检查：MASt3R 置信度均值超过阈值 $c_s$、深度一致性重叠率 $O_{lc} > \theta_o$（深度阈值 $\delta_d = 0.05$）、以及在活动地图上的 scan-to-model 跟踪误差 $e_t < \epsilon_t$。通过验证后，相对位姿作为约束加入位姿图，优化后每个高斯根据其关联关键帧的位姿增量 $\Delta T$ 同步更新位置和旋转：$\boldsymbol{\mu}_k' = \Delta T \, \boldsymbol{\mu}_k$，$R_k' = \Delta R \, R_k$，实现全地图的一致变形。

\textbf{活动/非活动双缓冲地图。}为解决跟踪漂移导致的新旧结构错位问题，2DGS-SLAM 维护活动子图 $M_A$（最近观测到的高斯）和非活动子图 $M_I$（历史高斯）。超过时间阈值未观测到的高斯自动转为非活动状态，回环校正后则重新激活被重新观测到的历史高斯，避免同区域重复创建。新增高斯时还通过体素哈希表（分辨率 $r_h$）检查空间占用，防止冗余。在 Replica 数据集上，2DGS-SLAM 的表面重建精度相比 S$^3$LAM 有进一步提升，且回环校正后全局一致性显著优于无回环的方案。

G2S-ICP SLAM（2025.07）~\cite{2507_18344} 沿着 CG-SLAM 的几何稳定性思路进一步推进——不再用正则化约束 3D 高斯的各向异性，而是直接采用 2D 高斯盘（2D Gaussian disk）表示场景，从表示层面消除深度不一致的根源。

\textbf{动机。}3D 椭球从不同视角相交于不同截面，导致多视角深度不一致；而 2D 盘在任意视角下投影一致，从根本上杜绝了深度模糊。具体而言，每个高斯基元的协方差矩阵 $C_k = R_k S_k S_k^\top R_k^\top$ 中缩放矩阵 $S_k = \text{diag}(s_1, s_2, 0)$——第三个尺度分量硬性置零，使高斯的空间支撑被约束在局部切平面上，从根本上杜绝了 3D 椭球沿法线方向的模糊。

在跟踪环节，G2S-ICP SLAM 将上述 2D 各向异性协方差嵌入广义 ICP（Generalized ICP, GICP）配准框架：源点云中每个点被建模为高斯分布 $\mathcal{N}(x_m^\text{src}, C_m^\text{src})$，其协方差矩阵的法线方向方差为零。这一表面约束先验使 GICP 的 Mahalanobis 距离最小化自然偏向几何一致的配准，无需修改底层配准公式。在建图环节，该方法引入距离感知缩放策略 $s_{1,2} \propto 1/z^p$（$z$ 为深度值，$p=0.333$），使远处高斯盘的初始尺度自动缩小以补偿投影稀疏性，提升视觉均匀性和位姿估计稳定性。G2S-ICP SLAM 在 NVIDIA RTX A5000 上以 30 FPS 实时运行（远超 CG-SLAM 的 8.5 Hz），在 Replica 数据集上平均深度 L1 误差为 0.74 cm，较 GS-ICP SLAM 降低 96\%。

沿着同一"2D Surfel 替代 3D 椭球"的技术路线，EGG-Fusion（2025.12）~\cite{2512_01296} 将 CG-SLAM 的 3D 高斯场替换为\textbf{2D 高斯 Surfel}（盘状基元），并将不确定性建模从启发式阈值筛选升级为\textbf{信息滤波器}（Information Filter）递归贝叶斯融合。

\textbf{动机。}CG-SLAM 的"两步淘汰"启发式（先降低透明度再移除）是硬决策，一旦误判就无法恢复；而信息滤波器通过递归贝叶斯更新逐步累积置信度，提供了更平滑、更鲁棒的不确定性建模方式。具体而言，每个 Surfel 的几何状态（位置+法线）维护一个 $6\times6$ 协方差矩阵 $\Omega^t$，每收到一帧新的 RGB-D 观测就执行一次信息矩阵更新 $\Omega^t = \Omega^{t-1} + H^\top \Sigma_z^{-1} H$（其中 $H$ 为观测雅可比矩阵，$\Sigma_z$ 为测量噪声协方差），逐步累积置信度而非依赖 CG-SLAM 的"两步淘汰"启发式。测量噪声 $\Sigma_z$ 与深度平方成正比，显式编码消费级传感器远距离噪声增大的物理特性。

此外，EGG-Fusion 引入两项与 CG-SLAM 不同的设计：一是自适应尺度初始化 $s = [\delta_s d/f_x,\, \delta_s d/f_y]$（$\delta_s=2.0$），远处的 Surfel 自动获得更大的世界空间尺度以在图像空间保持均匀覆盖，而非 CG-SLAM 的固定尺度；二是稀疏到稠密的位姿估计策略——先用稀疏特征对应（2D-3D 匹配 + Levenberg--Marquardt）提供对快速运动鲁棒的位姿初始化，再通过几何-光度联合稠密对齐精化。在 Replica 数据集上，EGG-Fusion 以 24 FPS 实时运行，表面重建误差仅 0.6 cm，比先前最佳 GS-SLAM 方法提升超过 20\%。

CG-SLAM（2024.03）~\cite{2403_16095} 修改了标准 3DGS 的\textbf{场表示与优化策略}——用一致性几何稳定的不确定性感知高斯场替代原始无约束的高斯场，$\alpha$-混合 渲染管线保持不变。CG-SLAM 系统性地解决了标准 3DGS 直接用于 SLAM 时的三重缺陷。在基线做法中，原始 3D 高斯场作为纯视角合成技术，其强各向异性使其倾向于过拟合输入图像而非准确建模 3D 表面。具体表现为高斯基元在优化后呈现出"箭形伪影"（arrow-shaped artifacts），即某个轴方向被极度拉伸以拟合个别视角的颜色，而非贴合真实表面几何；高斯基元可能不对齐环境表面，导致外推能力差并进一步退化相机跟踪；此外，建图过程中高斯基元数量的持续增长会不可避免地拖慢跟踪效率。CG-SLAM 的核心改进是构建一个\textbf{一致性、几何稳定且不确定性感知的 3D 高斯场}，具体包含三个层面的设计：

（1）\textbf{各向同性正则化}。在标准 3DGS 中，高斯基元的协方差矩阵 $\Sigma = RSS^\top R^\top$ 中缩放矩阵 $S = \text{diag}([s^x, s^y, s^z])$ 的三个分量可以独立优化，没有任何约束阻止某个方向极度拉伸。这种自由度在离线新视角合成中问题不大（多视角约束自然抑制极端形状），但在 SLAM 的在线增量建图场景中，早期帧数少、多视角约束不足，高斯极易退化为沿某一方向极度拉伸的"箭形"，既不贴合表面又产生误导梯度。为了抑制这种各向异性过拟合，CG-SLAM 引入软尺度正则化损失，鼓励高斯椭球趋向球形，在跟踪精度与渲染真实感之间取得平衡：

\begin{equation}
    \mathcal{L}_\text{iso} = \frac{1}{|G|} \sum_{i \in G} \max\!\Big(\frac{\max(\{s_i^x, s_i^y, s_i^z\})}{\min(\{s_i^x, s_i^y, s_i^z\})} - \gamma, 0\Big)
\end{equation}

直观地说，这个损失惩罚尺度因子的最大/最小比值超过阈值 $\gamma$ 的高斯——当三个尺度因子接近相等时比值为 1，越拉伸比值越大。其中 $s_i^{x,y,z}$ 为第 $i$ 个高斯的尺度因子，$\gamma=1.0$ 控制允许的各向异性程度。该损失是"软"的：各向异性比值低于 $\gamma$ 的高斯不受任何惩罚，不影响渲染质量；只有超过阈值的高斯被拉回球形。

（2）\textbf{深度不确定性模型}。CG-SLAM 从两个视角建立不确定性建模——渲染图像层面和高斯基元层面。

在\textbf{渲染图像层面}，$\alpha$-混合 深度本质上是对沿像素射线采样到的高斯椭球计算的期望值。在正态分布假设下，其不确定性与方差点点对应。CG-SLAM 利用 $\alpha$-混合 深度的方差渲染不确定性图：

\begin{equation}
    U = \sum_{i=1}^{N} \alpha_i T_i (d_i - D)^2
\end{equation}

其中 $D$ 为深度传感器观测值。不确定性值 $U$ 衡量的是渲染深度与传感器深度之间的加权离散程度——$U$ 大说明该像素附近的高斯分布分散，几何不收敛。为了抑制这种分散，CG-SLAM 进一步引入几何方差损失，将整张不确定性图的均值作为正则项：

\begin{equation}
    \mathcal{L}_\text{var} = \frac{1}{HW} \sum_{n=1}^{HW} |U_n|
\end{equation}

这个损失迫使高斯基元的位置在优化过程中不会远离真实深度，避免几何剧烈抖动。

在\textbf{高斯基元层面}，CG-SLAM 首先引入一个关键观察：单独使用 $\alpha$-混合 深度无法有效约束高斯基元的空间位置。原因是 $\alpha$-混合 深度是多个高斯的加权平均，某个高斯位置的偏移可能被其他高斯补偿，系统仍能产生正确的渲染深度。为此，CG-SLAM 同时渲染\textbf{中值深度} $D_\text{median}$——在像素的 splatting 过程中，累积透射率 $T$ 首次降至 0.5 以下时对应的高斯即为"中值高斯"，其深度记为中值深度。然后利用 $\alpha$-混合 深度与中值深度之间的对齐损失，确保两种深度估计一致：

\begin{equation}
    \mathcal{L}_\text{align} = \frac{1}{HW} \sum_{n=1}^{HW} |D_\alpha^n - D_\text{median}^n|
\end{equation}

对齐损失的直觉是：$\alpha$-混合 深度由权重最大的高斯主导，对齐损失迫使该"主导高斯"恰好出现在中值深度处，从而推动所有高斯基元聚集在场景表面附近，而非分散在空间中。

在此基础上，CG-SLAM 为每个高斯基元 $i$ 计算跨关键帧窗口 $F = \{f_1, f_2, ..., f_k\}$ 的深度不确定性 $v_i$。这里需要引入\textbf{支配像素}（dominated pixel）的概念：对于某个像素，其 $\alpha$-混合 深度由权重最大的高斯主导，该像素即为这个高斯的支配像素。高斯基元 $i$ 的不确定性由其在所有关键帧中支配像素处的深度差异聚合而成。

\begin{equation}
    v_i = \frac{1}{M_1 + ... + M_k} \sum_{f_k \in F} \sum_{p=1}^{M_k} \alpha_i^{k,p} T_i^{k,p} (D_p^k - d_i^k)^2
\end{equation}

其中 $M_k$ 为高斯 $i$ 在关键帧 $f_k$ 中的支配像素数目，$\alpha_i^{k,p}$ 和 $T_i^{k,p}$ 分别为高斯 $i$ 在像素 $p$ 上的不透明度贡献和透射率，$D_p^k$ 为传感器深度观测值，$d_i^k$ 为高斯 $i$ 在该视角下的深度值。$v_i$ 本质上度量高斯 $i$ 在多个视角下与真实表面的一致程度——值越大说明该高斯越不可靠。CG-SLAM 定期筛除不确定性超过阈值 $\tau=0.025$ 的不可靠基元：先将它们降为低透明度，给予重新优化的机会；若在后续迭代中仍无法恢复，则彻底移除。这种"两步淘汰"策略比直接删除更加自适应，持续提升高斯场中有效基元的比例。

（3）\textbf{位姿导数解析与高效跟踪}。在标准 3DGS 中，相机位姿（旋转和平移）并非可优化的参数——原始 3DGS 假设位姿已知，专注于优化高斯属性。CG-SLAM 首次对 EWA（椭圆加权平均，Elliptical Weighted Average）splatting 框架中相机位姿关于渲染损失的导数进行了完整的数学分析（详见其补充材料），并实验发现使用\textbf{李代数表示}（Lie algebra $\mathfrak{so}(3) | \mathbf{T}$）进行参数化比直接参数化旋转矩阵更有利于相机跟踪，尤其在旋转分量的优化上收敛更快。基于这一理论分析，CG-SLAM 开发了一套定制化 CUDA 框架，有效解耦跟踪与建图。

跟踪模块包含两部分。

\textbf{序列跟踪}：固定场景表示，通过恒速假设（将上一帧的位姿按上一帧的相对变换外推）初始化当前位姿，然后以光度损失 $\mathcal{L}_\text{color}$ 和几何损失 $\mathcal{L}_\text{geo}$ 的加权组合 $\mathcal{L}_\text{tracking} = \lambda_1 \mathcal{L}_\text{color} + \lambda_2 \mathcal{L}_\text{geo}$ 进行梯度下降优化位姿（Replica 数据集上 $\lambda_1=0.2, \lambda_2=1.0$，强调几何约束）。

\textbf{滑动窗口光束法平差（Bundle Adjustment, BA）}：为缓解 SLAM 中典型的累积漂移问题，CG-SLAM 设置包含 $k=4$ 个共视关键帧的滑动窗口，联合优化窗口内的相机外参和高斯属性。共视性判定方面，CG-SLAM 用预训练的 NetVLAD（网络向量局部聚合描述符，Network Vector of Locally Aggregated Descriptors）描述子编码关键帧，通过余弦相似度而非先前工作中的视锥重叠方法确定共视帧，更加高效准确。

\textbf{建图总损失}。CG-SLAM 的建图优化使用六项损失的加权组合：

\begin{equation}
    \mathcal{L}_\text{mapping} = \underbrace{\lambda_1 \mathcal{L}_\text{color} + \lambda_2 \mathcal{L}_\text{ssim} + \lambda_3 \mathcal{L}_\text{geo}}_{\text{标准渲染与几何约束}} + \lambda_4 \mathcal{L}_\text{align} + \lambda_5 \mathcal{L}_\text{iso} + \lambda_6 \mathcal{L}_\text{var}
\end{equation}

其中前三项（颜色 $\mathcal{L}_\text{color}$、SSIM $\mathcal{L}_\text{ssim}$、几何 $\mathcal{L}_\text{geo} = \frac{1}{HW}\sum|D_\alpha^n - D_n|$）是标准的光度-深度监督；后三项（对齐、各向同性、方差）是 CG-SLAM 针对高斯场一致性和稳定性新增的正则化项。在 Replica 数据集上，权重设置为 $\lambda_{1..6} = \{0.7, 0.1, 0.25, 0.25, 0.1, 0.15\}$，建图迭代 60 次。

\textbf{高斯管理}。初始化时基于第一帧的深度观测密集投影高斯基元到 3D 空间。后续建图中，设置阈值 $\eta=0.5$ 检测累积不透明度 $O < \eta$ 的像素，视为未观测或未充分构建区域，在这些位置利用颜色和深度信息生成新的高斯基元。同时继承了原始 3DGS 的分裂、克隆和移除策略进行高斯致密化。

就关键参数而言，CG-SLAM 在 NVIDIA RTX 4090 上以约 8.5 Hz 运行（轻量版使用半分辨率图像可进一步提升至约 15 Hz 跟踪速度），在 Replica 数据集上跟踪精度超过所有基线方法 25\%--75\%，同时在重建精度和运行效率上达到或超越当时的最佳 NeRF-SLAM 方法。

\subsubsection{抗锯齿与不确定性建模}

本路线介绍通过抗锯齿渲染和传感器噪声的不确定性建模来增强 3DGS-SLAM 鲁棒性的方法。

几何表示的改进解决了表面重建的精度问题，但 3DGS 的渲染分辨率敏感性和传感器噪声仍是影响 SLAM 鲁棒性的重要因素。本小节讨论如何通过抗锯齿和不确定性建模来增强系统的稳定性。

MipSLAM（2026.03）~\cite{2603_06989} 修改了标准 3DGS 的\textbf{像素级高斯贡献计算}——用椭圆自适应数值积分替代像素中心点采样——$\alpha$-混合 渲染管线和深度渲染保持不变。

\textbf{动机。}标准 3DGS 在像素中心做点采样来评估高斯贡献，当相机分辨率、焦距或缩放发生变化时，这种点采样违反 Nyquist-Shannon 采样定理，产生锯齿伪影；在 SLAM 的在线场景中，高斯基元持续投影且相机参数可能变化，锯齿会直接导致位姿估计漂移。该方法系统性地解决了 3DGS-SLAM 在相机参数（分辨率、焦距、缩放）变化时的锯齿与漂移问题，包含 EAA 抗锯齿和 SA-PGO 位姿图优化两个核心模块。

\textbf{（1）EAA：椭圆自适应抗锯齿}。

\textit{问题}。标准 3DGS 在计算像素 $u$ 处的高斯贡献时，直接在像素中心点处评估投影 2D 高斯 $G_i^{2D}(u)$，代入 $\alpha$-混合 公式得到颜色 $C(u) = \sum_{i=1}^{N} c_i \alpha_i G_i^{2D}(u) \prod_{j=1}^{i-1}(1 - \alpha_j G_j^{2D}(u))$ 和深度 $\hat{D}(u) = \sum_{i=1}^{N} z_i \alpha_i G_i^{2D}(u) \prod_{j=1}^{i-1}(1 - \alpha_j G_j^{2D}(u))$。这种点采样在相机分辨率变化时违反 Nyquist-Shannon 采样定理，产生锯齿伪影。直接引入 MipSplatting 的 3D 滤波或 Analytic-Splatting 的解析积分虽然可以缓解锯齿，但前者需要全局重训练高斯图、后者计算代价过高，均不适用于 SLAM 中数千高斯基元持续投影的在线场景。

\textit{方案}。EAA 的核心思想是：不在像素中心做一次点评估，而是在像素覆盖区域内沿高斯的等概率椭圆轮廓进行多次重要性加权采样，用数值积分逼近解析积分。具体分三步实现。

\textit{坐标系变换}。首先将积分域从像素矩形变换到投影高斯的固有椭圆坐标系。对投影 2D 协方差 $\Sigma^{2D}$ 做特征值分解 $\Sigma^{2D} = Q^T \Lambda Q$，其中 $\Lambda = \mathrm{diag}(\lambda_1, \lambda_2)$ 为特征值矩阵，$Q = [\mathbf{e}_1, \mathbf{e}_2]$ 为正交特征向量矩阵。直觉上，这相当于将像素空间"旋转+缩放"到与高斯椭圆主轴对齐的局部坐标系——在此坐标系下，等概率轮廓线就是同心椭圆，采样点沿这些椭圆分布最为均匀。

\textit{重要性加权数值积分}。在主轴坐标系下，对每个采样点 $\mathbf{x}_k$ 计算主轴坐标 $\mathbf{v}_k = Q^T(\mathbf{x}_k - \boldsymbol{\mu}^{2D})$，然后用重要性加权求和替代点采样：

\begin{equation}
    \tilde{\alpha}_i = \alpha_i \frac{\sum_{k=1}^{K} w_k \exp\!\left(-\tfrac{1}{2} \mathbf{v}_k^T \Lambda^{-1} \mathbf{v}_k\right)}{\sum_{k=1}^{K} w_k}
\end{equation}

其中 $K$ 为采样点数，$w_k$ 为重要性权重。这一公式直接替换标准 3DGS 中颜色和深度渲染公式里的 $\alpha_i G_i^{2D}(u)$，其余 $\alpha$-混合 结构不变。直觉上，$\tilde{\alpha}_i$ 度量的是高斯 $i$ 对像素 $u$ 覆盖区域的\textit{总体}贡献，而非仅在该像素中心点的瞬时值。

\textit{自适应采样密度}。各向异性高斯（一个方向极度拉伸而另一方向很窄的高斯）的条件数 $\kappa = \lambda_{\max}/\lambda_{\min}$ 很大，需要更密集的采样才能准确近似积分。EAA 自适应地增加高条件数高斯的采样密度，同时通过距离场检测像素边界，在边界附近增加采样权重以捕捉不连续性。增量 3D 滤波模块在滑动关键帧窗口内维护每个高斯的观测频率，确保抗锯齿效果随在线建图逐步完善。

\textbf{（2）SA-PGO：频谱感知位姿图优化}。

\textit{问题}。传统 PGO 在空间域独立优化各位姿节点，仅依赖几何残差约束相邻位姿，无法检测和抑制轨迹中的高频振荡（表现为相机位姿的快速抖动），这种振荡通常源于 SLAM 前端的累积漂移或传感器噪声。

\textit{方案}。SA-PGO 将轨迹视为时空信号，在频域中分析其行为，并利用图拉普拉斯谱分解引导优化。

\textit{频域轨迹分析}。对位姿序列 $\{p_1, \ldots, p_N\}$（$p_i = [\mathbf{t}_i^T, \boldsymbol{\theta}_i^T]^T \in \mathbb{R}^6$ 包含平移和轴角旋转），计算滑动窗口 DFT：

\begin{equation}
    \mathcal{F}_k^{(d)}(\omega) = \sum_{n=1}^{N_w} p_{k+n-1}^{(d)} \, w(n) \, e^{-j 2\pi \omega (n-1)/N_w}
\end{equation}

其中 $N_w$ 为窗口大小，$\omega$ 为离散频率索引，$w(n)$ 为 Hanning 窗函数。由此提取\textbf{频率签名}——加权频率质心 $S_k = \sum_{d,\omega} \omega \cdot |\mathcal{F}_k^{(d)}(\omega)|^2 / \sum_{d,\omega} |\mathcal{F}_k^{(d)}(\omega)|^2 + \epsilon_\text{reg}$。质心越大表示高频能量越多，即该窗口内位姿抖动越剧烈，越可能有漂移或噪声。

\textit{谱正则化位姿图优化}。SA-PGO 构建加权邻接矩阵 $A_{ij} = S_{ij}$（$S_{ij}$ 综合了频率相干性和几何正则性），计算归一化图拉普拉斯 $L_\text{norm} = D^{-1/2}(D - A)D^{-1/2}$，其 Fiedler 特征值 $\lambda_2^L$ 度量位姿图的代数连通性——$\lambda_2^L$ 越大，图越连通，优化越稳定。最终优化目标为：

\begin{equation}
    \mathbf{x}^* = \arg\min_{\mathbf{x}} \left(E_\text{geo}(\mathbf{x}) + \lambda_\text{spe} R_\text{spe}(\mathbf{x}) + \lambda_\text{smo} R_\text{smo}(\mathbf{x})\right)
\end{equation}

其中 $E_\text{geo}$ 为标准 PGO 几何残差（由谱置信度 $S_{ij}$ 加权），$R_\text{spe}$ 惩罚谱相似位姿间的不一致性，$R_\text{smo}$ 在时域惩罚高频振荡。$\lambda_\text{spe}$ 和 $\lambda_\text{smo}$ 根据 Fiedler 特征值自适应调整：$\lambda_2^L$ 大（连通性好）时增强谱正则化，小时退回到纯几何优化以保持稳定性。

\textit{整体损失}。跟踪与建图通过光度损失和深度损失联合优化：$\mathcal{L} = \lambda_\text{rgb}\|\hat{C}(u) - I_\text{gt}(u)\|_2^2 + \lambda_\text{depth}\|\hat{D}(u) - D_\text{gt}(u)\|_2^2$，其中颜色和深度渲染中的高斯贡献已替换为 EAA 的 $\tilde{\alpha}_i$。SA-PGO 在后端异步运行，周期性精化全局位姿图。MipSLAM 是首个支持任意相机参数重配置的频率感知 3DGS-SLAM 系统，在多个分辨率下保持高保真渲染，同时通过频域正则化提升跟踪精度。

VarSplat（2026.03）~\cite{2603_09673} 沿用 MipSLAM 所在的子地图式 3DGS-SLAM 管线（场景划分为多个子地图，每个子地图包含一组 3D 高斯基元），但切换了鲁棒性增强的切入角度。

\textbf{动机。}在 SLAM 的在线增量建图中，高斯基元的参数误差不可避免——低纹理区域、反射表面和深度不连续处的高斯尤其不可靠，但这些区域在标准 3DGS 中与其他区域被同等对待，不可靠观测的梯度会误导位姿优化和地图更新。MipSLAM 在像素级\textbf{贡献计算}上做文章（用椭圆自适应积分替代点采样），而 VarSplat 则在标准 $\alpha$-混合 渲染之上为每个高斯基元增加一个\textbf{外观方差}参数 $\sigma_i^2$，通过全方差定律（Law of Total Variance）将其合成为逐像素不确定性图 $V$，从而让系统能"知道自己不知道什么"——低纹理区域、反射表面、深度不连续处自然产生高不确定性。该方法系统性地解决了 3DGS-SLAM 在不可靠观测下的位姿漂移与建图退化问题。

\textbf{（1）外观方差的直觉}。

标准 3DGS 中每个高斯基元 $G_i$ 由位置 $\boldsymbol{\mu}_i$、不透明度 $\alpha_i$、尺度 $\mathbf{s}_i$、协方差矩阵 $\Sigma_i$ 和球谐系数 $c_i$（编码平均颜色）定义。VarSplat 在此基础上增加外观方差 $\sigma_i^2 \in \mathbb{R}^3$——注意它\textit{不是}定义几何范围的协方差 $\Sigma_i$，也\textit{不是}定义平均外观的球谐系数 $c_i$，而是建模围绕该平均颜色的观测离散程度。直觉上，即使球谐颜色正确建模了平均外观，不同视点下的实际颜色观测仍可能因遮挡边界变化、半透明表面混合等因素产生波动；$\sigma_i^2$ 正是量化这种波动的统计量。在深度不连续处和遮挡边界附近，微小的视角变化会改变重叠高斯的可见性和 $\alpha$ 权重，导致颜色观测不一致。因此这些区域自然学到更大的 $\sigma_i^2$。

\textbf{（2）方差渲染：全方差定律 $+$ $\alpha$ 合成}。

VarSplat 的核心创新在于：不依赖额外的预训练不确定性预测器，而是直接利用 3DGS 光栅化器已有的 $\alpha$-混合 管线，通过\textbf{全方差定律}将逐基元方差合成为逐像素方差。对于随机变量 $X$（像素颜色）和条件变量 $Z$（3D 高斯），全方差定律给出：

\begin{equation}
    \mathrm{Var}[X] = \underbrace{\mathbb{E}[\mathrm{Var}[X \mid Z]]}_{\text{期望逐基元方差}} + \underbrace{\mathrm{Var}(\mathbb{E}[X \mid Z])}_{\text{基元间颜色散度}}
\end{equation}

直觉上，逐像素不确定性来自两个来源：一是各高斯自身的观测波动（$\sigma_i^2$ 项），二是多个高斯在颜色上的不一致（$c_i$ 项）——即使每个高斯的方差为零，不同高斯在颜色上仍有差异，这种差异本身也贡献不确定性。

具体地，在标准 $\alpha$-混合 权重 $w_i = \alpha_i \prod_{j=1}^{i-1}(1-\alpha_j)$ 下，逐像素方差 $V$ 通过单遍光栅化渲染得到：

\begin{equation}
    V = \sum_{i} w_i(\sigma_i^2 + c_i^2) - \left(\sum_{i} w_i c_i\right)^2
\end{equation}

该公式与标准颜色渲染 $C = \sum_i w_i c_i$ 和深度渲染 $D = \sum_i w_i z_i$ 共享完全相同的光栅化通道，无需额外的网络推理或多次前向传播。这是 VarSplat 与依赖预训练不确定性预测器的方法（如 WildGS-SLAM 从 DINOv2 特征预测不确定性图）的根本区别：不确定性是从渲染过程本身涌现的，而非外部注入的。

\textbf{（3）端到端方差学习：从零优化，非预训练}。

VarSplat 将 $\sigma_i^2$ 作为每个高斯的可微参数，与位置、不透明度、尺度、球谐系数联合优化。与 ActiveNeRF 等通过神经网络预测方差的方法不同，VarSplat 直接在 3DGS 光栅化器中优化标量方差参数。建图阶段的损失函数为：

\begin{equation}
    \mathcal{L} = \lambda_\text{color} \mathcal{L}_\text{color} + \lambda_\text{depth} \mathcal{L}_\text{depth} + \lambda_\text{reg} \mathcal{L}_\text{reg} + \lambda_\text{var} \mathcal{L}_\text{var}
\end{equation}

其中颜色损失 $\mathcal{L}_\text{color}$ 采用 L1 与 SSIM 的加权组合，深度损失 $\mathcal{L}_\text{depth}$ 为 L1，正则化 $\mathcal{L}_\text{reg}$ 控制高斯尺度。方差损失 $\mathcal{L}_\text{var}$ 基于高斯负对数似然设计：

\begin{equation}
    \mathcal{L}_\text{var} = \frac{1}{2V}\left(\|\hat{I} - I\|_2^2 + \|\hat{D} - D\|_2^2\right) + \log(V)
\end{equation}

直觉上，当渲染残差（$\|\hat{I} - I\|_2^2$）大时，优化器会增大 $V$（即增大概率/方差），使第一项不至于过大；而 $\log(V)$ 项则防止方差无限增长——它惩罚过大的不确定性，迫使系统在观测充分时收敛到低方差。梯度通过链式法则回传到每个高斯的 $\sigma_i^2$：$\partial \mathcal{L}_\text{var}/\partial \sigma_i^2 = (\partial \mathcal{L}_\text{var}/\partial V) \cdot w_i$，即第 $i$ 个高斯的方差梯度正比于其 $\alpha$-混合 权重——对渲染贡献越大的高斯，其方差更新越积极。

\textbf{（4）不确定性引导的下游任务}。

渲染得到的逐像素不确定性图 $V$ 和逐基元方差 $\sigma_i^2$ 同时服务于三个下游模块。

\textit{跟踪}。对输入帧 $(I, D)$，跟踪损失为 $\mathcal{L}_\text{track} = \lambda_c \sum_p w_p \|\hat{I} - I\|_1 + (1-\lambda_c)\|\hat{D} - D\|_1$，其中 $w_p$ 是经对数缩放归一化的方差权重：高不确定性像素的权重被降低，避免低纹理、反射等不可靠区域主导梯度方向。跟踪时方差参数冻结，仅优化位姿。

\textit{子地图配准}。检测到回环后，VarSplat 将查询帧定位到数据库子地图中。配准损失使用与跟踪相同的方差权重 $w_p$ 调制光度残差，仅在可靠区域计算配准残差，避免深度不连续处的不可靠观测干扰位姿变换估计。

\textit{回环检测}。VarSplat 利用逐基元方差 $\sigma_i^2$ 计算子地图级可靠度比率 $r = \sum_j \{w_s \cdot j\} / \sum_j j$，即在应用方差加权前后不透明度之和的比值。该比率编码了子地图中"有多少外观信息是可靠的"，用于调制查询子地图与数据库子地图之间的相似度评分，使几何稳定的子地图获得更高的匹配优先级。

\textit{总结}。与 MipSLAM 在渲染层面修改高斯贡献计算不同，VarSplat 保留标准 $\alpha$-混合 管线，在其上叠加方差渲染层。两者互补：MipSLAM 消除分辨率变化带来的锯齿与漂移，VarSplat 量化观测不可靠性并引导下游优化聚焦于可靠区域。VarSplat 是首个在 3DGS-SLAM 中从零端到端学习逐基元外观方差、并通过全方差定律渲染不确定性图的系统，在 Replica、TUM-RGBD、ScanNet 和 ScanNet++ 上实现了鲁棒的跟踪和建图性能。

PointSLAM++（2026.01）~\cite{2601_11617} 与 VarSplat 虽然同属锚点范式，但切入角度截然不同。

\textbf{动机。}VarSplat 在标准 3D 高斯上叠加方差参数量化观测不确定性，但高斯数量的增长仍与像素数成正比，内存压力持续增大；而 Scaffold-GS 风格的锚点表示通过 MLP 解码生成子高斯属性，可以更少的锚点管理更大区域的高斯分布，天然适合内存受限的在线 SLAM 场景。VarSplat 在标准 3D 高斯上叠加方差参数以量化观测不确定性；而 PointSLAM++ 采用 Scaffold-GS 风格的\textbf{神经锚点}表示——每个锚点通过 MLP 解码生成 $k$ 个子高斯的属性（不透明度、颜色、四元数旋转、尺度），并将视角方向嵌入锚点外观特征以补偿视角相关的外观变化。

这一表示差异带来了建图策略的根本不同：PointSLAM++ 引入\textbf{层次化锚点}机制，将 ORB 特征生成的锚点标记为"主锚点"（不可分裂/删除，小学习率优化），以保持跟踪稳定性；同时在梯度超过阈值 $\bar{g} > \theta_g$ 的体素中从深度数据生成"次锚点"，以补充纹理稀缺区域的几何表达。在位姿估计方面，PointSLAM++ 提出级联式\textbf{渐进位姿优化}（Progressive Pose Optimization, PPO）：以 GICP 点云配准提供粗估计，再由 ORB 特征点云配准精化，最后通过 BA 联合优化位姿与路标点。这种多阶段设计使系统能在深度噪声下保持鲁棒——在 TUM-RGBD 上跟踪误差较纯 GICP 方法降低超过 50\%；在 ScanNet++（深度误差大、相机位姿间隔大）上，GS-ICP SLAM 等依赖 GICP 的方法无法完成定位，而 PointSLAM++ 通过 ORB 特征与多阶段精化维持稳定跟踪。

\subsection{鲁棒性与实时加速}
\label{ss:15_2}

本节介绍两条技术路线：（1）初始化与加速策略；（2）特征增强与鲁棒跟踪。

\subsubsection{初始化与加速策略}

本路线介绍提升高斯基元初始化质量和系统计算效率的方法，包括 SfM-free 初始化、像素采样策略和 GPU 加速方案。

抗锯齿和不确定性建模提升了系统鲁棒性，但高斯基元的初始化质量和计算效率仍然是实际部署的关键瓶颈。本小节讨论如何通过更好的初始化和算法-硬件协同设计加速收敛。

SGAD-SLAM（2026.03）~\cite{2603_21055} 在 SplaTAM 的简化球状高斯表示基础上做了两个关键改动。

\textbf{动机。}SplaTAM 的逐像素致密化机制在每像素仅使用一个各向同性球状高斯时几何表达能力不足，而放宽为可自由移动的 3D 高斯虽然表达力更强，但在在线 SLAM 的稀疏视角下容易过拟合。第一，引入\textbf{深度偏移量} $\delta_i \in \mathbb{R}^1$——每个像素对齐高斯的 3D 位置不再固定于观测深度 $D_i(j)$，而是沿射线偏移至 $D_i(j) + \delta_i$。这一自由度使像素对齐高斯在省略致密化（每像素仅一个高斯）的严格约束下仍能逼近真实几何，论文消融表明即使与可全场景自由移动的 3DGS 相比，偏移量方案的渲染质量也更高。第二，跟踪环节\textbf{完全绕过可微渲染}——将深度图下采样为 3D 点集，对每个点用邻域协方差建模局部几何（即高斯分布），再通过 GICP（Generalized ICP）将当前帧的几何分布与全局累积分布匹配，直接求解位姿。相比 SplaTAM 的轮廓感知光度-深度损失梯度下降，几何匹配对纹理缺失和大运动更鲁棒，且天然支持并行加速。

从"绕过可微渲染"回到可微渲染管线内部，RGS-SLAM（2026.01）~\cite{2601_00705} 保留 SplaTAM 的可微渲染管线与关键帧管理框架，但将核心致密化机制从"逐帧残差驱动增密"替换为"DINOv3 密集特征匹配 $+$ 多视角三角化的一次性初始化"。

\textbf{动机。}SplaTAM 的残差驱动致密化在稀疏视角下收敛缓慢且容易产生空间不一致——每帧仅根据当前渲染残差局部添加高斯，缺乏全局视野；而基于语义特征的三角化能为后续光度优化提供更稳定、覆盖更均匀的初始高斯分布。具体而言，每个关键帧触发时，用 DINOv3 描述子建立密集多视角对应，经置信度过滤后批量三角化出约 1000 个各向异性椭球高斯种子（而非 SplaTAM 的各向同性球状高斯），一次插入地图，后续仅做轻量合并与剪枝，不再迭代增密。这一设计的动机是：残差驱动增密在稀疏视角下收敛缓慢且容易产生空间不一致，而基于语义特征的三角化能为后续光度优化提供更稳定、覆盖更均匀的初始高斯分布。在跟踪端，RGS-SLAM 放弃了 SplaTAM 的轮廓掩码和深度损失，仅用光度损失（加仿射亮度模型补偿曝光变化），并采用解析 $SE(3)$ Jacobian 替代通用自动微分加速位姿更新。在 TUM RGB-D 上，训练时间从基线的 14.8 分钟降至 12 分钟（减少约 20\%），同时保持相当的定位精度和渲染质量。

实时加速方面，RTGS（2025.10）~\cite{2510_06644} 从 SplaTAM 的观察出发：跟踪过程中仅 14\% 的高斯贡献了绝大部分位姿梯度，且非关键帧以全分辨率处理带来大量冗余计算。

\textbf{动机。}3DGS-SLAM 在跟踪阶段对全部高斯执行完整渲染和反向传播，但绝大多数高斯对位姿梯度的贡献可忽略不计；同时，非关键帧与关键帧使用相同分辨率处理，造成大量冗余计算，严重限制了系统在边缘 GPU 上的实时性能。RTGS 在不改变 SplaTAM 高斯表示和渲染管线的前提下，引入两层冗余消除。

\textbf{自适应高斯剪枝}利用跟踪过程已有的位置梯度 $\nabla_{\boldsymbol{\mu}}L$ 与协方差梯度 $\nabla_{\Sigma}L$ 计算重要性分数 $s = \|\nabla_{\boldsymbol{\mu}}L\|_2 + \beta\|\nabla_{\Sigma}L\|_2$（复用跟踪阶段的梯度，不引入额外开销），采用 mask-then-prune 策略：先在 $K$ 次迭代中屏蔽低重要性高斯（排除渲染但仍保留以计算 tile-Gaussian 变化率），第 $K+1$ 次迭代时永久移除。剪枝间隔 $K$ 根据 tile-Gaussian 交叉变化率动态调整（变化率超过 5\% 则缩短至 $K_0/2$，否则延长至 $2K_0$）。剪枝上限设为 50\%，超过此值 ATE 急剧上升。

\textbf{动态下采样}将关键帧保持原始分辨率 $R_0$，非关键帧则根据距最近关键帧的时序距离逐步降采样：紧跟关键帧的非关键帧降至 $R_0/16$，后续每帧以因子 $m=2$ 递增，最高不超过 $R_0/4$，保证分辨率变化的平滑过渡以避免轨迹漂移。算法层面在三个基线 SLAM（GS-SLAM、MonoGS、Photo-SLAM）上均实现 2.5$\times$--3.6$\times$ 加速，ATE 和 PSNR 下降不超过 5\%。RTGS 还设计了 28\,nm 专用硬件插件（渲染引擎、梯度合并单元、预处理引擎、位姿计算单元），作为 GPU 旁路模块加速渲染和反向传播，在边缘 GPU（NVIDIA Jetson AGX Orin）上达到 $\geq$ 30 FPS 的实时性能，四数据集平均能效提升 32.7$\times$--73.0$\times$。

FlashSLAM（2024.12）~\cite{2412_00682} 修改了 3DGS-SLAM 的相机位姿估计环节——用预训练特征匹配 $+$ SVD 刚体变换替代恒速假设 $+$ 梯度下降优化——高斯表示和建图流程保持不变。

\textbf{动机。}SplaTAM 等基线方法通过恒速假设初始化位姿再用 100--200 次迭代梯度下降精化，这种优化范式的收敛域狭窄——当帧间运动量大时初始位姿偏离真实值过远，梯度容易陷入局部最优，且每次迭代都需要完整的渲染-反向传播过程，速度缓慢。

\textbf{问题：恒速假设 $+$ 梯度优化的固有局限。}在 SplaTAM、MonoGS 等基线做法中，相机位姿通过以下流程估计：先用恒速假设（将上一帧位姿按前一帧的相对变换外推）初始化当前帧位姿，再通过可微渲染的光度-深度联合损失 $\mathcal{L} = \mathcal{L}_\text{color} + (1-\lambda)\mathcal{L}_\text{depth}$ 进行 100--200 次迭代梯度下降优化。这一范式有两个根本缺陷。第一，梯度下降的收敛域狭窄——当帧间运动量大时（稀疏设置下帧间距可能跨越数十帧），初始位姿偏离真实值过远，梯度容易陷入局部最优。第二，每次迭代都需要完整的渲染-反向传播过程，稠密设置下单帧跟踪耗时约 1 秒。在稀疏设置下问题更加严重：恒速假设完全失效，现有方法甚至无法完成跟踪。

\textbf{方案：特征匹配 $+$ SVD 配准的两阶段跟踪。}FlashSLAM 的核心改进是将位姿估计从"优化"范式切换为"匹配 $+$ 求解"范式。直觉上，如果能在两帧之间建立可靠的像素对应关系，就可以直接通过几何求解得到位姿，无需任何迭代优化。前端使用 SuperPoint 检测关键点并提取局部描述子，通过 LightGlue 进行帧间特征匹配建立 2D-2D 对应 $(p_{t-1}, p_t)$，随后利用深度信息将 2D 匹配反投影为两组 3D 点集 $P_{t-1}$ 和 $P_t$。位姿估计转化为经典的最小二乘刚体变换问题——寻找旋转 $\mathbf{R}$ 和平移 $\mathbf{t}$，使两组对应点之间的距离最小：

\begin{equation}
    \min_{\mathbf{R}, \mathbf{t}} \sum_{i=1}^{N} \| \mathbf{p}_t^{(i)} - (\mathbf{R} \mathbf{p}_{t-1}^{(i)} + \mathbf{t}) \|^2
\end{equation}

直觉上，这等价于在 3D 空间中找到最佳旋转和平移使两组点云"对齐"。FlashSLAM 通过 SVD（Singular Value Decomposition，奇异值分解）分解闭式求解：先计算两组点集的质心 $\mathbf{c}_{t-1}$、$\mathbf{c}_t$，去质心后构建协方差矩阵 $\mathbf{H} = \mathbf{Q}_{t-1}^\top \mathbf{Q}_t$，对其做 SVD 分解 $\mathbf{H} = \mathbf{U}\boldsymbol{\Sigma}\mathbf{V}^\top$，则最优旋转为 $\mathbf{R} = \mathbf{V}\mathbf{U}^\top$，平移为 $\mathbf{t} = \mathbf{c}_t - \mathbf{R}\mathbf{c}_{t-1}$。这种基于特征匹配的跟踪方式天然具有更宽的收敛域——LightGlue 的特征匹配不受恒速假设约束，即使帧间运动量大也能建立可靠对应，并且 SVD 闭式求解无需迭代，速度远快于梯度优化。

\textbf{抗深度噪声设计。}消费级深度传感器（如 iPhone 的 LiDAR）存在显著噪声。FlashSLAM 引入\textbf{动态深度截断}策略：对每一帧，将深度分布的第 70 百分位作为截断阈值，过滤掉过远的不可靠深度点。阈值设为 70\% 而非更低的百分位，是为了保留足够多的有效匹配点，同时排除深度噪声导致的离群点。在建图阶段添加新高斯时，FlashSLAM 先对已映射点集 $P_o$ 与可见高斯集 $G_v$ 执行 ICP 对齐，仅当对齐适应度 $f > 0.2$ 且误差 $e < 0.1$ 时才接受变换 $(\mathbf{R}, \mathbf{t})$ 并将其应用于校正新点集 $P_n$ 的位置，防止深度噪声导致高斯放置偏差。

\textbf{可选梯度精化与 Loss-Weighted Sampling。}在特征匹配 $+$ SVD 配准的基础上，FlashSLAM 提供可选的梯度精化阶段，对渲染图像与输入图像计算光度损失，仅需 30--70 次迭代（而非基线的 100--200 次）。关键设计是 Loss-Weighted Sampling 采样策略：关键帧 $i$ 被选中的概率与其损失成正比，即 $P(i) = L_i / \sum_{j=1}^{N} L_j$，使得渲染质量最差的关键帧获得更多优化机会。为避免"遗忘"已经对齐良好的帧，该方法以概率 $p = 0.4$ 执行 Loss-Weighted Sampling，以概率 $1-p = 0.6$ 执行均匀随机采样，在局部精化与全局一致性之间取得平衡。

\textbf{效果。}消融实验显示：不加精化时跟踪速度 78 ms/帧、ATE 0.65 cm、PSNR 35.32 dB；50 次精化迭代时 ATE 降至 0.55 cm、PSNR 升至 39.21 dB，耗时 485 ms/帧。在 TUM fr3 场景上，FlashSLAM 的跟踪时间比 SplaTAM 减少 90\%；在 iPhone 13 Pro Max 自采集的噪声数据上，SplaTAM 和 MonoGS 完全跟踪失败，FlashSLAM 仍能准确重建场景，展示了其在消费级传感器上的鲁棒性。在 Replica 数据集上 FlashSLAM 达到平均 PSNR 38.45、SSIM 0.978、LPIPS（Learned Perceptual Image Patch Similarity）0.042，帧率达到 899 FPS。

\subsubsection{特征增强与鲁棒跟踪}

本路线介绍通过密集特征提取和鲁棒核函数增强 3DGS-SLAM 在恶劣条件下跟踪稳定性的方法。

初始化和加速策略解决了效率问题，但在纹理缺失、低光照等恶劣条件下仍需增强跟踪鲁棒性。本小节讨论通过密集特征和鲁棒核函数提升跟踪质量的方法。

FeatureSLAM（2026.01）~\cite{2601_05738} 沿着 FlashSLAM"特征辅助跟踪"的方向进一步深入。

\textbf{动机。}FlashSLAM 仅在前端用特征匹配建立对应关系，建图阶段仍然回归纯光度-深度优化，在低纹理或重复纹理区域仍可能因匹配不足而失败；同时，已有的语义增强 3DGS 方法依赖重型 ViT/DINO 骨干网络进行离线特征蒸馏，无法满足在线 SLAM 的实时性要求。FeatureSLAM 则将语义特征贯穿建图和跟踪的全流程，用 SAM2 多尺度金字塔特征替代手工设计的局部描述子，同时引入结构感知深度光栅化和几何正则化，在实现实时在线 SLAM 的同时支持开放集语义分割。

\textbf{动机：纯几何跟踪的脆弱性与特征蒸馏的离线瓶颈。}FlashSLAM 证明了特征匹配可以替代梯度优化实现鲁棒跟踪，但它使用的 SuperPoint+LightGlue 仅提供稀疏局部描述子，在低纹理或重复纹理区域仍可能因匹配不足而失败。另一方面，已有的语义增强 3DGS 方法（如 Feature3DGS、LangSplat）依赖 ViT/DINO 等重型骨干网络进行离线特征蒸馏，无法满足在线 SLAM 的实时性要求。FeatureSLAM 的核心思路是：选择速度足够快的基础模型编码器，将多尺度语义特征直接嵌入每个高斯，使特征在建图和跟踪中同时发挥作用，而非作为离线后处理步骤。

\textbf{SAM2 特征金字塔与在线蒸馏。}FeatureSLAM 采用 SAM2 的 Hiera-MAE 编码器来提取多尺度特征金字塔——三个层级分别产生 256/64/32 通道、分辨率为 $64{\times}64$/$128{\times}128$/$256{\times}256$ 的特征图。直觉上，高分辨率层级捕获精细边缘而低分辨率层级捕获语义上下文，两者形成互补的多尺度表达。对每个高斯 $G_i$，在投影位置对三层特征图进行双线性上采样后拼接为 $u_i$，再通过联合自编码器压缩为紧凑嵌入 $f_i = \phi(u_i)$，去除三层间的冗余。自编码器在 COCO 和 ScanNet 上预训练以获取通用特征表达能力；在线 SLAM 阶段冻结主体权重，仅在解码器头部插入 LoRA 模块（初始为零卷积，逐步解冻并预热学习率），在防止灾难性遗忘的同时适配体素渲染特有的伪影模式。每个高斯由此获得语义特征嵌入 $f_i^D$，建图损失中新增特征重建项 $\mathcal{L}_\text{feat} = \|F - F^\text{obs}\|_1$。

\textbf{结构感知深度光栅化。}与 FlashSLAM 继承传统单深度值假设不同，FeatureSLAM 提出更精细的深度模型。传统方法假设每个高斯在每个像素处产生单一深度值（期望深度 $D_\text{exp}(\mathbf{x}) = \sum_i T_i \alpha_i z_i$），这在表面不连续处和倾斜平面上会产生偏差。FeatureSLAM 推导出：在透视投影下，每个 3D 高斯在其投影椭圆上产生\textbf{局部平面深度场}——深度在高斯的 2D 投影区域内线性变化，而非保持常数。将这一逐高斯变化融入深度渲染，得到修正的期望深度 $D_\text{exp}(\mathbf{x}) = \sum_i T_i(\mathbf{x})\,\alpha_i(\mathbf{x})\,D_i(\mathbf{x})$，消除了 GS-ICP-SLAM 的远平面偏差，保留了倾斜表面，并产生更锐利的轮廓边缘。此外，FeatureSLAM 还计算中值深度 $D_\text{med}(\mathbf{x})$（沿射线累积透明度首次超过 0.5 处的深度）作为补充——期望深度平滑但有偏，中值深度保持锐利边缘，二者互为补充。

\textbf{深度-法线一致性与几何正则化。}基于结构感知深度模型，FeatureSLAM 引入三项几何约束。(1)~\textbf{分块 Pearson 相关深度损失}：将渲染深度与观测深度划分为 $p{\times}p$ 块，计算块内零均值归一化互相关——该损失对绝对尺度偏差不敏感，专注于保留相对深度变化（倾斜平面、不连续处、局部表面形状）。(2)~\textbf{深度-法线一致性}：分别从期望深度、中值深度和传感器深度通过有限差分计算法线 $\mathbf{n}_\text{exp}$、$\mathbf{n}_\text{med}$、$\mathbf{n}_\text{obs}$，对每对法线施加余弦相似度损失，鼓励深度场形成连贯平面和锐利边缘，抑制浮游高斯。(3)~\textbf{eRank 屏障正则化}：在线建图中，高斯容易退化为针状（尺度向量 $s_i$ 退化为秩-1），导致法线不稳定、深度过度自信。FeatureSLAM 计算尺度分布的归一化奇异值熵 $H_i = -\sum_k p_{i,k}\ln(p_{i,k})$，定义有效秩 $\text{erank}(s_i) = \exp(H_i) \in [1,3]$（1 对应针状退化，3 对应各向同性），施加屏障损失 $\mathcal{L}_\text{erank} = \lambda_e / \text{erank}(s_i)$，并配合薄盘惩罚 $\mathcal{L}_\text{thin} = \lambda_\text{thin} \exp(-s_{i,3}^2)$（$s_{i,3}$ 为最小尺度分量），防止零宽度盘状高斯干扰 GICP 跟踪和深度渲染。

\textbf{Feature-GICP 跟踪。}FlashSLAM 用 SuperPoint+LightGlue+SVD 做稀疏特征匹配；FeatureSLAM 采用完全不同的路线——将 GICP 从纯几何对应扩展为几何+语义联合对应。标准 GICP 通过最小化 Mahalanobis 距离对齐两组点集，但在无纹理、镜面反射或深度缺失区域，几何对应变得不可靠。FeatureSLAM 的关键观察是：当模型和观测都包含密集的、视角稳定的语义信息时，特征相似度可以辅助消歧几何对应。具体而言，Feature-GICP 在标准 GICP 能量函数中加入语义匹配项：

\begin{equation}
    T^* = \arg\min_{\mathbf{R}, \mathbf{t}} \sum_{i=1}^{N} \alpha_i \left( \mathbf{d}_i^\top \left( \mathbf{C}_t^{(i)} + \mathbf{R}\,\mathbf{C}_s^{(i)}\mathbf{R}^\top \right)^{-1} \mathbf{d}_i + d_f \right)
\end{equation}

其中 $(i, \sigma(i))$ 为几何对应对，$\mathbf{d}_i = \mathbf{p}_t^{(i)} - \mathbf{R}\mathbf{p}_s^{\sigma(i)} - \mathbf{t}$，$\alpha_i$ 为高斯不透明度（代理该位置体素组合的复杂度），$d_f = -\|f_i^{s,t} - f_{\sigma(i)}^{s,t}\|_2^2$ 为语义特征距离惩罚项。直觉上，即使在纹理贫乏区域几何对应模糊，语义特征仍可区分不同物体表面，使对应关系保持稳定。该问题经线性化后用 Gauss-Newton 求解。

\textbf{语义梯度剪枝。}与 FlashSLAM 的简单深度截断不同，FeatureSLAM 在剪枝中同时考虑颜色和语义重要性。对每个高斯 $G_i$，定义联合重要性分数：

\begin{equation}
    \tau_i = \sum_{\mathbf{x} \in \Omega_i} \left( \lambda_c \left| \frac{\partial \mathcal{L}_\text{rgb}}{\partial \alpha_i(\mathbf{x})} \right| + \lambda_f \left| \frac{\partial \mathcal{L}_\text{feat}}{\partial \alpha_i(\mathbf{x})} \right| \right)
\end{equation}

其中 $\lambda_c$ 和 $\lambda_f$ 平衡颜色与特征重要性。重要性分数低的高斯对外观和语义的贡献均可忽略，作为剪枝候选。跟踪关键帧时保守剪枝（$p=10\%$），建图关键帧时积极剪枝（$p=30\%$），在保持跟踪鲁棒性与控制地图规模之间取得平衡。

\textbf{效果。}在 Replica 数据集上，FeatureSLAM 在所有场景上达到 SOTA 跟踪精度（ATE RMSE），超越 GS-ICP-SLAM、Photo-SLAM、CartGS 等基线，同时以约 5 FPS 运行——无需预计算语义掩码、全局 BA 或回环闭合。在 TUM RGB-D 真实噪声数据上，GS-ICP-SLAM 因深度质量下降产生大量浮游高斯和边缘模糊，FeatureSLAM 则保持稳定的场景表示和重建质量。在开放集语义分割任务上，FeatureSLAM 匹配或超越 LERF、LangSplat 等离线蒸馏方法，而训练时间仅为后者的 $1/10$，且无需任何数据预处理。

DINO-SLAM（2025.07）~\cite{2507_19474} 沿着 FeatureSLAM"语义特征增强场景表示"的思路，从特征编码策略和表示通用性两个维度做出不同选择。

\textbf{动机。}FeatureSLAM 将 SAM2 多尺度特征直接嵌入高斯参数，但 DINO-SLAM 认为\textbf{原始 DINO 特征本身不足以充分编码场景结构}——DINO 特征仅捕获语义相似性，缺乏对空间布局和表面几何的显式建模，需要与几何信息交叉融合才能获得更丰富的场景表示。因此 DINO-SLAM 设计了场景结构编码器（Scene Structure Encoder, SSE），通过外观-几何交叉注意力将 DINO 特征蒸馏为更丰富的 Enhanced DINO（EDINO），同时支持 NeRF 和 3DGS 两种表示，而非仅限于高斯。

\textbf{SSE：外观-几何交叉注意力蒸馏。}DINO-SLAM 用 DINOv2 ViT-small（而非 FeatureSLAM 的 SAM2 Hiera-MAE）提取基础语义特征 $f_d$，随后通过两个并行 MLP 分别提取外观特征 $f_a = F_{ed}(f_d)$ 和几何特征 $f_g = F_{eg}(d)$（$d$ 为深度图）。核心创新在于两层交叉注意力模块：第一层以几何-外观联合信息引导 DINO 特征精炼 $f_{\text{atten-d}} = \text{softmax}(f_g f_a^\top / d_a)\,f_d$，使语义编码感知空间布局；第二层反向用精炼后的 DINO-几何信息增强外观表示 $f_{\text{atten-a}} = \text{softmax}(f_g f_{\text{atten-d}}^\top / d_{\text{atten-d}})\,f_a$。最终 EDINO 特征拼接为 $f_{ed} = [f_{\text{atten-a}}',\, f_{\text{atten-d}},\, f_g]$，同时编码语义、外观和几何三重信息。相比 FeatureSLAM 的自编码器压缩（去除层级冗余），SSE 的交叉注意力显式建模了三类信息间的互补关系。

\textbf{双管线设计与 EDINO 融合。}在 3DGS 管线中，DINO-SLAM 将 EDINO 特征直接赋给每个高斯 $G^i$ 的独立参数通道 $f_{ed}^i$，渲染时通过标准 $\alpha$-混合 获得特征图 $f_d' = \sum_i^N f_{ed}^i \alpha^i \prod_{j=1}^{i-1}(1-\alpha^j)$。监督采用双损失策略：编码特征损失 $\mathcal{L}_{ef} = \|\phi_d - f_{ed}\|_1$ 监督 tri-plane/高斯参数学习 EDINO 表示，DINO 特征损失 $\mathcal{L}_{df} = \|f_d' - f_d\|_1$ 确保渲染特征与原始 DINO 对齐——前者引导场景表示吸收结构化信息，后者保持与 DINO 语义空间的一致性。跟踪线程仅优化相机位姿（颜色+深度残差），不参与特征优化；建图线程执行全局 BA 以缓解遗忘和漂移，关键帧通过光心距离阈值 $d_l{=}0.7$ 和视差角阈值 $r_l{=}15^\circ$ 选择。在 Replica、ScanNet 和 TUM 上，DINO-SLAM 的双管线均达到 SOTA 水平。

CB-KNN（2025.11）~\cite{2511_23221} 沿着 SplaTAM 基于轮廓渲染的 SLAM 框架，从渲染端而非特征端提升跟踪鲁棒性。

\textbf{动机。}在线 SLAM 中高斯参数不可避免地存在误差（初始化偏差、密度不足、浮游噪声），这些误差会导致渲染结果产生尖锐伪影和局部偏差，进而误导位姿梯度——与其要求渲染越精确越好，不如让渲染过程本身对参数误差具备一定的容错能力。与 FeatureSLAM 为每个高斯嵌入语义特征以辅助对应匹配不同，CB-KNN 不引入任何特征编码器或额外参数。

\textbf{动机：渲染质量不等于跟踪质量。}传统 3DGS-SLAM 的跟踪优化依赖渲染图像与观测之间的逐像素残差，但参数误差会导致渲染结果产生尖锐伪影和局部偏差，进而误导位姿梯度。CB-KNN 挑战了"渲染越精确、跟踪越好"的直觉假设：适度模糊反而能抑制异常高斯的局部干扰，为位姿优化提供更平滑、更稳定的梯度景观。关键约束是——所有修正仅作用于当前帧的临时渲染，\textbf{绝不回写原始高斯参数}，保持"临时修正$\to$关键帧平滑$\to$原始地图保真"的三级原则。

\textbf{CB-KNN 临时修正机制。}对每个像素 $p$，CB-KNN 在标准 alpha-compositing 之前插入一步邻域平滑：首先从可见高斯中选出对该像素贡献最大的 $K$ 个最近邻 $G_p = \{g_{p1}, \ldots, g_{pK}\}$，然后对每个 $g_{pk}$ 执行位置和颜色的双重修正。位置修正是将其 2D 投影向 $K$ 邻域的质心 $C_p = \frac{1}{K}\sum_{g_{pj} \in G_p} \pi(u_{pj}, E_t, K_\text{cam})$ 收缩：$\pi'(g_{pk}) = \pi(g_{pk}) + \lambda\,\frac{C_p - \pi(g_{pk})}{\|C_p - \pi(g_{pk})\| + \epsilon}$，其中 $\lambda \in [0.1, 0.3]$ 为偏移系数，$\epsilon = 10^{-6}$ 防止除零。颜色修正则以各邻域高斯对像素 $p$ 的衰减值为权重做加权平均 $c'_{pk} = \sum_{g_{pj} \in G_p} w_{pj}\, c_{pj}$，抑制异常高斯的颜色噪声。经修正后得到临时高斯集 $G'_p$，后续 RGB、深度、轮廓渲染均基于 $G'_p$ 而非原始 $G$ 执行。

\textbf{自适应 $K$ 值选择。}$K$ 值由局部高斯密度和场景复杂度联合决定。将关键帧划分为 $8{\times}8$ 像素网格，统计每个网格内的高斯投影数 $D$，定义局部密度 $\rho = D/64$；同时将相邻关键帧间的运动幅度归一化为 $\mu \in [0,1]$。$K$ 值计算为 $K = K_0 + (0.5 \cdot \rho - (1 - \mu)) \cdot \gamma$，其中基准 $K_0 = 8$，调节系数 $\gamma = 0.3$。密度低或运动剧烈时 $K$ 增大以增强平滑，密度高或运动平稳时 $K$ 减小以保留细节。在 Replica 合成场景上取 $K_0 = 5$（深度质量高、运动微小），在 TUM-RGBD 和 ScanNet 真实噪声场景上取 $K_0 = 8$。

\textbf{关键帧与非关键帧的差异化处理。}CB-KNN 修正仅施加于关键帧——非关键帧直接用原始高斯 $G$ 渲染并基于恒速模型初始化位姿，保证跟踪速度；关键帧则经 CB-KNN 修正后渲染平滑 RGB、深度和轮廓图，仅对轮廓值 $S(p) > 0.99$ 的像素（完整覆盖区域）计算跟踪损失 $\mathcal{L}_t = \sum_{p: S(p)>0.99} (\mathcal{L}_1(D(p) - D_\text{GT}(p)) + 0.6\,\mathcal{L}_1(C(p) - C_\text{GT}(p)))$，其中颜色权重 $0.6$ 由实验确定。关键帧选择本身也受益于 CB-KNN：平滑后的 3D 点云用于计算帧间重叠度，使关键帧选取标准更稳定。

\textbf{效果。}以 SplaTAM 为基线，CB-KNN 在 Replica 上将 ATE RMSE 从 0.39 cm 降至 0.32 cm，在 TUM-RGBD 上从 3.31 cm 降至 3.18 cm，在 ScanNet 上从 13.43 cm 降至 12.48 cm；渲染方面，TUM-RGBD 的 PSNR 从 22.80 dB 提升至 23.87 dB，Replica 上 LPIPS 达到 0.09（所有方法中最优）。CB-KNN 不引入任何新参数或特征网络，计算开销仅为关键帧上的邻域查询和加权平均，在 NVIDIA A40 上保持与 SplaTAM 相近的帧率。

RoGER-SLAM（2025.10）~\cite{2510_22600} 转向一个不同的鲁棒性维度——FeatureSLAM 通过语义特征增强对应质量来对抗纹理贫乏，RoGER-SLAM 则直面\textbf{输入信号本身退化}的问题：传感器噪声和低光照导致原始 RGB 图像和深度图都不再可靠，此时无论特征多么丰富，基于退化输入的优化都会偏离正确解。

\textbf{动机。}当输入图像被强噪声淹没或光照严重不足时，标准的逐像素光度损失本身就是不可信的监督信号——以错误的目标做优化只会产生更错误的结果。

\textbf{SP-RoFusion：结构保持鲁棒融合。}标准 3DGS SLAM 用原始 RGB 图像作为光度监督信号，但在强噪声下该信号本身就不可信。RoGER-SLAM 的思路是构建一个\textbf{伪监督目标}，融合多模态信息以补偿单模态退化。具体而言，对渲染图像 $I_\text{render}$、观测深度 $D_\text{gt}$ 和结构边缘图 $G$，生成融合伪图像 $I_\text{fuse} = \beta_r I_\text{render} + \beta_d \text{Norm}(\text{Rep}_3(D_\text{gt})) + \beta_g \text{Norm}(\text{Rep}_3(G))$，其中 $G$ 由对渲染灰度图施加 Sobel 算子提取。直觉上，当颜色信号被噪声淹没时，深度和边缘仍提供可靠的几何骨架；当深度缺失时，颜色和边缘仍约束外观。建图损失 $L_\text{map} = 0.5 L_\text{color} + L_\text{depth} + \lambda_\text{illum} L_\text{illum}$ 中，照明项权重 $\lambda_\text{illum}$ 根据当前颜色残差自适应调节——残差越大（说明退化越严重），结构融合项的权重越高，形成动态防御机制。

\textbf{自适应跟踪目标。}与 FeatureSLAM 用 Feature-GICP 引入语义对应不同，RoGER-SLAM 保持直接法跟踪框架，但解决\textbf{固定颜色-深度权重不适应场景变化}的问题。在纹理丰富场景中颜色残差更可靠，在低纹理场景中深度残差更可靠，静态权重无法兼顾。RoGER-SLAM 设计了动态重加权机制：$w_\text{im} = \sigma_\text{im}/(L_\text{color} + \epsilon_\text{im})$，$w_\text{depth} = \sigma_\text{depth}/(L_\text{depth} + \epsilon_\text{depth})$，并通过正则项 $\lambda_R(\sigma_\text{depth} w_\text{im} - \eta)^2$ 将权重比约束在标定值 $\eta$ 附近，防止任一权重发散至极端值。跟踪损失仅使用高不透明度像素（$O_p > 0.99$）作为监督，排除低置信度区域的干扰。

\textbf{双条件触发的 CLIP（对比语言-图像预训练，Contrastive Language-Image Pre-training）增强模块。}SP-RoFusion 能应对中等噪声，但当噪声和低光照同时出现时，伪监督目标本身也会退化。RoGER-SLAM 引入一个按需激活的 CLIP 增强管线：利用 CLIP 冻结编码器对噪声不变的特征提取能力，通过多级可学习解码器去噪；低光照增强则先用神经表示归一化（Neural Representation Normalization, NRN）将不同退化程度映射到统一水平，再用 CLIP 文本-图像对齐损失引导增强——鼓励增强后的图像特征与"明亮"文本提示相似、与"暗淡"文本提示远离。关键在于\textbf{选择性激活}：仅当全局亮度 $\mu_L < 80$（低光照）且残差噪声方差 $\sigma_R^2 > 30$（高噪声）同时满足时才触发，正常帧仅用 SP-RoFusion 以保持效率。这种设计使系统在良好条件下保持标准 SLAM 的速度，仅在真正恶劣的帧上付出额外计算代价。

GSORB-SLAM（2024.10）~\cite{2410_11356} 与 Photo-SLAM 共享"ORB 特征 + 3DGS 渲染"的基本架构，但将两者的耦合从 Photo-SLAM 的松散分工（前端 ORB 做定位、后端 3DGS 做渲染）推进为\textbf{紧密联合优化}，并在深度表示、致密化和视点选择三个环节分别做了针对性改进。

\textbf{动机。}Photo-SLAM 的松散分工导致两个问题：前端 ORB 的重投影误差信息未传递给高斯优化，浪费了几何约束；后端 3DGS 的标准 $\alpha$-混合深度受噪声高斯干扰而不稳定，影响表面重建质量。

\textbf{表面深度替代重渲染深度。}Photo-SLAM 及大多数 3DGS-SLAM 用标准 $\alpha$-混合渲染深度 $D = \sum z_i \alpha_i T_i$，但射线沿途的每个高斯（包括噪声伪影）都会贡献权重，导致渲染深度不稳定。GSORB-SLAM 借鉴 2DGS 和 RTG-SLAM 的思路，用累积透射率中值定位最接近几何表面的高斯作为表面深度 $D_s = \{z_i \mid T_i > 0.5\}$，并引入透射率完备性估计 $\hat{T} = \sum 2\alpha_i T_i$ 以判断该深度是否可信——若 $\hat{T}$ 未收敛则 $D_s$ 本身不确定。跟踪时仅保留透射率接近 1（$T > 0.99$）的高斯作为几何特征，排除低置信区域对位姿优化的干扰。

\textbf{紧密耦合跟踪：联合光度-重投影误差。}Photo-SLAM 将定位完全交给 ORB 因子图优化，GSORB-SLAM 则在跟踪阶段将 ORB 重投影误差 $\mathcal{L}_{\text{rpj}} = \sum_{j \in M}\sum_{i \in P}\rho(\|p_i - \pi(W_{cj}, P^j_i)\|_2)$ 与基于表面深度的光度误差 $\mathcal{L}_{\text{rgb}} + \mathcal{L}_{\text{sur}}$ 加权联合优化：$\mathcal{L}_{\text{tracking}} = \mathbb{1}(T > 0.99)(\lambda_1^t \mathcal{L}_{\text{rgb}} + \lambda_2^t \mathcal{L}_{\text{sur}}) + \lambda_3^t \mathcal{L}_{\text{rpj}}$。为防止错误特征匹配主导损失下降方向，跟踪前先移除 ORB 外点，使重渲染损失作为主分量。后端则沿用 ORB-SLAM2 的稀疏特征图优化（bundle adjustment），与 Photo-SLAM 的回环闭合策略形成互补。

\textbf{自适应高斯扩展与正则化。}Photo-SLAM 利用非活跃 ORB 特征点引导高斯增生（几何致密化），GSORB-SLAM 采用不同策略：由外观掩码 $M_c$（灰度重渲染）和几何掩码 $M_d$（深度误差 $< 0.05$ 且有效）识别欠重建区域，并通过自适应深度误差阈值 $\hat{\epsilon}_d = \text{median}\{E_d\} + 40 \cdot \text{std}\{E_d\}$ 抑制冗余添加。扩展条件为累积透射率不足（$\hat{T} < 0.8$）或几何-外观联合条件满足。建图损失中额外引入各向同性正则化 $\mathcal{L}_{\text{iso}}$ 和标量正则化 $\mathcal{L}_s$，专门抑制 Photo-SLAM 中未处理的条状高斯伪影。

\textbf{混合图视点选择。}Photo-SLAM 未显式设计训练视点选择机制，GSORB-SLAM 观察到仅依赖共视关系图（如 TANBRIDGE）可能导致训练视角覆盖面不足。其混合图机制分三步构建渲染帧集 $R$：首先加入邻接帧（增强当前视角优化）和已优化帧（适配后端位姿更新）；其次对共视关系图内的候选帧要求重叠度 $OG < \epsilon_1$（排除冗余），对图外候选帧要求 $OG > \epsilon_2$（补充共视特征未覆盖的视角）；最后随机加入 $n_r$ 帧以缓解灾难性遗忘。这一机制在保证视角多样性的同时避免了过拟合。

\textbf{效果。}相比 ORB-SLAM2 基线 RMSE 降低 16.2\%，相比 3DGS-SLAM 基线 PSNR 提升 3.93 dB，在跟踪精度和重建质量上同时达到 SOTA。

\subsection{恶劣环境与特殊配置}
\label{ss:15_3}

本节介绍两条技术路线：（1）深度先验与稀疏传感器；（2）其他 RGB-D SLAM 进展。

\subsubsection{深度先验与稀疏传感器}

并非所有深度传感器都提供稠密可靠的深度测量。本小节讨论如何利用深度先验和稀疏传感器实现有效的 3DGS SLAM。

RGB-D 传感器的深度测量并不总是稠密可靠的。ToF-Splatting（2025.04）~\cite{2504_16545} 修改了 3DGS SLAM 的深度获取环节——用\textbf{多帧积分模块}替代单帧深度补全——3DGS 的渲染管线和高斯表示保持不变。

\textbf{动机。}ToF（飞行时间，Time-of-Flight）传感器在功耗上极具优势——低成本传感器仅需约 $200$ mW、成本低于 \$3——但其极低分辨率（$8 \times 8$ 即 64 个深度点）产生的深度极度稀疏，仅占一幅 $640 \times 480$ 图像的 $0.02\%$，无法直接用于高斯基元初始化或深度监督。

\paragraph{Baseline 做法。}在标准 3DGS SLAM（如 MonoGS、SplaTAM）中，稠密深度图承担三个关键角色：(1) 高斯基元的位置初始化——每个深度像素反投影为 3D 点作为高斯中心 $\mu \in \mathbb{R}^3$；(2) 跟踪阶段的深度监督——光度损失 $\mathcal{L}_{\text{rgb}}$ 与深度损失 $\mathcal{L}_{\text{depth}}$ 联合优化位姿；(3) 建图阶段的几何约束。然而这些方法均假设深度输入是稠密的（$100\%$ 密度）。面对极稀疏 ToF 输入，先前的 ToF-SLAM 利用 DELTAR 深度补全网络在单帧上从稀疏 ToF 测量恢复稠密深度图，但深度补全高度依赖输入稀疏深度的质量——噪声或缺失点会迅速传播并恶化稠密深度的精度，且这类模型跨域泛化能力较差。

\paragraph{核心改进：多帧积分模块。}ToF-Splatting 的核心改进是用\textbf{多帧积分模块}替代单帧深度补全。直觉上，如果只有一个稀疏深度图，补全网络的"想象力"容易出错；但如果同时拥有多个视角的信息，就可以利用视角间的视差关系交叉验证、互相补充，从而得到更可靠的稠密深度。具体实现基于 Depth on Demand（DoD）框架并扩展至多视角处理：从关键帧图中选取 $N = 4$ 个局部帧（基线距离接近 $b = 15$ 以确保足够视差），在每次迭代中使用不同的源视角交替更新深度预测，从而平滑融合多视角线索且避免顺序依赖。

该模块融合三路线索联合生成稠密深度图 $D^k$：(a) 稀疏 ToF 深度提供精确但稀少的绝对深度锚点；(b) 单目颜色图像提供边缘和纹理信息引导深度不连续处的边界锐化；(c) 多视角几何线索通过帧间视差恢复稠密相对深度。三个信息源互补：ToF 提供绝对尺度但极稀疏，多视角几何提供稠密但仅相对尺度的深度，颜色引导边界精度。此外，单目深度先验（Depth Anything v2 预测的归一化深度 $\hat{D}^k \in [0,1]^{H \times W}$）被注入 DoD 的单目编码器，提供对无纹理区域和视场角受限场景的鲁棒偏置。

\paragraph{异常值过滤。}廉价 ToF 传感器在特定表面上会产生严重异常值（某些深度点偏差极大），这些异常值若直接输入 DoD 会严重恶化预测精度。为此，ToF-Splatting 设计了一个简单有效的两步过滤策略：首先，DoD 在不使用 ToF 测量的条件下预测一个纯多视角深度图 $\tilde{D}^k$（此时深度仍有度量尺度，因为尺度由相对位姿 $P^{jk}$ 约束）；然后计算 $\tilde{D}^k$ 与 ToF 测量 $H^k$ 之间的 $L_1$ 误差，剔除误差高于分位数阈值 $q = 0.75$ 的 ToF 点，得到过滤后的稀疏深度 $\hat{H}^k$；最后将 $\hat{H}^k$ 重新输入 DoD 进行最终预测。这种"多视角先过滤，再融合"的策略巧妙地利用了多视角几何作为异常值检测器。

\paragraph{跟踪与建图。}整个系统采用帧到模型（frame-to-model）架构，跟踪损失为光度与深度的加权和：

\begin{equation}
\mathcal{L}_{\text{track}} = \alpha_{\text{track}} \underbrace{\|I_k - \hat{I}(F^k)\|_1}_{\mathcal{L}_{\text{rgb}}} + (1 - \alpha_{\text{track}}) \underbrace{\|D^k - \hat{D}(F^k)\|_1}_{\mathcal{L}_{\text{depth}}}
\end{equation}

其中 $\alpha_{\text{track}} = 0.9$，$\hat{I}$ 和 $\hat{D}$ 分别为渲染 RGB 和深度图。该公式直观地将"图像看起来对不对"和"深度测量准不准"两个信号按 $9:1$ 的比例加权——光度信号占主导，深度信号作为辅助约束。跟踪分两阶段进行：先用纯光度损失优化 $\tau_{\text{rgb}}^T = 30$ 步获得初始位姿，再执行多帧积分生成稠密深度图 $D^k$，最后用完整跟踪损失精化 $\tau_{\text{rgbd}}^T = 70$ 步。关键帧选择基于渲染不透明度（opacity）定义的新颖度因子 $\eta^k$：像素 $i$ 的不透明度 $O(KF^k_i) < \tau_{\text{th}}$ 时标记为不确定区域，$\eta^k$ 为不确定区域与已覆盖区域的比值，当 $\eta^k > \eta_{\text{th}} = 0.1$ 时选择为关键帧。建图损失在跟踪损失基础上增加 SSIM 损失 $\mathcal{L}_{\text{ssim}}$、法向量一致性损失 $\mathcal{L}_{\text{normal}}$（通过渲染深度和多帧预测深度的 3D 梯度计算法向量）和各向同性正则化 $\mathcal{L}_{\text{iso}}$（惩罚过度拉长的高斯，促进等比例缩放），每轮建图迭代 $\tau_{\text{rgbd}}^M = 60$ 步，每 $\tau_{\text{rgbd}}^M / 2$ 步剪枝不透明度低于 $0.5$ 的高斯。

\paragraph{系统耦合与实验表现。}跟踪、多帧积分和建图三个模块紧密耦合——跟踪质量影响关键帧选择和位姿精度（影响多帧积分），多帧积分的深度质量影响高斯初始化和建图质量，建图的渲染不透明度又反馈到关键帧选择策略，形成自我强化的循环。在 ZJUL5 真实稀疏 ToF 数据集上，ToF-Splatting 的轨迹误差（ATE）低于 3 cm，在多个场景上跟踪和建图精度均显著优于 ToF-SLAM、KinectFusion、ElasticFusion 等方法（其中部分方法在最具挑战性的场景上完全失败）。消融实验表明：仅用 $0.04\%$ 密度（约 128 个点）的稀疏深度，ToF-Splatting 即可接近使用 $100\%$ 稠密深度的 MonoGS 的性能。单次深度估计耗时约 90 ms（含 Depth Anything v2 推理），整体管线在 RTX 3090 上以约 1.5 FPS 运行。

MemGS（2025.09）~\cite{2509_13536} 将关注点从深度获取质量转向嵌入式平台的 GPU 内存效率。

\textbf{动机。}3DGS-SLAM 的高斯基元数量在运行过程中持续增长，在移动机器人（MAV）等资源受限设备上很快耗尽宝贵的 GPU 内存，限制了系统的运行时长和可重建场景规模。与 ToF-Splatting 采用端到端 3DGS SLAM 架构不同，MemGS 以 ORB-SLAM3 为前端提供位姿和稀疏地图点，3DGS 映射作为后端——这意味着其高斯基元管理和初始化逻辑与锚点方法完全不同。

\paragraph{Patch-Grid（PG）采样初始化。}锚点方法中高斯基元由稠密深度图的像素反投影初始化。MemGS 面对的是 ORB-SLAM3 前端提供的稀疏关键点——即便经过后续的克隆/分裂操作，这些局部集中、全局稀疏的特征点仍不足以建模整个场景。为此，MemGS 引入 Patch-Grid（PG）采样：将图像划分为 2D 网格补丁，当某个补丁内的关键点数量低于阈值时，在补丁内均匀采样额外点并用最近邻关键点的深度赋值，从而获得更均匀的空间覆盖。

\paragraph{基于几何相似性的体素合并。}ToF-Splatting 通过多帧积分提升深度质量，但不涉及高斯基元的冗余管理。MemGS 的核心改进是在体素空间中合并几何相似的冗余高斯，直接针对 SLAM 过程中高斯基元数量持续增长所导致的内存膨胀问题。具体地，首先用梯度掩码 $M(G)$ 筛选梯度低于阈值 $\epsilon = 0.001$ 的稳定高斯（即几何和外观已充分收敛的基元），然后在体素空间中对落入同一体素的候选高斯对 $\{G_i, G_j\}$ 计算平方马氏距离 $d_M^2 = (\mu_j - \mu_i)^T \Sigma_i^{-1}(\mu_j - \mu_i)$。仅当 $d_M^2$ 低于卡方分布临界值 $\chi^2_{0.05} = 7.815$（3 自由度，95\% 置信水平）时才执行合并。

合并后高斯的位置 $\mu_k^*$ 通过最大化两个高斯的联合对数似然解析求解：$\mu_k^* = (\Sigma_i^{-1} + \Sigma_j^{-1})^{-1}(\Sigma_i^{-1}\mu_i + \Sigma_j^{-1}\mu_j)$，本质上是协方差加权的最优位置。尺度和旋转则通过最小化 Wasserstein-2 距离（度量两个高斯分布的概率分布差异）用 L-BFGS 迭代优化，初始旋转由球面线性插值（slerp）给出，初始尺度取两者的均值。整个合并过程在自定义 CUDA 核上执行，单次耗时约 3 ms。在不影响系统运行时性能的前提下，MemGS 显著降低了 GPU 内存使用，使 3DGS SLAM 在 MAV 等资源受限设备上成为可行方案。

\subsubsection{其他 RGB-D SLAM 进展}

本路线介绍频域分析、结构化初始化和 TSDF 融合等推动 RGB-D SLAM 进步的补充技术路线。

在频域分析方向，FGS-SLAM（2025.03）~\cite{2503_01109} 针对 SplaTAM 框架中高斯初始化的不确定性问题——空间域信息无法判别场景各区域需要"多少、多密"的高斯——提出用傅里叶频域分析驱动自适应致密化。

\textbf{动机。}标准 3DGS-SLAM 的致密化基于轮廓值和深度误差等空间域指标，但这些指标无法区分"纹理丰富的细节区域需要密集高斯"和"平坦表面只需少量高斯"，导致高斯分配不均匀。与 SplaTAM 的核心差异体现在两方面。

\textbf{跟踪架构}：SplaTAM 依赖可微渲染 + 轮廓掩码在图像空间优化位姿，需要维护稠密地图才能跟踪；FGS-SLAM 则构建独立的稀疏高斯地图，利用 GICP（广义迭代最近点）直接在 3D 空间中估计变换矩阵 $T$，将源点云与稀疏地图对齐。这一设计借鉴了 GICP 与 3DGS 共享高斯分布假设的数学基础，使跟踪不再受制于稠密地图质量。

\textbf{频域致密化}：SplaTAM 基于轮廓值和深度误差判断"哪里缺高斯"；FGS-SLAM 先对当前帧做 2D 傅里叶变换，用高斯高通滤波器 $H(u,v) = 1 - e^{-D^2(u,v)/2D_0^2}$ 分离高/低频区域，再以三角阈值法自适应确定采样间隔——高频区域（纹理丰富的细节）以间隔 $m$ 密集采样并赋予小半径 $\lambda_h$，低频区域（平坦表面）以间隔 $n$（$n > m$）稀疏采样并赋予大半径 $\lambda_l$。此外，系统引入高斯缺失区域检测掩码 $M_m = M_i \cup M_d \cup M_c$（联合不透明度掩码、深度掩码和颜色掩码），仅在缺失区域新增高斯，避免盲目添加导致的爆炸式增长。稀疏地图与稠密地图共享高斯属性，通过联合渲染保持一致性。建图与跟踪并行运行于独立线程。在 Replica 数据集上，FGS-SLAM 的建图迭代次数平均不到 3 次——比 SplaTAM（30 次）、MonoGS（60 次）等方法少 5--100 倍——频域先验使高斯参数收敛极快。系统整体以 36 FPS 运行，比 SplaTAM（约 4 FPS）快近一个数量级。

SEGS-SLAM（2025.01）~\cite{2501_05242} 沿用 Photo-SLAM 的跟踪--建图解耦架构（ORB-SLAM3 前端定位 + 3DGS 后端渲染），但识别出一个被先前工作忽视的关键问题：3DGS-SLAM 系统在优化过程中高斯的底层结构会逐渐退化——即使 Photo-SLAM 用 ORB-SLAM3 点云初始化高斯，优化过程中高斯仍会偏离场景的真实几何结构，导致边缘模糊和结构失真。

\textbf{动机。}标准 3DGS 的自适应致密化以光度损失梯度为唯一驱动，没有显式的结构保持机制，导致高斯在优化中逐渐偏离 ORB 特征点捕获的锐利边缘和几何结构。

\textbf{结构增强锚点（SEPM）。}SEGS-SLAM 的核心观察是 Scaffold-GS 中基于锚点的 3DGS 组织方式能有效保持场景结构：锚点在训练中位置固定，其下属高斯的分布始终围绕锚点展开。SEGS-SLAM 将这一思想引入 SLAM——对每个关键帧的 ORB-SLAM3 点云 $P_k$ 做增量体素化 $V_k = \lfloor P_k / l \rfloor$（$l$ 为体素尺寸），将体素中心作为锚点。锚点继承点云的强结构先验，且因位置固定而持续保持场景结构，有效防止优化过程中的结构退化。高斯参数（位置偏移、颜色、旋转、缩放）通过轻量 MLP 从锚点特征解码，颜色分支额外融入外观嵌入。

\textbf{Appearance-from-Motion 嵌入（AfME）。}此前 NeRF-W 等方法用图像索引作为外观嵌入，但 SLAM 中测试集占比高达 80\%（远多于 NeRF 的 12.5\%），图像索引方式难以泛化到未见视角。AfME 改用相机位姿 $(R, t)$ 作为外观编码器的输入，通过轻量 MLP $M_a$ 将位姿映射为嵌入向量 $\phi^{(a)}_{R,t} = M_a(R, t)$，再送入颜色解码器。位姿作为输入有三个优势：每帧唯一、天然编码空间信息、比离散索引更连续。训练后固定高斯属性仅变化位姿输入时，渲染图像的光照条件随位姿平滑变化，验证了 AfME 能有效将外观编码到连续位姿空间中。

\textbf{频率金字塔正则化（FPR）。}SEGS-SLAM 进一步观察到结构化高斯在渲染高频细节（物体边缘、复杂纹理区域）时表现不佳。FPR 在多个尺度 $s \in \{s_0, s_1, \ldots, s_n\}$ 上对渲染图像和真实图像分别做 2D FFT，通过高通滤波器提取高频分量后计算频域损失 $\mathcal{L}_{hf}$。FPR 主要用于引导锚点致密化：当某体素内所有高斯的平均梯度超过阈值 $\tau_g$ 时，在该体素中心新增锚点——高频区域的渲染-真值差异会增大 $\mathcal{L}_{hf}$，推动梯度超过阈值，从而在高频区域增加更多锚点。系统整体损失为 $\mathcal{L} = (1-\lambda)\mathcal{L}_1 + \lambda(1-\text{SSIM}) + \lambda_{vol}\mathcal{L}_{vol} + \lambda_{hf}\mathcal{L}_{hf}$。

\textbf{实验表现。}SEGS-SLAM 支持单目、双目和 RGB-D 三种输入。消融实验中，仅引入 SEPM 即已超越现有 SOTA 方法；在此基础上添加 AfME 和 FPR 进一步提升。在 TUM RGB-D 数据集上（光照变化大、场景高度杂乱），SEGS-SLAM 在 PSNR 上显著优于所有对比方法（比 MonoGS 提升约 19.86\%）；在 Replica 数据集上同样取得最高平均渲染质量。跟踪精度与 Photo-SLAM、ORB-SLAM3 相当，实时跟踪约 17 FPS、渲染约 400 FPS。

CaRtGS（2024.10）~\cite{2410_00486} 系统性地分析了 GS-SLAM 中在线渲染优化的\textbf{计算错位}（computational misalignment）现象，并提出三阶段自适应对齐策略。

\textbf{动机。}在线 SLAM 受实时约束仅有数千次优化迭代（远少于离线 3DGS 的数万次），且其像素级反向传播因多线程竞争而效率低下；同时，最早加入的关键帧被过度优化而最新关键帧优化不足，导致渲染质量随时间退化。CaRtGS 是一个即插即用模块，可集成到 Photo-SLAM、GS-ICP-SLAM 等解耦架构的后端渲染管线中。

\textbf{优化不充分问题}：标准 3DGS 的离线训练可达数万次迭代，但在线 SLAM 受实时约束仅有数千次，且其像素级反向传播（pixel-wise backpropagation）因多线程竞争共享高斯基元而触发串行原子操作，严重限制并行效率。CaRtGS 引入 splat-wise 反向传播——将线程映射从像素切换到 splat，每个线程独立维护一个 splat 的状态，通过前向阶段预存透射率 $T$ 和累积颜色实现高效的像素状态交换。在相同运行时间内将总优化迭代数从基线的 4.6 提升至 15.4（3 倍），平均 PSNR 从 32.1 dB 提升至 33.8 dB。

\textbf{长尾优化问题}：为防灾难性遗忘，GS-SLAM 通常随机从关键帧池中选取旧帧再优化，但池的持续扩张导致最早加入的关键帧被反复优化、最新关键帧则优化不足，PSNR 呈下降趋势。CaRtGS 改为基于训练 loss $L$ 的自适应选取：为每个关键帧维护剩余迭代数 $r_i$ 和最近 loss $l_i$，当所有关键帧的迭代配额耗尽后，按 loss 排序将新配额优先分配给 loss 最高的前 $d_k$ 个关键帧。这一策略在 Replica Room2 上将 PSNR 从 34.9 dB 提升至 36.4 dB。

\textbf{弱约束致密化问题}：标准 3DGS 仅通过低不透明度剪枝控制模型规模，约束不足。CaRtGS 引入不透明度正则化损失 $\mathcal{L}_o = \frac{1}{N}\sum_i o_i$，鼓励高斯基元学习低不透明度，使对渲染贡献小的基元更易被剪枝淘汰。完整损失为 $\mathcal{L} = \mathcal{L}_\text{rendered} + \lambda_o \mathcal{L}_o$，其中 $\mathcal{L}_\text{rendered} = (1-\lambda_\text{ssim})\mathcal{L}_1 + \lambda_\text{ssim}\mathcal{L}_\text{ssim}$。在 Replica 单目设置下，正则化系数 $\lambda_o=0.001$ 即可将高斯基元数量减半，PSNR 损失可忽略；在 VECtor 数据集上，CaRtGS 以仅十分之一的高斯基元取得超过 3 dB 的 PSNR 提升。系统整体在 22 FPS 以上的实时约束下运行。

GSFusion（2024.08）~\cite{2408_12677} 的核心转变在于地图表示层面。

\textbf{动机。}SplaTAM 将 3DGS 作为唯一地图表示，但纯高斯地图难以提供精确的几何约束（如用于机器人导航的障碍物检测）；同时，SplaTAM 的逐像素致密化会产生大量冗余高斯，导致内存持续膨胀。SplaTAM 将 3DGS 作为\emph{唯一}地图表示，GSFusion 则构建\textbf{双地图系统}——在 Supereight2 体积映射引擎上同时维护一张 TSDF 体积地图和一张 3DGS 地图。TSDF 地图提供准确的几何结构，3DGS 地图负责照片级渲染，两者互补而非替代。这一设计带来的核心收益是：可以利用 TSDF 体素的权重信息来\textbf{去重}——查询待初始化位置最近的体素，若权重为 1 表示首次观测、可安全添加高斯；若权重已超过 1 表示该位置已被访问过、无需重复初始化。相比 SplaTAM 基于渲染轮廓和深度残差的逐像素致密化，GSFusion 的初始化策略进一步引入了\textbf{四叉树对比度分解}：对每帧 RGB 图像按对比度递归划分，直到每个叶节点区域的对比度低于阈值——平坦区域产生大的叶节点（初始化少量大高斯），边缘细节区域产生小的叶节点（初始化多个小高斯）。每个叶节点中心仅生成一个高斯，初始旋转为单位矩阵，缩放由叶节点中心到其角的反投影长度确定。这一策略使高斯数量从逐像素采样降低一个数量级，且紧凑的初始分布直接消除了 SplaTAM 所需的后续致密化和剪枝过程。

\textbf{在线优化策略。}SplaTAM 对所有帧统一执行固定次数的建图优化。GSFusion 则利用"新初始化高斯数量"作为信息增益指标来区分关键帧与非关键帧——新高斯数超过阈值时标记为关键帧。关键帧执行 $m$ 次优化迭代，非关键帧仅执行 $n$ 次（$n < m$），节省的计算量被随机分配给历史关键帧列表中的帧，以缓解在线优化中的遗忘问题（catastrophic forgetting）。在 Replica 和 TUM RGB-D 数据集上，GSFusion 以比 SplaTAM 少一个数量级的高斯基元达到了可比的渲染质量，同时输出 TSDF 几何地图以支持机器人导航等下游任务。

HoloGS（2024.05）~\cite{2405_02005} 的出发点是：SplaTAM 等 3DGS-SLAM 系统的相机位姿和初始化点云仍依赖离线运动恢复结构（Structure from Motion, SfM）预处理（如 COLMAP），这一步骤耗时且不适用于需要即时重建的场景。

\textbf{动机。}COLMAP 等 SfM 工具需要完整图像集且计算耗时，无法满足 AR/VR 等需要即时获取 3D 重建结果的应用需求；而商用 MR 头显（如 HoloLens 2）本身配备高质量的视觉-惯性追踪系统和深度传感器，可直接提供位姿和稠密点云，从而绕过 SfM 预处理步骤。HoloGS 利用 Microsoft HoloLens 2 的原生传感器数据\textbf{完全绕过 SfM}——三个核心数据流直接从硬件获取：(1) RGB 图像（$1920 \times 1080$）及对应相机位姿（来自 HoloLens 内建视觉-惯性追踪系统）；(2) ToF 深度图及对应位姿（来自 HoloLens 深度传感器，截断阈值 3\,m）；(3) 将深度图逐像素按内参反投影、再通过位姿变换至 RGB 相机坐标系后合并得到的稠密点云。

与 SplaTAM 从单帧深度图反投影初始化高斯不同，HoloGS 将多帧深度点云融合为统一点云（约 62 万点，远多于 SfM 稀疏点云的约 4 万点）作为高斯中心的初始位置，其余参数（球谐系数、不透明度、尺度、旋转）沿用标准 3DGS 的默认初始化。优化管线也与 SplaTAM 不同：SplaTAM 在在线帧序设定下逐帧优化位姿和高斯参数，而 HoloGS 假设位姿由传感器直接提供且固定不变，沿用了标准 3DGS 的离线训练流程（30,000 次迭代，$\mathcal{L} = (1-\lambda)L_1 + \lambda L_{\text{D-SSIM}}$），不做位姿精化。实验在两个自采 HoloLens 数据集（户外雕像"Denker"和室内植物"Ficus"）上对比内部传感器数据与外部 SfM 数据的效果：内部数据的 PSNR 约 20\,dB，低于外部 SfM 数据约 27\,dB，渲染图像边缘模糊且存在浮游伪影（floater artifacts），致密化点云的 Chamfer 距离几何精度平均弱约 10 倍。

论文分析其主要原因在于 HoloLens RGB 相机位姿精度不足（而非深度点云质量），并提出在训练过程中联合优化相机位姿（类似 NeRF 中的位姿精化策略）作为改进方向。此外，HoloGS 展示了从训练后高斯中心提取致密化点云的后处理能力——这一"点云填充"效果可类比 MVS 稠密重建，但在低纹理均匀区域（如花盆、雕像基座）仍存在空洞。

Qu et al.（2024.08）~\cite{2408_05635} 在 SplaTAM 的各向同性高斯表示和轮廓感知渲染管线基础上，做了两项核心改动。

\textbf{动机。}SplaTAM 的恒速假设在大运动或剧烈旋转下初始位姿偏差较大，同时优化全部 6 自由度容易陷入局部最优；此外，标准 3DGS 高斯的多视角不一致性会导致表面表示不准，影响位姿估计和重建精度。第一，将 SplaTAM 的恒速模型初始化 + 联合位姿优化替换为\textbf{旋转-平移解耦策略}：交替迭代优化旋转和平移分量，而非同时优化全部 6 自由度。第二，引入\textbf{深度先验几何正则化}——3DGS 高斯的多视角不一致性会导致表面表示不准，进而降低位姿估计和重建精度，深度先验约束弥补了这一缺陷。在关键帧管理上，Qu et al. 也改用基于平均视差（parallax）的选择策略，替代 SplaTAM 的视锥重叠度方案。WildfireX-SLAM（2025.10）~\cite{2510_27133} 构建大规模森林火灾 SLAM 数据集，揭示了 3DGS-SLAM 在森林环境中的独特挑战。

\textbf{动机。}森林火灾场景具有极端的视觉特征——浓烟遮挡、火焰的动态光照明、植被的高频纹理重复——这些条件远超标准室内 SLAM 数据集的难度范围，需要专门的数据集来评估和推动 3DGS-SLAM 在极端环境下的鲁棒性。

\begin{takeaway}
RGB-D SLAM 是 3DGS 在线重建的起点：SplaTAM 首次实现 3DGS 同时跟踪与建图，GSFusion 的双地图系统用 TSDF 补偿纯高斯地图的几何缺陷，四叉树采样和深度先验正则化则持续提升重建效率和精度。
\end{takeaway}

% ======================================================================


\section{单目与RGB SLAM}
\label{sec:15_2}

与 RGB-D SLAM 相比，单目和纯 RGB SLAM 面临更大的挑战：缺少直接深度测量意味着系统必须从单张或多帧图像中推断几何信息。高斯基元的初始化、位姿估计的收敛域和尺度一致性都变得更加困难。本节梳理 28 篇论文，从早期开拓到近期的前馈范式和多视角基础模型融合，展现这一方向的快速演进。

\subsection{经典优化范式}

本节介绍三条技术路线：（1）早期开拓：从零开始的单目高斯 SLAM；（2）深度先验引导的 RGB SLAM；（3）解耦跟踪与建图的架构。

\subsubsection{早期开拓：从零开始的单目高斯 SLAM}

本路线介绍不依赖深度传感器的单目 3DGS-SLAM 开拓性工作，它们从零开始解决了高斯初始化和位姿优化两大核心困难。

单目 3DGS-SLAM 的核心困难在于：没有深度传感器，高斯无法直接初始化，位姿优化的收敛域也大幅缩小。本小节介绍两个奠定单目高斯 SLAM 基础的开拓性工作。

Gaussian Splatting SLAM（MonoGS）（2023.12）~\cite{2312_06741} 是首个将 3DGS 应用于单目 SLAM 的工作。

\textbf{动机。}在此之前的 NeRF-SLAM 方法需要逐射线采样和体渲染，单帧渲染耗时数百毫秒，无法实时运行；而标准 3DGS 虽然渲染速度快，但其自适应致密化机制和位姿假设是为离线新视角合成设计的，不适合帧序到达的在线 SLAM 场景。本方法修改了 3DGS 的\textbf{相机位姿优化和建图}环节——用解析雅可比矩阵替代自动微分进行位姿梯度计算，用各向同性正则化约束高斯形状，用基于共视性的关键帧管理替代固定间隔策略——3DGS 的光栅化渲染管线和高斯基元表示保持不变。其背景挑战在于：没有深度传感器时，高斯基元无法从深度图直接初始化位置，且位姿优化的收敛域大幅缩小。

\textbf{跟踪与解析雅可比。}标准 3DGS 的 CUDA 光栅化仅对高斯参数（位置、颜色、不透明度、协方差）计算了显式梯度，但并未推导相机位姿 $\mathbf{T} \in \mathrm{SE}(3)$ 关于渲染结果的梯度。在基线做法中，若要端到端地通过可微渲染优化相机位姿，只能依赖自动微分（autodiff）逐层回传梯度，这在每帧需要多次迭代的 SLAM 场景下引入了显著的计算开销。MonoGS 的核心创新是推导了 3DGS 中相机位姿的\textbf{解析雅可比矩阵}——这是首次在 EWA（椭圆加权平均）Splatting 框架下给出 $\mathrm{SE}(3)$ 位姿关于高斯参数的完整解析梯度。

直观地理解，3D 高斯的世界坐标 $\boldsymbol{\mu}^w$ 通过相机位姿 $\mathbf{T}$ 投影到图像平面得到相机坐标 $\boldsymbol{\mu}^c$，而这个投影过程关于位姿的导数决定了渲染误差如何反向传播到位姿参数。利用李代数（Lie algebra）在流形上求导，可以推导出最小参数化雅可比——所谓最小参数化，是指雅可比矩阵的维度严格匹配自由度数，消除冗余计算：

\begin{equation*}
\frac{\partial \boldsymbol{\mu}^c}{\partial \boldsymbol{\xi}} = -[\boldsymbol{\mu}^c]_{\wedge}, \quad
\frac{\partial \boldsymbol{\mu}^c}{\partial \mathbf{t}} = \mathbf{R},
\end{equation*}

其中 $\boldsymbol{\xi} \in \mathbb{R}^3$ 为旋转向量（对应 $\mathrm{SE}(3)$ 中旋转分量的李代数），$[\cdot]_{\wedge}$ 将三维向量映射为 $3 \times 3$ 反对称矩阵，$\mathbf{t}$ 为平移向量，$\mathbf{R}$ 为旋转矩阵。直观地，旋转导数 $-[\boldsymbol{\mu}^c]_{\wedge}$ 表示相机绕原点旋转时场景点的瞬时速度方向（叉积关系），而平移导数 $\mathbf{R}$ 则直接是旋转矩阵——因为平移对相机坐标的影响就是刚体旋转后的平移量。在跟踪阶段，仅优化当前帧的相机位姿，不更新高斯参数。单目模式下最小化光度残差：

\begin{equation*}
\mathcal{E}_{\mathrm{pho}} = \left\| I(\mathbf{T}, \mathcal{G}) - \mathbf{I} \right\|_1,
\end{equation*}

其中 $I(\mathbf{T}, \mathcal{G})$ 从位姿 $\mathbf{T}$ 渲染高斯集合 $\mathcal{G}$ 得到合成图像，$\mathbf{I}$ 为观测图像。RGB-D 模式下额外加入几何残差 $\mathcal{E}_{\mathrm{geo}} = \| D(\mathbf{T}, \mathcal{G}) - \mathbf{D} \|_1$，其中深度图 $D_p = \sum_i z_i \alpha_i \prod_{j=1}^{i-1}(1 - \alpha_j)$ 通过 $\alpha$-混合 渲染，$z_i$ 为高斯 $i$ 沿相机射线的距离。两个损失通过权重 $\lambda_{\mathrm{pho}} = 0.9$ 加权联合（RGB-D 模式）。相机位姿使用 Adam 优化器，旋转学习率 $0.003$、平移学习率 $0.001$，每帧执行 100 次迭代（当位姿更新幅度低于 $10^{-4}$ 时提前终止）。这一解析推导消除了自动微分开销，使跟踪阶段在每帧仅需约 50 次梯度下降迭代即可收敛，实现 3 FPS 实时运行。

\textbf{单目深度初始化。}在单目模式下，MonoGS 面临深度初始化难题——没有深度传感器，高斯基元无法从深度图直接反投影初始化位置。MonoGS 的策略是利用渲染深度（即当前地图已有的深度估计）来引导新高斯的初始化，但通过不同的方差设计区分可靠与不可靠区域。对于有渲染深度估计的像素 $p$，高斯中心以低方差初始化 $\boldsymbol{\mu} \sim \mathcal{N}(D_p, 0.2\,\sigma_D)$——低方差意味着初始位置紧密围绕估计深度，优化器只需微调；对于无深度估计的未观测区域像素，高斯中心以高方差初始化 $\boldsymbol{\mu} \sim \mathcal{N}(\tilde{D}, 0.5\,\sigma_D)$——其中 $\tilde{D}$ 为渲染深度的中位数，高方差使这些高斯的初始位置覆盖更大的深度范围。这一设计的直觉是：高方差初始化的多数错误高斯在后续多帧优化中，因无法在多个视角下产生一致的渲染贡献，其不透明度会迅速趋近于零而自然消失。

\textbf{高斯剪枝。}仅靠不透明度衰减不足以清理所有错误高斯。MonoGS 进一步引入基于可见性的剪枝：仅在关键帧窗口 $\mathcal{W}$ 已满时执行，对于最近 3 个关键帧内插入的高斯，若未被至少 3 个其他帧观测到，则判定为几何不稳定予以剪枝。此外，所有不透明度低于 $0.7$ 的高斯也被清除。

\textbf{关键帧管理与高斯共视性。}SLAM 需要从连续视频流中选取最具信息量的帧联合优化地图，而非使用所有帧。MonoGS 的关键帧选择基于\textbf{高斯共视性}（Gaussian Covisibility）——一个利用 3DGS 光栅化天然深度排序特性定义的共视度量。具体地，一个高斯被判定为"从某视角可见"，当且仅当：(a) 它参与了该视角的光栅化渲染；(b) 射线在该高斯处的累积不透明度尚未达到 $0.5$（即未被前方高斯完全遮挡）。这一判据利用了 3DGS 沿射线排序的渲染顺序，无需额外的遮挡检测启发式。基于此，两帧 $i$, $j$ 之间的共视性通过 IOU（交集/并集）和重叠系数（Overlap Coefficient, OC）度量：

\begin{equation*}
\mathrm{IOU}_{\mathrm{cov}}(i, j) = \frac{|\mathcal{V}^i \cap \mathcal{V}^j|}{|\mathcal{V}^i \cup \mathcal{V}^j|}, \quad
\mathrm{OC}_{\mathrm{cov}}(i, j) = \frac{|\mathcal{V}^i \cap \mathcal{V}^j|}{\min(|\mathcal{V}^i|, |\mathcal{V}^j|)},
\end{equation*}

其中 $\mathcal{V}^i$ 为从帧 $i$ 可见的高斯集合。关键帧 $i$ 被注册的条件为：与上一关键帧 $j$ 的 $\mathrm{IOU}_{\mathrm{cov}}$ 低于阈值（Replica 数据集为 $0.95$，TUM 为 $0.90$），或相对平移 $\|t_{ij}\|$ 超过帧 $i$ 中位深度的 $\text{kf\_m}$ 倍（Replica 为 $0.04$，TUM 为 $0.08$）。关键帧窗口 $\mathcal{W}$ 的大小限制为 $8$--$10$ 帧，旧关键帧在与最新关键帧的 $\mathrm{OC}_{\mathrm{cov}}$ 低于 $0.3$ 时被移除。

\textbf{建图与各向同性正则化。}建图阶段联合优化当前窗口 $\mathcal{W}$ 内的关键帧和随机选取的 2 个历史关键帧（防止全局地图遗忘），优化高斯参数的同时固定相机位姿。MonoGS 引入\textbf{各向同性正则化}：

\begin{equation*}
\mathcal{E}_{\mathrm{iso}} = \sum_{i=1}^{N} \| \mathbf{s}_i - \bar{\mathbf{s}}_i \mathbf{1} \|_1,
\end{equation*}

其中 $\mathbf{s}_i \in \mathbb{R}^3$ 为第 $i$ 个高斯的三个缩放参数，$\bar{\mathbf{s}}_i = \frac{1}{3}\sum_k s_i^k$ 为其均值。这一公式惩罚缩放参数偏离其均值——即鼓励三个轴方向的缩放趋于一致，使椭球趋向球体。设计动机在于：3DGS 光栅化在视线方向上没有约束，即使在有深度观测的情况下，高斯仍然可以沿射线方向自由拉伸。在标准新视角合成中这不是问题（因为输入视角经过精心选取），但在连续 SLAM 中，单目或稀疏视角下缺乏足够的多视角约束，高斯容易退化为沿视线方向的细长椭球（"针状高斯"），产生渲染伪影并干扰后续帧的跟踪。完整的建图优化目标为：

\begin{equation*}
\min_{\{\mathcal{G}^k\}} \sum_k \mathcal{E}^k_{\mathrm{pho}} + \lambda_{\mathrm{iso}} \mathcal{E}_{\mathrm{iso}},
\end{equation*}

其中 $\lambda_{\mathrm{iso}} = 10$，建图阶段每关键帧执行 150 次迭代。消融实验表明，各向同性正则化在 Replica 数据集（包含大量弱纹理区域）上显著提升跟踪精度，在 TUM 数据集上改善较小。

该方法可无缝扩展到 RGB-D SLAM，并能重建微小甚至透明物体。在 TUM 单目评估中，MonoGS 在不使用任何深度先验网络的情况下超越了 DSO、DROID-SLAM 等经典和学习方法。

Splat-SLAM（2024.05）~\cite{2405_16544} 是首个实现全局优化的 RGB-only 高斯 SLAM 系统。

\textbf{动机。}MonoGS 采用 frame-to-model 跟踪，没有全局位姿和地图优化，长期运行会产生不可逆的累积漂移；同时 MonoGS 仅依赖渲染深度进行高斯初始化，在多视角约束不足的区域容易产生几何退化。Splat-SLAM 的核心思路是将\textbf{frame-to-frame 光流跟踪}与\textbf{可变形高斯地图}结合，通过回环检测和全局光束法平差（BA）消除漂移，并通过在线地图变形将全局优化结果传播到高斯场景表示中。在建图环节，Splat-SLAM 沿用了 MonoGS 的高斯共视性关键帧选择、基于可见性的剪枝和增密策略以及各向同性正则化，但在跟踪和全局一致性上做出了根本性的架构改变。

\textbf{基于光流的稠密光束法平差跟踪。}与 MonoGS 通过解析雅可比直接优化相机位姿的光度误差不同，Splat-SLAM 采用预训练的循环光流模型 RAFT（2020）~\cite{teed2020raft} 预测帧间稠密光流，再结合 DSPO（Disparity, Scale and Pose Optimization）层（2024）~\cite{zhang2024glorie} 在因子图 $G(V, E)$ 上联合优化关键帧位姿和逐像素视差。因子图中，节点 $V$ 存储关键帧位姿和视差（视差的倒数即为深度），边 $E$ 存储帧间光流对应关系。

这一设计切换的根本动机是收敛域：光流跟踪是 frame-to-frame 的——当前帧仅与相邻关键帧建立对应关系，而非依赖已有地图的渲染结果——这使得跟踪不受地图质量的制约，在场景探索初期地图尚不完善时尤为关键。DSPO 层包含两个交替优化的目标。第一个目标是稠密光束法平差（Dense BA），在滑动窗口内的局部因子图上联合优化位姿和视差。直觉上，Dense BA 的目标是最小化光流预测的帧间对应关系与当前位姿/深度估计之间的不一致程度——如果位姿和深度都完美，光流对应点应精确落在重投影位置上。

\begin{equation*}
\mathcal{E}_{\mathrm{DBA}} = \sum_{(i,j) \in E} \left\| \hat{\mathbf{x}}_{ij} - K_j^{-1}\left(T_j T_i^{-1}(1/d_i) K^{-1}[\mathbf{p}_i, 1]^\top\right) \right\|^2_{\Sigma_{ij}},
\end{equation*}

其中 $\hat{\mathbf{x}}_{ij}$ 为光流预测的对应像素坐标，$K$ 为相机内参，$T_i$, $T_j$ 为关键帧 $i$, $j$ 的 camera-to-world 外参，$d_i$ 为像素 $\mathbf{p}_i$ 的视差，$\|\cdot\|_{\Sigma_{ij}}$ 为以光流置信度为权重的 Mahalanobis 距离。第二个目标引入单目深度先验 $D^{\mathrm{mono}}$（由预训练的 DPT 模型在线预测），通过可学习的逐帧尺度 $\theta_i$ 和偏移 $\beta_i$ 将相对深度对齐到多视点深度：

\begin{equation*}
\mathcal{E}_{\mathrm{mono}} = \lambda_1 \sum_{\mathbf{p}_i \in V} \left\| d_i^h - (\theta_i / D^{\mathrm{mono}}_i + \beta_i) \right\|^2 + \lambda_2 \sum_{\mathbf{p}_i \in V} \left\| d_i^l - (\theta_i / D^{\mathrm{mono}}_i + \beta_i) \right\|^2.
\end{equation*}

这里的关键设计是将视差分为高误差部分 $d^h$ 和低误差部分 $d^l$：仅对多视角优化不准确的像素（高误差部分）使用单目深度修正，保留已准确的多视角视差（低误差部分）用于校准单目深度的尺度和偏移——通过设置 $\lambda_1 < \lambda_2$ 确保尺度参数主要从可靠的低误差视差中学习。高低误差的分类基于双视角深度一致性：对于关键帧 $i$ 中的像素 $(u,v)$，将其投影到所有其他关键帧 $k$ 中检查 3D 距离是否低于阈值，统计全局一致性计数 $n_i(u,v)$，低于阈值则标记为高误差。这两个目标交替优化以避免尺度歧义——若同时联合优化视差、尺度和偏移，尺度自由度会使优化退化。

\textbf{回环检测与全局光束法平差。}MonoGS 仅在滑动窗口内优化，没有回环机制。Splat-SLAM 在局部跟踪之上引入了完整的全局一致性后端。回环检测通过计算当前活跃关键帧与所有历史关键帧之间的平均光流幅值实现——当光流低于阈值 $\tau_{\mathrm{loop}}$（表明视角重叠足够）且帧间时间间隔超过阈值 $\tau_t$（避免冗余边）时，添加单向回环边。回环优化仅优化活跃关键帧及其相连的回环节点，控制计算量。全局 BA 则构建包含所有历史关键帧的完整因子图，每 20 个关键帧执行一次，每次优化前对视差和位姿平移进行归一化以维持数值稳定性。

\textbf{代理深度图与高斯初始化。}MonoGS 在单目模式下使用渲染深度的中位数和高方差初始化新高斯。Splat-SLAM 构造了更精确的\textbf{代理深度图}（Proxy Depth Map）作为深度监督来源：对于多视点一致性验证通过的像素（低误差），直接使用多视点深度 $\check{D}$；对于未通过验证的像素（高误差），使用经尺度对齐后的单目深度 $D^{\mathrm{mono}} + \beta$。这一混合策略确保即使在纹理稀疏或视角覆盖不足的区域，也能从单目深度先验中获得合理的深度估计。新高斯的初始化沿用 MonoGS 的 RGB-D 模式策略——将代理深度图反投影为 3D 点云再下采样，每个点成为一个高斯的中心——区别在于深度来源从真实传感器变为代理深度图。

\textbf{可变形高斯地图。}这是 Splat-SLAM 最核心的架构创新，也是 MonoGS 完全不具备的能力。当后端优化（回环 BA 或全局 BA）更新了关键帧位姿和深度后，高斯地图需要相应变形以保持一致。每个高斯 $g_i$ 记录了其被创建时锚定的关键帧。当该关键帧的位姿从 $T$ 更新为 $T'$、代理深度从 $D$ 更新为 $D'$ 时，关联高斯的属性按如下规则变形：
\begin{itemize}[itemsep=0pt,topsep=2pt,leftmargin=10pt,label=$\bullet$]
\item \textbf{位置}：将高斯中心 $\boldsymbol{\mu}_i$ 投影到原位姿 $T$ 下找到像素对应 $(u,v)$，然后沿光轴方向平移深度差 $\Delta d = D'(u,v) - D(u,v)$——注意不是重新锚定到新深度（因为高斯不一定精确位于表面上），而是保持其与表面的相对偏移；
\item \textbf{缩放}：$s'_i = \left(1 + \frac{\Delta d}{(T^{-1}\boldsymbol{\mu}_i)_z}\right) s_i$，其中 $(T^{-1}\boldsymbol{\mu}_i)_z$ 为高斯在旧相机坐标系下的 $z$ 轴深度；
\item \textbf{旋转}：$R'_i = \Delta R \, R_i$，其中 $\Delta R = R'R^{-1}$ 为位姿旋转更新量。
\end{itemize}
对于投影到缺失深度或视野范围外的高斯，仅施加刚体变换（不调整缩放）。论文发现，主动变形比让优化器被动适应新位姿更有效——特别是渲染质量在主动变形后显著提升。这一设计使得 3D 高斯表示在保持高保真渲染的同时，获得了神经隐式表示（如 GlORIE-SLAM 的神经点云）才具备的全局变形能力，且无需昂贵的重新积分操作。

\textbf{曝光补偿。}Splat-SLAM 还引入了 MonoGS 未涉及的曝光补偿机制：为每个关键帧优化一个仿射变换 $(a_k, b_k)$，将渲染颜色校正为 $a_k C_k + b_k$，使系统对真实视频中常见的亮度变化具有鲁棒性。完整的建图优化目标为：

\begin{equation*}
\min_{\mathcal{G}, a, b} \sum_{k \in \mathrm{KFs}} \left\| (a_k C_k + b_k) - C^{\mathrm{gt}}_k \right\|_1 + \lambda_d \left\| D^r_k - D_k \right\|_1 + \lambda_{\mathrm{reg}} \sum_{i=1}^{|G|} \| \mathbf{s}_i - \bar{\mathbf{s}}_i \mathbf{1} \|_1,
\end{equation*}

其中光度损失和几何损失通过权重 $\lambda_d$ 平衡，$\lambda_{\mathrm{reg}}$ 控制各向同性正则化强度。建图同样在滑动窗口内的关键帧上批量进行，与 MonoGS 一致。

\textbf{实验表现。}Splat-SLAM 是首个同时具备回环闭合、全局 BA、代理深度和在线高斯变形的 RGB-only SLAM 系统。在 Replica、TUM-RGBD 和 ScanNet 上，Splat-SLAM 在重建和渲染精度上均达到或超越已有 RGB-only 方法，在跟踪精度上与使用其他 3D 表示的竞争方法持平，同时保持较小的地图尺寸和运行时间。

\subsubsection{深度先验引导的 RGB SLAM}

本路线介绍利用单目深度估计和视觉里程计为 RGB SLAM 提供几何先验的方法，弥补缺失的深度信息。

深度估计网络的进步为单目 SLAM 提供了关键的几何先验。本小节讨论如何利用单目深度估计和视觉里程计弥补缺失的深度信息。

MonoGS 实现了首个单目 3DGS SLAM，但其跟踪模块深度耦合于高斯地图的渲染质量——解析雅可比需要通过已有高斯渲染图像来计算位姿梯度，场景探索初期地图尚未收敛时跟踪极易退化；同时，MonoGS 的单目深度初始化依赖渲染深度的中位数和高方差"猜测"，初始几何精度受限。MonoGS++（2025.04）~\cite{2504_02437} 保留 MonoGS 的光栅化渲染管线和高斯基元表示不变，但对跟踪前端、地图初始化、高斯管理和正则化四个环节做了全面替换，从根本上将跟踪与建图解耦。

\textbf{动机。}MonoGS 的跟踪深度耦合于高斯地图的渲染质量——场景探索初期地图尚未收敛时跟踪极易退化；同时，MonoGS 的单目深度初始化依赖渲染深度的中位数和高方差"猜测"，初始几何精度受限。

\textbf{DPVO 替代解析雅可比跟踪。}与 MonoGS 通过解析雅可比在可微渲染中联合优化位姿和高斯参数不同，MonoGS++ 将前端完全替换为预训练的 Deep Patch Visual Odometry（DPVO）——一种受 DSO 启发的学习型稀疏补丁里程计。DPVO 对每个 RGB 关键帧 $I_i$ 采样 $K$ 个方形补丁，每个补丁以齐次坐标 $P_k^i = [u, v, 1, d]^\top$ 参数化（$d$ 为逆深度，补丁内假设为常数），构建补丁图并在局部窗口 $N(i)$ 内通过可微光束法平差联合优化位姿和补丁逆深度。核心优化目标为最小化 Mahalanobis 距离：

\begin{equation*}
\sum_{(i,j,k) \in E} \left\| K T_j T_i^{-1} K^{-1} \hat{P}_k^i - \hat{P}_k^{j \leftarrow i} \right\|^2_{\Sigma_{kj}},
\end{equation*}

其中循环网络为每条边预测光流修正量 $\delta_{kj} \in \mathbb{R}^2$ 和置信度权重 $\Sigma_{kj}$。这一架构性改变使跟踪完全独立于高斯地图——即使地图尚未收敛或存在退化区域，DPVO 仍能通过帧间特征匹配获得准确位姿，且无需在每帧执行 MonoGS 那样密集的光度误差迭代（MonoGS 每帧 100 次梯度下降）。

\textbf{稀疏补丁反投影初始化。}MonoGS 在单目模式下用渲染深度中位数和高方差策略初始化新高斯（本质是"猜"深度）。MonoGS++ 改为等待 DPVO 完成 $N$ 帧初始化后，将所有已优化补丁的中心通过相机外参反投影到世界坐标系，直接生成全局稀疏点云：

\begin{equation*}
P = \{T_i^{-1}K^{-1}\hat{P}_k^i \mid i \leq N, k \leq K\},
\end{equation*}

每个点作为一个高斯的中心，初始高斯总数为 $|G| = N \times K$。深度来源从"渲染估计"变为"VO 优化结果"——DPVO 的可微光束法平差已对补丁逆深度进行了多视角约束优化，初始几何精度显著高于 MonoGS 基于渲染深度中位数的随机采样。

\textbf{动态高斯插入。}MonoGS 在每个关键帧直接为所有无高斯覆盖的像素创建新高斯。MonoGS++ 引入冗余过滤：每个新关键帧的补丁反投影更新点云 $P' = \{T_i^{-1}K^{-1}\hat{P}_k^i \mid i \leq j, k \leq K\}$ 后，对每个候选点 $q$ 计算其与现有高斯中心的最小距离 $d(q, P) = \min_{p \in P} \|q - p\|_2$，仅当距离超过阈值 $\tau$ 时才创建新高斯——距已有高斯太近的点判定为已重建区域的冗余点，直接跳过。这一机制避免在已重建区域重复添加高斯，将地图规模的增长控制在仅覆盖新观测到的空间区域。

\textbf{清晰度增强增密。}MonoGS 沿用标准 3DGS 的梯度引导增密（梯度大的高斯分裂或克隆），但该策略在弱纹理区域失效——梯度信号不足以触发增密，且渲染像素被 $\alpha$-混合 权重最大的高斯主导。MonoGS++ 提出\textbf{清晰度驱动}增密作为替代：对每个像素 $p$，找到 $\alpha$-混合 权重最大的主导高斯 $G_m(p)$，然后检查同像素其他高斯 $G_n(p)$ 与 $G_m(p)$ 的投影重叠——重叠掩码 $M(n) = \mathbb{1}(|O_i(n) \cap O_i(m)| > \tau)$ 为真时，将 $G_n$ 分裂。这一策略直接以渲染清晰度为目标（重叠过大意味着多个高斯竞争同一像素导致模糊），不依赖梯度信号，在墙壁、地毯等弱纹理区域尤其有效。

\textbf{平面正则化。}MonoGS 的各向同性正则化鼓励高斯趋向球体（三个缩放轴一致，$\mathcal{E}_{\mathrm{iso}} = \sum_i \| \mathbf{s}_i - \bar{\mathbf{s}}_i \mathbf{1} \|_1$）。MonoGS++ 换用\textbf{平面正则化}，鼓励高斯趋向扁平盘状，与各向同性的目标恰好相反：

\begin{equation*}
\mathcal{L}_{\mathrm{reg}} = |\min(0.01, \min(s))|,
\end{equation*}

通过最小化三个缩放分量中的最小值 $\min(s)$ 来压扁高斯。设计动机在于：墙壁、桌面、地板等平面结构在室内场景中占据大部分面积，扁平高斯比球形高斯更适合表示这些表面——3D 高斯在优化中本身就有扁平化的趋势（如 GaussianShader 和 GaussianPro 所观察到的），平面正则化进一步强化了这一特性。完整的建图目标为：

\begin{equation*}
\mathcal{L} = \lambda_{\mathrm{color}}\mathcal{L}_{\mathrm{color}} + \lambda_{\mathrm{reg}}\mathcal{L}_{\mathrm{reg}},
\end{equation*}

其中颜色损失 $\mathcal{L}_{\mathrm{color}} = (1-\lambda_{\mathrm{photo}})\mathcal{L}_{\mathrm{photo}} + \lambda_{\mathrm{photo}}\mathcal{L}_{\mathrm{SSIM}}$，$\lambda_{\mathrm{photo}} = 0.2$。

\textbf{效果。}跟踪与建图解耦带来了显著的效率提升：MonoGS++ 的 FPS 比 MonoGS 快 5.57 倍（相比使用 RGB 输入的 MonoGS 快 16.32 倍）。在 Replica 数据集上，MonoGS++ 的跟踪精度（ATE RMSE）和渲染质量（PSNR）均优于 MonoGS（RGB），在 TUM-RGBD 上同样在 PSNR 和 LPIPS 上全面领先。消融实验验证了三个新组件的独立贡献：移除动态插入导致地图冗余膨胀，移除清晰度增密导致弱纹理区域模糊，移除平面正则化导致平面细节退化。

IG-SLAM（2024.08）~\cite{2408_01126} 在 MonoGS 的 3DGS 场景表示和渲染管线上做了根本性的跟踪替换——将 MonoGS 的解析雅可比位姿优化替换为 DROID-SLAM 的稠密光束法平差（Dense BA），从而获得稠密深度图、深度协方差和位姿三个输出作为建图监督。

\textbf{动机。}MonoGS 的解析雅可比跟踪深度耦合于高斯地图的渲染质量——场景探索初期地图尚未收敛时跟踪极易退化；而 DROID-SLAM 的稠密 BA 提供了独立于地图质量的深度和位姿估计，且能从 Hessian 矩阵推导出逐像素深度协方差，为建图提供不确定性信息。跟踪模块维护帧图 $(V, E)$，基于共视性选关键帧并建立帧间边；循环门控循环单元（Gated Recurrent Unit, GRU）网络在帧间特征相关性的驱动下预测光流修正量 $r_{ij}$ 和置信度 $w_{ij}$，经 Gauss-Newton 线性化求解 Hessian 系统 $\mathbf{Hx} = \mathbf{b}$ 联合更新位姿和逆深度。跟踪收敛后，IG-SLAM 利用同一 Hessian 结构推导\textbf{深度协方差} $\Sigma_d = \mathbf{P}^{-1} + \mathbf{P}^{-1}\mathbf{E}^T\boldsymbol{\Sigma}_G^{-1}\mathbf{E}\mathbf{P}^{-1}$——这一信息来源是 MonoGS 完全不具备的。

\textbf{深度不确定性引导的建图。}深度协方差在两个环节发挥关键作用。

\textbf{高斯初始化掩码}：将协方差归一化至 $[0, 1]$ 后，以阈值 $0.2$ 生成二值掩码（高协方差区域被屏蔽），再经最大值滤波和多数投票滤波平滑，确保高斯仅在深度可靠的像素处初始化——这与 MonoGS 用方差区分可靠/不可靠区域的策略有相似直觉，但 IG-SLAM 的不确定性直接来自稠密 BA 的 Hessian，而非手工设计的方差参数。

\textbf{加权深度损失}：建图损失 $\mathcal{L} = \mathcal{L}_c + (1-\lambda)\mathcal{L}_d$ 中的深度分量以逆协方差加权 $\mathcal{L}_d = \|\hat{D} - D\|_{\Sigma_d^{-1}}^1$，使不确定区域的深度误差对优化的影响降低。建图的高斯增密与剪枝沿用 MonoGS 的策略（每 150 次迭代克隆小高斯、分裂大高斯、剪枝不透明度低于 $0.1$ 的高斯）。

\textbf{层级优化与学习率衰减。}与 MonoGS 固定分辨率的建图训练不同，IG-SLAM 采用 3 级粗到细金字塔（下采样因子 $s=0.8$，比 Photo-SLAM 的 $s=0.5$、2 级更平滑），每级从对应分辨率的深度图反投影初始化新高斯。此外，由于高斯从稠密深度图初始化时彼此间距很近，过高的位置学习率会导致高斯在目标位置附近来回抖动（"TV 静态噪声"效应），因此 IG-SLAM 引入线性衰减策略 $\eta(t) = (1-t)\eta_i + t\eta_f$，其中 $t = n/\tau$，从覆盖粗到细全范围的初始学习率平滑过渡至精细调优的终值。

\textbf{全局光束法平差。}当关键帧总数超出滑动窗口长度后，IG-SLAM 每 10 个关键帧对所有历史关键帧执行全局 Dense BA（基于距离度量构建稀疏连接图），跟踪结束时再做最后一次全局优化。全局 BA 的位姿和深度更新同步传播到建图模块。系统整体以 10 FPS 单进程运行，在 EuRoC 数据集的大规模序列上实现了照片级 3D 重建。

与 IG-SLAM 同样基于 MonoGS 框架但聚焦于深度初始化环节，UDGS-SLAM（2024.09）~\cite{2409_00362} 针对 MonoGS 单目深度初始化的弱点——以渲染深度的中位数作为初始估计、高方差覆盖的多数错误高斯需后续多帧才能自然消失——引入了更可靠的深度来源：\textbf{UniDepth 网络}直接从单帧 RGB 图像预测度量深度（metric depth），而非仅提供相对深度排序。

\textbf{动机。}MonoGS 在单目模式下以渲染深度的中位数和高方差初始化新高斯，本质是"猜测"深度——大部分错误高斯需在后续多帧中因不透明度衰减而自然消失，这一过程既缓慢又浪费计算资源；如果能在初始化时就获得几何上合理的深度，将大幅缩短收敛时间。这一改进使高斯基元在初始化时就拥有几何上合理的 3D 位置，大幅缩短了后续优化的收敛时间。

然而 UniDepth 的逐像素估计存在局部不一致性（相邻像素深度跳变），UDGS-SLAM 通过\textbf{统计滤波}（基于 IQR 的异常值抑制）在像素邻域内剔除偏离四分位距的深度离群值，确保滤波后的深度场在局部区域平滑连续。在跟踪与建图阶段，系统联合优化相机轨迹 $\mathbf{T}_{cw}$ 和高斯场景参数 $\mathcal{G}$，损失函数结合光度误差 $E_{\mathrm{pho}}$ 与几何误差 $E_{\mathrm{geo}}$：$L(\mathcal{G}, \mathbf{T}_{cw}) = \lambda \cdot E_{\mathrm{pho}} + (1 - \lambda) \cdot E_{\mathrm{geo}}$，其中几何项由 UniDepth 滤波深度图充当监督信号——这与 MonoGS 仅依赖渲染深度的方式形成对比。消融实验还系统分析了 UniDepth 不同骨干编码器（backbone encoder）对最终跟踪精度和渲染质量的影响，揭示了深度估计网络架构选择与 SLAM 端到端性能之间的耦合关系。

\subsubsection{解耦跟踪与建图的架构}

本路线介绍将跟踪与建图解耦的多种架构设计，让专门的模块分别负责快速位姿估计和高质量场景重建。

将跟踪与建图解耦是提升效率和鲁棒性的有效策略——让专门的跟踪模块负责快速位姿估计，让建图模块专注高质量场景重建。本小节讨论多种解耦架构设计。

GSO-SLAM（2026.02）~\cite{2602_11714} 选择 Direct Sparse Odometry（DSO）作为跟踪前端，以 2DGS 作为场景表示，通过 Expectation-Maximization（EM）框架实现跟踪与建图的双向耦合。

\textbf{动机。}松耦合方法（如 Photo-SLAM 将 ORB-SLAM 与 3DGS 各自独立运行）中，VO 的深度估计与 GS 的场景表示互不通信，VO 的误差无法被 GS 修正，GS 的精细几何也无法反馈给 VO 精化位姿——两者都只能各自为战。E-step 固定位姿与深度、优化 2DGS 参数（RGB 渲染损失 $L_1 + D\text{-}SSIM$、半稠密深度一致性损失 $L_1(D_r, D)$、法线一致性先验）；M-step 固定 GS、通过滑动窗口 BA 优化关键帧位姿与半稠密逆深度。M-step 的深度正则化采用加权平均策略——将 DSO 的半稠密深度与 E-step 渲染的稠密深度融合作为 BA 初始值。这既保留了直接法深度估计的可靠性，又利用了稠密重建的完整性，无需额外的正则化项即可在高频几何细节与收敛稳定性之间取得平衡。

\paragraph{基于图像梯度的 Gaussian Splat 初始化。}
GSO-SLAM 的另一项核心贡献是一种利用 DSO 已有计算结果的零开销初始化策略。DSO 在优化过程中自然计算图像梯度，GSO-SLAM 将这些梯度视为 2D 高斯投影分布的概率估计：假设图像强度服从高斯分布，通过对数概率密度的梯度建立方程组，以最小二乘求解每个像素的 2D 协方差 $\Sigma^{2D}$。由于同一 3D 点可被多个关键帧观测到，系统将来自不同关键帧 $obs(p)$ 的 2D 协方差通过 $\Sigma^{2D}_i = J_i W_i \Sigma^{3D} W_i^T J_i^T$ 关联到同一 3D 协方差，线性化后堆叠所有观测的最小二乘估计 $\Sigma^{3D}$。为保证协方差矩阵的正定性，系统施加正则化（加 $\epsilon I$）和特征值裁剪两步修正，再通过特征值分解提取旋转矩阵 $R$ 和缩放矩阵 $S$（最小特征值置零，适配 2DGS 的扁平表示），直接获得接近最终收敛状态的初始高斯基元。消融实验表明，该初始化在 Replica room1 上比 KNN 初始化减少 2700 次额外迭代（64.6 秒）即可达到相同 PSNR（34.45\,dB）。

\paragraph{前端/后端并行架构与实验表现。}
系统在双线程上并行运行：主线程处理前端跟踪（两帧直接对齐估计新帧位姿、关键帧选择、滑动窗口局部 BA 对应 M-step），2DGS 优化线程在后台持续执行 E-step。新增关键帧时，通过上述初始化策略生成新高斯基元插入场景，并优先在最近关键帧的窗口内进行精化训练，待活跃窗口充分收敛后采样更多视角以防止过拟合。场景表示选用 2DGS 而非 3DGS，因为 2DGS 的射线--splat 2D 交点渲染产生多视角一致的深度估计，在几何精度上优于 3DGS 的像素--射线 3D 交点方式。系统无需回环闭合模块和离线精化阶段，纯增量式 VO 即可在 INS 大规模走廊数据集上以 23.11\,dB 的平均 PSNR 超越 MonoGS（16.76\,dB）和 Photo-SLAM（17.62\,dB），运行速度比 MonoGS 快 36 倍、比 SplaTAM 快 141 倍。

同一技术路线中更早的探索是 DROID-Splat（2024.11）~\cite{2411_17660}，同样以 DROID-SLAM 的端到端稠密光流网络为跟踪核心，但采用了更简洁的``直耦合''架构——跟踪器与渲染器通过共享位姿状态直接耦合，无需 HI-SLAM2 那样的四模块分离与离线精化流程。

\textbf{动机。}纯高斯渲染优化的跟踪精度始终落后于传统方法和端到端 SLAM，而端到端跟踪器虽鲁棒但无法产生照片级场景表示；DROID-Splat 旨在弥合这一差距，通过"稠密跟踪 + 稠密渲染"的统一框架同时实现高精度定位和高质量重建。

\paragraph{端到端稠密跟踪与帧图优化。}
跟踪模块维护一个动态帧图 $(V, E)$，关键帧节点存储视差 $d_i$ 和位姿 $g_i \in SE(3)$，边 $(i,j)$ 表示帧间共视关系。与 HI-SLAM2 直接使用 DROID-SLAM 的光流输出作为深度代理不同，DROID-Splat 完整继承了 DROID-SLAM 的可微 BA 框架：循环 GRU 网络预测残差场 $r_{ij}$ 和置信度 $w_{ij}$，驱动对应点更新 $p_{ij}^* = r_{ij} + p_{ij}$，通过基于重投影的可微 BA 联合优化位姿和视差。系统在单 GPU 上同步运行局部前端（小窗口 BA）和全局后端（全图 BA），前者快速处理输入流，后者优化大规模长期连接——这一前端/后端同步设计与 HI-SLAM2 的四模块流水线形成对比，后者将前端跟踪、回环闭合、建图和离线精化作为独立阶段顺序执行。

\paragraph{Pseudo-RGBD 模式与单目深度先验对齐。}
在单目设置下，DROID-Splat 引入 Pseudo-RGBD（P-RGBD）模式整合单目深度先验。系统选择 Metric3D 作为默认先验网络（消融实验表明其时序一致性优于 ZoeDepth 和 DepthAnything），并为每帧估计全局仿射参数（尺度 $s_i'$ 和偏移 $o_i'$）对齐先验深度，正则化项为 $E_{\mathrm{reg},m} = \sum_i \| d_i^* - (s_i' d_i' + o_i') \|^2$。这一策略与 HI-SLAM2 形成直接对比：HI-SLAM2 正是发现全局仿射假设无法校正空间非均匀的深度扭曲，才引入了空间自适应的 $2{\times}2$ 尺度网格；DROID-Splat 保留了更简单的全局仿射对齐，依赖 Metric3D 本身较好的时序一致性来弥补空间精度的不足。此外，仿射参数 $(s, o)$ 与位姿 $G$ 之间存在优化歧义——DROID-Splat 采用与 HI-SLAM 相同的块坐标下降策略：先固定尺度/偏移优化位姿，再固定位姿图优化结构、尺度和偏移，交替执行。

\paragraph{未标定相机的两阶段处理。}
DROID-Splat 的一个突出特点是支持完全未标定的 in-the-wild 视频（如手机拍摄）。系统观察到内参 $\xi$ 与单目先验变量之间存在类似位姿--尺度的歧义，因此设计了两阶段处理流程：第一阶段固定先验参数，联合优化相机内参（使用 DROID-SLAM 的标定扩展）；第二阶段锁定标定结果，以标准 P-RGBD 模式运行。这一设计使系统能够在无需任何相机先验信息的条件下完成重建，而 HI-SLAM2 假设内参已知，专注于在线跟踪精度。

\paragraph{基于视觉相似性的回环检测。}
回环检测策略与 HI-SLAM2 有本质区别。HI-SLAM2 基于光流距离检测回环（要求共视区域足够大以保证收敛），而 DROID-Splat 采用端到端视觉特征：使用 EigenPlaces 网络提取每帧特征向量，插入 CPU 端的 FAISS 数据库进行近邻检索。回环候选需同时满足三个条件——特征距离 $d_{f,ij} < \tau_f$、朝向差异 $d(R_i, R_j) < \tau_r$、时间间隔 $|t_i - t_j| > \tau_t$。检测到候选后，将双向边加入后端帧图，通过可微 BA 在已有的密集对应框架内优化回环约束——这一做法比 Go-SLAM 的回环 BA 更有效。作者特别指出，采用更保守的帧图构建策略本身就能显著提升跟踪精度，这一发现与 HI-SLAM2 的"局部 BA + 全局 BA 分离"理念一致。

\paragraph{MCMC 引导的高斯致密化。}
在渲染端，DROID-Splat 对标准 3DGS 的致密化策略进行了重要改进。原始 3DGS 基于梯度阈值进行分裂和克隆（SplaTAM、MonoGS 等均沿用此策略），但这一启发式策略缺乏全局视野，可能产生冗余或不足的高斯基元。DROID-Splat 引入 MCMC 引导的致密化：将高斯采样建模为马尔可夫链蒙特卡洛过程，通过随机游走采样自适应控制高斯基元分布，同时预设基元总数上限。消融实验表明，在匹配总基元数量（TUM-RGBD 上平均 12 万个高斯）的条件下，MCMC 策略在渲染指标上一致优于标准梯度策略。此外，系统还测试了 2D 高斯泼溅（2DGS）替代 3DGS 的效果，发现 2DGS 在户外无界场景中因生成平滑表面而受浮动高斯影响较小，但在 TUM-RGBD 等室内杂乱场景中容易过平滑细节或构建错误表面。

\paragraph{渲染反馈与跟踪耦合。}
DROID-Splat 还探索了一个前瞻性问题：渲染目标能否反过来改善跟踪？系统设计了两种反馈机制：一是将渲染损失的梯度反向传播到位姿图，在渲染更新时微调位姿；二是利用渲染深度致密化跟踪器的视差状态。实验表明，当存在完美的 RGBD 真值监督时，反馈机制确实有效；但使用单目先验或在真实数据上时，跟踪稳定性反而受到严重影响——这解释了为什么后续的 HI-SLAM2 选择在离线精化阶段才进行位姿--地图联合优化，而非在线阶段直接反馈。

\paragraph{实验表现。}
DROID-Splat 在 Replica 和 TUM-RGBD 上均取得竞争性结果。在单目模式下，由于稠密表示同时应用于跟踪器和渲染器，渲染指标一致排名最高；在 RGBD 模式下 Replica 基准已趋饱和（PSNR ${\sim}39$\,dB，深度 L1 ${\sim}0.6$\,cm）。系统在 RTX 4090 上运行，单目模式下受限于 Metric3D 的推理速度。消融实验系统性地验证了各组件的贡献：帧图构建策略最为关键，单目先验的尺度优化至关重要，MCMC 致密化是最有效的渲染改进，而 2DGS 在室内场景中收益有限。在 in-the-wild 手机视频上，系统展示了从完全未标定输入进行照片级重建的能力。作者指出 Replica RGBD 模式的评估协议存在过拟合问题（奖励在训练帧上的性能），建议关注非训练帧的泛化指标。

HI-SLAM2（2024.11）~\cite{2411_17982} 是 HI-SLAM 的后继系统，从隐式 NeRF 表示转向显式 3DGS，以解决前一版本在尺度对齐精度、地图更新效率和场景边界预设三方面的局限。

\textbf{动机。}HI-SLAM 使用隐式 NeRF 表示，地图更新需要重新训练 MLP 且无法高效执行回环校正后的全局地图变形；3DGS 的显式表示允许通过解析闭式变换即时更新地图，天然适配回环闭合后的全局一致性维护。其核心设计思想是``混合解耦''：用 DROID-SLAM 作为稠密 SLAM 前端提供深度代理（而非直接从地图优化中估计深度），同时以 3DGS 作为紧凑的显式地图表示，在跟踪与建图解耦后通过联合优化重新耦合，兼顾效率与精度。系统由四个模块组成：基于 DROID-SLAM 的在线跟踪器（估计关键帧位姿和深度图）、在线回环闭合模块（检测回环并执行 $Sim(3)$ 位姿图 BA）、连续建图器（基于 3DGS 的在线地图构建与更新）、以及离线精化阶段（全量 BA 与位姿--地图联合优化）。

\paragraph{空间自适应尺度网格对齐。}
在深度估计方面，HI-SLAM2 修改了 3DGS-SLAM 中单目深度先验的对齐策略——用空间自适应的 $2{\times}2$ 尺度网格（scale-grid）替代 HI-SLAM 的单一全局尺度参数——3DGS 的 $\alpha$-混合渲染管线保持不变。问题的根源在于：单目深度先验（如 OmniData）并非度量深度，且其尺度扭曲在图像空间内呈非线性分布——例如画面左侧的深度可能被系统性低估而右侧被高估。HI-SLAM 为每帧估计一个全局尺度 $s$ 和偏移 $b$，本质上是假设整幅深度图仅存在仿射变换偏差。但实际扭曲是空间非均匀的，单一参数无法同时校正所有区域。

HI-SLAM2 的解决方案是为每帧深度先验 $\hat{d}_i$ 建立一个系数维度为 $(m,n)$（默认 $m{=}n{=}2$）的 2D 尺度网格 $s_i$。直觉上，这相当于在深度图上铺设若干``控制点''，每个控制点拥有独立的尺度系数，像素位置处的尺度通过双线性插值 $\mathrm{Bi}(p_i, s_i)$ 从周围四个网格系数获得。对齐后的深度先验为 $\hat{d}_i \cdot \mathrm{Bi}(p_i, s_i)$，深度先验因子定义为：

\begin{equation}
r_d = \| \hat{d}_i \cdot \mathrm{Bi}(p_i, s_i) - d_i \|^2
\end{equation}

其中 $d_i$ 是 BA 估计的逆深度。直观地说，这个因子迫使对齐后的深度先验``贴近''BA 的深度估计，而网格系数自动学习如何校正空间非均匀扭曲。网格分辨率设为 $2{\times}2$ 而非更高，是因为低纹理区域的光流预测置信度低，导致重投影误差对网格分块的权重不足，过高的分辨率反而会引入优化不稳定。

然而，将深度先验因子直接纳入 BA 联合优化位姿、深度和尺度系数，会导致尺度漂移并阻碍收敛（HI-SLAM 已发现此问题）。因此 HI-SLAM2 引入联合深度与尺度对齐（JDSA）模块，将尺度估计从主 BA 循环中分离出来。JDSA 的优化目标为：

\begin{equation}
\min_{s, d} \sum_{(i,j) \in \mathcal{E}} \| \mathbf{q}_{ij} - \pi(\mathbf{T}_{ij} \pi^{-1}(\mathbf{p}_i, d_i)) \|_{\Sigma_{ij}}^2 + \sum_{i \in \mathcal{V}} \| \hat{d}_i \cdot \mathrm{Bi}(\mathbf{p}_i, s_i) - d_i \|^2
\end{equation}

第一项是标准重投影误差（$\mathbf{q}_{ij}$ 为光流预测的目标点，$\mathbf{T}_{ij} = \mathbf{T}_j\mathbf{T}_i^{-1}$ 为帧间刚体变换，$\Sigma_{ij}$ 为光流置信度权重矩阵），第二项是深度先验因子。JDSA 与局部 BA 交替执行：局部 BA 固定尺度系数仅优化位姿和深度，JDSA 固定位姿联合优化尺度和深度，两者交替确保系统尺度稳定且深度先验被良好对齐。

在求解效率方面，JDSA 利用尺度变量维度远小于深度变量的结构特点，将 Hessian 矩阵分块为：

\begin{equation}
\begin{bmatrix} \mathbf{B} & \mathbf{E} \\ \mathbf{E}^T & \mathbf{C} \end{bmatrix} \begin{bmatrix} \Delta s \\ \Delta d \end{bmatrix} = \begin{bmatrix} v \\ w \end{bmatrix}
\end{equation}

其中 $\mathbf{B}$ 对应尺度变量（维度很小），$\mathbf{C}$ 对应深度变量。由于深度先验因子作用于先验而非深度变量本身，$\mathbf{C}$ 保持对角结构（即深度像素间相互独立），可高效求逆为 $\mathbf{C}^{-1} = \mathrm{diag}(1/c_1, \ldots, 1/c_n)$。利用 Schur 补先求解尺度增量 $\Delta s$，再回代求解深度增量 $\Delta d$，避免了大规模稠密矩阵求逆。这一改进在 Replica 数据集上将重建精度提升了 1.54\,cm。

\paragraph{无偏深度渲染与几何监督。}
在地图表示方面，系统引入无偏深度渲染，修改了 3DGS 的深度计算方式——用射线--高斯交点处的平面方程精确交点深度替代直接取高斯均值深度——$\alpha$-混合的颜色渲染保持不变。此前的工作（SplaTAM、Gaussian SLAM）直接使用高斯基元均值位置的深度作为渲染深度，但当射线在高斯基元边缘处穿过而非经过其中心时，均值深度会引入系统性偏差。直觉上，一个``扁平''高斯的均值可能在表面后方，而射线与其椭球面的实际交点在前方，两者之差即为偏差。

HI-SLAM2 的无偏深度计算方法为：沿射线方向求解射线与高斯椭球面的平面方程，确定实际交点深度。由于从同一视点出发穿过同一高斯的所有射线共面，平面方程仅需对每个高斯求解一次，不影响 splat 光栅化的计算效率。这一设计使深度渲染更忠实地反映表面几何位置，显著改善了表面拟合精度。

在建图优化损失方面，系统组合了四项损失函数。光度损失 $L_c$ 计算曝光补偿后渲染图像与观测图像的 L1 差异；深度损失 $L_d$ 计算渲染深度与 BA/JDSA 估计深度的 L1 差异；法线先验损失 $L_n$ 采用余弦嵌入形式 $L_n = \sum_{k \in \mathcal{K}} |1 - \hat{\mathbf{n}}_k^T \mathbf{n}_k|$（其中 $\hat{\mathbf{n}}_k$ 为渲染深度图的法线、$\mathbf{n}_k$ 为 OmniData 预测的法线先验），通过余弦相似度监督低纹理区域的几何一致性——低纹理区域光流信号弱，法线先验提供了关键的形状约束。尺度正则化项 $L_s = \sum_{i \in \mathcal{G}} \| s_i - \bar{s}_i \|$ 惩罚高斯基元沿某一方向的过度拉伸（$\bar{s}_i$ 为三个尺度分量的均值），防止椭球退化为极端细长形状。总损失为 $L = \lambda_c L_c + \lambda_d L_d + \lambda_n L_n + \lambda_s L_s$，使用 Adam 优化器每新增关键帧执行 10 次迭代。

此外，系统简化了 3DGS 的颜色表示：直接使用 RGB 值而非球谐（SH）系数，降低优化复杂度。视角相关的颜色变化通过离线精化阶段的逐帧曝光补偿处理——每帧优化一个 $3{\times}4$ 仿射变换矩阵（$3{\times}3$ 颜色变换矩阵 $\mathbf{A}$ 加 $3{\times}1$ 偏移向量 $\mathbf{b}$），将渲染图像校正为 $\hat{I}' = \mathbf{A}\hat{I} + \mathbf{b}$。高斯基元的初始化通过将关键帧深度图反投影到 3D 空间完成：每个像素的深度值反投影为 3D 点作为高斯均值，朝向初始化为单位朝向，尺度取最近 3 个邻居的平均距离以适应局部点密度，不透明度初始化为 0.5，颜色取对应像素 RGB 值。为控制地图增长，系统实施随机下采样和低不透明度高斯剪枝策略，每 500 次迭代重置不透明度，每 150 次迭代交替执行致密化与剪枝。

\paragraph{$Sim(3)$ 回环闭合与显式高斯变形。}
HI-SLAM2 的回环闭合模块在线并行运行，采用基于 $Sim(3)$ 的位姿图 BA（PGBA）同时校正位姿漂移和单目系统固有的尺度漂移。回环检测基于三个准则：光流距离低于阈值 $d_{\mathrm{flow}}$（确保足够共视区域以可靠收敛）、当前位姿估计的朝向差异在阈值 $\theta_{\mathrm{ori}}$ 内、帧索引差超过最小阈值 $t_{\mathrm{temp}}$（排除局部 BA 窗口内的近邻帧）。检测到回环后，系统将 $SE(3)$ 位姿转换为 $Sim(3)$ 表示（每帧附加独立尺度因子，初始为 1），在位姿图上执行 BA。

PGBA 的优化目标同时最小化两类因子：回环边的重投影因子和从局部 BA 滑动窗口中保留的``非活跃''重投影边推导出的相对位姿因子：

\begin{equation}
\min_{\mathbf{T}, d} \sum_{(i,j) \in \mathcal{E}^*} \| \mathbf{q}_{ij} - \pi(\mathbf{T}_{ij} \pi^{-1}(\mathbf{p}_i, d_i)) \|_{\Sigma_{ij}}^2 + \sum_{(i,j) \in \mathcal{E}^+} \| \log(\hat{\mathbf{T}}_{ij} \mathbf{T}_i \mathbf{T}_j^{-1}) \|_{\Sigma^{\mathrm{rel}}_{ij}}^2
\end{equation}

其中 $\mathcal{E}^*$ 为检测到的回环闭合边，$\mathcal{E}^+$ 为相对位姿因子集合，$\hat{\mathbf{T}}_{ij}$ 为从稠密对应推导的相对位姿，$\Sigma^{\mathrm{rel}}_{ij}$ 为基于调整理论估计的相对位姿方差权重。直观地说，第一项让回环帧对重新对齐，第二项保持已优化局部结构的一致性。为保证收敛，PGBA 使用阻尼 Gauss-Newton 算法，阻尼因子 $\lambda = 10^{-4}$，正则化因子 $\delta = 10^{-1}$，兼顾避免局部极小和数值稳定性。

回环闭合后，地图更新通过显式高斯变形完成——这是 3DGS 显式表示相比隐式 NeRF 的独特优势。每个高斯基元锚定于其初始化关键帧 $i$，当关键帧位姿从 $\mathbf{T}_i$ 更新为 $\mathbf{T}_i'$ 时，高斯的均值、朝向和缩放分别按如下公式变换：

\begin{equation}
\boldsymbol{\mu}_j' = (\mathbf{T}_i'^{-1} \mathbf{T}_i \boldsymbol{\mu}_j) / s_i, \quad R_j' = R_i'^{-1} R_i R_j, \quad \mathbf{s}_j' = s_i \mathbf{s}_j
\end{equation}

其中 $s_i$ 是 $Sim(3)$ 位姿的尺度因子。直觉上，均值先通过旧位姿变换到世界坐标系、再通过新位姿的逆变换回来（并除以尺度因子），等价于将高斯随锚定帧一起``搬运''到校正后的位置。朝向做类似旋转校正，缩放乘以尺度因子以匹配全局尺度变化。由于变换是解析的闭式操作，无需重新训练即可即时反映全局一致性校正。

\paragraph{离线精化。}
离线精化阶段依次执行三个步骤，进一步提升几何与外观质量。第一步是补充分类帧插入：在线阶段的关键帧选择依赖光流距离，无法预知未来轨迹，可能导致视野边缘区域观测不足（例如相机前进后回头的场景）。离线阶段通过将每帧像素投影到相邻关键帧，统计落入邻帧视野之外的像素比例，当超过阈值时标记为观测不足并插入补充关键帧。第二步是全量 BA：在 PGBA 的基础上对所有重叠帧对重新计算重投影因子（PGBA 仅对回环边计算），以更细的粒度优化全局一致性。第三步是位姿--地图联合优化：在光栅化渲染过程中计算位姿 Jacobian，使用 Adam 优化器联合优化高斯参数和相机位姿（同时优化逐帧曝光补偿参数），利用前两步的精化结果作为初始值。

\paragraph{实验表现。}
实验表明，HI-SLAM2 在 Replica 数据集上的 ATE 比仅在线跟踪降低 29.3\%（主要来自回环闭合和离线精化的累积贡献），在 ETH3D SLAM 的 61 条序列上取得最高 AUC（含极端光照和完全黑暗场景），在 Waymo 户外数据集上跟踪精度亦为最佳。其几何与渲染质量不仅大幅领先 RGB-only 方法（ScanNet++ 上完整重建包括低纹理地面、墙壁甚至玻璃窗），甚至在多个指标上超越 RGB-D 方法。系统在线运行 22\,FPS（Replica），离线精化仅需 12 秒。语义扩展实验中，高斯基元增加语义通道后 mIoU 比 HI-SLAM 提升 14.93\%，证明 3DGS 显式表示在语义边界清晰度上优于隐式 NeRF。

GEVO（2024.09）~\cite{2409_09295} 同样利用 3DGS 地图渲染替代原始图像存储，但其关注点与 Splat-SLAM 的全局一致性不同——GEVO 聚焦于\textbf{内存效率与灾难性遗忘}问题。

\textbf{动机。}此前的方法通过存储所有历史关键帧反复重训高斯来缓解遗忘，但这些图像占据总内存的 96\%--99\%，在 Replica 上高达数百 MB，严重制约了 SLAM 系统在资源受限设备上的部署。GEVO 的核心策略是仅保留当前滑动窗口内的原始图像（8--10 帧），用已有地图渲染的图像替代历史关键帧——但由于渲染图像的保真度会随地图退化而逐渐下降，单纯替换并不足够，因此 GEVO 提出了三个互补机制来维持渲染图像的可用性。

\textbf{占用保持初始化。}与 Splat-SLAM 的代理深度图（融合多视点深度与单目先验）不同，GEVO 采用基于多视立体（MVS）代价体的初始化策略：在滑动窗口的 8--10 个关键帧上构建光度一致性代价体 $V(u,d)$（64 个深度假设，图像 4 倍下采样至仅需 4.7 MB），经信念传播提取深度图后上采样至全分辨率，再通过 SPGF$^*$ 算法生成高斯。关键创新在于高斯同时表示\textbf{障碍物}和\textbf{自由区域}——自由区域高斯编码了无障碍空间，用于后续识别和剪枝错误遮挡已有高斯的新高斯（即缓解 RO）。为提升颜色保真度，GEVO 修改了 SPGF$^*$ 使每个高斯仅表示颜色相似的面片。这一初始化过程与定位方法解耦，可搭配 MonoGS 的直接法或 ORB-SLAM 的特征法。

\textbf{两阶段一致性感知优化。}GEVO 将优化分为局部和全局两个阶段。局部阶段仅用当前滑动窗口的图像优化新初始化高斯与可见的旧高斯（目标函数为光度损失 $E_{\mathrm{pho}}$ 加各向同性正则化 $E_{\mathrm{iso}}$），由于不扰动全局地图，旧高斯渲染的历史关键帧保持高保真。全局阶段将局部地图合并入全局地图，但通过三个步骤控制遗忘：（1）插入时将新高斯的不透明度降至 0.2，防止其立即遮挡旧高斯；（2）仅优化一个活跃子集 $A_t = G_t \cup E_t$，其中 $E_t$ 为在当前窗口中渲染误差高于阈值 $\tau$ 的旧高斯——这些高斯既不一致（误差高）又充分可见（贡献大），优化它们不会扰动观察不足的旧高斯，从而缓解 IRO；同时引入光度一致性损失 $E_{\mathrm{pc}}$，从 4 个历史视角渲染图像并与全局地图的渲染结果对齐，阻止地图保真度的累积退化；（3）剪枝不透明度低于 0.7 或占用概率低于 0.9 的高斯（占用概率由自由区域高斯通过高斯混合回归计算），进一步消除 RO。

\textbf{实验表现。}GEVO 将额外内存开销降至 58 MB（TUM）和 83.5 MB（Replica），比 MonoGS 和 Photo-SLAM 低 8--145 倍，比 GO-SLAM 低 53--168 倍——额外内存几乎完全由当前滑动窗口的帧占据。计算开销仅比 MonoGS 增加 8\%--10\%（主要来自仅优化少于 15\% 的活跃高斯）。渲染精度在超过半数序列上优于存储 300 个关键帧的 MonoGS，ATE 亦因基于 MVS 的几何验证初始化而更稳定。消融实验表明，三项技术叠加可将最终 PSNR 恢复至 19 dB，仅比存储全部原始图像低 0.4 dB。

与 Photo-SLAM 用 ORB 特征前端做位姿估计、后端 3DGS 做渲染的解耦思路一脉相承，MGSO（2024.09）~\cite{2409_13055} 同样采用跟踪与建图解耦的架构——但将前端替换为\textbf{直接法}光度 SLAM（DSO），核心动机在于：ORB-SLAM3 仅跟踪角点与边缘等稀疏特征点（约数百个），而 DSO 的梯度优先像素选择策略能产生密度远高于特征点的稠密点云，且其空间分布天然适配 3DGS 的最终高斯分布——复杂区域密集、简单区域稀疏但仍有覆盖。

\textbf{动机。}Photo-SLAM 用 ORB 特征点初始化高斯，但 ORB 特征仅覆盖约 30\% 的图像像素，大量弱纹理区域（墙壁、地板）缺乏初始高斯，需要后续昂贵的致密化迭代来弥补。MGSO 的出发点正是这一经验观察：直接将 DSO 输出的稠密点云作为 3DGS 的初始化来源，省去从零随机初始化和大量增密/剪枝迭代，从而以更紧凑的地图获得更高质量的重建。

\textbf{针对 3DGS 初始化的密度增强与平坦区域插值。}DSO 的像素选择策略将图像划分为不同大小的块，在每个块中选取梯度最高的像素，这种分层次选择使得点云密度天然适配场景复杂度。但 DSO 的点云密度是为跟踪最优而设计的，对 3DGS 初始化而言仍显不足。MGSO 对此做了两项针对性增强。第一，在 DSO 原始跟踪像素之外增加不参与位姿估计的额外像素——使系统同时拥有最优的跟踪像素密度（不过多）和最优的 3DGS 初始化点云密度（不过少）。第二，对梯度接近零的平坦区域（DSO 的光度跟踪要求至少存在一定梯度，故此类区域天然缺乏跟踪点），MGSO 用 Delaunay 三角剖分将已跟踪像素连接为三角网格，在三角面片内插值生成新点，深度取三角面片顶点的平均值——平坦表面上的线性插值精度较高，因此这一策略在最适合插值的区域发挥作用。与 Photo-SLAM 将非活跃 ORB 特征点（约 70\% 的 2D 特征）转化为超基元的方式相比，MGSO 的插值策略初始点数量更多、且聚焦于平坦区域，插值精度更高。

\textbf{3DGS 优化与加速策略。}MGSO 沿用标准 3DGS 的可微光栅化渲染和 $\mathrm{L}_1 + \lambda \cdot \mathrm{SSIM}$ 光度损失，颜色表示采用 MonoGS 的 RGB 方式替代球谐函数以加速训练。在加速策略上，MGSO 直接采用 Photo-SLAM 引入的高斯金字塔渐进训练（对真实图像逐层平滑下采样构建金字塔，优化初期使用低分辨率目标，随训练推进逐步提高分辨率），以及 AbsGS 的同向视角空间位置梯度（homo-directional view-space position gradient）策略替代原始 3DGS 的增密判断标准，并周期性剪枝低不透明度高斯以消除透明浮点伪影。与 Photo-SLAM 针对稀疏 ORB 特征点设计几何致密化策略不同，MGSO 不需要额外的致密化逻辑——SLAM 模块输出的稠密点云本身已提供充足的高斯初始化来源。当新关键帧到来时，MGSO 直接将 SLAM 模块输出的点云位置和颜色作为新高斯的初始化参数，增密/剪枝频率设为每 1000 次迭代。

\textbf{实验表现。}MGSO 仅使用单目 RGB 输入，在桌面 GPU 上以 30\,FPS 以上运行，甚至在笔记本电脑（RTX 3080 移动版）上也能完成实时重建（Replica 上 PSNR 达 31.9\,dB）。在 Replica 上平均 PSNR 为 31.4\,dB，在 EuRoC 上为 22.1\,dB（Photo-SLAM 为 19.7\,dB），在 TUM 上同样优于 Photo-SLAM。最显著的优势在地图紧凑性：EuRoC 上地图仅 8.3\,MB（Photo-SLAM 为 110\,MB，缩小约 13 倍），TUM 上仅 2.9\,MB（Photo-SLAM 为 17\,MB）——原因在于高质量的初始化点云大幅减少了冗余高斯的创建。消融实验证实点云密度对重建质量至关重要：将点云密度减半后重建质量显著下降；而增密频率的提高不仅无益反而增大地图并降低质量——说明 MGSO 的结构化初始点云已足够准确，不需要 3DGS 的增密/剪枝做大幅度修正。MGSO 目前没有回环闭合机制，全局一致性是未来改进方向。

\subsection{前馈范式与基础模型融合}

传统 3DGS-SLAM 依赖逐帧梯度优化，收敛速度受限。以 MonoGS（15.2.1 节）为代表的典型 baseline 中，每关键帧需要随机初始化高斯属性，然后通过数百次可微渲染--反向传播迭代来联合优化高斯参数和相机位姿，单次迭代约 20\,ms，一个关键帧的总优化耗时约 1 秒，根本无法满足实时需求。最近的工作开始探索用前馈网络直接预测高斯属性，替代传统的逐帧优化——将耗时的梯度下降优化替换为单次前向推理，同时利用大规模预训练获得更强的几何先验。

% [C.6a] Flash-Mono 五要素展开：三大瓶颈动机 → 有状态transformer循环前馈架构(ViT+交叉注意力+DPT头) → 子地图划分与隐藏状态回环闭合/Sim(3)优化、2D surfels替代3DGS、轻量后端20次迭代 → 三阶段课程训练+DL3DV/ScanNet++ → 10+ FPS实时SOTA跟踪渲染
Flash-Mono（2026.04）~\cite{2604_03092} 在上述 baseline 的基础上，系统分析了阻碍实时单目 GS-SLAM 的三大瓶颈：\textbf{动机。}传统的 Train-from-Scratch 范式要求每关键帧执行数百次优化迭代（单次约 20 ms），总计约 1 秒/关键帧，根本无法满足实时需求；增量前馈方法虽然快速但累积位姿与尺度漂移；标准 3DGS 的三维椭球体表示缺乏显式表面约束。为此，Flash-Mono 提出 Predict-and-Refine 范式——用前端前馈网络直接预测高质量高斯属性，后端仅需极少迭代精化，核心设计思想是``用强预测替代弱初始化''。

\paragraph{循环前馈前端模型}Flash-Mono 的前端核心是一个有状态 transformer 循环前馈模型，直觉上可以理解为：模型不仅要``看''当前帧，还要``记住''之前看过的场景——这个``记忆''就是持久隐藏状态 $\mathbf{M}_{t-1}$，它逐步积累局部多帧的几何与视觉上下文，作为跨帧信息传递的桥梁。具体来说，每帧 RGB 图像 $I_t \in \mathbb{R}^{H \times W \times 3}$ 经 ViT 编码器提取视觉 token $\mathbf{F}_t \in \mathbb{R}^{K \times C}$，随后两个互连解码器通过交叉注意力在当前视觉 token 与持久隐藏状态 $\mathbf{M}_{t-1}$ 之间双向交换信息（初始状态 $\mathbf{M}_0$ 为零向量），联合输出三项预测：

\begin{equation}
\hat{\mathbf{T}}_t, \hat{\mathcal{G}}_t, \mathbf{M}_t = f(I_t, \mathbf{M}_{t-1})
\end{equation}

其中：(a) 相机位姿 $\hat{\mathbf{T}}_t \in SE(3)$（表示从当前相机坐标系到初始帧坐标系的变换），由与视觉 token 拼接的可学习 pose token $\mathbf{z}_t$ 经 MLP 头提取——pose token 的作用是让模型在交叉注意力中汇聚场景几何线索用于位姿推理；(b) 逐像素 2D Gaussian surfel（即平面二维高斯基元，每个 surfel 定义在当前相机坐标系下，属性包括均值 $\hat{\mu}$、置信度 $\hat{w}$、颜色 $c$、不透明度 $\alpha$、旋转 $r$ 和 2D 缩放 $s$），由两个 DPT 头（Dense Prediction Transformer，一种将 transformer 特征解码为逐像素密集预测的经典架构）分别解码均值/置信度和其他属性；(c) 更新后的隐藏状态 $\mathbf{M}_t$，携带着聚合的多帧几何与视觉上下文向前传播。

模型的融合过程可形式化为：

\begin{align}
\mathbf{F}_t &= \mathrm{Encoder}(I_t) \\
\tilde{\mathbf{F}}_t, \tilde{\mathbf{z}}_t, \mathbf{M}_t &= \mathrm{Decoders}((\mathbf{F}_t, \mathbf{z}_t), \mathbf{M}_{t-1})
\end{align}

解码器通过交叉注意力实现双向信息交换：当前帧的视觉 token 可以查询历史隐藏状态中的几何上下文，同时隐藏状态也吸收当前帧的新信息。这一设计使模型既能利用历史帧的累积知识（类似 SLAM 中的地图），又保持流式处理能力（不同于 VGGT 等 batch 方法需要一次性处理所有帧）。

\paragraph{训练目标与三阶段课程}模型在 DL3DV 和 ScanNet++ 上以三阶段课程训练，总损失函数为三项加权和：

\begin{equation}
\mathcal{L}_{\mathrm{total}} = \lambda_{\mathrm{pose}}\mathcal{L}_{\mathrm{pose}} + \lambda_{\mathrm{geo}}\mathcal{L}_{\mathrm{geo}} + \mathcal{L}_{\mathrm{render}}
\end{equation}

其中位姿损失 $\mathcal{L}_{\mathrm{pose}} = \sum_{t=1}^{L}(\|\hat{q}_t - q_t\|_2 + \|\hat{t}_t - t_t\|_2)$ 监督四元数旋转和平移；几何损失 $\mathcal{L}_{\mathrm{geo}} = \sum_{t=1}^{L}\sum_{n=1}^{HW}\hat{w}_{t,n}\|\hat{p}_{t,n} - p_{t,n}\|_2^2 - \log(\hat{w}_{t,n})$ 以置信度加权的 $L_2$ 距离监督逐像素点云预测（$p_{t,n}$ 由 GT 深度图经相机内参反投影得到）；渲染损失 $\mathcal{L}_{\mathrm{render}}$ 包含 MSE、LPIPS 和深度 $L_2$ 三项，且分为逐帧渲染监督和全局合并渲染监督两部分——后者将整个序列的 2DGS 预测合并为全局地图后渲染，避免模型为规避帧间冲突而过度缩小单个高斯尺度。三阶段课程策略为：第一阶段冻结骨干网络（从 CUT3R 初始化），仅训练 2DGS 属性预测头，学习率 $2 \times 10^{-4}$，序列长度 1--4 帧（5,000 步热身）；第二阶段解冻结解码器、位姿头和均值预测头，学习率降至 $1 \times 10^{-5}$，序列长度固定为 4 帧，使模型适应任务数据分布并缓解循环架构中的梯度爆炸；第三阶段进一步降低学习率至 $5 \times 10^{-6}$，序列长度扩展至 32 帧，训练模型处理更长的时序依赖。

Flash-Mono 的三项关键技术贡献使其在速度、精度和一致性上同时取得突破。

\paragraph{子地图划分与隐藏状态回环闭合}
循环模型虽然天然适合流式处理，但存在一个根本问题——灾难性遗忘：随着序列增长，累积漂移不可避免地增大。为缓解这一问题，输入流被划分为短子地图（每个子地图重新初始化隐藏状态），子地图间保留一帧重叠以链接局部轨迹——通过重叠帧的相对变换可将局部位姿估计链接为连续轨迹。每个子地图的最终隐藏状态被缓存为紧凑的子地图描述符，存入``隐藏状态袋''（bag of hidden states）。

回环闭合的直觉是：隐藏状态 $\mathbf{M}$ 实际上是场景局部几何和视觉信息的紧凑摘要——如果相机回到了曾经到过的地方，加载对应的``历史记忆''就应该能让模型在那个旧坐标系下重新定位当前帧。当外观方法检测到回环候选（当前关键帧 $I_j$ 属于子地图 $C_b$，历史关键帧 $I_i$ 属于子地图 $C_a$）时，系统取出历史子地图的缓存隐藏状态 $\mathbf{M}_a$，与当前帧做单次前向推理 $f(I_j, \mathbf{M}_a)$——在 $\mathbf{M}_a$ 的条件下，模型将 $I_j$ 解读到子地图 $C_a$ 的坐标系下，输出重定位位姿 $\hat{\mathbf{T}}_j^a \in SE(3)$ 和对应的点云 $\mathcal{P}_j^a = \{\hat{p}_k^a\}$。相对位姿计算为 $\hat{\mathbf{T}}_{ji} = (\hat{\mathbf{T}}_j^a)^{-1}\hat{\mathbf{T}}_i^a$。

为解决单目系统的尺度歧义，Flash-Mono 对同一帧 $I_j$ 在两个隐藏状态下生成的点云进行比较：$\mathcal{P}_j^a$（历史隐藏状态）和 $\mathcal{P}_j^b = \{\hat{p}_k^b\}$（当前隐藏状态）。关键洞察是：两个点云都来自同一帧图像、处于相机坐标系下，因此仅相差一个尺度因子。相对尺度通过最小二乘求解：

\begin{equation}
s^* = \arg\min_s \sum_k \| \hat{p}_k^b - s \cdot \hat{p}_k^a \|^2
\end{equation}

最终，尺度与相对位姿组合为完整的 $\mathrm{Sim}(3)$ 回环约束：

\begin{equation}
\mathbf{H}_{ji} = \begin{bmatrix} s^*\mathbf{R}_{ji} & \mathbf{t}_{ji} \\ \mathbf{0}^T & 1 \end{bmatrix}
\end{equation}

全局位姿图在 $\mathrm{Sim}(3)$ 空间优化，节点为关键帧位姿 $\mathbf{T}_k^W \in \mathrm{Sim}(3)$，边包含三类约束：帧间约束（子地图内连续帧的相对变换）、子地图间对齐约束（相邻子地图共享帧的相对尺度估计）和回环闭合约束。全局最优位姿通过最小化非线性最小二乘代价函数求解：

\begin{equation}
\mathbf{T}^{W*} = \arg\min_{\mathbf{T}^W} \sum_{(i,j) \in \mathcal{E}} \| \log(\mathbf{H}_{ji}^{-1}((\mathbf{T}_i^W)^{-1}\mathbf{T}_j^W)) \|^2_{\Sigma}
\end{equation}

其中残差在李代数 $\mathfrak{sim}(3)$ 中计算，通过 GTSAM 高效求解。

\paragraph{2D Gaussian surfels 替代 3DGS}
Flash-Mono 采用平面 2D Gaussian surfel（二维高斯基元）作为地图表示。每个 surfel 定义在相机坐标系下，属性包括位置 $\mu$、颜色 $c$、不透明度 $\alpha$、旋转 $r$ 和 2D 缩放 $s$。渲染通过体积 $\alpha$-混合实现：

\begin{equation}
w_i(\mathbf{p}) = \alpha_i \exp\left(-\frac{1}{2}(\mathbf{p}-\mu_i)^T\boldsymbol{\Sigma}_i^{-1}(\mathbf{p}-\mu_i)\right), \quad (\hat{C}, \hat{D}, \hat{A}) = \sum_{i=1}^{N}(c_i, z_i, 1)w_i \prod_{j=1}^{i-1}(1-w_j)
\end{equation}

其中 $\mathbf{p}$ 为像素坐标，$\boldsymbol{\Sigma}_i \in \mathbb{R}^{2 \times 2}$ 是由旋转 $r_i$ 和 2D 缩放 $s_i$ 诱导的屏幕空间协方差，$z_i$ 为 surfel 沿光轴的深度。相比 3D 高斯椭球体，平面 2D 表示提供更强的表面先验——surfel 天然贴合表面而非悬浮在空间中——有效抑制浮游伪影（floaters）并提升几何保真度。这在 SLAM 场景中尤为关键，因为微小的几何不一致会通过位姿估计误差快速累积为漂移。

\paragraph{轻量级后端精化}
前端的高质量预测为后端提供了强几何与外观先验，因此每个关键帧仅需 20 次优化迭代（对比 MonoGS 和 S3PO-GS 的 250 次），实现 10 倍加速。后端在独立线程中运行，管线包含四个阶段：
(1)~自适应体素化——前端预测的逐像素 2DGS 可能过于密集，为降低内存开销，将地图划分为 $2 \times 2$ 的 2DGS 基元块，块内基元通过属性平均合并为一个基元：

\begin{equation}
\mu_{\mathrm{merged}} = \frac{1}{N}\sum_{n=1}^{N}\mu_n, \quad r_{\mathrm{merged}} = \frac{\sum_{n=1}^{N}(r_n, r_1)}{\|\sum_{n=1}^{N}(r_n, r_1)\|}
\end{equation}

其中四元数平均采用标准化内积确保旋转一致性。深度变异超过阈值 $\tau_d$ 的块被保留以维持几何细节。该模块将基元数量从 1.35M 减少至 0.56M（减少超过 58\%），PSNR 仅从 19.70 降至 19.44。
(2)~全局地图融合——先渲染现有全局地图，剪除 RGB 或深度重建误差高的基元，再将体素化后的新基元从相机坐标系变换到世界坐标系（$\mu_{\mathrm{world}} = s_k\mathbf{R}_k\mu_{\mathrm{cam}} + \mathbf{t}_k$），仅在累积图显示未充分重建的区域添加新基元。
(3)~轻量精化——仅精化最近 $K$ 个关键帧关联的局部地图区域，20 次迭代即可收敛。
(4)~回环后地图校正——将高斯刚体绑定到其原始关键帧，当关键帧位姿从 $\mathbf{T}_{\mathrm{old}}$ 更新为 $\mathbf{T}_{\mathrm{new}}$ 时，通过增量变换 $\Delta \mathbf{T} = \mathbf{T}_{\mathrm{new}}\mathbf{T}_{\mathrm{old}}^{-1}$ 即时变形地图，避免了耗时的重新渲染优化。

实验结果表明，Flash-Mono 在 ScanNet（in-domain）、BundleFusion（out-of-domain）和 KITTI（大规模户外）三个挑战性数据集上均取得 SOTA 跟踪精度和渲染质量，同时以 10+ FPS 实时运行。系统总参数量 795.7M，推理时显存占用约 3GB（RTX 4090），前端单帧推理 62\,ms，后端 77.5\,ms/帧，FPS 计算包含前端推理、后端建图、回环检测和位姿图优化等所有开销。消融实验验证了各组件的有效性：无后端精化时 PSNR 为 20.14，10 次精化即提升至 22.41，证明前端的强初始预测；子地图长度为 8 帧时跟踪误差最低（ATE RMSE 0.106），过短缺乏时序上下文，过长则因 RNN 遗忘导致子地图内漂移累积（16 帧时误差已明显上升）；隐藏状态回环闭合在跟踪精度上大幅优于 PnP+RANSAC 基线和无回环配置，验证了将隐藏状态作为紧凑场景描述符的有效性。模型加速实验表明，通过 float16 量化和 CUDA Graphs 优化，前端推理延迟从 283\,ms 降至 85\,ms（3.33$\times$ 加速），在笔记本级 RTX 4060（8GB）上的推理时间（83\,ms）已接近桌面级 RTX 4090（62\,ms）。

M$^3$（2026.03）~\cite{2603_16844} 同样利用多视角基础模型增强单目 SLAM，但走了不同的技术路线。

\textbf{动机。}Flash-Mono 通过循环隐藏状态积累跨帧几何上下文，但循环架构存在固有的灾难性遗忘问题，需要子地图划分来缓解；M$^3$ 则从另一个角度解决多帧一致性问题——通过为核心骨干 Pi3X 增加密集匹配头，使基础模型直接输出逐像素稠密特征描述符，从而在不同帧之间建立精确的逐像素对应关系。

\paragraph{密集匹配头与训练策略}匹配头的结构为 DPT 解码块（$DPT_\mathrm{desc}$）串联一个两层 MLP，中间以 GELU 激活，输出经 L2 归一化为单位向量以确保匹配稳定性。训练时仅微调匹配头，冻结骨干的编码器、解码器和已有预测头——$DPT_\mathrm{desc}$ 以点预测头的预训练参数初始化，利用几何预测与密集匹配之间的结构知识共享加速收敛。损失函数采用对称 InfoNCE 目标，在批次内以第一帧为参考，通过 GT 三维坐标重合建立逐像素对应真值，以双向 softmax 交叉熵鼓励互描述符一致性，附加匹配置信度正则项。训练 200K 迭代（AdamW + 余弦学习率，峰值 $1 \times 10^{-4}$），每批 8 帧。

\paragraph{由粗到细的稠密对应与动态区域抑制}SLAM 推理时，密集匹配头为每帧输出描述符图和点图。对帧对 $(I_i, I_j)$，M$^3$ 先利用基础模型预测的位姿将 $I_i$ 的局部点图变换到 $I_j$ 的坐标系下并投影至图像平面，获得粗略对应位置；再在投影点半径 $r=4$ 的局部邻域内搜索描述符余弦相似度最高的像素，精化对应关系。这种由粗到细的策略将匹配计算从全局二次搜索降为局部线性精化，兼顾精度和效率。此外，M$^3$ 引入基于描述符一致性的动态区域检测：将参考关键帧的描述符图投影到当前帧后，比较投影描述符与当前帧预测描述符的像素级点积相似度，乘以历史运动图实现时序一致性调制，生成运动图 $M_i \in [0,1]^{H \times W \times 1}$。低一致性像素（即动态物体或遮挡区域）在优化中被降权，抑制重影伪影和轨迹漂移。

\paragraph{紧耦合流式 SLAM 框架}与 Flash-Mono 的子地图划分策略不同，M$^3$ 采用滑动窗口（长度 $L=8$，含最多 4 个历史关键帧和 4 个新帧）一次性送入增强版 Pi3X，单次前向推理同时产出前端跟踪和后端全局优化所需的全部信息——历史帧的点图和描述符用于更新全局因子图，新帧的预测提供初始位姿、点图和描述符用于实时跟踪与关键帧选择。帧分类沿用 ARTDECO 的三分策略（关键帧/建图帧/普通帧），通过对应比率阈值 $\tau_k$ 和像素位移百分位数阈值 $\tau_m$ 自动区分。前端跟踪以密集对应驱动的重投影误差最小化实现：相对位姿 $\mathbf{T}_{kf}$ 通过最小化 Huber 核加权的重投影残差 $\sum_{(m,n)} \|p_{k,n} - \pi(\mathbf{T}_{kf}\mathbf{X}^f_{f,m})\|w(q_{m,n})$ 估计，其中密集对应 $M(I_f, I_k)$ 来自匹配头的逐像素预测。为解决基础模型不同推理间因尺度歧义导致的内参不一致问题，M$^3$ 以首次推理的 RANSAC 内参为参考，后续推理按焦距比对齐内参并重缩放点图。全局优化在 $\mathrm{Sim}(3)$ 流形上联合优化所有关键帧位姿（与 Flash-Mono 类似处理尺度漂移），以因子图维护帧间约束边，最小化集体重投影误差。3DGS 地图采用 Neural Gaussian（含个体特征 $f_l$ 和区域特征 $f_g$），从关键帧和建图帧通过 LoG 概率图初始化，优先关注高频细节和重建薄弱区域，同时严格排除动态像素。序列完成后执行 10K 次全局精化，联合优化位姿和所有高斯参数。

在实验方面，M$^3$ 在 ScanNet++、ScanNetV2、VR-NeRF（室内场景）和 KITTI、Waymo、FAST-LIVO2（户外场景）共 60 个场景上取得了 SOTA 位姿精度和渲染质量。ATE RMSE 相比 VGGT-SLAM 2.0 降低 64.3\%，相比 MASt3R-SLAM 大幅领先。消融实验验证了各设计选择：搜索半径 $r=0$ 时精度明显下降，证明局部精化的必要性；去掉内参对齐后粗匹配退化；去掉全局优化后定位精度降低；用 Pi3 替换 Pi3X 有小幅性能损失；前端后端解耦引入冗余推理、降低稳定性；窗口大小 $L=8$ 在时序稳定性和计算开销间取得最佳平衡，验证了整体架构的合理性。

MonoEM-GS（2026.04）~\cite{2604_10593} 同样面向单目 RGB SLAM，但与 GaussianFlow SLAM 的端到端可微渲染优化范式不同，它将 3DGS 同时视为高斯混合模型（GMM，Gaussian Mixture Model）和泼溅渲染基元，用\textbf{增量 EM} 代替基于梯度的联合优化。其设计动机在于：前向几何预测模型（如 MapAnything）虽然能从单张 RGB 图像估计点云和相机位姿，但逐帧累积预测不可避免地引入噪声和歧义结构；EM 框架通过概率化数据关联（而非硬匹配）将新测量稳健地融合到已有地图中，从而获得更几何一致的重建。具体而言，每帧的 MapAnything 前向推理产出点云、位姿和 DINOv3 语义特征，位姿经子地图裁剪和彩色 ICP 对齐后固定，地图更新则由两阶段完成。

第一阶段为 EM 融合：对新预测点找到 $K$ 个最近高斯，以 DINOv3 特征余弦相似度 $d_\text{cos}$、法线一致性 $d_n$ 和 Mahalanobis 距离 $d_\sigma$ 的组合计算责任权重 $r^{ik} \propto \exp(-\tfrac{1}{2}d_\sigma^2 + d_\text{det} + \lambda_1(d_n-1) + \lambda_2(d_\text{cos}-1))$，只保留最佳解释置信度超过阈值 $\tau$ 的对应，随后以加权运行平均 $\mathcal{M}^*_t \leftarrow (1-\gamma_t)\mathcal{M}^*_t + \gamma_t \mathcal{M}^*_{t,\text{new}}$ 更新高斯参数，其中等向性分数 $\gamma_k$ 衡量测量点在法线切平面上是否均匀分布于高斯两侧，防止单向观测导致退化。第二阶段为光栅化精化：将 EM 未融合的预测点初始化为新高斯 $G_t$，对 $G_t$ 和子地图 $\mathcal{M}^*_t$ 分别执行高斯泼溅渲染，利用渲染深度检测未映射区域并合并同面高斯，确保地图覆盖完整。此外，每个高斯在初始化时即绑定 DINOv3 语义特征并冻结，支持即位开放集 3D 语义分割——在 Replica 数据集上以 f-mIoU 指标取得最优。在 TUM RGB-D 和 7-Scenes 上，MonoEM-GS 的位姿精度略逊于 MASt3R-SLAM 等优化型基线，但法线一致性和重建几何质量更优，地图更稀疏却结构更清晰。

GaussianFlow SLAM（2026.04）~\cite{2604_15612} 在 MonoGS 的光度损失跟踪框架上引入光流作为额外的几何监督信号，解决单目设定下缺乏深度约束时地图与位姿互相退化的自增强循环问题。

\textbf{动机。}在单目设定下，没有深度传感器的约束，高斯位置和相机位姿的优化是高度耦合的——不准确的位姿导致高斯放置错误，错误的高斯又反过来恶化位姿估计，形成互相退化的自增强循环。核心概念\textbf{GaussianFlow}定义了高斯在连续帧间的投影运动——将帧 $I_t$ 中渲染某像素的高斯集合归一化到规范空间，再反归一化到帧 $I_{t+1}$ 的图像空间，所得位移即为该像素的 GaussianFlow $\boldsymbol{\zeta}_t$。系统鼓励 $\boldsymbol{\zeta}_t$ 与预训练网络估计的光流 $\mathbf{X}_t$ 对齐，从而在不依赖深度传感器的情况下为高斯位置和相机位姿提供一致的几何约束。

为实现可扩展的 kernel 级优化，作者推导了 GaussianFlow 损失关于高斯参数（均值、协方差）和 $\mathrm{SE}(3)$ 位姿的\textbf{闭环解析梯度}——对投影协方差 $\boldsymbol{\Sigma}'$ 进行特征分解，得到 $\boldsymbol{\Sigma}'^{\pm 1/2}$ 关于特征向量和特征值的显式导数，可直接嵌入 3DGS 的 GPU 光栅化 kernel 进行反传，避免了自动微分的计算图开销。在跟踪阶段，初始位姿由光度损失与前后两帧的 GaussianFlow 损失联合估计，随后通过卷积门控循环单元（Convolutional GRU, ConvGRU）循环更新模块执行基于光流的稠密光束平差（DBA）进行二阶位姿精化——GaussianFlow $\boldsymbol{\zeta}_t$ 作为 ConvGRU 输入预测光流 $\mathbf{X}_t$，后者引导位姿更新，更新后的地图又重新生成 $\boldsymbol{\zeta}_t$，形成互补收敛循环。在建图阶段，损失函数联合光度损失、GaussianFlow 损失、各向同性正则化和不透明度熵损失，位姿固定仅优化高斯参数。

此外，针对 GaussianFlow 持续几何校正可能引入的欠致密化或浮点伪影，作者提出\textbf{归一化误差}致密化与剪枝策略：将每个高斯的逐像素误差按其屏幕空间密度贡献归一化（消除大高斯天然高误差的偏差），对归一化 DSSIM 误差或光流误差超阈值的高斯分别进行分裂或剪枝，比 MonoGS 的不透明度阈值剪枝更精细地清理不稳定基元。在 TUM RGB-D 和 EuRoC 数据集上，GaussianFlow SLAM 的跟踪精度和渲染质量均优于已有单目 3DGS-SLAM 方法。

\subsubsection{户外与多相机 RGB SLAM}

本路线介绍面向户外真实环境和多相机配置的 RGB SLAM 方法，应对尺度漂移、剧烈光照变化和无界场景等挑战。

从实验室室内场景扩展到户外真实环境，需要应对尺度漂移、剧烈光照变化和无界场景等新挑战。本小节讨论面向户外场景和多相机配置的 RGB SLAM 方法。

S3PO-GS（2025.07）~\cite{2507_03737} 针对户外单目 RGB SLAM 的尺度漂移问题，提出自洽的 3DGS 点图追踪与 patch 尺度对齐建图。

\textbf{动机。}引入外部预训练模型补充几何信息虽可增强鲁棒性，却带来外部模块与 3DGS 地图之间的尺度对齐难题——大旋转大位移场景下，对齐误差累积导致尺度漂移，使得位姿估计偏离真实轨迹。S3PO-GS 的核心思路是：让位姿估计直接锚定在 3DGS 地图自身的尺度上，预训练模型仅充当像素对应的桥梁，不参与位姿计算，从根本上绕过尺度对齐问题。

\textbf{自洽 3DGS 点图追踪。}与 MonoGS 用解析雅可比直接优化位姿不同，S3PO-GS 将跟踪分解为两步：先用 PnP+RANSAC 估计初始位姿，再用可微渲染精化。初始位姿估计的关键是建立当前帧与相邻关键帧之间的 2D-3D 对应。具体地，从相邻关键帧视点 $T_{\mathrm{ak}}$ 出发，利用 3DGS 的 $\alpha$-混合 光栅化渲染深度图 $D_{\mathrm{ak}}$，再通过相机内参逆变换得到渲染点图 $X_{\mathrm{ak}}^r$——该点图的 3D 坐标完全来自 3DGS 地图，尺度与重建场景严格一致。与此同时，将当前帧和相邻关键帧输入预训练的点图模型（MASt3R/DUSt3R），获得两组归一化尺度的点图 $X_{\mathrm{ak}}^p$, $X_n^p$ 及置信度分数。通过最小化点图距离建立像素对应，过滤低置信度点后，利用已建立的对应关系将渲染点图 $X_{\mathrm{ak}}^r$（3D 坐标来自地图）与当前帧（2D 像素）连接起来，形成 2D-3D 对应。由于 $X_{\mathrm{ak}}^r$ 的尺度直接源自 3DGS 场景，PnP 解算的位姿天然保持正确尺度，无需额外对齐。初始位姿随后通过 MonoGS 式的解析雅可比光度优化精化——但得益于点图追踪提供的精确初始值，仅需 MonoGS 迭代次数的 10\% 即可收敛。

\textbf{Patch 尺度对齐与点图替换建图。}单目 SLAM 缺乏几何监督，常见做法是引入单目深度先验，但跨帧深度估计存在尺度不一致。S3PO-GS 的建图模块用预训练点图作为几何先验，但需要将点图的归一化尺度对齐到当前 3DGS 场景尺度。直接全局归一化会因渲染点图中的离群值和分布不一致导致对齐失败。S3PO-GS 为此设计了 patch 级局部对齐：将渲染点图 $X^r$ 和预训练点图 $X^p$ 切分为 $P \times P$ 的 patch，计算每个 patch 的均值 $\mu$ 和标准差 $\sigma$，选取分布相似的 patch（满足 $|\mu_r - \mu_p| < \beta\,\sigma_p$）进行归一化。在归一化空间中筛选"正确点"集合 $\mathcal{CP}$（两图差异小于阈值 $\tau_r$），由此计算尺度因子 $\gamma' = \overline{X^r[\mathcal{CP}]} / \overline{X^p[\mathcal{CP}]}$，迭代至稳定。在建图时，对齐后的预训练点图用于两方面：一是点图替换——当渲染点图中某点与对齐预训练点图的偏差超过阈值时，用后者替换前者，再经随机稀疏下采样初始化新高斯，减少错误点的引入；二是几何监督损失 $\mathcal{L}_{\mathrm{geo}} = \|X^r - \hat{X}^p\|_1$，与 MonoGS 的光度损失和各向同性正则化联合优化，在局部关键帧窗口内同时精化位姿和高斯参数。

该方法在 Waymo、KITTI 和 DL3DV 数据集上取得了 3DGS SLAM 框架下的最优跟踪精度和新视角合成质量。OpenGS-SLAM（2025.02）~\cite{2502_15633} 同样利用预训练点图回归网络（DUSt3R（Dense Unconstrained Stereo 3D Reconstruction））处理户外单目 RGB SLAM，但与 S3PO-GS 的"渲染点图桥接 + PnP"跟踪范式不同，它采用更直接的 pairwise 策略：将相邻两帧输入共享参数的双分支 transformer 网络，跨分支的信息共享使输出的两组点图天然对齐，从中提取 2D-3D 对应后直接用 RANSAC+PnP 解算帧间相对位姿 $T^k_{\mathrm{trans}}$，再级联得到绝对位姿 $T^k = T^k_{\mathrm{trans}} T^{k-1}$。这种设计的代价是位姿尺度完全依赖点图回归的尺度一致性，而非像 S3PO-GS 那样锚定在 3DGS 地图尺度上。

为此，OpenGS-SLAM 提出\textbf{自适应尺度映射器}：利用连续三帧 $(k{-}1, k, k{+}1)$ 生成的两对交叉视角点图，匹配相同 3D 点并计算逐对尺度比 $\sigma_{ij} = \|X_i'^{k} - X_j'^{k+1}\| / \|X_i^{k-1} - X_j^{k}\|$，以首帧为参考累积乘积得到全局尺度因子 $S_k = S_{k-1} \cdot \frac{1}{N}\sum_{ij} \sigma_{ij}$，在插入新高斯前对点图坐标进行尺度校正。这一策略与 S3PO-GS 的 patch 级局部对齐形成对比——后者在渲染点图与预训练点图之间做局部归一化，而 OpenGS-SLAM 在纯预训练点图之间做跨帧累积对齐，无需渲染点图参与。此外，针对户外车辆前视相机不复访场景的特点（不同于室内 SLAM 的学习率单调衰减），OpenGS-SLAM 设计了\textbf{自适应学习率调整}：根据相邻关键帧的旋转角 $\theta_{\mathrm{rad}} = \arccos\!\bigl(\frac{\mathrm{trace}(R_{\mathrm{diff}}) - 1}{2}\bigr)$ 动态重置累积迭代数 $N_{\mathrm{iter}}^{\mathrm{new}} = N_{\mathrm{iter}} \cdot (1 - \theta / 90^\circ)$，转弯时学习率回升以快速适应新场景，直行时持续衰减以精化已有区域。在 Waymo 数据集上，OpenGS-SLAM 的平均 PSNR 比 MonoGS 提升 10\%，跟踪精度接近采用全局 BA 的 GlORIE-SLAM，但无需后端优化。消融实验表明，移除自适应尺度映射器导致跟踪和渲染同时退化，移除自适应学习率调整则引起转弯处轨迹漂移。

TVG-SLAM（2025.06）~\cite{2506_23207} 同样基于三帧结构进行几何推理，但与 Splat-SLAM 的光流-DSPO 跟踪范式截然不同，其核心贡献是将\textbf{三视角几何}引入跟踪和建图的全流程。

\textbf{动机。}双视角几何（如光流或极线约束）在某些退化配置下（如共线运动）会失去可观性，而三视角约束（如三焦点张量）在这些退化配置下仍能提供有效的几何约束，特别适合户外低视差场景。在匹配层面，TVG-SLAM 使用 DUST3R（2024）~\cite{wang2024dust3r} 生成密集点图，以中间关键帧为桥梁将两对帧间匹配 $(I_{k-1}, I_k)$ 和 $(I_k, I_t)$ 桥接为三视角对应 $(p_{k-1}, p_k, p_t)$，通过共享像素的交叉验证过滤不一致对应，比 Splat-SLAM 的光流对应多了几何冗余。

在跟踪层面，TVG-SLAM 构建混合损失 $\mathcal{L}_{\mathrm{track}} = \lambda_p \mathcal{L}_{\mathrm{photo}} + \lambda_{2D} \mathcal{L}_{2D} + \lambda_{3D} \mathcal{L}_{3D}$，其中光度项与 Splat-SLAM 类似，但两个几何项是全新设计：$\mathcal{L}_{2D}$ 基于三焦点张量（trifocal tensor）计算当前帧中观测点到另外两视角确定的极线的代数残差，提供纯 2D 几何约束而无需显式三维重建——三视角约束在共线运动等退化配置下仍保持可观性，特别适合户外低视差场景；$\mathcal{L}_{3D}$ 利用 DUST3R 点图提供的两个独立 3D 估计，通过 Procrustes 对齐后计算 3D-3D 对应残差。为应对异步跟踪-建图中建图线程滞后导致渲染图过时的问题，TVG-SLAM 提出 DART（Dynamic Attenuation of Rendering Trust）机制：以上次关键帧以来的帧数 $\Delta N_f$ 作为地图新鲜度代理，通过 sigmoid 函数 $\lambda_p = w_{\min} + (w_{\max} - w_{\min})\sigma(k(\Delta N_f - N_m))$ 平滑衰减光度权重——地图新鲜时 $\lambda_p$ 保持高值，地图过时时权重降低，系统自动依赖不受地图质量影响的三视角几何约束。

在建图层面，TVG-SLAM 提出 TUGI（Tri-view Uncertainty-guided Gaussian Initialization）：利用三视角匹配中同一场景点的多个独立 3D 估计计算方差 $\sigma_g^2$，以此编码高斯的初始缩放（$\propto \sigma_g^2$）和不透明度（$\alpha_{\mathrm{new}} = \sigma^{-1}(a(1 - k\sigma_g^2))$，不确定性越大不透明度越低），使不确定区域的高斯以更宽的支撑开始优化，避免渲染伪影。实验在 Waymo、Small City 和 Cambridge Landmarks 三个户外数据集上进行，ATE 相比 OpenGS-SLAM 分别降低 28.2\%、65.7\% 和 69.0\%，渲染 PSNR 亦全面领先。消融实验表明，移除 TUGI 导致跟踪精度下降最大（ATE 增加 38.3\%），移除 DART 使 ATE 增加 21.0\%，两个几何约束互补缺失任一均显著降低精度。

MCGS-SLAM（2025.09）~\cite{2509_14191} 将 S3PO-GS 的单目设定推广至多相机 rig（3--5 个已标定相机），核心挑战从"单目尺度漂移"转为"多相机间的跨视角深度与尺度一致性"。

\textbf{动机。}多相机 rig 中每个相机独立运行单目深度估计，各相机的深度预测存在独立的空间变化尺度偏移，直接融合会导致跨视角深度不一致，进而引起几何退化。，MCGS-SLAM 设计了多相机光束调整（MCBA）模块，将时序对（同相机跨时间）和跨视角对（同时间不同相机）的投影残差统一到一个光束调整框架中。时序对通过已知外参 $\mathbf{T}_C^B$ 将体位姿变换映射到相机坐标系，跨视角对则直接利用预标定的相机间外参 $\mathbf{T}_{C_i C_j}$，使位姿更新 $\Delta\boldsymbol{\xi}$ 和深度更新 $\Delta\mathbf{d}$ 通过 Schur complement 联合求解。在尺度对齐方面，S3PO-GS 的 patch 级局部对齐仅在单帧内操作，而多相机设定中各相机的单目深度预测（Metric3Dv2）存在独立的空间变化尺度偏移。

MCGS-SLAM 为此引入联合深度-尺度对齐（JDSA），为每个关键帧分配一个可学习的 $m \times n$ 尺度网格，通过双线性插值生成逐像素尺度因子，将预测深度重缩放为优化深度；JDSA 与 MCBA 交替运行，既利用多视角光度-几何一致性约束尺度，又用尺度对齐损失 $\|\tilde{d}_i \cdot B_i - d_i\|^2$ 保证对齐稳定，避免了直接联合优化尺度与位姿导致的不收敛问题。建图阶段，当 MCBA 更新关键帧位姿 $\Delta\mathbf{T}_t$ 时，所有锚定于该帧的高斯的均值和协方差同步传播（$\boldsymbol{\mu}_i \leftarrow \Delta\mathbf{T}_t \cdot \boldsymbol{\mu}_i$，$\Sigma_i \leftarrow \mathbf{R}\Sigma_i\mathbf{R}^\top$），避免重新初始化或重新渲染。在 Waymo 数据集上，MCGS-SLAM 平均 ATE 1.298\,m（HI-SLAM2 为 1.473\,m），平均 PSNR 26.50\,dB（DROID-Splat 为 25.62\,dB）；在 AirSim 低视差 Garden 场景中 PSNR 达 29.36\,dB，较所有单相机基线高约 4\,dB，验证了多视角互补的增益。

BGS-SLAM（2025.07）~\cite{2507_23677} 同样面向户外场景，但与 S3PO-GS 的单目 + 点图追踪范式截然不同，从三个层面重新设计了系统。

\textbf{动机。}单目 SLAM 存在固有的尺度歧义，即使引入预训练深度估计也难以获得可靠的绝对尺度；双目立体相机通过已知基线直接提供度量深度，从根本上消除了尺度漂移问题，使系统更适合户外大尺度场景。

\textbf{传感器层面}，S3PO-GS 仅用单目 RGB，需预训练点图模型补充几何先验并解决尺度对齐问题；BGS-SLAM 直接采用双目立体相机，通过预训练深度立体网络（MonSter）生成稠密视差图，经标准三角化公式 $\hat{D} = bf/d$ 转换为度量深度——立体几何天然提供绝对尺度，从根本上消除了单目 SLAM 的尺度漂移和 patch 对齐需求。相比 LiDAR，立体深度在空间密度和场景覆盖上更具优势（LiDAR 稀疏但精确，立体稠密但有噪声）。

\textbf{跟踪层面}，S3PO-GS 将追踪嵌入 3DGS 可微渲染管线，用 PnP+RANSAC + 光度优化精化位姿；BGS-SLAM 则完全解耦跟踪与建图，采用 ORB-SLAM2 作为外部追踪器——提取左右图的 ORB 特征进行立体匹配，通过帧级 BA 和滑窗局部 BA 联合优化位姿与 3D 地图点。这一设计牺牲了追踪与建图之间的紧耦合，但换来了户外大尺度场景中远比可微渲染更鲁棒的位姿估计（SplaTAM、MonoGS 等方法在 KITTI 完整序列上因跟踪发散而失败）。

\textbf{建图层面}，BGS-SLAM 的核心贡献是针对户外深度图高维 $L_1$ 优化中局部极小值过多的问题，提出压缩采样策略：利用户外深度图在差分域的稀疏性（$K$-sparse），通过满足 RIP 条件的感知矩阵 $P \in \{0,1\}^{m \times d}$（$m = O((K+S)\log d)$，实验中均匀采样约 25\% 像素），将高维优化转化为凸优化问题 $\min_{d,E}\|d\|_1 + \lambda\|E\|_1$，同时去除大像素级误差并保留结构信息，使局部极小值数量从 $O(n^d)$ 降至 $O(n \cdot d)$。

映射损失组合了五项：RGB 损失（$L_1$ + SSIM，$\lambda_{\mathrm{rgb}}=0.8$, $\lambda_{\mathrm{ssim}}=0.2$）、加权几何损失（以 RGB 梯度 $g_{\mathrm{rgb}} = \|\nabla I\|$ 作为权重，强调边缘区域）、法线一致性损失（从立体深度图和渲染深度图分别计算法线并约束一致）、平滑损失（惩罚法线突变），以及天空分割模块过滤无效深度监督。实验在 KITTI 和 KITTI-360 上进行，PSNR 较 MonoGS 等方法提升超过 6\,dB（KITTI）和 10\,dB（KITTI-360），且在完整长序列上成功运行而 SplaTAM 因内存溢出在 A40 48GB GPU 上失败。ConfidentSplat（2025.09）~\cite{2509_16863} 同样面临单目深度先验的可靠性问题，但与 S3PO-GS 用预训练点图做尺度对齐的路线不同，它直接在深度图层面融合两种互补来源：\textbf{动机。}多视角几何约束产生的稠密深度在纹理丰富区域可靠但在弱纹理区域退化，而单目深度先验虽然缺乏度量精度但全局一致性更好——两种深度源在不同区域具有互补的可靠性，逐像素自适应混合可以取长补短。

ConfidentSplat 的关键创新是为每个像素计算显式的可靠性权重，逐像素自适应混合两个深度源。具体地，多视角置信度 $w_{\mathrm{mv}}(\mathbf{p})$ 来自几何一致性分数 $n_i(\mathbf{p})$——统计相邻 $N_{\mathrm{key}}=30$ 个关键帧中，有多少帧的 3D 重建点与当前帧点的距离小于阈值 $\eta \cdot \mathrm{mean}(\hat{D}_i)$（$\eta=0.01$），归一化为 $w_{\mathrm{mv}} = n_i / N_{\mathrm{key}}$；单目置信度取补 $w_{\mathrm{mono}} = 1 - w_{\mathrm{mv}}$。融合深度 $D_i(\mathbf{p}) = w_{\mathrm{mv}} \hat{D}_i + w_{\mathrm{mono}} D^{\mathrm{mono}}_{\mathrm{scaled},i}$ 作为代理深度监督可变形 3DGS 地图的初始化和优化。跟踪前端基于 DROID-SLAM 式的 RAFT 光流 + DSPO 因子图联合优化位姿、逆深度和尺度/偏移参数；后端通过回环检测和全局 BA 修正轨迹，并将位姿更新传播至高斯参数（主要是均值 $\bm{\mu}$）进行非刚性变形，保持全局一致性。在 Replica 上深度 L1 误差降至 2.37\,cm（MonoGS 为 27.24\,cm），ScanNet 上平均 PSNR 达 27.78\,dB（GlORIE-SLAM 为 22.68\,dB）；在自定义移动端数据上对比 Splat-SLAM 基线，PSNR 从 28.82 提升至 32.74\,dB，深度 L1 降低 9.6\%。

MotionGS（2024.05）~\cite{2405_11129} 同样基于可微渲染位姿优化，但不同于 S3PO-GS 依赖预训练点图建立 2D-3D 对应，MotionGS 从跟踪管线的帧选择机制切入，提出\textbf{双关键帧策略}以减少冗余帧参与优化。

\textbf{动机。}标准 3DGS-SLAM 对所有帧统一处理，但大量帧之间的视角变化极小，参与联合优化时仅带来冗余梯度而无新增信息，不仅浪费计算还可能导致过拟合。双关键帧策略分为两层：运动滤波器提取每帧的密集视觉特征（6 层残差模块 + 3 层下采样，输出 256 维 1/8 分辨率特征图），通过 4 层相关性金字塔计算帧间运动向量，当运动向量范数超过阈值或帧间距超过上限时选为\emph{运动关键帧}；信息滤波器则进一步从运动关键帧中筛选——当前帧相对前一信息关键帧的场景几何相对补集 $\mathrm{RC}(i,j) = (G_i - G_i \cap G_j)/(G_i \cup G_j)$ 超过阈值时，才加入滑动窗口成为\emph{信息关键帧}，窗口满时按重叠系数 $\mathrm{OC}_{i,j}$ 淘汰最冗余帧。这种两级筛选使参与联合优化的帧数大幅减少，同时保持几何覆盖的完整性。

在场景表示方面，MotionGS 借鉴 Compact 3DGS 的 mask 思路，为每个高斯的尺度 $s_i$ 和不透明度 $o_i$ 引入可学习二值掩码 $B_i = \sigma(b_i)$，在优化过程中自动抑制对渲染贡献极小的"微型高斯"，配合 mask 损失 $\mathcal{L}_M = \frac{1}{N}\sum_{i}\sigma(b_i)$ 鼓励稀疏化，优化结束后丢弃掩码。这使得 MotionGS 在 Replica 上地图存储小于 50 MB（对比 SplaTAM 的约 500 MB，减少约 90\%），同时 PSNR 和 LPIPS 均优于同期 3DGS SLAM 方法。

MGS-SLAM（2024.05）~\cite{2405_06241} 同样采用解耦式架构（前端 VO + 后端 3DGS 并行），但与 MGSO 的核心差异在于三点：前端从传统光度法（DSO）替换为学习式稀疏 VO（DPVO），通过 patch 级光流校正进行 BA 获取粗位姿；深度几何不再依赖 VO 直接输出的点云，而是引入预训练 MVS 网络——以当前跟踪帧为参考图像、前 $N$ 个关键帧及位姿为输入，经 FPN 提取三层特征并通过可微分 warp 构建 cost volume，由粗到精解码得到稠密先验深度图；后端不再与前端完全独立，而是以 VO 粗位姿为初始值，通过可微渲染同时优化高斯地图和位姿。

\textbf{深度平滑正则化与 SDAR 尺度校正。}MVS 网络估计的先验深度图存在误差，直接用于几何监督会引入错误引导。MGS-SLAM 对此设计了两个互补机制。深度平滑损失 $\mathcal{L}_{\mathrm{smooth}} = \|d_{i,j-1}-d_{i,j}\|_2 + \|d_{i+1,j}-d_{i,j}\|_2$ 对渲染深度图的相邻像素施加平滑约束，促使高斯在几何上对齐到正确位置——与 NeSLAM 对 RGB-D 传感器噪声深度做去噪不同，这里正则化的对象是从高斯地图渲染出的深度图，本质是引导高斯获得更准确的几何位置。稀疏-稠密调整环（SDAR）分三步闭合前后端尺度：第一步，利用 DPVO 稀疏点云（几何精度更高）对 MVS 先验深度图做线性方差校正 $\bar{D}_d = \frac{\sigma_s}{\hat{\sigma}_d}\hat{D}_d + \mu_d(\frac{\mu_s}{\hat{\mu}_d}-\frac{\sigma_s}{\hat{\sigma}_d})$，对齐均值和方差；第二步，将校正后的深度图反投影为点云，下采样后初始化新高斯加入地图；第三步，利用 3DGS 实时渲染能力，为前端跟踪帧渲染深度图并初始化其点云深度，确保前端跟踪尺度与后端高斯地图一致。此外，关键帧选择基于高斯共视性 IOU（而非简单的帧间距），当两帧可见高斯集合的 IOU 低于阈值时才创建新关键帧，在保证几何覆盖的同时减少冗余。

\textbf{实验表现。}MGS-SLAM 仅使用单目 RGB 输入，在 Replica 上 ATE 达到 0.32\,cm，超越包括 RGB-D 方法在内的所有对比方法（SplaTAM 为 0.36\,cm，DPVO 为 0.54\,cm）；TUM 上 ATE 为 2.93\,cm，比 MonoGS（5.66\,cm）低 48\%，比前端 DPVO（4.49\,cm）低 35\%——后者直接验证了稀疏 VO 与高斯建图联合优化的有效性。Replica 上渲染 PSNR 为 29.78\,dB，优于其他单目方法（Photo-SLAM 为 28.01\,dB）。消融实验表明，移除深度平滑损失导致 PSNR 下降约 1.2\,dB，移除 SDAR 导致 ATE 增加约 70\%，两者缺一不可。

从在线 SLAM 转向离线建图场景，3DGM（2024.05）~\cite{2405_17187} 则从离线建图角度处理环境-物体分离问题。

\textbf{动机。}户外驾驶场景中，移动物体（车辆、行人）在不同遍历中频繁变化，若不加处理会污染静态环境重建，导致"幽灵车辆"等伪影；但人工标注瞬态物体的成本极高，需要一种无监督的自动分割方案。3DGM 将这一先验形式化为鲁棒估计问题——环境像素为内点（inlier），瞬态物体像素为离群点（outlier）——从而无需任何人工标注即可联合完成 2D 瞬态物体分割和 3D 环境建图。

\textbf{三阶段管线：EmerSeg 与 EnvGS。}3DGM 的处理流程分为三步。首先，多条遍历的 RGB 图像联合送入 SfM 获取相机位姿和稀疏点云，初始化 3DGS 场景——多遍历提供比单次遍历更丰富的视角覆盖和更强的三角化约束。第二步 EmerSeg 是分割的关键：将 DINOv2 语义特征蒸馏到每个高斯的 $d$ 维特征向量 $f_i$ 中，联合优化 RGB 渲染损失 $\mathcal{L}_{\mathrm{rgb}}$ 和特征渲染损失 $\mathcal{L}_{\mathrm{feat}}$。由于内点像素（环境）在数量上远超离群像素（瞬态物体），模型梯度主要由一致性内点驱动，使得训练收敛后特征残差图 $|F_t(\xi_t; G) - F_t|$ 的高残差区域恰好对应瞬态物体。3DGM 进一步利用残差图的空间信息：归一化后用轮廓检测提取高残差区域，结合空间先验（过滤过小区域和天空区域）并合并相邻轮廓取凸包，生成瞬态物体掩码 $M_t$，全程无需额外训练。第三步 EnvGS 使用鲁棒优化精化环境高斯：损失函数 $\mathcal{L} = \sum_t \mathcal{L}_{\mathrm{rgb}}(M_t \odot \hat{I}_t(\xi_t; G),\, M_t \odot I_t)$ 通过掩码 $M_t$ 屏蔽离群像素，使瞬态物体对 3D 地图优化的影响最小化。

在 Mapverse 基准（基于 Ithaca365 和 nuPlan 数据集）上，3DGM 的无监督瞬态物体分割 IoU 达 60--70\%，接近有监督方法水平；去除瞬态物体后渲染 PSNR 约 25--28\,dB，相比未去除提升 1--2\,dB。

SplatMAP（2025.01）~\cite{2501_07015} 同样面向单目 RGB 场景，但与 S3PO-GS 用预训练点图模型补充几何先验的路线不同，它直接采用 DROID-SLAM 的稠密光流前端（RAFT 光流 + ConvGRU 迭代精化位姿与深度），以 factor graph 捕获帧间共视关系，无需任何外部预训练模型。

\textbf{动机。}单目 SLAM 在建图初期因观测不足、三角化基线短，生成的点云存在系统性误差，这些误差沿管线传播后在墙壁和家具上产生"重影"伪影；传统 3DGS 的 clone/split 增密策略对稠密 SLAM 点云既低效又无法修正这类早期误差。SplatMAP 为此提出 SLAM 驱动自适应增密（SIAD），核心是用 DROID-SLAM factor graph 优化过程中产生的动态可靠性掩码 $m_i(p)$ 替代传统增密：掩码综合三项判据——深度一致性（深度值 $>0$ 且跨帧几何一致）、帧间重投影距离（$\|T_{ij}\cdot\mathbf{x}_i - \hat{\mathbf{x}}_j\| < \mathrm{thresh}$）、以及 ConvGRU 置信度权重（$\mathrm{mean}(w_{ij}(p)) \geq \epsilon$）。

每当 SLAM 更新位姿增量 $\Delta\xi^{(k)}$ 和深度增量 $\Delta d^{(k)}$ 后，SIAD 对三类点分别处理：持续可靠点（$m^{(k-1)}=1, m^{(k)}=1$）按 $g_p^{(k)} = \exp(\Delta\xi^{(k)})\circ T_{ij}\circ\Pi_c^{-1}(p, d^{(k-1)}+\Delta d^{(k)})$ 更新位置；失效点（$m^{(k-1)}=1, m^{(k)}=0$）直接剪枝；新增可靠点（$m^{(k-1)}=0, m^{(k)}=1$）初始化新高斯。这一过程完全绕过了传统 3DGS 的 clone/split 操作，避免了稠密点云场景下的冗余增密开销。在建图优化方面，SplatMAP 的损失函数 $\mathcal{L} = \mathcal{L}_{\mathrm{rgb}} + \lambda_{\mathrm{geo}}\mathcal{L}_{\mathrm{geo}}$ 同样联合光度与几何约束，但与 S3PO-GS 用渲染点图与预训练点图的 $L_1$ 距离做几何监督不同，SplatMAP 设计了边缘感知法线损失：从渲染深度图的梯度计算法线 $\mathbf{N} = \nabla_x\mathbf{D}\times\nabla_y\mathbf{D}/\|\nabla_x\mathbf{D}\times\nabla_y\mathbf{D}\|$，惩罚法线梯度 $\mathcal{L}_{\mathrm{geo}} = \frac{1}{HW}\sum|\nabla\mathbf{N}|\otimes\omega(|\nabla\mathbf{I}|)$，其中权重函数 $\omega(x) = \exp(-|x-1|^2/\sigma^2)$ 采用高斯调制——在低梯度区域平滑法线，在高梯度边缘保留细节，避免了原始多项式权重 $\omega(x)=(x-1)^q$ 对锐利边缘的过度惩罚。光度项则用多尺度 SSIM（MS-SSIM）替代 3DGS 的单尺度 SSIM，在三个下采样尺度上计算 SSIM 后取平均。

在 Replica 上 SplatMAP 取得 36.864\,dB PSNR、0.985 SSIM 和 0.040 LPIPS，相比此前单目 SOTA（Photo-SLAM 的 33.30\,dB）分别提升 10.7\%、6.4\% 和 49.4\%；在 TUM-RGBD 上 PSNR 为 23.12\,dB，较 GlORIE-SLAM（20.99\,dB）提升 10.2\%，且超越了部分 RGB-D 方法（NICE-SLAM、ESLAM 的 Replica PSNR）。

从追求重建质量转向追求存储效率，RP-SLAM（2024.12）~\cite{2412_09868} 同样采用解耦架构（特征 SLAM 跟踪位姿，3DGS 优化场景），但与上述方法追求重建质量不同，其核心目标是实现\textbf{紧凑且高效的场景表示}。

\textbf{动机。}已有解耦式方法在增量建图时采用逐像素密集采样初始化新高斯，导致大量冗余基元和存储开销；同时使用固定关键帧窗口，连续迭代中旧区域的观测逐渐被排除，引发灾难性遗忘。RP-SLAM 针对这两个瓶颈分别提出改进。

\textbf{四叉树自适应采样与 KNN 过滤。}RP-SLAM 不对关键帧逐像素采样，而是以图像局部梯度方差引导的四叉树递归细分自适应分配采样密度：对每个区格计算梯度方差 $V(C) = \frac{1}{|C|}\sum_{p \in C}(\nabla I(p) - \mu_C)^2$，当 $V(C) > \tau_{\mathrm{th}}$ 且区格尺寸大于自适应最小值 $c_{\mathrm{th}}$ 时继续细分，使采样集中在纹理丰富区域。以 $640 \times 480$ 图像为例，最小区格 $c=8$ 时仅生成 2872 个采样点（$c=4$ 则达 9898 个），显著减少平坦区域的冗余基元。采样像素经深度反投影生成候选高斯后，进一步通过 KNN 过滤去除空间冗余：对每个候选高斯 $G_{\mathrm{new}}$ 搜索 3 个最近邻 $\mathcal{G}=\{G_1, G_2, G_3\}$，若其到所有邻居的距离均满足 $\mathrm{dis}(G_{\mathrm{new}}, G_i) < \lambda r_i$（$r_i$ 为邻居的最小缩放分量），则判定为冗余并丢弃。两阶段筛选确保地图始终保持紧凑表示。

\textbf{动态关键帧窗口。}固定窗口优化只关注与当前关键帧共视的局部区域，旧区域在迭代中逐渐"遗忘"。RP-SLAM 将关键帧按共视性分为共视集 $\mathcal{K}_1$ 和非共视集 $\mathcal{K}_2$，每次迭代从两组各随机采样 $k_1$, $k_2$ 帧与当前关键帧组成动态窗口 $\mathcal{W} = \{K_{\mathrm{new}}\} \cup \mathcal{S}_1 \cup \mathcal{S}_2$。共视帧保证局部一致性，非共视帧以概率形式被回顾，从而在保持局部精度的同时缓解遗忘。此外，针对单目模式缺乏深度信息的问题，RP-SLAM 直接利用特征 SLAM 提供的稀疏点云为新高斯提供初始 3D 位置，而非 MonoGS 的随机深度初始化——消融实验显示，该策略使单目场景的初始深度渲染即能描述场景的大致几何结构，为后续优化提供有效起点。在 TUM、Replica 和 ScanNet++ 数据集上，RP-SLAM 在保持 SOTA 渲染精度的同时，高斯基元数量仅为对比方法的 30\%--50\%，存储开销显著降低。

从"各自提出方法"转向"系统性基准评估"，Schmidt et al.（2024.12）~\cite{2412_03263} 提供了首个面向户外自然场景的系统性对比基准。

\textbf{动机。}此前 3DGS-SLAM 的评测主要在室内合成数据集（Replica）或有限真实场景（TUM-RGBD）上进行，缺乏在季节变化和光照变化两个关键维度上的系统性评估，无法回答"3DGS-SLAM 在户外真实环境中究竟有多鲁棒"这一关键问题。——一个植被茂密、表面弱纹理的户外环境——在季节（夏/秋/冬/春）和光照（白天/黄昏/夜晚/外灯补光）两个维度上评估了 11 种 SLAM 方法，涵盖传统方法（ORB-SLAM3）、深度学习方法（DROID-SLAM, DPV-SLAM）、NeRF-based 方法（GO-SLAM, GlORIE-SLAM, Co-SLAM, ORBEEZ-SLAM）和 3DGS-based 方法（MonoGS, Photo-SLAM, Splat-SLAM, GS-ICP-SLAM）。

核心发现揭示了三类方法在户外场景中的互补性：季节变化对跟踪精度影响较小（NeRF 和 3DGS 方法跨季节 ATE 波动不大），但光照条件构成决定性挑战——夜晚场景下多数单目方法完全失效。在单目模式中，3DGS 方法表现两极分化：Splat-SLAM 跨季节平均 ATE 2.12\,m，与 DROID-SLAM（2.08\,m）竞争力相当，而 MonoGS 在四个季节中两次失败、Photo-SLAM 则在夜晚完全崩溃，说明 Frame-to-Model 式位姿优化（MonoGS）在弱纹理户外场景中远不如 Frame-to-Frame 式（Splat-SLAM 借助 DROID-SLAM 前端）鲁棒。RGB-D 模式下这一差距显著缩小：MonoGS 平均 ATE 降至 0.90\,m，且在夜晚场景中仅 1.26\,m——远优于传统方法 ORB-SLAM3 的 4.99\,m——表明深度传感器补充的几何信息对 NeRF/3DGS 方法在挑战条件下的鲁棒性至关重要。

传统方法跨季节一致性最好（ORB-SLAM3 RGB-D ATE 标准差仅 0.06\,m），但对光照变化高度敏感（黄昏时 ATE 从 0.59\,m 飙升至 5.32\,m）。NeRF 方法 GO-SLAM 在光照维度上最鲁棒（单目 ATE 3.32\,m，所有方法中唯一在四种光照下均有结果的单目方法），但 GPU 占用达 12--14\,GB，运行时间约为实时的 4 倍。总体而言，该基准为户外 3DGS SLAM 的方法选择提供了实用参考：当深度信息可用且低光照是主要挑战时，3DGS 方法（MonoGS RGB-D）优于传统方法；当计算资源受限时，Photo-SLAM 在中等光照条件下提供了精度-效率的最佳平衡。

\begin{takeaway}
单目与 RGB SLAM 的核心挑战是缺少深度信息带来的初始化和尺度歧义：解耦式跟踪-建图架构、四叉树自适应采样和动态关键帧窗口有效缓解了冗余基元和灾难性遗忘问题，而户外基准测试揭示了光照变化是决定性挑战。
\end{takeaway}

% ======================================================================

\subsection{近期方法脉络：在线单目3DGS、结构增强3DGS-SLAM、显式几何SLAM等}

这一组近期工作延续 online 3DGS、3DGS-SLAM、3D Gaussian Splatting，主要把改动集中在扩展、增量改进、规模扩展等环节。下面只保留相对既有路线的变化，便于和前文方法直接对照。

\paragraph{MoonSplat}
MoonSplat（2026.06）~\cite{2606_17935} 延续 online 3DGS 在“在线单目3DGS”上的基本设定，主要改变是：MoonSplat：Monocular Online Gaussian Splatting，加入Sim(3)全局优化。3DGS 的显式高斯表示和可微渲染能力保持为基础。
其设计动机是：针对在线3DGS缺少全局优化导致的相机位姿脆弱和长序列效率低问题，提出鲁棒在线重建框架。
具体而言，方法在单目序列中增量维护高斯场，并用Sim(3)全局优化修正尺度、位姿和地图一致性。
适用边界是：单目尺度和回环仍具挑战；动态场景可能破坏全局优化假设。

\paragraph{MMD-SLAM}
MMD-SLAM（2026.06）~\cite{2606_19874} 延续 3DGS-SLAM 在“结构增强3DGS-SLAM”上的基本设定，主要改变是：MMD-SLAM：Atlanta World结构先验引导Multi-Meta Gaussian表示，并结合点线融合位姿优化和Gaussian evolution。3DGS 的显式高斯表示和可微渲染能力保持为基础。
其设计动机是：把Atlanta World主方向结构编码进高斯地图，并用点线融合约束改善跟踪与建图一致性。
具体而言，方法引入点线融合位姿优化以增强纹理弱区域的跟踪，再用带主方向的Multi-Meta Gaussian表示显式利用结构先验，最后通过Gaussian evolution自适应更新地图。
适用边界是：依赖结构化环境假设；非曼哈顿/非Atlanta场景中结构先验可能失效。

\paragraph{Triangle Splatting SLAM}
Triangle Splatting SLAM（2026.05）~\cite{2605_31419} 延续 3D Gaussian Splatting 在“显式几何SLAM”上的基本设定，主要改变是：Triangle Splatting SLAM：用triangle soup进行跟踪与建图，并在线通过受限Delaunay三角剖分转为连通mesh。3DGS 的显式高斯表示和可微渲染能力保持为基础。
其设计动机是：提出首个基于Triangle Splatting的稠密RGB-D SLAM系统，使地图可在线转为连通mesh，服务仿真、碰撞和编辑等下游任务。
具体而言，系统通过可微渲染优化未结构化三角形集合完成跟踪和建图，并借助Delaunay三角剖分在线生成可连接网格。
适用边界是：不是3DGS本体方法；三角形表示与高斯表示的效率/鲁棒性差异仍需更多比较。

\paragraph{MyGO-Splat}
MyGO-Splat（2026.06）~\cite{2606_29738} 可作为“Closed-loop monocular Gaussian SLAM”中的新锚点候选。相较于 3DGS-SLAM，主要改变是：用 Gaussian primitives 解析栅格化深度和法线反向监督位姿优化。3DGS 的显式高斯表示和可微渲染能力保持为基础。
其设计动机是：提出多目标闭环几何反馈，将 refined Gaussian map 的深度和法线用于 tracking 自校正。
具体而言，从 Gaussian primitives 解析渲染 pixel-wise depth 和 surface normals，并引入 scale-aware adaptive alignment 连接单目先验与尺度一致性。
这一锚点判断来自当前分类中的机制关系：候选锚点：扩展，且相关性 0.94；后续方法可围绕 加入闭环几何反馈以调节 RGB-only tracking drift。 继续比较和扩展。
适用边界是：闭环反馈依赖当前地图几何质量，早期错误或低纹理场景可能导致反馈偏差。

\paragraph{MoPe}
MoPe（2026.06）~\cite{2606_29237} 延续 3DGS-SLAM 在“Dynamic-scene robustness in Gaussian SLAM”上的基本设定，主要改变是：Motion Permanence：将动态性建模为跨时间的持久属性而非单帧观测。3DGS 的显式高斯表示和可微渲染能力保持为基础。
其设计动机是：指出动态性不是瞬时视觉属性，而应作为地图级记忆变量维护，从表示层避免动态内容被吸收到静态地图。
具体而言，在 monocular Gaussian Splatting SLAM 中记录和传播区域或高斯的 motion permanence，当行人暂停、遮挡后重现等情况发生时仍保持动态判断，减少错误建图。
适用边界是：依赖动态性估计和时间关联质量，对长期遮挡、缓慢移动或语义不确定对象仍可能误判。

\paragraph{Robust and Efficient Monocular 3D Gaussian SLAM for Kilometer-Scale Outdoor Scenes}
Robust and Efficient Monocular 3D Gaussian SLAM for Kilometer-Scale Outdoor Scenes（2026.06）~\cite{2606_30436} 延续 3DGS-SLAM 在“Kilometer-scale outdoor Gaussian SLAM”上的基本设定，主要改变是：运动自适应混合 tracking、按需地图管理和大尺度高效 Gaussian mapping。3DGS 的显式高斯表示和可微渲染能力保持为基础。
其设计动机是：提出 condition-triggered 三层求解 tracking pipeline，在 Essential matrix 和 PnP 等模型间动态切换以应对几何退化。
具体而言，在保证 drift-free camera poses 的基础上进行高保真 Gaussian mapping，并结合按需机制控制大尺度地图内存。
适用边界是：单目户外场景仍面临尺度漂移、动态交通和光照变化；系统复杂度较高。
\section{多模态SLAM}
\label{sec:15_3}

纯视觉 SLAM 在纹理缺失、剧烈运动和光照变化等条件下存在固有局限。多模态 SLAM 通过融合视觉与 IMU（惯性测量单元，Inertial Measurement Unit）、LiDAR、雷达、事件相机等互补传感器，显著提升鲁棒性和精度。3DGS 作为可微的场景表示，天然适合多模态融合——不同传感器的约束可以统一到光度、几何和运动损失框架中联合优化。本节梳理 25 篇论文，按传感器组合分类。

\subsection{惯性与LiDAR融合}
\label{ss:15_4}

本节介绍两条技术路线：（1）视觉-惯性融合；（2）LiDAR-视觉融合。

\subsubsection{视觉-惯性融合}

本路线介绍视觉与 IMU 紧耦合融合的 3DGS SLAM 方法，利用惯性测量提供运动先导和重力约束以提升跟踪鲁棒性。

IMU 提供高频率运动测量，为视觉 SLAM 提供关键的运动先导和重力约束。本节首先以 VIGS-SLAM 为锚点方法，深入解析视觉-惯性紧耦合在 3DGS SLAM 中的完整设计范式，再逐步展开其他方法的改进脉络。

\textbf{问题背景与设计动机。}尽管近年来基于 3DGS 的 SLAM 方法已实现稠密、照片级的地图构建，但纯视觉设计在运动模糊、弱纹理和曝光变化等挑战条件下显著退化。VIGS-SLAM（2025.12）~\cite{2512_02293} 的核心思路是将 IMU 的惯性约束紧密融入视觉跟踪和高斯建图的统一优化框架，使系统在上述挑战场景中保持鲁棒。具体而言，系统以 RGB 帧序列 $\{I_i\}_{i=1}^N$ 和原始 IMU 测量 $\{\mathbf{a}_k, \boldsymbol{\omega}_k\}_{k=1}^M$（加速度和角速度）为输入，通过两个并行模块协同工作：前端跟踪模块负责相机位姿估计与关键帧管理，高斯泼溅建图模块负责增量构建照片级场景地图。回环闭合检测在后台并行执行，以保证全局一致性。

\textbf{前端跟踪：视觉-惯性联合优化。}跟踪模块在线处理图像流，当光流幅度超过阈值时创建新关键帧；为限制 IMU 预积分在长时间间隔上的漂移，系统同时强制每 $t_{kf}$ 秒至少创建一个关键帧。对每个关键帧 $i$，系统联合优化相机位姿 $\mathbf{T}_i = (\mathbf{R}_i, \mathbf{t}_i)$、视差 $d_i$、速度 $\mathbf{v}_i$ 和 IMU 偏差 $\mathbf{b}_i = [\mathbf{b}_g^T, \mathbf{b}_a^T]^T$（分别为陀螺仪和加速度计偏差）——相比纯视觉方法仅优化位姿和深度，这里增加了对运动状态和传感器偏差的显式建模。

视觉残差项沿用 DROID-SLAM 的学习型 GRU 更新算子，迭代预测关键帧对 $(i, j)$ 之间的稠密对应 $\mathbf{p}^{*}_{ij}$ 及其置信度 $\mathbf{w}_{ij}$，通过最小化重投影误差联合优化位姿和深度：

\begin{equation}
    E_{\mathrm{vis}}(\mathbf{T}, d) = \sum_{(i,j) \in E} \sum \| \mathbf{p}^{*}_{ij} - \pi(\mathbf{T}_{ij}\,\pi^{-1}(\mathbf{u}_i, d_i)) \|_{\mathbf{W}_{ij}}^2,
\end{equation}

其中 $E$ 为局部帧图的共视边集合，$\pi$ 和 $\pi^{-1}$ 分别为投影和反投影函数，$\mathbf{u}_i$ 为帧 $i$ 中的 2D 像素坐标，$\| \cdot \|_{\mathbf{W}}$ 为按置信度加权的 Mahalanobis 距离——置信度越高（即对应越可靠）的像素对优化贡献越大。

惯性残差项则基于 IMU 预积分（2017）~\cite{forster2015imu}。直觉上，IMU 以数百赫兹的频率输出角速度和加速度，直接逐个采样点优化代价过高；预积分将连续关键帧时间戳 $t_i$ 到 $t_j$ 之间的所有 IMU 数据高效聚合为三项相对运动量——相对旋转 $\Delta \mathbf{R}_{ij}$、相对位置 $\Delta \mathbf{p}_{ij}$ 和相对速度 $\Delta \mathbf{v}_{ij}$——以及对应的 Jacobian 和协方差。惯性残差通过最小化当前相机相对运动与预积分测量之间的差异来约束位姿和运动状态：

\begin{equation}
    E_{\mathrm{iner}}(\mathbf{T}, \mathbf{v}, \mathbf{b}) = \sum_{(i, j=i+1)} \| [\mathbf{r}^{\mathrm{rot}}_{ij},\, \mathbf{r}^{\mathrm{pos}}_{ij},\, \mathbf{r}^{\mathrm{vel}}_{ij},\, \mathbf{r}^{\mathrm{bias}}_{ij}]^T \|_{\boldsymbol{\Sigma}_{ij}^{\mathrm{iner}}}^2,
\end{equation}

其中旋转残差 $\mathbf{r}^{\mathrm{rot}}_{ij}$ 衡量视觉估计的相对旋转与 IMU 预积分旋转（校正陀螺仪偏差后）的一致性；位置和速度残差 $\mathbf{r}^{\mathrm{pos}}_{ij}$、$\mathbf{r}^{\mathrm{vel}}_{ij}$ 检验从速度和重力推算的相对运动是否与预积分测量吻合；偏差平滑项 $\mathbf{r}^{\mathrm{bias}}_{ij} = \mathbf{b}_j - \mathbf{b}_i$ 正则化偏差的时域变化，防止过拟合。前端跟踪在局部帧图 $E$ 上联合最小化 $E_{\mathrm{vis}} + E_{\mathrm{iner}}$，使用 Levenberg-Marquardt 算法通过自定义 CUDA 核高效求解。

\textbf{IMU 初始化：三阶段渐进策略。}视觉-惯性紧耦合的一个关键难点在于系统启动时如何可靠地估计 IMU 参数（重力方向、度量尺度、初始偏差）。VIGS-SLAM 采用精心设计的三阶段渐进初始化：\emph{阶段一}（纯视觉初始化），利用前 $N^{\mathrm{vis}}_{\mathrm{init}}$ 个关键帧仅最小化视觉残差 $E_{\mathrm{vis}}$，估计位姿和视差（存在一个全局尺度歧义）；\emph{阶段二}（纯惯性优化），继续视觉跟踪至 $N^{\mathrm{iner}}_{\mathrm{init}}$ 个关键帧后，固定位姿不变，引入重力方向 $\mathbf{R}_{wg}$、逐关键帧速度 $\mathbf{v}_i$、IMU 偏差 $\mathbf{b}_i$ 和全局对数尺度参数，仅最小化惯性目标函数以恢复度量尺度和对齐重力方向；\emph{阶段三}（视觉-惯性联合优化），将前两阶段的结果作为初值，联合最小化视觉和惯性目标函数进行最终精调。这种渐进策略的核心动机是延迟视觉-惯性耦合——在 IMU 参数尚不可靠时过早引入惯性约束反而会引入噪声，渐进策略通过逐步增加约束自由度来确保鲁棒收敛。

\textbf{回环闭合：Pose Graph Bundle Adjustment（PGBA）。}为高效处理回环闭合，VIGS-SLAM 采用 PGBA 而非完整全局 BA，以少量精度换取显著的速度提升。一个并行的回环检测器基于光流差异构建回环边 $E^*$；同时，系统动态扩展一个相对位姿图 $E^+$，将局部帧图的相对位姿信息不断追加进去。PGBA 仅对回环边 $E^*$ 执行稠密视觉更新，对其余位姿图边 $E^+$ 仅施加轻量级的相对位姿约束，在 $\mathrm{Sim}(3)$ 空间优化以纠正长期尺度漂移。

\textbf{回环后的高斯一致性更新。}回环闭合后所有关键帧的位姿和尺度都会更新，如何在不重新初始化和重新优化所有高斯的前提下保持地图一致性？VIGS-SLAM 的策略是将位姿更新传播到每个高斯所锚定的关键帧上。设关键帧 $k$ 的 PGBA 前/后位姿分别为 $(\mathbf{R}_k^-, \mathbf{p}_k^-)$ 和 $(\mathbf{R}_k^+, \mathbf{p}_k^+)$，尺度变化为 $s_k$。对于从关键帧 $k$ 初始化的任何高斯 $g_i$（均值 $\boldsymbol{\mu}_i^-$，协方差 $\boldsymbol{\Sigma}_i^-$），更新分三步：先将高斯从旧世界坐标变换到旧相机局部坐标 $\mathbf{x}_{\mathrm{loc}}^- = (\mathbf{R}_k^-)^T(\boldsymbol{\mu}_i^- - \mathbf{p}_k^-)$，再施加尺度缩放 $\mathbf{x}_{\mathrm{loc}}^+ = s_k \mathbf{x}_{\mathrm{loc}}^-$，最后映射到新世界坐标：

\begin{equation}
    \boldsymbol{\mu}_i^+ = \mathbf{R}_k^+(s_k\, \mathbf{x}_{\mathrm{loc}}^+) + \mathbf{p}_k^+.
\end{equation}

协方差做类似更新 $\boldsymbol{\Sigma}_i^+ = \mathbf{R}_k^+(s_k^2 (\mathbf{R}_k^-)^T \boldsymbol{\Sigma}_i^- \mathbf{R}_k^-)(\mathbf{R}_k^+)^T$，而不透明度和颜色保持不变。整个过程以批量操作高效执行。

\textbf{高斯泼溅建图。}建图模块与跟踪并行运行。新关键帧到来时，系统将视差 $d_i$ 转换为深度后反投影到 3D 空间初始化高斯，颜色取对应像素值，不透明度初始化为 $0.5$。每帧执行 10 次建图迭代：从跟踪帧图 $E$ 中随机采样关键帧（外加两帧全局关键帧），渲染颜色和深度图，联合最小化颜色损失 $\mathcal{L}_c$、深度损失 $\mathcal{L}_d$ 和各向同性正则化 $\mathcal{L}_{\mathrm{iso}}$（防止稀疏观测区域高斯过度拉伸），总损失为 $\mathcal{L} = \lambda_d \mathcal{L}_d + \lambda_c \mathcal{L}_c + \lambda_{\mathrm{iso}} \mathcal{L}_{\mathrm{iso}}$。

\textbf{实验表现。}VIGS-SLAM 在五个挑战性数据集（EuRoC、RPNG AR Table、UTMM、FAST-LIVO2 和自采集数据集）上与 15 种方法进行了对比，涵盖经典视觉 SLAM（DSO、ORB-SLAM3）、学习型视觉 SLAM（DROID-SLAM）、经典视觉-惯性 SLAM（VINS-Mono、OKVIS）、学习型视觉-惯性 SLAM（DBA-Fusion）以及 3DGS SLAM（Splat-SLAM、HI-SLAM2）等。在 EuRoC 数据集上取得最佳总体 ATE，大多数序列排名第一或第二；在 RPNG AR Table 上跟踪误差减半于第二名 VINS-Mono，同时取得最高渲染质量；在户外低帧率、运动抖动和反射面的 FAST-LIVO2 数据集上，经典 VIO 系统频繁初始化失败，而 VIGS-SLAM 从第一帧鲁棒初始化并在所有序列上保持稳定跟踪。消融实验验证了六个设计选择（IMU 偏差估计、IMU 融合、惯性初始化阶段、IMU 位姿初始化、回环闭合、逐关键帧偏差建模）各自的必要性——移除任一组件均导致跟踪精度和鲁棒性下降。

GI-SLAM（2025.03）~\cite{2503_18275} 采取了与 VIGS-SLAM 不同的技术路线来实现 IMU 集成：VIGS-SLAM 采用紧耦合的 IMU 预积分框架，将惯性测量聚合为相对运动量并参与完整的 BA 优化；GI-SLAM 则将 IMU 约束以损失函数的形式直接嵌入梯度下降优化流程，实现了更轻量的视觉-惯性融合，同时支持单目、双目和 RGB-D 多种传感器配置（有/无 IMU 均可运行）。

\textbf{与 VIGS-SLAM 的核心差异：轻量 IMU 损失。}VIGS-SLAM 的惯性残差基于成熟的 IMU 预积分理论（2017）~\cite{forster2015imu}，将连续关键帧间的数百个 IMU 采样聚合为相对旋转、位置、速度三项预积分量及其协方差，需要显式估计速度、偏差等运动状态变量。GI-SLAM 的设计动机是降低惯性融合的工程复杂度：它不使用预积分，而是直接利用帧间的 6-DOF IMU 测量构建平移和旋转两项约束。平移约束从上一帧线速度 $\mathbf{v}_{t-1}$ 和当前帧加速度 $\mathbf{a}_t$ 出发，积分得到 IMU 预测位移 $\Delta\mathbf{t}_{\mathrm{imu}} = \mathbf{v}_{t-1}\Delta t + \frac{1}{2}\mathbf{a}_t\Delta t^2$，与优化变量 $\Delta\mathbf{t}_{\mathrm{opt}}$ 之间计算 $\ell_2$ 损失；旋转约束类似地从角速度 $\boldsymbol{\omega}_t$ 积分得到 IMU 预测相对旋转 $\Delta\mathbf{r}_{\mathrm{imu}} = \boldsymbol{\omega}_t\Delta t$，与优化变量 $\Delta\mathbf{r}_{\mathrm{opt}}$ 之间计算 $\ell_2$ 损失。最终 IMU 损失 $\mathcal{L}_{\mathrm{imu}} = w_t\mathcal{L}_{\mathrm{trans}} + w_r\mathcal{L}_{\mathrm{rot}}$ 通过超参数 $w_t, w_r$ 平衡两项约束。这种设计省去了预积分的 Jacobian 计算和协方差传播，代价是在长时间间隔上不如预积分精确。

\textbf{定位框架：SE(3) 参数化的梯度下降。}与 VIGS-SLAM 采用学习型 GRU 迭代更新稠密光流对应不同，GI-SLAM 将相机位姿参数化为 $SE(3)$ 齐次变换矩阵 $\mathbf{P}_t$，以增量形式 $\mathbf{P}_t = \mathbf{P}_{t-1}\Delta\mathbf{P}$ 从上一帧递推当前位姿。优化目标是最小化传感器数据（RGB、深度、IMU）与从当前 3DGS 地图渲染图像之间的对齐误差，即加权组合 RGB 损失 $\mathcal{L}_{\mathrm{rgb}}$（光度一致性 $\ell_1$ 损失）、深度损失 $\mathcal{L}_{\mathrm{depth}}$（几何一致性 $\ell_1$ 损失）和上述 IMU 损失 $\mathcal{L}_{\mathrm{imu}}$。由于 3DGS 的泼溅渲染过程天然可微，上述三部分损失均可通过反向传播直接计算梯度并更新 $\Delta\mathbf{P}$。

\textbf{关键帧管理：共视性评分机制。}GI-SLAM 的关键帧选择也不同于 VIGS-SLAM 的光流阈值策略。系统维护一个动态关键帧窗口 $W_k$，通过统一评分函数 $s_i$ 综合三个标准来决定是否将帧 $i$ 提升为关键帧：（1）共视性——基于 3DGS 渲染过程中的可见高斯重叠度 $\sigma_G = |G_i \cap G_j| / |G_i \cup G_j|$，自然适应场景复杂度而无需显式几何重叠检查；（2）基线跨度——归一化平移距离 $\|\mathbf{t}_{ij}\|/d_{\mathrm{med}}$ 衡量视角多样性，最大化多视约束强度；（3）运动稳定性——通过硬阈值排除线速度或角速度过大的帧（运动模糊风险）。评分超过阈值 $\sigma_{\mathrm{kf}}$ 且窗口有容量时选为关键帧，与最新关键帧重叠超过 $\sigma_{\mathrm{overlap}}$ 时剔除冗余帧。这种设计在保持多视约束的同时最小化计算开销。

\textbf{建图与传感器灵活性。}建图模块采用与 MonoGS 类似的 3DGS 场景表示，通过可微泼溅渲染实现梯度驱动的地图更新——包括高斯位置、协方差、颜色和不透明度的联合优化，以及自适应密度控制（克隆欠重建区域的高斯、剪枝不透明度低于阈值或空间范围过小的高斯）。与 VIGS-SLAM 相比，GI-SLAM 不包含回环闭合模块，因此在长序列上的全局一致性依赖前端的增量跟踪精度。但系统的传感器配置灵活性显著更高：IMU 损失是可选的附加项，移除后系统退化为纯视觉 SLAM，仍可正常工作。

\textbf{实验表现。}GI-SLAM 在 EuRoC 和 TUM-RGBD 数据集上进行了评估。IMU 增强模块有效提升了相机跟踪的准确性、鲁棒性和效率——消融实验表明移除 IMU 损失后跟踪精度显著下降，尤其在运动模糊和弱纹理序列上退化明显。系统在多个配置（单目、双目、RGB-D，有/无 IMU）下均达到与最佳实时方法相当的性能，验证了轻量 IMU 损失在 3DGS SLAM 框架中的有效性。

\textbf{VIGS SLAM（2025.01）~\cite{2501_13402}：松耦合 IMU 预积分驱动的 GICP 跟踪。}如果说 VIGS-SLAM 代表了视觉-惯性紧耦合的成熟范式，VIGS SLAM 则探索了一条截然不同的技术路线——将 IMU 以松耦合方式注入点云匹配跟踪，以实现大规模场景的 SLAM。两者的核心差异在于 IMU 的角色：VIGS-SLAM 中 IMU 预积分残差直接参与 BA 联合优化（紧耦合），而 VIGS SLAM 中 IMU 预积分仅作为 GICP 点云配准的初始位姿估计，不进入优化目标函数。

\textbf{GICP 跟踪与 IMU 初始估计。}VIGS SLAM 以 GS-ICP SLAM 为基线，采用广义 ICP（GICP）进行前端跟踪。具体而言，对每帧 RGB-D 图像，从深度图生成点云并为每个点计算 $k$-近邻协方差矩阵，从而将点云建模为一组高斯分布 $\mathcal{G}_k = \{(\mathbf{x}_m, \boldsymbol{\Sigma}_m)\}$。跟踪目标是在源点云（当前帧）与目标点云（已建地图）之间估计最优刚体变换 $T_k^*$，通过最大化对应点对高斯分布的似然函数实现——这与 VIGS-SLAM 基于稠密光流的 GRU 迭代更新策略完全不同，几何意义更直观，且天然适配深度传感器。关键创新在于 IMU 预积分如何发挥作用：VIGS-SLAM 将预积分量作为残差参与联合优化，而 VIGS SLAM 仅利用帧间 IMU 预积分值（相对旋转 $\Delta\mathbf{R}^k_{k+1}$、相对位移 $\Delta\mathbf{p}^k_{k+1}$）为 GICP 提供初始对应关系——准确的初始猜测使 GICP 即使在快速运动下也能快速收敛，大幅降低纯视觉跟踪的计算负载。

\textbf{预积分参数的在线校正。}由于 IMU 持续积分加速度和角速度，漂移在低速运动时逐渐累积。为防止积分误差恶化跟踪性能，VIGS SLAM 在每帧 GICP 跟踪完成后，用优化后的位姿 $T_k^* = [\mathbf{R}_k^*|\mathbf{p}_k^*]$ 反向更新 IMU 预积分的初始参数：位置 $\mathbf{p}_t^k$ 和旋转 $\mathbf{R}_t^k$ 直接用 GICP 结果覆盖，速度 $\mathbf{v}_t^k$ 则通过有限差分 $\mathbf{v}_t^k = (\mathbf{p}_k^* - \mathbf{p}_{k-1}^*)/\Delta t$ 估计。这种"跟踪-校正"循环确保 IMU 预积分不致长期漂移，但代价是松耦合架构无法像 VIGS-SLAM 的紧耦合 BA 那样显式建模速度和偏差的不确定性。

\textbf{建图：复用跟踪信息与尺度归一化。}建图模块采用标准 3DGS 框架，联合优化光度损失 $\mathcal{L}_{\mathrm{photo}}$（$\ell_1$）、SSIM 损失 $\mathcal{L}_{\mathrm{SSIM}}$ 和深度损失 $\mathcal{L}_{\mathrm{depth}}$。一个工程亮点是 GICP 跟踪阶段已为每个点计算了协方差矩阵，这些值在建图阶段直接复用为高斯的初始协方差，避免了冗余计算。此外，SLAM 场景中高斯是增量引入的，早期和后期添加的高斯可能因场景覆盖度不同而产生尺度失衡，因此系统对新高斯施加尺度归一化以维持一致性。

\textbf{面向大规模场景的设计。}VIGS SLAM 的设计目标是突破现有 3DGS SLAM 的房间尺度限制。核心论点是：直接法（稠密光度损失跟踪）要求关键帧间距很小以保持光度一致性，导致大量关键帧和高内存消耗；而 GICP + IMU 初始估计的组合使关键帧间距可以大幅增大，因为点云匹配对视角变化的容忍度远高于稠密光度损失。系统在大规模视觉-惯性数据集 uHumansV1（$65\text{m} \times 65\text{m}$ 办公空间，含动态人员）和 uHumansV2（公寓至大型办公场景）上进行评估。在 uHumansV1 上，VIGS SLAM 的跟踪精度与经典 VIO 方法（ORB-SLAM3、VINS-Mono）相当，远超所有 3DGS SLAM 对比方法（MonoGS、PhotoSLAM）；相比基线 GS-ICP SLAM，跟踪误差从数米降至数十厘米级别。在 uHumansV2 上，其他 3DGS SLAM 方法在从公寓扩展到大型办公场景时性能显著退化，而 VIGS SLAM 在两种尺度上均保持一致的跟踪和渲染质量。

VBGS-SLAM（2026.04）~\cite{2604_02696} 则从根本不同的范式出发：VIGS-SLAM 和 GI-SLAM 都沿袭确定性优化思路——前者通过 LM 求解器联合最小化视觉-惯性残差，后者通过梯度下降最小化光度-深度-IMU 损失；VBGS-SLAM 将整个 SLAM 问题重新建模为概率生成模型的变分推理。

\textbf{与 VIGS-SLAM 的核心差异：从确定性优化到变分推理。}VIGS-SLAM 将高斯泼溅建图和位姿跟踪视为两个独立模块，分别通过可微渲染反向传播和 BA 优化求解。VBGS-SLAM 的关键创新是将两者耦合在统一的生成概率框架中：场景被建模为 $K$ 个高斯分量组成的高斯混合模型，每个分量的空间位置和颜色各自服从多元正态分布，参数以 Normal-Inverse-Wishart（NIW）共轭先验描述；相机位姿 $\mathbf{T}_t$ 被提升为 $\mathfrak{se}(3)$ 切空间上的随机变量 $\boldsymbol{\xi}_t \sim \mathcal{N}(\boldsymbol{\mu}_{\xi,t}, \boldsymbol{\Sigma}_{\xi,t})$，通过指数映射映射到 $SE(3)$。这意味着位姿不再是确定性优化变量，而是带有完整后验不确定性（协方差矩阵 $\boldsymbol{\Sigma}_{\xi,t}$）的分布。

\textbf{闭式变分更新：免反向传播的位姿-地图联合推断。}在均值场近似下，变分后验被因式分解为点分配 $q(\mathbf{z})$、空间参数 $q(\boldsymbol{\mu}_{k,s}, \boldsymbol{\Sigma}_{k,s})$、颜色参数 $q(\boldsymbol{\mu}_{k,c}, \boldsymbol{\Sigma}_{k,c})$、混合权重 $q(\boldsymbol{\pi})$ 和位姿 $q(\mathbf{T}_t)$ 五组独立因子。通过坐标上升变分推理（CAVI）和指数分布族的共轭性质，每组因子均可推导出闭式更新规则。位姿更新是核心创新点——为处理 $SE(3)$ 流形上的非线性，系统在当前位姿估计的切空间上做一阶 Taylor 展开，得到 Jacobian $\mathbf{G}_{k,s} = [\bar{\mathbf{R}}_t, -\bar{\mathbf{R}}_t[\boldsymbol{\mu}_{k,s}]_\times]$，进而推导出位姿后验的闭式更新：精度矩阵 $\boldsymbol{\Sigma}_{\xi,t+1}^{-1}$ 由所有观测点对 Jacobian 的加权累加加上先验精度构成，均值 $\boldsymbol{\mu}_{\xi,t+1}$ 则综合残差项和先验均值加权求和。这一过程完全无需反向传播通过可微渲染器——位姿和地图参数在计算充分统计量的单次遍历中同步更新，显著降低了计算开销。

\textbf{位姿不确定性向地图的传播。}VIGS-SLAM 的建图模块假设跟踪输出的位姿是确定值，不确定性仅通过 IMU 协方差在 BA 中传播。VBGS-SLAM 则在变分框架中实现了更紧密的耦合：点分配 $\gamma_{nk}$（第 $n$ 个点属于第 $k$ 个高斯分量的概率）的计算包含对位姿分布 $q(\mathbf{T}_t)$ 的期望，使分配结果自然吸收位姿不确定性；空间参数的充分统计量 $\mathcal{T}(\mathbf{s}_n)$ 同样依赖位姿后验。这种设计使得位姿估计的不确定性能够传播进地图参数的更新中，而非事后校正。

\textbf{关键帧管理与自适应高斯控制。}关键帧选择沿用 MonoGS 的双标准策略——共视性触发（观察足够多未见高斯时插入新关键帧）和时间约束（防止慢速运动中过久未插入关键帧导致漂移）。选中的关键帧存入固定大小缓冲区，以滑动窗口变分推理联合优化窗口内所有关键帧位姿和相关高斯参数，旧关键帧被边缘化。自适应高斯管理针对两个方向：对欠覆盖区域，基于软曼哈顿距离分析检测未重建区域并从当前 RGB-D 帧反投影插入新高斯；对冗余高斯，采用基于混合权重 $\pi$ 不变性的剪枝策略——若某分量经过多轮更新后先验权重 $\pi$ 未发生变化，说明该分量未获得有效观测支持，将其重新初始化而非简单删除。

\textbf{实验表现。}VBGS-SLAM 在 Replica、TUM-RGBD 和 AR-TABLE 三个数据集上与 SplaTAM、MonoGS 等方法进行了对比。在 Replica 上取得最佳平均 ATE $0.33\,\mathrm{cm}$，较 SplaTAM 改善 $8.3\%$；在长序列 AR-TABLE 上平均 ATE $5.94\,\mathrm{cm}$，优于 SplaTAM 的 $7.75\,\mathrm{cm}$——后者因突发跟踪跳变触发不受控的高斯插入而在多个序列上失败，而 VBGS-SLAM 在所有序列上保持稳定。渲染质量同样具有竞争力：Replica 上 PSNR $37.94\,\mathrm{dB}$，TUM-RGBD 上 PSNR $23.46\,\mathrm{dB}$（最佳），AR-TABLE 上 $35.24\,\mathrm{dB}$ 且 GPU 用量最低。IMU 消融实验表明核心 RGB-D 流水线已提供稳定的高精度定位，IMU 传播仅带来边际提升（ATE 从 $5.94\,\mathrm{cm}$ 降至 $5.71\,\mathrm{cm}$）。

\subsubsection{LiDAR-视觉融合}

本路线介绍 LiDAR 点云与视觉图像在 3DGS SLAM 中的紧耦合方法，以 Gaussian-LIC 系列为代表。LiDAR 提供精确的几何测量和远距离感知能力，与视觉的丰富纹理信息互补。Gaussian-LIC（2024.04）~\cite{2404_06926} 是首个实时照片级 LiDAR-惯性-相机 SLAM 方法，旨在解决纯视觉 SLAM 在无界户外场景中的三大痛点：弱纹理导致几何退化、天空和变化曝光导致渲染伪影、以及大规模场景下实时性能的瓶颈。其核心思路是利用 LiDAR 的精确点云为高斯基元提供比视觉反投影更精确的初始化先验，从而加速收敛并提升重建质量。

\textbf{多模态里程计与数据准备。}系统由两个并行模块组成：LiDAR-惯性-相机里程计（负责位姿估计和建图数据准备）以及 3DGS 在线建图（负责照片级重建）。里程计模块基于先前工作 Coco-LIC（一种连续时间因子图优化框架），每 0.1 秒增量优化一段轨迹，融合三种因子：点到地图的 LiDAR 因子、帧到地图的视觉因子和惯性因子，即使在剧烈运动和光照变化等挑战条件下也能输出鲁棒、实时的位姿。系统支持固态 LiDAR 和机械旋转 LiDAR 两种类型。对每个时间窗口内的 LiDAR 点，系统以每 $N_l$ 个点随机保留一个的比例降采样后，投影到最新图像上着色并剔除相机视场外的点，与位姿标注的图像一起作为一帧发送给建图线程。

\textbf{LiDAR-视觉联合高斯初始化。}尽管 LiDAR 点提供了精确的几何先验，但其覆盖范围远小于相机——特别是机械旋转 LiDAR 的重复扫描模式。为此，系统维护一个类似 VINS-Mono 的视觉滑窗来补偿 LiDAR 盲区：滑窗包含基于共视性从最近图像中精选的 $N_k$ 个关键帧，每个关键帧包含由 KLT 稀疏光流跟踪的 Shi-Tomasi 角点；角点在连续被跟踪 $N_t$ 个关键帧后通过三角化生成 SfM 点，与 LiDAR 点合并用于高斯初始化。在初始化策略上，系统对每个 3D 点创建一个以该点位置为中心的高斯：球谐零阶系数填充其 RGB 颜色，不透明度设为 $o_a$，旋转设为单位矩阵。为抑制混叠伪影，尺度按深度自适应分配——离相机近的高斯分配较小尺度、远的分配较大尺度（直觉上，近处物体需要更精细的高斯来表达细节），形式化为 $s = \frac{d}{f} \cdot \mathbf{e}$，其中 $d$ 为点在当前相机坐标系下的深度、$f$ 为焦距、$\mathbf{e}$ 为全 1 向量。

\textbf{增量地图扩展与优化。}首帧之后，系统每 5 帧取一个关键帧用于地图扩展。收到关键帧时，先合并最近 5 帧的全部点云，再渲染当前高斯地图的不透明度图 $O$，生成掩码 $M$ 筛选当前地图尚未可靠覆盖的新区域（$M = O < \tau$）——只有投影到掩码像素上的点才被用于初始化新高斯，从而避免在已重建区域产生冗余基元。优化阶段，系统从所有历史关键帧中随机选取 $K$ 帧进行联合优化，避免灾难性遗忘并维护全局几何一致性。损失函数结合 $\ell_1$ 范数和 D-SSIM 项：

\begin{equation}
\mathcal{L} = (1-\lambda)\|E[C] - I\|_1 + \lambda \mathcal{L}_{\text{D-SSIM}},
\end{equation}

其中 $I$ 为真实图像，$C$ 为渲染图像，$\lambda$ 为权重系数。为处理户外场景中常见的变化曝光，系统额外优化一个 $3 \times 4$ 仿射矩阵 $E$ 作用于渲染颜色——前 3 列用于颜色通道的逐通道缩放，第 4 列用于偏移，使高斯的球谐颜色能在不同曝光条件下自适应校正。

\textbf{天空建模。}在无界户外场景中，天空若不加建模，表面高斯会被迫拟合天空外观，产生严重渲染伪影。系统在上半球以世界坐标系中心为原点、大半径 $R$ 为范围随机采样 $N_s$ 个位置初始化天空高斯，球谐零阶设为白色、初始不透明度设为 $o_b$（高于普通高斯的 $o_a$），尺度由最近邻天空高斯之间的距离决定，确保背景天空高斯始终位于前景之后。

\textbf{CUDA 加速策略。}系统纯 C++ 和 CUDA 实现，跟踪和建图模块基于 ROS 并行运行。论文精心设计了三项加速策略以突破实时性能瓶颈：
(1)~\textbf{逐高斯反向传播}（Per-Gaussian Backpropagation）：标准 3DGS 反向传播按像素分配 CUDA 线程，但每个像素对应的高斯数量差异极大，导致最慢像素成为瓶颈、其他线程空闲等待，且多像素反向传播到同一高斯时产生原子梯度加法冲突。逐高斯策略改为将线程分配给瓦片内所有高斯，提高并行度的同时减少冲突。
(2)~\textbf{瓦片剔除}（Tile Culling）：标准 3DGS 将投影的 2D 高斯近似为圆形来确定关联瓦片，但在 SLAM 增量建图中常见的细长 2D 高斯（近处被极度拉伸）对大部分瓦片的贡献微弱，白白浪费计算。系统通过负载均衡策略高效剔除这些弱贡献瓦片。
(3)~\textbf{稀疏 Adam}：仅优化当前相机视锥内的高斯而非全部，将优化器更新范围与场景总规模解耦。

\textbf{实验验证。}系统在三个公开 LiDAR-惯性-相机数据集上评估，涵盖固态 LiDAR（FAST-LIVO、R3LIVE）和机械旋转 LiDAR（MCD），共 9 个序列，包含 LiDAR 退化、视觉挑战和大规模校园等场景。所有序列使用相同超参数。在渲染质量上，Gaussian-LIC 在训练视角和新视角上均显著优于 NeRF-SLAM、MonoGS 和适配 LiDAR 的 SplaTAM——其 LiDAR 点和视觉 SfM 点的联合初始化提供了远超纯视觉方法的几何先验。在跟踪鲁棒性上，所有对比方法（基于单模态输入或简单重渲染损失跟踪）在复杂户外场景中均出现严重漂移甚至失败，而 Gaussian-LIC 通过紧密耦合三模态数据在连续时间因子图中持续获得鲁棒位姿。在实时性方面，Gaussian-LIC 是唯一满足实时要求的方法，在传感器数据采集时间窗口内同时完成跟踪和建图。消融实验证实天空建模、曝光处理和视觉 SfM 点三个组件各自不可或缺——移除任一组件均导致渲染质量下降。

Gaussian-LIC2（2025.07）~\cite{2507_04004} 继承了 Gaussian-LIC 的连续时间因子图里程计框架和增量高斯建图管线，针对其三项不足做出针对性改进：Gaussian-LIC 的视觉 SfM 点补偿方案在弱纹理区域和帧间视角变化较大时不够鲁棒，高斯地图的几何精度因缺少深度监督而欠佳，以及高斯地图的光度信息尚未反馈到里程计中辅助位姿估计。

\textbf{零样本深度补全替代视觉 SfM 点。}Gaussian-LIC 通过维护视觉滑窗三角化 SfM 点来补偿 LiDAR 盲区，但光流跟踪在弱纹理区域容易失败。Gaussian-LIC2 改用零样本深度补全模型 SPNet（基于 TensorRT 部署），直接以稀疏 LiDAR 深度图和 RGB 图像为输入预测稠密深度图。论文详细论证了选择深度补全而非单目深度估计或 MVS 的动机：单目相对深度估计存在尺度模糊，即使与 LiDAR 最小二乘对齐也难以在全局恢复准确度量；单目度量深度估计在大规模户外场景精度不足；MVS 对视差敏感且泛化能力有限。深度补全则以稀疏 LiDAR 为度量锚点，天然具备正确尺度且泛化性强。具体流程为：对每个关键帧的稀疏 LiDAR 深度图 $D_s$ 和 RGB 图像 $C$，SPNet 输出稠密深度图 $D_c$；通过深度变化阈值 $\tau_1$ 检测补全失败（罕见），再过滤负深度和高梯度区域得到 $D_c^f$；然后将 $D_s$ 划分为 $30\times 30$ 块，对无 LiDAR 覆盖的块选取 $D_c^f$ 中深度最小且小于 $\tau_2$ 的像素，反投影为补充点云 $P_s$，与着色 LiDAR 点合并用于高斯初始化。这一策略从根本上解决了稀疏 LiDAR（如 VLP-32C）和 LiDAR 与相机视场不匹配导致的 LiDAR 盲区欠重建问题。

\textbf{深度正则化提升几何精度。}Gaussian-LIC 的建图损失仅包含 RGB 渲染的 $\ell_1$ + D-SSIM 项，高斯仅被光度信号驱动优化，导致深度渲染质量差——论文指出这是现有 3DGS SLAM 方法的共性缺陷：过度关注视觉质量而忽视几何精度。Gaussian-LIC2 在损失函数中引入深度渲染损失项：

\begin{equation}
\mathcal{L} = \mathcal{L}_c + \beta_d\,\mathcal{L}_d, \quad \mathcal{L}_c = (1-\lambda)\|\hat{C} - C\|_1 + \lambda \mathcal{L}_{\text{D-SSIM}}, \quad \mathcal{L}_d = \|(\tfrac{\hat{D}}{\hat{O}} - D_s)\odot M_d\|_1,
\end{equation}

其中 $\mathcal{L}_c$ 与 Gaussian-LIC 的 RGB 损失形式相同，$\mathcal{L}_d$ 为深度损失，$\hat{D}$ 为渲染深度图，$\hat{O}$ 为渲染不透明度图（用于归一化，避免增量建图中因遮挡导致深度偏差），$D_s$ 为稀疏 LiDAR 深度图作为监督信号，$M_d$ 为有效深度掩码。权重 $\beta_d = 0.005$ 保持较小以确保 RGB 质量不退化。这一正则化不仅提升了深度渲染精度，还通过抑制过拟合改善了新视角合成质量——消融实验证实去除深度监督后定位精度下降，因为基于高斯地图的定位本质上需要从未知视角连续渲染，而过拟合的高斯地图在新视角下会产生伪影。

\textbf{高斯地图光度约束集成到里程计。}这是 Gaussian-LIC2 探索的最具前瞻性的改进方向。Gaussian-LIC 的里程计仅使用 LiDAR 点到地图因子、视觉重投影因子和惯性因子三种约束。Gaussian-LIC2 新增第二种相机因子选项：利用当前高斯地图从估计位姿渲染 RGB 图像 $\hat{C}$，与观测图像 $C$ 最小化光度残差（$\ell_1$ + D-SSIM），通过 Adam 迭代 $N_t=30$ 步优化相机位姿，然后将优化后的位姿作为因子图中的先验约束 $r_{C_2}$ 参与后续 LiDAR-惯性-相机联合优化。这一设计利用了连续时间轨迹的天然优势：因子图可在任意时间戳查询位姿，因此高斯地图渲染与原始传感器数据在时间轴上无缝对齐。在 LiDAR 退化场景（如固态 LiDAR 对准地面或墙壁）中，传统的 LiDAR 点到平面因子失效，而高斯地图的光度约束提供了替代定位信号——论文在 Visual\_Challenge 序列上的实验表明，光度约束方案（Ours-2）在该弱纹理序列上甚至优于基于光流重投影的默认方案（Ours-1），因为光度约束直接比较图像结构而非依赖特征匹配。

\textbf{CUDA 加速的继承与扩展。}Gaussian-LIC2 完整继承了 Gaussian-LIC 的逐高斯反传、快速瓦片剔除和稀疏 Adam 三项策略，并保留完整的三阶球谐各向异性高斯表示（不做简化）。在此基础上新增四项优化：分离球谐系数（低阶和高阶分开处理，消除耗时的拼接操作）、高效可分离高斯卷积 SSIM（源自 Taming-3DGS）、warp-level predicate masking（统一线程退出决策避免分支发散，并在全线程标记返回时提前终止整个 warp）、以及 pinned memory 非阻塞 CPU-GPU 传输（数据传输与 GPU 计算并行执行）。这些策略确保系统在大规模户外场景中维持实时性能。

\textbf{实验验证。}系统在五类数据集上评估，涵盖固态 LiDAR（Livox Avia、Mid-360）和机械旋转 LiDAR（OS1-64、VLP-32C），包括公开数据集（R3LIVE、FAST-LIVO、MCD、M2DGR）和自采数据集（提供精确轨迹真值和深度真值，支持序列外新视角评估）。在定位上，连续时间因子图紧密耦合方案在 LiDAR 退化和大规模场景中均取得最优精度，渲染式方法（MonoGS、SplaTAM、Co-SLAM）在户外大规模场景中因视场变化过大而频繁失败。在渲染质量上，Gaussian-LIC2 在序列内和序列外新视角合成上均显著优于 Gaussian-LIC——深度补全消除了 LiDAR 盲区的欠重建，深度正则化抑制了过拟合带来的新视角伪影，而 Gaussian-LIC 因忽视几何优化导致深度渲染质量差且无法重建 LiDAR 未扫描区域（如天花板）。论文还与离线方法 LetsGo 对比：在相同输入条件下，Gaussian-LIC2 在数据采集时间内完成重建即达到可比质量，额外迭代 10,000 次仅需 24 秒便显著超越 LetsGo 的 30,000 次迭代结果，证明增量建图管线的高效性。系统还展示了基于连续时间轨迹和 Gaussian 地图的下游应用：视频帧插值（利用轨迹任意时间戳查询 + 高斯地图任意视角渲染实现帧率倍增）和快速网格提取（增大深度权重后 TSDF 融合 + Marching Cubes 提取精确 mesh）。

GS-LIVO（2025.01）~\cite{2501_08672} 在上述 Gaussian-LIC 的 LiDAR-视觉融合路线上做了根本性的架构升级——将里程计框架从连续时间因子图替换为迭代误差状态卡尔曼滤波器（IESKF），并用全局八叉树 + GPU 滑动窗口的双层架构取代了 Gaussian-LIC 的单层全量高斯地图，从而首次将实时高斯 SLAM 部署到资源受限的嵌入式平台。其动机直指 Gaussian-LIC 等现有方法的两项瓶颈：一是地图构建线程远慢于位姿估计线程，机器人无法快速感知环境变化；二是高斯基元数量随场景增长导致的 GPU 显存膨胀，使得现有方法在嵌入式平台完全无法运行。

\textbf{全局八叉树地图与双层存储架构。}与 Gaussian-LIC 将所有高斯直接存储在 GPU 显存中不同，GS-LIVO 引入了空间哈希索引的递归八叉树作为全局地图的 CPU 端存储结构。根体素通过空间哈希键 $k = \lfloor \mathbf{p} / v_s \rfloor$ 索引，其中 $v_s$ 为根体素边长。八叉树可根据环境复杂度递归细分，天然支持多级细节（LoD）自适应：室内场景采用 0.03\,m 细粒度根体素、户外大规模场景采用 1.0\,m 粗粒度根体素，均支持 2 级细分。这一设计的核心考量在于：CPU 内存（RAM）容量远大于 GPU 显存（VRAM），且可通过交换空间轻松扩展，因此将全局八叉树驻留在 CPU 端，仅在 GPU 端维护当前视场（FoV）内的高斯子集，从根本上将显存占用与场景总规模解耦。

\textbf{LiDAR-相机联合初始化：从深度自适应到表面贴合。}Gaussian-LIC 对每个 3D 点创建各向同性高斯（旋转设为单位矩阵），尺度按深度自适应分配。GS-LIVO 则更进一步——利用 LiDAR 点的法向量（由 LiDAR-惯性 SLAM 系统估计）设定高斯旋转矩阵，将每个高斯建模为贴合物体表面的 2D 平面片元：缩放矩阵 $\mathbf{S}_i$ 在法向量正交方向上取八叉树层级决定的尺度，在法向量方向上取超参数 $\epsilon$（一个很小的切片厚度），使高斯变成紧贴表面的薄片。直觉上，表面几何主要沿法线方向变化剧烈，而沿切线方向平缓，因此将高斯压扁为表面片元比各向同性高斯更高效地编码了几何。在颜色初始化上，与 Gaussian-LIC 的球谐零阶直接填充不同，GS-LIVO 通过双线性插值从相机图像获取球谐零阶值，高阶初始化为零，随后用光度梯度迭代优化。

\textbf{增量式高斯滑动窗口：五步增量更新协议。}这是 GS-LIVO 区别于 Gaussian-LIC 稀疏 Adam 优化策略的核心创新。Gaussian-LIC 通过稀疏 Adam 仅优化当前视锥内的高斯，但高斯数据仍全量存储在 GPU 显存中；GS-LIVO 则设计了专门的滑动窗口，仅在 FoV 内维护连续内存布局的高斯子集。滑动窗口由三个缓冲区协同工作：空间哈希表（SHT，空间坐标到 CPU 内存指针的映射）、CPU 高斯缓冲区（CGB，连续内存存放活跃体素的高斯参数）、GPU 高斯缓冲区（GGB，GPU 端镜像用于并行优化）。每帧通过五步增量更新协议维护：(1)~\textbf{回写全局地图}：将上一帧滑动窗口中仍位于当前 FoV 内的体素标记为 OVERLAP，移出 FoV 的体素将优化后的参数回写全局八叉树并标记 DELETE；(2)~\textbf{删除与压缩}：将标记 DELETE 的体素与缓冲区尾部交换后移除，保持内存连续性；(3)~\textbf{重叠与新增加}：用当前 LiDAR 帧的空间哈希键识别与上一帧滑动窗口的重叠体素及需要新增的区域；(4)~\textbf{追加新叶体素}：将 ADD 体素追加到 CPU 缓冲区尾部并更新哈希表；(5)~\textbf{CPU/GPU 同步}：将 CPU 缓冲区数据直接传输到 GPU 缓冲区，确保高斯可立即被 GPU 优化和渲染。这一增量策略通过帧间重叠复用将内存拷贝降至最低，同时避免了原始 3DGS 中基于瓦片的深度排序导致前景污染背景的混叠问题。

\textbf{IESKF 紧耦合多传感器融合里程计。}在里程计框架上，GS-LIVO 与 Gaussian-LIC 采用了截然不同的路线。Gaussian-LIC 基于 Coco-LIC 的连续时间因子图，每 0.1 秒增量优化一段轨迹；GS-LIVO 则以迭代误差状态卡尔曼滤波器（IESKF）为框架，借鉴 FAST-LIVO2 的顺序更新结构，但用高斯地图的光度渲染替代了 FAST-LIVO2 中基于稀疏 patch 的光度误差。具体而言，视觉测量模型利用高斯地图渲染当前帧视角的图像 $\hat{I}_k$，通过最小化 $\hat{I}_k$ 与观测图像 $I_k$ 的光度残差来优化位姿 $\boldsymbol{\xi}$（相机位姿的李代数表示）。GS-LIVO 推导了光度损失到相机位姿的雅可比矩阵，并进一步系统性地将更新从相机位姿传播到 IMU 位姿，嵌入 IESKF 框架。这一设计的关键优势在于：现有大多数高斯基 SLAM 方法仅依赖优化器更新相机位姿，不评估更新后的协方差；而 GS-LIVO 的 IESKF 框架能将视觉更新的协方差传播到后续的 IMU 和 LiDAR 更新中，形成真正紧密耦合的多传感器融合系统，显著提升融合一致性。论文对比了 Gaussian 渲染与 FAST-LIVO2 的 patch warp 方法——随着 patch 尺寸增大，传统 warp 方法出现明显的拼接缝隙，而高斯渲染能无缝合成完整图像，包括非朗伯表面的照片级渲染。

\textbf{实验验证}覆盖桌面端（Intel i9-13900KF + RTX 4090，128\,GB RAM）和嵌入式端（Jetson Orin NX 16G），使用 FAST-LIVO2 数据集（CBD03、HKU01 大规模校园场景）、MARS-LVIG 数据集（HKairport01、HKisland03 无人机山海环境，含 D-RTK 精确真值）以及自采室内序列（MoCap 真值）。在渲染质量上，与 S3Gaussian 和 LetsGo 等离线 LiDAR 辅助方法相比，GS-LIVO 的固定层级八叉树在降低计算开销的同时保持可比渲染质量。在定位精度上，户外大规模场景轨迹 RMSE 达 0.042\,m，显著优于 R$^3$LIVE（1.465\,m）和 LVI-SAM（4.665\,m），与 FAST-LIVO 精度可比。在实时性上，桌面端室内场景全管线（含滑动窗口维护和高斯优化）处理时间低于 90\,ms，PSNR 稳定在 25--30\,dB；消融实验证实滑动窗口策略在保持 PSNR 可比的同时，显著降低 GPU 显存消耗和计算时间——无滑动窗口时显存和耗时随场景增长持续膨胀，有滑动窗口时两者保持稳定。在嵌入式端，根体素尺寸 0.5\,m、图像分辨率 $256 \times 216$、滑动窗口 20\,000 个高斯的配置下，优化耗时 15.3\,ms、地图维护 18.9\,ms、全管线 48.3\,ms，PSNR 23.52\,dB，并成功集成到基于 A* 全局规划与 LQR 轨迹跟踪的完整自主导航系统中——据作者所知，这是首个在 ARM 架构嵌入式平台上部署的实时高斯基 SLAM 系统。

\textbf{局限与展望}：当前八叉树层级和体素尺寸需手动设置，未来可探索依据观察距离、结构复杂度和纹理丰富度自适应调整 LoD；对颜色均匀区域，可合并相似高斯以进一步压缩存储和计算。

LIVE-GS（2025.07）~\cite{2507_23273} 指出上述 Gaussian-LIC 系列方法的两项不足：一是仅依赖稀疏位姿图优化实现回环闭合，缺乏照片级渲染所要求的表面级稠密几何对齐——位姿图仅在关键帧层面对齐轨迹，残余位姿不一致会直接导致高斯地图中出现鬼影伪影；二是沿用原始 3DGS 的 kNN 初始化策略，在 LiDAR 数据非均匀密度下（远处点稀疏）产生大量畸形的过大高斯，膨胀内存并拖慢渲染。为此，LIVE-GS 提出三项针对性改进。

\textbf{双地图架构与 Surfel Bundle Adjustment。}与 Gaussian-LIC 的单一高斯地图不同，LIVE-GS 维护稀疏 surfel 地图（用于高效定位）和稠密高斯地图（用于照片级渲染）两种表示，灵感源自 Elastic LiDAR Fusion。前端处理 LiDAR-相机-IMU 同步数据流，通过联合优化几何约束（基于 surfel 的点到平面 ICP 残差）和光度约束（稀疏特征点直接亮度比较）估计帧间相对位姿。当平移或旋转超过阈值时创建关键帧，并初始化和局部精炼高斯。后端执行全局优化：在位姿图中建立里程计约束边和 surfel-to-surfel 几何对齐边，前者保证连续性，后者引入远程结构约束。Surfel 对齐误差采用 G-ICP 框架建模，通过最近邻搜索建立对应关系并用马氏距离阈值剔除离群点，联合协方差由匹配 surfel 对各自的协方差合成。全局优化使用 iSAM2 增量求解，仅更新受影响的部分而非每步重新线性化整个系统。优化收敛后，仅对位姿被修正的关键帧关联的高斯施加刚体变换（更新位置和旋转），再以固定位姿对受影响高斯的尺度、不透明度和球谐系数进行光度精炼以恢复渲染质量——这一全局精修机制使 LIVE-GS 在大规模户外场景中维持长期一致性，消融实验证实 SBA 在所有测试数据集上均提升 PSNR 并降低最大高斯尺度。

\textbf{深度不变的高斯初始化。}Gaussian-LIC 按深度自适应分配高斯尺度（近小远大），这在 LiDAR 非均匀密度下产生问题——远处稀疏点被分配过大的高斯，渲染时呈现模糊。LIVE-GS 提出一种基于像素投影足迹的初始化策略：对每个 LiDAR 点先通过 kNN + PCA 获取旋转矩阵，建立各向异性的三维协方差；然后通过相机投影的雅可比矩阵计算该高斯在图像平面上的二维投影面积；再计算目标面积（直径 $n$ 像素的圆形足迹，论文评估 $n{=}1$ 和 $n{=}5$ 两种设定）与原始投影面积的比值 $k$，将三维尺度各分量乘以 $\sqrt{k}$ 使投影面积恰好等于目标值。直觉上，这确保每个高斯在图像平面上占据预设的固定像素面积，无论其距相机远近——近处高斯三维尺度小但占据合理像素，远处高斯三维尺度虽大但同样在像素空间受限，从而避免了传统方法中远处高斯因密度稀疏而过度膨胀的问题。消融实验表明 5 像素变体在 CBD02 上将地图从 235\,MB 压缩至 55\,MB（4.3 倍缩减），PSNR 仅下降 0.5\,dB。

\textbf{有界 sigmoid 重参数化。}在天空等观测不足区域，高斯在优化过程中容易无限膨胀、遮挡后方已优化好的高斯，降低渲染质量。LIVE-GS 对可学习尺度参数 $s$ 施加有界 sigmoid 映射 $\sigma_{\text{bounded}} = \sigma_{\min} + (\sigma_{\max} - \sigma_{\min}) \cdot \text{sigmoid}(s)$，将尺度严格约束在 $[\sigma_{\min}, \sigma_{\max}]$ 区间内。上界 $\sigma_{\max}$ 通过自适应机制动态调整：若超过 15\% 的高斯尺度接近上界（$>0.95\sigma_{\max}$），说明约束过紧，将 $\sigma_{\max}$ 放大 1.2 倍并重新映射；若超过 95\% 的高斯尺度接近下界，则缩小 $\sigma_{\max}$ 至原来的 0.8 倍。这一机制在保持高斯紧凑的同时允许场景中确实需要的大尺度高斯存在。消融实验显示有界 sigmoid 将最大尺度从数百米（无约束时）降至亚米级，与 SBA 联合使用时在所有数据集上取得最佳结果。

\textbf{实验验证。}系统在六个数据集上评估，涵盖 FAST-LIVO2 数据集（城市街景 CBD02、Retail、SYSU01）、MARS-LVIG 数据集（无人机海岸 airport03）、自采室内数据集（Reception、Demo room），以及 Botanic Garden 数据集（厘米级 RTK-GPS 真值用于轨迹评估）。在轨迹精度上，LIVE-GS 在 Botanic Garden 三个序列上均取得最优 ATE（0.528/0.529/0.498\,m），优于 Gaussian-LIC 和 GS-LIVM，证实 surfel bundle adjustment 对抑制长程漂移的有效性。在渲染质量上，1 像素变体在 CBD02、Retail 和 airport03 上取得最高 PSNR；5 像素变体以极小的质量损失换取大幅地图压缩。在鲁棒性上，GS-ICP-SLAM 和 GS-SDF 在 airport03 大规模户外场景中完全失败，而 LIVE-GS 正常运行。PINGS（2025.02）~\cite{2502_05752} 在统一的弹性点基隐式神经地图内融合有符号距离场（SDF）和高斯泼溅辐射场，通过几何一致性约束实现两种模态的相互增强——SDF 提供精确的表面几何，辐射场提供照片级外观，两者的联合表示支持同时满足导航和渲染的下游需求。

上述方法均围绕 LiDAR-视觉-惯性紧耦合展开。从 LiDAR 数据的初始化策略出发，LVI-GS（2024.11）~\cite{2411_02703} 利用着色 LiDAR 点初始化 3DGS——每个 LiDAR 点通过投影获取对应像素的 RGB 颜色，直接生成带颜色的初始高斯集合，加速收敛。金字塔训练方法在多级分辨率上渐进学习特征，避免高频噪声导致的优化不稳定。

从初始化策略转向大规模场景对齐，LiV-GS（2024.11）~\cite{2411_12185} 首次在大规模户外场景中直接对齐离散稀疏 LiDAR 数据和连续可微高斯地图，克服了传统 LiDAR 建图固定分辨率的局限——高斯地图的可微性使得 LiDAR 点到地图的对齐可以通过梯度优化连续精化。与 LiV-GS 同样处理 LiDAR 稀疏性问题但路径不同，GS-LIVM（2024.10）~\cite{2410_17084} 利用高斯过程回归（GPR）从统计建模角度缓解稀疏不均匀 LiDAR 观测带来的问题——GPR 在观测稀疏区域提供概率性的深度预测和不确定性估计，使高斯基元的初始化在 LiDAR 扫描间隙仍能获得合理的位置估计。

与前述方法均依赖视觉+LiDAR 融合不同，Splat-LOAM（2025.03）~\cite{2503_17491} 构建纯 LiDAR 的 Gaussian Splatting 里程计和建图管线，利用球面投影驱动高斯基元的优化——不同于视觉域的透视投影，球面投影天然适配 360$^\circ$ LiDAR 的全方位扫描模式，将 LiDAR 点直接作为高斯中心进行可微渲染和位姿优化。

回到多传感器融合方向，LiGSM（2025.03）~\cite{2503_05425} 用 LiDAR 点云初始化 3DGS 提供更稠密的起点，联合图像和点云构建损失估计位姿并优化传感器外参，实现动态自适应的传感器标定。

从通用场景转向特定应用领域，AgriGS-SLAM（2025.10）~\cite{2510_26358} 面向果园场景，耦合 LiDAR 里程计和多相机 3DGS 渲染，利用 LiDAR 回环闭合纠正长程漂移，多视角高斯泼溅的批量渲染从互补视角恢复果园的重复行结构和因风力/季节变化的外观差异，实现跨季节的鲁棒建图。

最后，上述所有方法都假设传感器已预先标定。MM3DGS（2024.04）~\cite{2404_00923} 挑战了这一前提，首次展示用未标定相机图像和惯性测量实现准确 SLAM——无需预先标定相机内参和相机-IMU 外参，系统在运行过程中在线估计所有标定参数，跟踪精度比当时最佳提升 3 倍。

\subsection{非传统传感器融合}
\label{ss:15_5}

本节讨论雷达、事件相机与热红外等非传统传感器与 3DGS SLAM 的融合方法。

\subsubsection{雷达、事件相机与热红外 SLAM}

本路线介绍基于雷达、事件相机和热红外等非传统传感器的 3DGS SLAM 方法，聚焦恶劣环境和特殊感知模态下的建图与定位。

雷达在恶劣天气和光照条件下具有独特优势。毫米波雷达对光照变化、雨雾雪等恶劣天气以及非金属材料遮挡具有天然鲁棒性，且成本低、功耗小，使其成为视觉和 LiDAR 在恶劣环境下失效时的关键互补传感器。然而，雷达 SLAM 长期面临一个核心瓶颈：现有雷达里程计（RO/RIO）方法通过逐帧估计速度或帧间匹配雷达点/图像来推算位姿，本质上遵循航迹推算（dead-reckoning）范式，漂移随时间不可避免地累积。虽然回环闭合和地点识别可以在 revisit 同一地点时纠正漂移，但这种纠正需要重返已观测区域，限制了在线探索场景下的适用性。视觉里程计中的多帧光束调整（BA）已被证明能在不依赖回环的情况下显著降低漂移——它联合优化滑动窗口内的相机位姿和场景结构。LiDAR 领域也已出现 BA 式优化（如 BALM、LONER）。然而，将 BA 扩展到雷达领域长期未被探索，主要挑战在于雷达缺乏稠密场景表示，且雷达感知涉及距离、方位角和 Doppler 三个维度的复杂信号模型。

RadarSplat-RIO（2026.04）~\cite{2604_13492} 首次将高斯泼溅的光束调整范式引入雷达 SLAM，填补了这一空白。

\textbf{动机与问题定义}：雷达里程计（如 MRIO）的典型管线是：从雷达原始数据经 FFT 生成方位-距离（RA）图像，通过 CFAR 检测提取特征点，利用 Doppler 速度估计雷达自运动速度（将每个特征点的 Doppler 速度建模为自运动速度在视线方向上的投影，通过最小二乘求解），再结合 IMU 预积分进行位姿图优化。这一管线的根本问题在于帧间里程计缺乏全局一致性校正——长序列运行中漂移不可避免，而传统回环闭合需要 revisit 同一地点，限制了纠正时机。

\textbf{核心方法}由前端和后端两部分构成。前端沿用 MRIO 的雷达-惯性里程计获取初始位姿，后端执行\textbf{位姿精化、建图和光束调整}三个过程，三者均基于 RadarSplat++ 场景表示的可微渲染。

RadarSplat++ 是对 RadarSplat 的关键扩展。每个高斯基元存储 3D 位置 $\mu$、旋转四元数 $q$、缩放向量 $S$ 和\textbf{雷达散射截面（RCS）}属性 $\sigma$（替代传统 3DGS 的 RGB 颜色）。基于雷达方程的渲染过程将 3D 高斯投影到极坐标空间，沿仰角轴进行正交加权平均，根据 RCS 值 $\sigma$ 渲染出 RA 图像。RadarSplat-RIO 在此基础上提出了\textbf{可微 Range-Doppler 渲染}——利用雷达自运动速度 $v$ 计算每个高斯基元的 Doppler 速度：

\begin{equation*}
v_i^{\text{Doppler}} = -\mathbf{r}_i^\top \mathbf{v},
\end{equation*}

其中 $\mathbf{r}_i$ 是从雷达到第 $i$ 个高斯的单位视线方向向量。随后，通过天线增益加权和 Doppler 速度的高斯核软分箱，将 RA 空间的功率和 Doppler 图像转换为 RD 图像——这一转换对 RA 图像和 Doppler 速度完全可微，使梯度能够回传至高斯参数和位姿。RadarSplat++ 因此成为首个支持雷达全维度（距离-方位-Doppler）测量的可微渲染管线。

后端三个过程依次运行。

\textbf{位姿精化}仅优化最新关键帧的位姿 $T_i$，损失函数联合 RA 和 RD 渲染误差：$\mathcal{L}_{\text{pose}} = \sum_{i \in W} \left[ L\!\left(\hat{I}_{RA}, I_{RA}\right) + L\!\left(\hat{I}_{RD}, I_{RD}\right) \right]$，其中自运动速度 $v_i$ 由前端初始化并在优化过程中通过相邻位姿差递归更新。

\textbf{建图}使用滑动窗口（$N=10$ 个关键帧），仅利用 RA 图像的几何匹配损失加上高斯尺度正则项 $\mathcal{L}_s(S)$ 防止高斯基元过度膨胀。

\textbf{光束调整}联合优化位姿和高斯参数，结合 RA 与 RD 渲染损失以及尺度正则项：$\mathcal{L}_{BA} = \sum_{i \in W} \left[ L\!\left(\hat{I}_{RA}, I_{RA}\right) + L\!\left(\hat{I}_{RD}, I_{RD}\right) \right] + \lambda_s \mathcal{L}_s(S)$，其中 $L$ 为 $L_1$ 和 SSIM 的组合损失。BA 采用半径空间窗口（$r_{BA} = 10$\,m）而非时间滑动窗口——当轨迹 revisits 已建图区域时，空间窗口自动选择几何相关的关键帧参与优化，无需显式回环检测即可改善几何一致性。IMU 预积分因子通过 GTSAM 库融合雷达平移速度和 IMU 旋转速度，为前端提供惯性约束。

\textbf{关键创新}体现在三个方面：一是首次将 BA 的联合位姿-场景优化能力从视觉域和 LiDAR 域扩展到雷达域，使雷达 SLAM 能够在不依赖回环闭合的情况下显著抑制漂移；二是 RadarSplat++ 的可微 RD 渲染使 Doppler 测量直接参与位姿和场景优化，Doppler 速度对线性运动的直接约束在特征缺失的长走廊等场景中尤为关键（RA-only 对齐在这些场景下病态）；三是半径空间窗口的 BA 策略，通过空间邻近性而非时间邻近性选择关键帧，天然适配 revisit 场景的漂移校正，无需显式回环检测模块。

\textbf{实验验证}使用配备 TI MMWCAS-RF-EVM 毫米波雷达和 Intel RealSense D435i 的自定义数据集（5 条室内序列，每条含至少 1--2 个不同规模的回环），以立体视觉惯性 SLAM（RTAB-Map）作为伪真值。与 MRIO 基线对比，RadarSplat-RIO 的平均绝对平移误差和旋转误差分别降低超过 90\% 和 80\%，显著改善了位姿估计鲁棒性并有效抑制了长期漂移。消融实验证实：移除整个后端导致轨迹迅速漂移；禁用 BA 模块精度显著下降；省去 RIO 前端初始化则因缺乏可靠位姿初值导致最差性能——三者缺一不可。RD 渲染损失主要提升平移精度，因为 Doppler 测量直接约束线性运动，在特征缺失走廊中收益最为显著。半径窗口（$r_{BA} = 10$\,m）优于其他半径值和时间滑动窗口配置。

\textbf{局限与展望}：当前方法仅在室内环境验证，数据集尚未公开；伪真值来自视觉惯性 SLAM 而非运动捕获系统，精度上限受限于视觉 SLAM 本身的误差；雷达与 IMU 之间的硬件时钟同步不完善可能影响前端精度（但后端优化对此表现出良好鲁棒性）。未来可扩展至户外大规模场景，探索与回环闭合的协同，以及利用高斯场景表示实现雷达语义建图。

Rad-GS（2025.11）~\cite{2511_16091} 正是将雷达高斯 SLAM 推向户外公里级场景的工作。上一节 RadarSplat-RIO 在室内环境验证了雷达高斯表示的 BA 范式，但受限于室内毫米波雷达+IMU 的纯惯性架构，且未涉及动态物体。Rad-GS 将问题域从室内扩展到户外公里级场景，同时引入两个 RadarSplat-RIO 完全未涉及的新挑战：\textbf{动态物体干扰}和\textbf{大规模内存管理}。在传感器配置上，Rad-GS 采用 4D 毫米波雷达（Oculii G7）配合单目相机，而非 RadarSplat-RIO 的纯雷达-惯性组合。

\textbf{数据增强与 Doppler 动态去除}构成了 Rad-GS 的第一个核心模块，也是其与 RadarSplat-RIO 的根本差异所在。RadarSplat-RIO 利用 Doppler 速度 $v_i^{\text{Doppler}} = -\mathbf{r}_i^\top \mathbf{v}$ 构建 RD 渲染损失来约束位姿——Doppler 在那里的角色是\textbf{位姿优化的观测信号}。Rad-GS 则将 Doppler 的角色完全转向\textbf{动态物体识别}。具体而言，首先通过 CMDF 跨模态扩散融合模块将原始稀疏雷达点云增强为稠密点云（保留几何连续性但丢失 Doppler 信息），然后利用 Doppler 速度残差在原始稀疏雷达检测中标记动态索引：$|v_{\text{dop},i} - \hat{v}_{\text{dop},i}| > \delta_v$，其中 $\hat{v}_{\text{dop},i} = \mathbf{r}_i^\top \hat{\mathbf{v}}_m$ 是基于最小二乘估计的平台速度投影。接下来通过自适应八叉树将动态标签从稀疏原始点传播到增强稠密点云——对每个增强点 $\hat{p}_i$，取其 $K$ 个最近叶节点 $\{n_k\}_{k=1}^K$，当超过半数叶节点具有动态标签时标记该点为动态。动态标签进一步向上传播至最近公共祖先以保证空间连续性，最终投影到图像平面生成 2D 边界框，引导 EfficientSAM 在框内提取精确物体轮廓。这一单帧动态去除策略的关键优势在于：无需先验位姿估计或多帧时序累积，直接在 3D 空间利用 Doppler 物理量识别运动物体，从而避免了视觉残差方法和语义方法在户外光照变化下失效的问题。

\textbf{前端跟踪}采用\textbf{协方差引导的雷达-高斯匹配}——这一策略借鉴自 LiV-GS 的点云-高斯对齐思想。根据局部几何各向异性（通过 SVD 估计）将增强雷达点分为三种形状类型：平面型（$\sigma_3/\sigma_2 \leq \tau$，法向量约束）、线型（$\sigma_1/\sigma_2 \geq 1/\tau$，主方向约束）和混合型（加权组合），为每种类型构造形状自适应特征矩阵 $\mathbf{v} \in \mathbb{R}^{3 \times 3}$。位姿通过迭代优化对齐输入点云与最近高斯质心：$E = \sum_{\mathbf{x}_p \in P_R} w(\mathbf{x}_p)(\mathbf{v}_{\mathbf{x}_p} \cdot (\mathbf{T} \mathbf{x}_p - \boldsymbol{\mu}_g))^2 + R(\mathbf{v}_{\mathbf{x}_p}, \mathbf{v}_g)$，其中密度加权 $w(\cdot)$ 和方向一致性正则项 $R(\cdot)$ 共同约束匹配质量。这与 RadarSplat-RIO 的前端（基于 MRIO 的雷达-惯性里程计获取初始位姿）形成对比——Rad-GS 不依赖 IMU，而是通过协方差引导的点-高斯匹配直接估计位姿。

\textbf{后端精化}利用与雷达帧非同步的图像帧提升渲染质量。RadarSplat-RIO 的后端执行位姿精化-建图-BA 三阶段联合优化，全部基于 RadarSplat++ 的可微 RA/RD 渲染。Rad-GS 的后端则更加简洁：利用三次 Hermite 插值（相邻关键帧位姿 $P_0, P_1$ 和平台速度 $V_0, V_1$ 作为边界条件）获取非同步图像的中间位姿，然后对 3DGS 表示进行全局渲染精化。损失函数由三部分组成：$\mathcal{L} = (1-\lambda_1)E_{pho} + \lambda_1 E_{geo} + \lambda_2 E_{rou}$，其中光度误差 $E_{pho}$ 比较渲染图像与真实图像，几何误差 $E_{geo}$ 比较渲染深度与增强雷达深度。

\textbf{粗糙度感知损失} $E_{rou}$ 是 Rad-GS 针对户外场景的专门设计——户外环境中从粗糙树干到光滑建筑面的表面粗糙度差异极大，单一约束策略无法兼顾。RadarSplat-RIO 仅使用尺度正则项 $\mathcal{L}_s(S)$ 防止高斯过度膨胀，而 Rad-GS 根据高斯基元尺度比 $\sigma_{mean}/\sigma_3$（$\sigma_{mean}$ 为两个较接近尺度的均值）与阈值 $\tau$ 的关系自适应切换：当 $\tau \leq \sigma_{mean}/\sigma_3 \leq 1/\tau$ 时仅约束尺度差异 $\|\delta_\sigma\|^2$，否则额外增加最小尺度惩罚 $\gamma\|\sigma_{min}\|^2$。在 Nyl2 序列上，这一损失比各向同性损失提升 PSNR 2.96\,dB，比法线引导损失提升 2.88\,dB——各向同性损失过度平滑导致建筑边缘色彩模糊和树干几何漂移，法线引导损失虽锐化平面区域但在粗糙表面产生波纹伪影并模糊细长物体。

为应对公里级场景的内存挑战，Rad-GS 专门设计了\textbf{全局八叉树管理}策略。与 RadarSplat-RIO 的滑动窗口建图（$N=10$ 个关键帧）和半径空间窗口 BA 的策略不同，Rad-GS 采用增量式全局管理方案：新点通过贝叶斯融合合并到已有高斯基元之中——将雷达距离观测建模为零均值高斯噪声 $d = d_{true} + \varepsilon$, $\varepsilon \sim \mathcal{N}(0, \sigma^2)$，以贝叶斯方式更新高斯基元的均值和协方差（$\bm{\Sigma}'_i = (\bm{\Sigma}_i^{-1} + \sigma_\varepsilon^{-2}\mathbf{I}_3 + \bm{\Sigma}_p^{-1})^{-1}$），既降噪又避免冗余高斯创建。分裂时通过八叉树条件约束确保不增加节点数——新点从最近高斯按正态分布采样，空间分裂严格遵循八叉树结构。在 NTU4DRadLM 数据集上，这一策略使 GPU 内存在约 1200 帧后趋于稳定（约 11 GB），在 2400 帧时比 LiV-GS（LiDAR 输入）降低约 45\% 内存消耗，比 MonoGS（LiDAR 输入）降低约 39\%。

\textbf{实验验证}在两个数据集上进行。NTU4DRadLM 数据集（4 条序列，含 4D 雷达+单目相机+LiDAR+GNSS 真值）评估渲染质量和动态去除：Rad-GS 的动态去除模块在三个序列上平均 PSNR 比 T-3DGS 高约 1.7\,dB，比 SLS 高约 2.7\,dB，有效恢复了道路标线和建筑立面等背景细节。自定义 Campus 数据集（1.2\,km 闭环+0.5\,km 开环，RTK-GNSS 真值）评估定位精度：在 Campus Loop 上，Rad-GS 的绝对轨迹误差为 6.454\,m，旋转漂移 1.130\,$^\circ$/100\,m，优于 4DRadarSLAM（11.542\,m, 2.163\,$^\circ$/100\,m）和 LiV-GS（20.139\,m, 5.542\,$^\circ$/100\,m）。值得注意的是 Rad-GS 使用里程计位姿与 GT 位姿的渲染质量差距很小（PSNR 差约 0.5--1.1\,dB），表明增量建图过程中的位姿估计误差累积得到有效抑制。在含透明玻璃的 Cp 序列中，Rad-GS 因图像与雷达点云的不匹配而出现漂移——这是雷达-视觉融合的固有限制，因为雷达可穿透玻璃而相机不能。

4D 雷达里程计（2024.12）~\cite{2412_13639} 利用高斯建模的概率密度函数进行配准，多假设扫描匹配提供更丰富的场景模型。GS-GVINS（2025.02）~\cite{2502_10975} 则从另一个方向拓展了高斯泼溅在导航中的应用——将 3DGS 的可微渲染与 GNSS 紧耦合，面向公里级户外驾驶场景。与 RadarSplat-RIO 的雷达渲染管线不同，GS-GVINS 构建的是 GNSS-Visual-Inertial 因子图优化（FGO）框架：在已有的 GNSS RTK、INS 和特征点视觉因子基础上，新增\textbf{3DGS 光度因子}，通过渲染当前视角的图像并与观测图像比较 $L_1$ 光度误差来约束位姿。值得注意的是，3DGS 建图与 FGO 解耦——高斯参数仅在建图线程中独立优化，FGO 仅优化相机位姿和仿射亮度参数，避免高斯参数与导航状态耦合导致优化规模爆炸。

该方法的\textbf{设计动机}在于：在 GNSS 信号恶化的城市场景（卫星遮挡、多径效应），惯性导航漂移加速，而传统特征点视觉因子对弱纹理和极端曝光场景鲁棒性不足——3DGS 的稠密渲染在这些场景中仍能保持高保真度，提供额外的位姿约束。三个核心技术创新支撑了这一思路：

（1）\textbf{SE(3) 解析雅可比}：GS-GVINS 是首个将相机位姿对 3D 高斯的解析 Jacobian 直接融入 FGO 的系统（而非使用自动微分），利用李群李代数推导 $SE(3)$ 上的最小参数化雅可比，显著降低计算开销。

（2）\textbf{基于建图损失的自适应加权}：3DGS 因子的可靠性取决于高斯地图质量。GS-GVINS 设计方差模型 $\sigma^2_{3DGS,i} = \sum_k E^k_{\text{photometric}} + \lambda_{\text{isotropic}} E_{\text{isotropic}}$，将当前关键帧窗口内所有帧的光度误差和各向同性正则化误差汇聚为因子方差——建图质量差的区域自动降低 3DGS 因子的权重，避免低质量渲染引入错误约束。

（3）\textbf{运动感知高斯剪枝}：户外驾驶中车辆频繁经历急减速或近乎静止等极端动态状态，此时高斯地图更新不足且单目深度初始化粗略，导致渲染出现模糊、空洞或漂浮伪影。剪枝机制根据相邻帧相对平移量 $t_i$ 判定极端运动（$t_i < t_{\min}$ 或 $t_i > t_{\max}$），移除沿相机射线累积不透明度 $\alpha$ 在到达最新插入高斯前已达到 0.5 的旧高斯，降低前方冗余密度。

\textbf{实验验证}使用自采集数据（u-blox F9P GNSS + ICM-20948 IMU + FLIR Blackfly-S 相机）和公开 UrbanNav 数据集，覆盖 open-sky、sub-urban、urban 三类环境。与 GICI-LIB 和 IC-GVINS 对比，在 GNSS 友好的 open-sky 场景中三者性能接近（RTK 绝对定位占主导）；在 sub-urban B 数据集上，GS-GVINS 的 3D 平移 RMSE 分别比 GICI 和 IC-GVINS 降低 17.26\% 和 71.93\%，3D 旋转 RMSE 降低 14.62\%；在 ultra-dense urban canyon 中 IC-GVINS 发散，GS-GVINS 的 3D 平移 RMSE 仍比 GICI 降低 13.71\%。消融实验表明：在地下隧道等极端曝光场景中，特征点跟踪失效但 3DGS 渲染仍能准确约束垂直方向平移；运动感知剪枝在车辆减速至近乎停止时维持了渲染质量。

EGS-SLAM（2025.08）~\cite{2508_07003} 与 Gaussian-LIC 同属多传感器融合 GS-SLAM，但切入点截然不同：Gaussian-LIC 融合 LiDAR-惯性-相机以解决户外弱纹理和大规模重建，而 EGS-SLAM 引入事件相机以解决运动模糊下跟踪和建图同时退化的问题。事件相机以微秒级异步感知亮度变化，天然免于运动模糊，但其输出是稀疏、离散的事件流而非密集图像帧——如何将事件流与 RGB-D 帧在统一的 3DGS 表示下对齐并联合优化，是 EGS-SLAM 的核心挑战。

\textbf{曝光期间的连续轨迹建模与模糊感知渲染。}与 Gaussian-LIC 在连续时间因子图中优化轨迹不同，EGS-SLAM 采用更简洁的线性插值方案：对每帧曝光区间 $[0, \tau]$，通过球面线性插值（Slerp）连接起始位姿 $\mathbf{T}_0$ 和终止位姿 $\mathbf{T}_\tau$，以 $K$ 个等距采样点近似连续轨迹，每个采样点从 3DGS 地图渲染一幅瞬时清晰图像。模糊图像则建模为这些瞬时清晰图像在曝光时间上的积分近似：$\tilde{I}(\bm{u}) \approx \frac{1}{K}\sum_{k=0}^{K-1} \mathcal{I}(\mathbf{T}(\eta_k), \bm{u})$。这一物理成像模型的引入使得跟踪和建图都能在模糊输入下工作——跟踪通过最小化渲染模糊图像与观测模糊图像的残差来估计连续轨迹的两个端点位姿，建图则在固定高斯地图的条件下联合优化高斯参数和近期关键帧位姿。

\textbf{事件-图像双模态对齐：可学习相机响应函数。}事件数据具有高动态范围（HDR），而标准图像受限于低动态范围（LDR），二者的动态范围不一致阻碍了在统一 3DGS 表示下的联合优化。EGS-SLAM 引入可学习相机响应函数（CRF），将渲染的 HDR 图像映射到 LDR 空间，以 $N$ 个可训练控制节点 $\bm{Q}$ 通过线性插值逼近可微 CRF。CRF 在负值、正常值和过曝值三个区间分别采用线性、插值和 $-\alpha/\sqrt{\cdot}$ 衰减策略（$\alpha=0.01$），确保全范围可微且梯度稳定。

\textbf{事件渲染损失与无事件损失。}事件图无法直接从高斯场景渲染，EGS-SLAM 利用事件生成模型间接获取：对连续两个轨迹采样点渲染的灰度图像（通过 BT.601 亮度变换获取）取对数差分，模拟事件触发 $\hat{E}_k(\bm{u}) = \log(\hat{B}(\mathbf{T}(\eta_k), \bm{u})) - \log(\hat{B}(\mathbf{T}(\eta_{k-1}), \bm{u}))$。事件损失 $L_{HE}$ 对比累积事件图与渲染事件图，而无事件损失 $L_{NE}$ 则惩罚无事件发生区域的虚假亮度变化，有效抑制去模糊中常见的振铃伪影（ringing artifacts），权重 $\lambda_{NE}=0.4$。深度损失取曝光期间所有采样点渲染深度与观测深度的最小 $L_1$ 距离，以匹配运动模糊帧中深度与瞬时位姿的不确定性。

\textbf{实验表现。}EGS-SLAM 构建了合成数据集 EventReplica（基于 Replica 扩展的运动模糊-事件数据）和真实数据集 DEVD（DAVIS346 事件相机 + RealSense D435i 深度相机 + FZMotion 动捕真值）。在 EventReplica 上，EGS-SLAM 的平均 PSNR 达 27.62\,dB，较 MonoGS 的 23.32\,dB 提升 4.20\,dB；平均 ATE 从 MonoGS 的 9.22\,cm 降至 5.32\,cm（改善约 42\%）。ORB-SLAM2 和 Photo-SLAM 在多数序列中因运动模糊导致跟踪丢失。消融实验证实事件跟踪的贡献最为关键——仅启用事件跟踪即可大幅提升位姿精度，再加入事件建图则进一步提升渲染质量；无事件损失在深度质量受限的真实场景中显著抑制振铃伪影；CRF 对齐 HDR-LDR 动态范围差异带来额外精度提升。系统在 RTX 4080 上运行约 0.55 FPS，慢于 MonoGS 的 1.75 FPS，主要因为需要在 $K$ 个轨迹采样点上重复渲染以融合事件和图像信号。

IncEventGS（2024.10）~\cite{2410_08107} 在 EGS-SLAM 的基础上进一步去掉了 RGB-D 帧和深度传感器，仅凭单个事件相机完成增量式 3DGS 重建与位姿估计。去掉 RGB-D 后，IncEventGS 面临双重难题：既无密集图像帧可提供绝对亮度监督，也无深度观测可约束几何——跟踪和建图必须完全依赖事件流的亮度变化信号驱动。

\textbf{事件 chunk 与对数亮度变化监督。}与 EGS-SLAM 利用模糊图像建模（曝光区间内 $K$ 个采样点渲染积分）不同，IncEventGS 将事件流按 50\,ms 固定窗口切分为 chunk，每个 chunk 仅用两个端点 SE(3) 位姿 $\mathbf{T}_\text{start}$、$\mathbf{T}_\text{end}$ 通过 $\mathfrak{se}(3)$ 空间线性插值参数化连续轨迹。在每个 chunk 内随机采样两个邻近时刻 $t_k$ 和 $t_k{+}\Delta t$，从 3DGS 渲染对应亮度图像 $\hat{I}(t_k)$、$\hat{I}(t_k{+}\Delta t)$，合成对数亮度变化 $\hat{\mathbf{E}}(\mathbf{x}) = \log \hat{I}(t_k{+}\Delta t, \mathbf{x}) - \log \hat{I}(t_k, \mathbf{x})$，与事件流累积测量的亮度变化 $\mathbf{E}(\mathbf{x}) = C \cdot \sum_{t_k < t_i < t_k{+}\Delta t} p_i$（$C$ 为对比度阈值）求 $\ell_2$ 距离作为监督。这里的关键差异是：EGS-SLAM 的亮度变化来自灰度图像的 BT.601 变换后对数差分（事件图与模糊帧图的混合监督），而 IncEventGS 的亮度变化完全从 3DGS 渲染的灰度图直接取对数差分，无需任何帧式输入。损失函数联合事件 $\ell_2$ 损失与结构不相似性（DSSIM）损失：$\mathcal{L} = (1{-}\lambda)\mathcal{L}_\text{event} + \lambda \mathcal{L}_\text{SSIM}$，权重 $\lambda{=}0.05$。

\textbf{增量跟踪-建图与滑动窗口 BA。}IncEventGS 采用了经典 SLAM 的 tracking-mapping 交替范式，而非 EGS-SLAM 所使用的跟踪（模糊帧匹配）+ 建图（联合优化高斯与关键帧位姿）的分离流程。跟踪阶段仅优化最新 chunk 的两个端点位姿（200 步迭代），目标函数为 $\|\mathbf{E}(\mathbf{x}) - \hat{\mathbf{E}}(\mathbf{x})\|_2$；建图阶段将跟踪所得轨迹作为初始值，在最近 $n_w{=}20$ 个 chunk 的滑动窗口内联合优化所有轨迹参数与 3DGS 属性（1500 步迭代）。与 EGS-SLAM 的建图联合优化高斯参数与近期关键帧位姿相比，IncEventGS 的 BA 完全由渲染亮度变化驱动，没有任何 RGB 图像、深度图或 CRF 校正参与。

\textbf{从随机初始化到深度引导 bootstrap。}帧式 3DGS 方法通常依赖 COLMAP 提供初始点云和位姿，但仅有事件相机时这些不可获得。IncEventGS 在包围盒内随机采样点云初始化 3DGS，用前 $m{=}3$ 个 chunk 联合优化高斯属性和随机初始化的相机位姿（4500 步）。这一过程能产生合理的亮度渲染，但因初始基线过短，三维结构质量不佳。论文因此引入单目深度估计网络（RepVG）从渲染亮度图预测稠密深度图，反投影重新初始化高斯中心位置后再次优化——消融实验证实深度 bootstrap 将 ATE 从 1.534\,cm 降至 0.064\,cm（改善约 96\%），PSNR 从 17.80\,dB 提升至 21.74\,dB。

\textbf{增量高斯增长与可见性掩码。}跟踪获得新 chunk 的精确位姿后，通过渲染深度图反投影在已探索区域添加新高斯。为避免对已建图区域的冗余重建，系统渲染不透明度图 $V$，仅在不透明度低于阈值 $\lambda_V{=}0.8$ 的像素处添加新高斯——这些像素对应尚未被充分覆盖的未知区域。

\textbf{实验验证。}系统在 Replica 合成数据集（5 个场景，1000\,Hz 仿真事件流）和 TUM-VIE 真实数据集上评估。尽管事件 NeRF 方法（E-NeRF、EventNeRF、Robust e-NeRF）使用真值位姿训练，IncEventGS 在 PSNR/SSIM/LPIPS 上仍全面超越它们（如 office2 场景 PSNR 21.74\,dB vs.\ Robust e-NeRF 的 16.81\,dB），论文将此归因于 3DGS 显式表示相比 NeRF 隐式表示的建模效率优势。相比两阶段基线（E2VID 重建图像 → COLMAP 估计位姿 → 3DGS 训练），IncEventGS 的端到端 BA 在位姿精度和新视角合成上均显著更优，因为两阶段方法的误差会逐级累积。在轨迹估计上，IncEventGS 在 Replica 上的 ATE 为 0.046--0.067\,cm，优于事件视觉里程计 DEVO（0.138--0.289\,cm）和 ESVO2（0.251--2.506\,cm）。论文还在快速运动场景下与 MonoGS 和 SplaTAM 对比，事件相机的高时间分辨率使其在运动模糊下性能退化远小于帧式方法。

TRGS-SLAM（2026.03）~\cite{2603_20443} 则将上述连续时间轨迹思想从 LiDAR-视觉-惯性（Gaussian-LIC 系列的因子图框架）转向了一个此前鲜有 GS SLAM 涉足的全新模态组合——微测辐射热计热红外相机 + IMU。热红外相机能在完全黑暗和烟雾环境中成像，但微测辐射热计在剧烈运动下面临三重退化：运动模糊（热时间常数 $\tau$ 导致像素对入射功率做指数衰减积分）、滚动快门畸变（逐像素读出导致同一帧内不同像素对应不同时刻的位姿）、以及缓慢时变的固定模式噪声（FPN，每个像素的偏置不同且随时间漂移）。TRGS-SLAM 的核心贡献是在 3DGS 渲染管线中内嵌了这三类退化的物理模型，使得模糊、畸变和噪声不再是跟踪的障碍，反而成为可微渲染的一部分被联合优化。

\textbf{B 样条连续时间轨迹与两阶段 IMU 融合。}与 Gaussian-LIC 通过因子图每隔固定时间窗口增量优化轨迹不同，TRGS-SLAM 用两条均匀 B 样条分别参数化位置 $\mathbf{p}(t)$ 和旋转 $R(t)$（4 阶，保证 $C^2$ 连续即加速度连续），轨迹值由控制点确定，可在任意时刻查询。这一选择的关键动机是：运动模糊和滚动快门建模需要在每个像素的读出时刻 $t_{x_p,y_p,l} = t_{0,0,l} + x_p\,\Delta t_\text{pix} + y_p\,w\,\Delta t_\text{pix}$ 估计相机位姿——B 样条的闭式导数天然支持这种逐像素时间采样。IMU 融合分为两阶段：初始化前仅用陀螺仪残差 $\mathbf{r}_{g,m} = \boldsymbol{\omega}_C^C(t_m) - (\tilde{\boldsymbol{\omega}}_{I,m}^I - \mathbf{b}_g)$ 约束角速度，轨迹处于相对尺度；初始化后通过解析方法求解重力方向和尺度因子，切换到绝对尺度的 IMU 轨迹，同时引入加速度计残差 $\mathbf{r}_{a,m}$，构成完整的惯性约束。IMU 损失形式为 $\mathcal{L}_{\text{IMU},\mathcal{I}} = \sum_m \lambda_g\|\mathbf{r}_{g,m}\|_2 + \lambda_a\|\mathbf{r}_{a,m}\|_2$，辅以偏置正则 $\mathcal{L}_\text{bias} = \lambda_\text{bias}(\|\mathbf{b}_g\|_2 + \|\mathbf{b}_a\|_2)$ 防止异常残差导致过拟合。

\textbf{微测辐射热计感知渲染。}这是 TRGS-SLAM 区别于所有前述方法的核心差异。标准 3DGS 假设输入是锐利的全局快门图像，直接在单一位姿处光栅化。TRGS-SLAM 则将渲染管线改造为三步：(1) 在曝光时间窗口 $[t_\text{max,l}-T_b,\, t_\text{max,l}]$ 内均匀采样 $N_r$ 个时刻，通过 B 样条查询每个时刻的相机位姿并光栅化得到 $N_r$ 幅图像 $I_G(t_r)$；(2) 对这 $N_r$ 幅图像按预计算的权重 $\{W_r\}$ 加权求和，模拟微测辐射热计的热时间常数积分效应，得到模糊渲染图 $M_l = \sum_{r=0}^{N_r-1} I_G(t_r) \odot W_r$；(3) 叠加两个 B 样条建模的 FPN——像素级偏置 $O_\text{pix}(t)$（以零均值损失 $\mathcal{L}_\text{FPN}$ 约束避免与信号混淆）和全局偏置 $o_\text{global}(t)$——得到最终渲染 $\hat{I}_l = M_l + O_\text{pix}(t'_{0,0,l}) + o_\text{global}(t'_{0,0,l})$。图像损失为 $\ell_1$ 范数。直觉上，运动模糊不再是跟踪的"噪声"，而是由渲染管线显式建模的"信号"——跟踪优化的目标就是找到使模糊渲染与模糊观测最一致的轨迹。

\textbf{视角多样性不透明度重置与漂移校正。}TRGS-SLAM 对标准 3DGS 的全局不透明度重置策略做了 SLAM 场景下的针对性改造。全局重置在 SLAM 中既低效又破坏稳定性（已建立的高斯可能被误删），因此 TRGS-SLAM 引入视角多样性指标：对每个高斯 $g$ 累积观测方向的二阶矩矩阵 $V_g = V_{g,\text{prev}} + \sum_{r \in R_g} \mathbf{d}_{r,g}\,\mathbf{d}_{r,g}^\top$，仅对条件数（最大特征值/最小特征值）超过阈值的高斯执行不透明度重置——这些高斯仅从单一方向被观测，几何约束不足。此外，建图优化中使用 $\lambda_s > \lambda_w$ 对历史关键帧赋予比当前窗口更高的权重，防止地图覆盖已观测区域导致位姿漂移；当 SSIM 低于阈值时自动回滚系统状态并暂停建图，等待跟踪重定位。

\textbf{实验验证。}TRGS-SLAM 在 TRNeRF 数据集（含剧烈运动的热红外序列，真值位姿由立体视觉 SfM 提供）上进行评估。这是目前唯一能在严重热红外退化条件下准确跟踪的 GS SLAM 系统——对比方法（MonoGS、SplaTAM 等视觉 GS SLAM）在运动模糊和低对比度的热红外图像上跟踪频繁失败。离线精炼后，TRGS-SLAM 的热红外图像恢复质量与使用真值位姿的 TRNeRF 可比，实现了无需真值位姿的竞争性图像去模糊与去噪。

如果说 TRGS-SLAM 走的是"端到端物理建模"路线——将热红外传感器的退化模型嵌入可微渲染管线联合优化——那么 TOM-GS（2026.02）~\cite{2602_07493} 则选择了一条截然相反的路径：将跟踪与建图完全解耦为两个独立模块，各司其职。这一设计选择并非随意，而是源于一个实践观察：几何方法（ORB-SLAM3、LDSO 等）在热红外数据上频繁失败，而学习型视觉里程计（DROID-SLAM 等）在 RGB 域已展现出对低纹理和剧烈运动的强鲁棒性。TOM-GS 的核心问题由此转化为：能否将学习型里程计的成功迁移到热红外域，并以 3DGS 作为稠密建图后端？

\textbf{高位深适配与深度增强里程计。}TRGS-SLAM 在渲染端处理传感器退化，TOM-GS 则在前端就解决输入适配问题。热红外相机输出 14 或 16 位原始图像，而 DROID-SLAM 的网络仅在 8 位 RGB 上训练——这是一个此前无人处理的域迁移缺口。TOM-GS 采用 Fieldscale 自适应增强策略将高位深转换为 8 位灰度图；消融实验证实该策略一致优于 OpenCV CLAHE 和 Shin 方案，因为 Fieldscale 能提供均匀增强且无过曝伪影。里程计模块保留了 DROID-SLAM 的 ConvGRU 光流预测 + DBA 稠密光束调整框架，但新增了 DSO 层（Disparity and Scale Optimization）来融合预训练单目深度网络（评估了 Depth Anything、ZoeDepth、Metric3D、LOTUS 四种）的预测：通过最小二乘估计仿射参数 $\theta_i$、$\gamma_i$ 对齐单目深度与 DROID 深度的尺度，再对多视角一致性检测出的高误差像素用对齐后的单目深度替换（权重 $\alpha_1{=}0.01$, $\alpha_2{=}0.1$）。这一机制的动机很直接——热红外图像的低纹理导致 DROID 的光流估计退化，而单目深度先验恰好能弥补这一弱点。

\textbf{深度引导的 GS 建图。}建图模块接收里程计输出的关键帧位姿和稠密深度图，以 3DGS 作为场景表示。TRGS-SLAM 的 3DGS 直接在联合优化中端到端更新，TOM-GS 则是跟踪-建图分离架构——建图模块不反哺位姿，但需要从纯视觉深度先验（而非 LiDAR 点云）引导高斯生成。系统将 DROID 深度与对齐后的单目深度融合为 proxy depth：低误差像素保留 DROID 深度，高误差像素用单目深度替代，从 proxy depth 采样生成高斯基元。考虑到热红外图像无色彩信息，GS 表示省略球谐而直接用灰度强度作为颜色 $c_i$；损失函数包含颜色 $\ell_1$、D-SSIM 和深度 $\ell_1$ 三项（权重 $\alpha{=}0.2$, $\beta{=}0.2$）。增密采用 AbsGS 策略（基于 2D 高斯位置的梯度绝对值），剪枝移除不透明度低于 0.05 的高斯及未在至少 3 帧中观测到的新增高斯。

\textbf{实验验证。}系统在 RRXIO（9 条序列，含室内运动捕捉真值和户外场景）和 VIVID（14 条序列，覆盖室内外及多种光照条件）两个热红外 SLAM 基准上评估。定位方面，TOM-GS 在 RRXIO 热红外序列的平均 ATE（0.35\,m）显著优于 DROID-SLAM（1.35\,m）和几何方法 ORB-SLAM3（多数序列跟踪失败），尤其在 mocap\_dark 场景中热红外 ATE（0.08\,m）甚至优于光照良好条件下的视觉 ATE——验证了热红外在弱光场景的天然优势以及深度增强里程计的有效性。在 VIVID 上，热红外里程计整体优于 RGB，在暗光和手电筒场景中差距尤为显著。渲染方面，Metric3D + Fieldscale 组合在 VIVID 上取得最优 PSNR（27.22\,dB），领先 Depth Anything + Fieldscale 组合（26.59\,dB）；GS 建图的渲染质量全面超越 NeRF-VO 和 GLORIE-SLAM，且 GS 天然对原始热图噪声的平滑效果使重建图像视觉更干净。

如果说前述方法的目标都是在给定传感器配置下最大化 SLAM 系统自身的精度与鲁棒性，那么 Sapling-NeRF（2026.02）~\cite{2602_22731} 则代表了另一种思路：将 NeRF 重建与 LiDAR SLAM 松耦合，不以端到端联合优化为目标，而是让各模块各司其职——为森林生态监测这一具体应用场景提供可重复、可地理定位的定量测量。这一设计选择的动机源于一个实践瓶颈：地面激光扫描（TLS）和移动激光扫描（MLS）的点云过于稀疏，无法重建树苗（sapling，0.5--2\,m 高）的细枝和叶片等精细结构，而纯视觉 NeRF/3DGS 方法又缺乏真实尺度且无法在不同季节间保持一致的坐标框架。

\textbf{三级松耦合表示。}与 TOM-GS 将跟踪与建图解耦为两个独立模块不同，Sapling-NeRF 构建了三级抽象表示并松耦合：(1) Level 3 为 2\,Hz GNSS 提供米级精度的地球框架定位；(2) Level 2 基于自研 LiDAR SLAM（VILENS-SLAM，LiDAR-惯性里程计 + 回环检测）构建公顷级森林点云地图，支持多会话（multi-session）离线合并——通过在位姿图层面融合不同日期的采集轨迹，使跨季节地图保持局部一致性；(3) Level 1 为每棵目标树苗训练独立的 NeRF 模型，实现物体级密集重建和新视角合成。三级之间的连接方式是松耦合而非端到端：LiDAR SLAM 提供的轨迹仅作为 NeRF 的初始位姿和尺度锚定，NeRF 的实际训练位姿由 COLMAP SfM 精化后通过 Umeyama 方法对齐到 SLAM 坐标系并恢复真实尺度。

\textbf{物体级采集策略与尺度恢复。}数据采集采用两种模式：割草式遍历整个林分以建立 Level 2 地图；到达目标树苗时切换为围绕树苗的"穹顶"模式（dome pattern），从各个方向密集拍摄图像。这一物体级采集策略与 Block-NeRF 等将大场景空间划分为子地图的方案根本不同——Sapling-NeRF 只在感兴趣区域（树苗）上投入 NeRF 的重建精度，背景由 LiDAR 地图承担。NeRF 选择 NeRF 而非 3DGS 的原因也出于应用考量：森林场景中需要重建大量树苗，NeRF 的内存效率优于 3DGS，尽管渲染速度更慢。

\textbf{骨架提取与叶木分离。}从 NeRF 导出的稠密点云出发，系统使用 PC-Skeletor 提取树苗分支骨架（基于 k-最近邻图的迭代收缩，保留拓扑一致性），进一步通过过骨架化（overskeletonization）产生的末梢分叉点定位叶片区域，将点云分离为叶片 $\mathbf{L}_{ij}$ 和木质 $\mathbf{W}_{ij}$ 两部分。由此可导出茎干高度、分支拓扑、叶片空间分布和叶木比（LWR $= N_l / N_w$）等生态指标。

\textbf{实验验证。}系统在英国 Wytham Woods（23 棵树苗，7/8/12 月三次采集）和芬兰 Evo 森林（3 棵针叶树苗）上评估。NeRF 重建的几何精度全面优于 TLS：TLS 无法重建 1\,cm 以下的细枝和叶片，而 NeRF 导出点云的茎干高度与 TLS 测量偏差仅 1--2\,cm，叶木比则显著更合理（TLS 因骨架化失败导致大量叶片被误分类为木质，LWR 普遍低于 0.3，而 NeRF 的 LWR 在 12--33 之间）。跨季节对比展示了系统的长期监测能力：Sapling 01 在 7 月与 8 月间仅有叶片分布的微小变化，而到 12 月叶片基本脱落；Sapling 02 在 7 至 8 月间出现了明显的枝条断裂（论文以红色箭头标注）。新视角合成在小树苗（$<$1\,m）上表现最佳（PSNR 最高 21.12\,dB），因为围绕小树苗更容易获得完整的全方位视角覆盖。

\begin{takeaway}
多模态 SLAM 通过融合 LiDAR、IMU、事件相机等互补传感器突破纯视觉 SLAM 的固有局限，LiDAR 提供精确几何约束，事件相机弥补运动模糊，GNSS 实现全局定位，三级松耦合架构已在森林生态监测等实际场景中验证了多模态融合的有效性。
\end{takeaway}

% ======================================================================

\subsection{近期方法脉络：多模态高斯建图、鲁棒跟踪、LiDAR-constrained Gaussian mapping等}

这一组近期工作延续 Gaussian Splatting mapping、3DGS-SLAM、3DGS，主要把改动集中在扩展、增量改进等环节。下面只保留相对既有路线的变化，便于和前文方法直接对照。

\paragraph{LIT-GS}
LIT-GS（2026.06）~\cite{2606_20424} 延续 Gaussian Splatting mapping 在“多模态高斯建图”上的基本设定，主要改变是：LIT-GS：LiDAR-Inertial-Thermal Gaussian Splatting，引入LiDAR平面几何约束、热-LiDAR关联和点到平面残差。3DGS 的显式高斯表示和可微渲染能力保持为基础。
其设计动机是：将热成像引入LIV Gaussian mapping，并用LiDAR平面几何显式约束pose/structure refinement和Gaussian优化。
具体而言，方法以LIV视觉地图点作为置信跨模态锚点建立thermal-LiDAR关联，并把加权LiDAR point-to-plane residual注入优化。
适用边界是：需要多传感器同步和标定；热-LiDAR关联失败会影响稳定性。

\paragraph{Spectral GS-SLAM}
Spectral GS-SLAM（2026.06）~\cite{2606_21258} 延续 3DGS-SLAM 在“鲁棒跟踪”上的基本设定，主要改变是：Spectral GS-SLAM：观测性/退化感知跟踪，融合ICP和互补特征约束并自适应补偿欠约束方向。3DGS 的显式高斯表示和可微渲染能力保持为基础。
其设计动机是：提出面向实时3DGS-SLAM的谱分析跟踪框架，在结构退化时识别并补偿不稳定方向。
具体而言，方法结合ICP和feature-based constraints，通过谱/观测性分析判断优化病态方向，并自适应调整约束权重。
适用边界是：仍依赖足够互补传感/特征；极端动态或重复纹理环境可能失败。

\paragraph{Structured-Li-GS}
Structured-Li-GS（2026.06）~\cite{2606_27509} 延续 3DGS 在“LiDAR-constrained Gaussian mapping”上的基本设定，主要改变是：基于 LiDAR-inertial-visual SLAM 的点云锚定 Gaussian 初始化和多损失训练。3DGS 的显式高斯表示和可微渲染能力保持为基础。
其设计动机是：把 LiDAR-inertial-visual SLAM 输出的准确稠密点云用于 Gaussian 锚定、局部几何初始化和训练损失约束。
具体而言，从子采样点云初始化高斯位置和椭球参数，并结合 photometric、flattening、offset、depth、normal 等损失实现轻量化结构化 3DGS。
适用边界是：依赖高质量 LiDAR/SLAM 点云；对低成本传感器、动态物体和大尺度在线更新的鲁棒性需要进一步验证。

\paragraph{DL-SLAM}
DL-SLAM（2026.07）~\cite{2607_01860} 延续 3DGS-SLAM 在“Probabilistic dynamic-object handling in Gaussian SLAM”上的基本设定，主要改变是：Dual-Level Probability：同时处理瞬时静态对象的位姿约束价值和静态地图污染风险。3DGS 的显式高斯表示和可微渲染能力保持为基础。
其设计动机是：指出 transiently static objects 可为位姿估计提供几何约束，但不应被集成进静态地图。
具体而言，结合 per-pixel uncertainty 和对象/区域级概率，既利用暂时静止对象增强 pose estimation，又避免它们在 Gaussian map 中留下持久 artifact。
适用边界是：依赖概率估计和对象运动判别；纯几何/语义歧义下仍可能误分类。
\section{视觉定位与重定位}
\label{sec:15_4}

视觉定位（Visual Localization）是在已知场景中估计查询图像的 6-DoF 位姿，是 AR/VR、自动驾驶和机器人导航的关键能力。3DGS 作为场景表示为视觉定位带来了新机遇：可微渲染使基于分析-合成（analysis-by-synthesis）的位姿优化成为可能，显式高斯基元支持特征场构建和 2D-3D 对应。本节梳理 29 篇论文，从基于渲染的位姿优化到开放词汇定位，覆盖这一方向的技术全貌。

\subsection{渲染驱动与特征场定位}
\label{ss:15_6}

本节介绍两条技术路线：（1）基于渲染的位姿优化；（2）基于特征场的 2D-3D 对应。

\subsubsection{基于渲染的位姿优化}

本路线介绍利用 3DGS 的可微渲染直接在像素级对齐渲染图像与查询图像进行位姿优化的方法。

利用 3DGS 的可微渲染进行位姿优化是 3DGS 定位最直觉的范式。与传统稀疏特征定位管线（SuperPoint 提取局部特征 + SuperGlue 匹配 + PnP RANSAC 求解）依赖 2D-3D 特征对应不同，基于渲染的方法直接在像素级对齐渲染图像与查询图像，天然适合低纹理、弱特征等稀疏方法失效的场景。GSLoc（2024.10）~\cite{2410_06165} 是这一范式的开创性工作，通过推导完整的位姿解析梯度、设计粗到细优化策略和扩展检索数据库三个核心贡献，证明了 3DGS 场景表示用于视觉定位的可行性。

\paragraph{问题建模与位姿梯度推导。}
GSLoc 的核心思路直观而优雅：给定已训练的 3DGS 场景模型 $\mathcal{G}$，对于查询图像 $\mathbf{I}$，通过最小化渲染图像与查询图像的光度差异来搜索 6-DoF 位姿 $\mathbf{T}_w^c \in \mathrm{SE}(3)$。目标函数为：

\begin{equation}
    \mathbf{T}_w^{c*} = \arg\min_{\mathbf{T}_w^c \in \mathrm{SE}(3)} \, L_1 \bigl( \hat{\mathbf{I}}(\mathbf{T}_w^c, \mathcal{G}), \, \mathbf{I} \bigr),
\end{equation}

其中 $\hat{\mathbf{I}}(\mathbf{T}_w^c, \mathcal{G})$ 是从位姿 $\mathbf{T}_w^c$ 渲染得到的图像。然而，3DGS 的高性能渲染基于 CUDA 光栅化，不支持自动微分，因此无法直接对相机位姿求梯度。GSLoc 解决这一难题的方式是手动推导相机位姿关于高斯参数的全部解析梯度。具体而言，作者将位姿参数化在 Riemannian 流形上——即位姿所在的特殊欧氏群 $\mathrm{SE}(3)$——并用对应的李代数 $\mathfrak{se}(3)$ 获得最小 6-DoF 表示。例如，旋转分量的梯度通过反对称矩阵（skew-symmetric matrix）$[\cdot]^\wedge$ 构造，平移分量则利用旋转矩阵 $\mathbf{R}_w^c$ 的列向量传播。这种流形参数化保证了优化始终在合法位姿空间中进行，相比常见的四元数-平移参数化（9 个参数带归一化约束），仅用 6 个自由度且无需归一化，消融实验证实其收敛性明显更优。求解采用 Adam 优化器，初始学习率 $10^{-2}$ 并指数衰减至 $10^{-5}$，最多迭代 2000 步或直到损失变化小于 $10^{-5}$ 连续 3 次时停止。

\paragraph{粗到细渲染调度。}
即使有了正确的梯度，光度 $L_1$ 损失在 6-DoF 位姿空间中仍具有高度非凸性——高频图像细节（纹理、边缘）在位姿偏移时产生剧烈的损失变化，容易将一阶优化方法困在局部最优。GSLoc 对此提出的解决方案直观而有效：如果在优化初期先把图像"模糊化"，低频结构信息会主导梯度方向，引导优化走出局部最优；随着优化推进，逐步恢复高频细节以精确定位。具体实现为对渲染图像和查询图像同时施加逐渐衰减的 2D 高斯模糊核 $N_{2d}(\boldsymbol{\sigma}_j)$，将目标函数修改为：

\begin{equation}
    \mathcal{L}(\mathbf{T}_w^c) = L_1 \bigl( N_{2d}(\boldsymbol{\sigma}_j) * \hat{\mathbf{I}}(\mathbf{T}_w^c, \mathcal{G}), \; N_{2d}(\boldsymbol{\sigma}_j) * \mathbf{I} \bigr),
\end{equation}

其中 $*$ 表示卷积，核的协方差 $\boldsymbol{\sigma}_j$ 在前 1000 步内从 $10^{-1}$ 线性衰减到 $10^{-4}$（固定核尺寸 $L=200$ 像素）。实际采用两步策略以兼顾效率：第一步运行标准位姿优化（无模糊），若渲染 PSNR 未达到 25 dB 的收敛阈值则判定失败，重启并启用粗到细策略。成功收敛的案例通常仅需 100--300 次迭代（约 5--15 秒），在 Replica 合成场景上以 5$^\circ$ 旋转/5 cm 平移为成功标准时，当初始位姿的 3D 交并比（IoU）达到 0.4 以上，定位成功率即可超越稀疏特征基线。

\paragraph{渲染增强的检索数据库扩展。}
基于渲染优化的定位高度依赖初始位姿的质量——初始位姿与真值越接近，优化越容易收敛到正确解。GSLoc 利用 3DGS 的新视角合成能力巧妙地缓解了这一"冷启动"问题：从有限训练图像出发，利用已训练的 3DGS 模型在任意新视角渲染逼真图像，将这些渲染图像连同其已知位姿加入 NetVLAD 图像检索数据库。扩展后的数据库提供更密集的场景覆盖，使查询图像能匹配到更近的初始位姿候选。实验显示，扩展后初始位姿的中位 3D IoU 从 0.3 提升至 0.4，GSLoc 总体定位成功率提高约 10\%，在低纹理和空旷场景中优于稀疏特征方法。在 Deep Blending 真实场景上同样验证了这一结论的有效性。

GSplatLoc（2024.12）~\cite{2412_20056} 沿用 GSLoc 的梯度优化范式，但做了三项关键改变：将优化信号从 RGB 光度差异替换为深度图差异，将位姿参数化从 $\mathfrak{se}(3)$ 李代数改为四元数-平移表示，并借助 gsplat 库的全可微渲染管线直接反向传播梯度，无需像 GSLoc 那样手动推导解析梯度。目标函数设计体现了深度优化的特点：除了标准的深度 $L_1$ 损失 $\mathcal{L}_d$ 外，额外引入等高线损失 $\mathcal{L}_c$——用 Sobel 算子提取渲染深度图与观测深度图的梯度（边缘），约束二者的结构对齐——两项以 $\lambda_1\!=\!0.8$、$\lambda_2\!=\!0.2$ 加权组合。等高线损失专注于深度不连续处的对齐，对绝对深度尺度的偏差不敏感，在纹理缺失但几何结构清晰的区域提供关键梯度信号。场景初始化阶段，GSplatLoc 对参考深度图生成的点云施加 PCA 对齐，将点云主轴与坐标轴对齐以统一优化起点，显著提高了不同数据集间优化损失的收敛稳定性。求解采用 Adam 优化器，旋转学习率 $5\!\times\!10^{-4}$、平移学习率 $10^{-3}$，通常约 100 次迭代即收敛——远快于 GSLoc 的数百至数千步。

代价是 GSplatLoc 需要 RGB-D 传感器提供观测深度图（GSLoc 仅需 RGB 图像），且在 TUM RGB-D 等真实场景中传感器噪声导致精度从 Replica 上的平均 0.016 cm 降至 0.81 cm。HGSLoc（2024.09）~\cite{2409_10925} 同样利用 3DGS 渲染进行位姿精化，但走了一条与 GSLoc 完全不同的技术路线：GSLoc 需要手动推导位姿对高斯参数的解析梯度并做数百至数千步迭代优化，HGSLoc 则完全绕过梯度计算，将位姿搜索建模为经典的图搜索问题。具体而言，每个候选位姿 $(q_w, q_x, q_y, q_z, t_x, t_y, t_z)$ 作为搜索图中的一个节点，通过设定旋转和平移变化量 $\delta_q$、$\delta_t$ 向邻居节点扩展，用类 A* 启发式搜索高效定位最优位姿。启发评估函数为 $f(n) = g(n) + h(n)$，其中实际代价 $g(n_{\text{child}}) = g(n_{\text{current}}) + 1$ 记录搜索深度，启发代价 $h(n) = \Sigma |I_q - \hat{I}_n|$ 即渲染图像与查询图像的像素级绝对差之和。这一设计使得 HGSLoc 无需对 3DGS 渲染管线做任何可微性改造，即可作为轻量级即插即用模块附加在任意粗位姿估计器（如 APR 的 Marepo、SCR 的 ACE）之后。

为兼顾不同误差量级，作者设置分级优化步长：大误差场景用较大步长快速逼近真值区域，微小误差场景切换为精细步长进一步提升精度。消融实验表明，像素级绝对差作为启发函数明显优于 PSNR 和 SSIM——以 Playroom 场景为例，前者精化后误差为 0.29 cm / 2.72$^\circ$，而 PSNR 和 SSIM 分别为 0.76 cm / 6.29$^\circ$ 和 0.87 cm / 6.79$^\circ$。在 7Scenes 上，ACE+HGSLoc 的平均位姿误差降至 1.0 cm / 0.29$^\circ$，误差在 1 cm / 1$^\circ$ 以内的比例从 ACE 单独的 53.7\% 提升至 59.2\%。对人工注入噪声的 Deep Blending 位姿，平移误差最高可降低 98\%（Boats 户外场景），旋转误差最高降低 99.7\%，展现出对噪声的强鲁棒性，优于同时期基于神经网络联合优化的策略。方法的局限在于每个节点扩展需约 300--400 次 3DGS 渲染，实时性受限。

GS-CPR（2024.08）~\cite{2408_11085} 同样利用 3DGS 的渲染能力进行位姿精化，但其技术路线与 GSLoc 的梯度优化范式有本质区别。GSLoc 需要手动推导位姿对高斯参数的解析梯度并迭代优化数百至数千步，这一过程虽然精确但对 3DGS 的渲染管线提出了严格的可微性要求。GS-CPR 则从另一个角度切入：\textbf{既然 3DGS 可以高质量地渲染图像和深度图，为什么不直接利用 3D 几何信息，将位姿精化转化为经典的 2D-3D 对应问题？}这一动机使得 GS-CPR 完全绕过了梯度反传播——不需要对 3DGS 模型或位姿估计器求导，仅通过一次前向渲染和匹配即可完成精化。

\paragraph{从 2D-2D 匹配到 2D-3D 对应。}
具体流程如下：给定预训练位姿估计器（如 APR 或 SCR 模型）输出的粗初始位姿 $\hat{\xi}$，3DGS 模型从该位姿渲染合成图像 $\hat{\mathbf{I}}_r$ 和深度图 $\hat{\mathbf{D}}_r$。随后，利用 3D 视觉基础模型 MASt3R 在查询图像 $\mathbf{I}_q$ 与渲染图像 $\hat{\mathbf{I}}_r$ 之间建立稠密 2D-2D 像素对应关系 $\mathcal{C}_{q,r}$。关键步骤在于：利用渲染深度图 $\hat{\mathbf{D}}_r$、相机内参 $\mathbf{K}$ 和已知位姿 $\hat{\xi}$，将 $\mathcal{C}_{q,r}$ 中渲染图像一侧的 2D 像素反投影为 3D 坐标，从而将 2D-2D 匹配升级为 2D-3D 对应。最终，将这些 2D-3D 对应关系输入经典 PnP 求解器配合 RANSAC 鲁棒估计，一次性得到精化位姿。整个过程无需训练任何场景特定的特征提取器或描述子——MASt3R 作为预训练基础模型，直接在 RGB 图像对上操作，在合成-真实（sim-to-real）域差异下展现出强鲁棒性。相比之下，GSLoc 仅利用 2D 光度信息（渲染图像与查询图像的 $L_1$ 差异）进行迭代优化，而 GS-CPR 通过深度渲染额外引入了 3D 几何约束，这是其仅用一次前向过程即超越 NeRF-based 方法多步迭代精化的根本原因。

\paragraph{曝光自适应渲染。}
户外场景中，建图序列与查询序列常因时间、天气和曝光参数不同而产生显著的外观差异，直接匹配渲染图像与查询图像会严重损害对应质量。GS-CPR 为此在 3DGS 渲染管线中引入曝光自适应仿射颜色变换（ACT）模块：一个 4 层 MLP 以查询图像的亮度直方图（YUV 颜色空间，亮度通道分为 10 个 bin）为输入，预测 $3 \times 3$ 仿射矩阵 $\mathbf{Q}$ 和 3 维偏置向量 $\mathbf{b}$，对渲染像素施加变换 $\hat{\mathbf{I}}(r) = \mathbf{Q} \cdot \hat{\mathbf{I}}_{\text{rend}}(r) + \mathbf{b}$，使渲染图像自适应地逼近查询图像的曝光条件。该模块在 Cambridge Landmarks 户外数据集上显著降低了平移和旋转误差。

\paragraph{快速变体与实验表现。}
除基于 2D-3D 对应的精化管线外，GS-CPR 还探索了一个速度优先的替代方案：利用 MASt3R 的点图注册能力直接估计查询图像与渲染图像间的相对位姿，通过将 MASt3R 预测的点图与 3DGS 渲染深度图对齐来恢复场景尺度，最终以 $\hat{\xi}' = [R_\Delta \hat{R} \mid R_\Delta \hat{\mathbf{t}} + s\mathbf{t}_\Delta]$ 计算精化位姿。该变体仅需 71 ms 额外开销（完整管线约 180 ms），为 APR 方法的时间敏感场景提供了精度-速度折中方案。在 7Scenes 和 12Scenes 室内数据集上，GS-CPR 与 APR/SCR 方法结合后均大幅降低位姿误差，在两个室内数据集上达到新的最优精度，其中 ACE + GS-CPR 甚至超越了传统 HLoc（SuperPoint + SuperGlue）管线，表明 3DGS 有潜力替代经典点云作为视觉定位的场景表示。LoGS（2024.10）~\cite{2410_11505} 在 GSLoc 的梯度优化位姿精化基础上，将问题从"给定 3DGS 模型做位姿优化"拓展为从建图到定位的完整管线，核心动机是应对少样本（few-shot）场景——训练图像极少时，随机初始化的高斯难以优化到理想位置，重建伪影直接干扰后续定位。建图阶段的两项设计正是为解决这一问题：首先用 SuperPoint/SuperGlue 进行 SfM 三角化，获取精确稀疏点云作为高斯初始化，使 3DGS 从一个合理的空间分布出发而非随机猜测；再用预训练 DPT（Dense Prediction Transformer）估计单目深度图，通过 Pearson 相关系数约束渲染深度与伪深度的分布一致性（而非逐像素对齐），弥补深度真值的缺失。针对少样本覆盖不全导致的过拟合，LoGS 在相邻训练视角间用 SLERP 插值生成伪视角，将深度正则化损失也施加于这些虚拟观测，以零额外采集成本增加场景覆盖。

定位阶段沿用了 GSLoc 的分析-合成范式——通过最小化渲染图像与查询图像的光度差异迭代优化位姿——但做了三项关键改进。第一，在位姿精化之前增加了基于 SfM 点云的经典定位管线（NetVLAD 检索 → 共视聚类 → SuperGlue 局部匹配 → PnP-RANSAC），为梯度优化提供一个比 GSLoc 的纯检索初始化更可靠的粗位姿起点。第二，GSLoc 对全部像素施加光度 $L_1$ 损失，LoGS 设计了三重掩码 $M = (M_{\text{grad}} \cap M_{\text{fea}}) \cap M_{\text{occ}}$ 精确筛选信息丰富像素：梯度掩码 $M_{\text{grad}}$ 用 Scharr 算子选取边缘以上的像素以捕捉结构信息；特征点掩码 $M_{\text{fea}}$ 以 SfM 阶段 SuperPoint 提取的关键点为中心取邻域，聚焦于角点等高判别性区域；不透明度掩码 $M_{\text{occ}}$ 过滤掉未覆盖高斯的空像素。这一掩码策略有效减少了重建伪影和无信息区域的干扰，使优化器免于陷入由噪声像素引起的局部最优。第三，位姿优化目标函数中引入仿射亮度参数 $a$、$b$ 联合求解，以 $\hat{\mathbf{I}}(p) = e^a \mathbf{C}(p) + b$ 调整渲染图像的曝光条件，应对户外场景中查询图像与建图序列间的光照差异。

在四个大规模数据集上，LoGS 达到当时最优位姿估计精度，且在训练图像数量极少的条件下仍保持鲁棒。上述方法均假设训练图像与查询图像间满足光度一致性，但 in-the-wild 图像集合（如 Phototourism 的旅游照片）中外观差异极大——同一场景在不同时间、天气、季节下拍摄，仅靠 LoGS 的仿射亮度参数无法消除这种根本性的外观不一致。R3GS（2025.05）~\cite{2505_15294} 的核心洞察是：\textbf{如果渲染模型能根据查询图像的外观自适应调整，光度一致性假设就不再需要场景外观恒定这一前提。}具体而言，R3GS 沿用 GSLoc 的光度 $L_1$ 损失优化位姿，但将固定 3DGS 模型替换为外观条件化的混合表示——每个锚点将 CNN 提取的全局外观特征与多分辨率哈希网格的局部空间特征拼接后，通过浅层 MLP 动态预测高斯属性。重定位时，先对查询图像提取外观特征并据此生成一组与查询外观匹配的高斯参数，再从初始位姿渲染并与查询图像对齐，使光度优化真正在"同类外观"下进行。

初始位姿通过 CosPlace 场景识别检索 3 个最近邻训练图像，以其位姿的反距离加权平均作为起点（有别于 GSLoc 的 NetVLAD 检索 + 扩展渲染数据库策略），随后固定高斯仅优化位姿，Adam 迭代 500 次即收敛。此外，R3GS 的重建阶段为重定位提供了两个关键支撑：微调轻量级 LRASPP-MobileNet 人体检测网络生成可见性掩码，泛化至海报、横幅、车辆等瞬态物体，消除其对渲染质量的干扰；天空区域以 Fibonacci 采样在大半径球面上初始化 20,000 个高斯，固定位置仅优化外观，通过深度正则约束防止前景高斯溢出至天空区域，有效抑制浮游物伪影。在 Phototourism 三个场景上，R3GS 的重定位 ATE 从 3DGS 基线的 0.401/0.381/1.377 m 降至 0.101/0.021/0.367 m，同时达到当时最优渲染保真度，且因哈希网格紧凑存储和外观嵌入后可导出标准 3DGS 点云，存储需求也显著降低。

\subsubsection{基于特征场的 2D-3D 对应}

本路线介绍将高维语义特征嵌入 3DGS 构建特征场，通过 2D-3D 特征对应实现精确定位的方法。

将高维特征嵌入 3DGS 构建特征场，是实现 2D-3D 对应和精确匹配的关键策略。直觉上，如果每个 3D 高斯不仅携带颜色，还携带一组描述其在场景中视觉身份的高维特征向量，那么查询图像中的 2D 特征就可以直接与 3D 高斯的特征进行匹配，从而建立 2D-3D 对应关系——这正是特征场定位的基本思路。

STDLoc（2025.03）~\cite{2503_19358} 是该路线的开创性工作，提出基于 Feature Gaussian 的从稀疏到稠密定位范式，打破传统先图像检索再特征匹配的粗到精流水线，建立了后续一系列方法共同依赖的技术基线。

\textbf{场景表示——Feature Gaussian。}STDLoc 继承 Feature 3DGS 的特征场蒸馏思路，为每个高斯基元 $\theta_i = \{x_i, q_i, s_i, \alpha_i, c_i, f_i\}$ 额外附加可训练特征属性 $f_i$。训练时，先用 SuperPoint 等通用特征提取器从训练图像中提取稠密特征图 $F_t(I)$ 作为真值，再通过 3DGS 的 $\alpha$-混合光栅化将特征属性渲染为特征图 $\hat{F}_s$，以 L1 距离优化特征场损失 $L_f = \|F_t(I) - \hat{F}_s\|_1$，与辐射场损失 $L_{rgb}$ 联合训练。这样得到的 Feature Gaussian 场景 $\mathcal{G}$ 既能渲染图像，又能渲染特征图，为定位提供了统一的 3D 场景表示。

\textbf{匹配导向采样——从百万高斯中筛选场景地标。}一个典型场景包含数百万高斯基元，直接与全部高斯匹配在计算上不可行，且大量模糊高斯会干扰匹配质量。STDLoc 的核心洞察是：好的定位地标需要同时满足两个条件——\emph{多视角可识别}（从不同视角看都能正确匹配）和\emph{空间均匀分布}（覆盖场景各个区域）。为此设计评分策略：对每个高斯 $g_i$，将其中心投影到所有可见训练图像上，提取对应 2D 特征与 $f_i$ 计算余弦相似度并取均值，得到匹配评分 $s(g_i)$——评分越高说明该高斯的特征越适合匹配。为避免纹理丰富区域高斯密集导致采样不均，先通过最远点采样选取空间锚点，再在每个锚点的 $k$ 个空间近邻中选取评分最高的高斯。最终仅需几千个高斯即可构成场景地标集 $\mathcal{L}$，在保持定位效果的同时大幅降低匹配开销。

\textbf{场景专用检测器——自适应发现查询关键点。}有了地标集 $\mathcal{L}$，还需要确定查询图像中哪些位置应该参与匹配。通用检测器（如 SuperPoint）提取的关键点是场景无关的，与 $\mathcal{L}$ 中地标的 2D 投影未必对齐。STDLoc 为此训练场景专用检测器 $D_\phi$：一个轻量 4 层 CNN，输入查询特征图 $F_t(I_q)$，输出地标概率热图 $\mathcal{H}$。训练信号以自监督方式获得——将 $\mathcal{L}$ 中每个高斯中心投影到训练图像上，对应像素标记为正样本，用二元交叉熵损失优化。推理时对热图做非极大值抑制（NMS）确保关键点均匀分布。该检测器使查询图像的关键点与场景地标天然对齐，显著提高匹配成功率。

\textbf{稀疏到稠密定位流水线。}有了上述三个组件，STDLoc 构建了完整的两阶段定位管线。\emph{稀疏阶段}：对查询图像提取特征图 $F_t(I_q)$，经场景专用检测器获得稀疏关键点，与 $\mathcal{L}$ 中所有地标特征计算余弦相似度进行最近邻匹配，建立 2D-3D 对应集 $M_{sparse}$，通过 PnP+RANSAC 求解初始位姿 $\xi_{sparse}$——整个过程无需任何位姿先验。\emph{稠密阶段}：以 $\xi_{sparse}$ 为初值，从 $\mathcal{G}$ 渲染稠密特征图和深度图，与查询特征图进行粗到精的稠密匹配（先在 $1/8$ 分辨率建立粗对应，再在原分辨率精化），通过渲染深度图将 2D 匹配提升为 2D-3D 对应，再次用 PnP+RANSAC 求解精化位姿 $\xi_{dense}$。稠密阶段支持迭代执行以进一步提升精度。

\textbf{实验表现。}在室内 7-Scenes 数据集上，STDLoc 相比最优 SCR 方法 ACE 平移和旋转误差均降低约 30\%，在 5cm/5$^\circ$ 阈值下召回率达 99.1\%，2cm/2$^\circ$ 阈值下达 90.9\%。在室外 Cambridge Landmarks 数据集上同样超越所有对比方法，包括传统结构化方法和 NeRF-based 方法。每个场景的训练时间约 90 分钟（Feature Gaussian）加 30 分钟（检测器），在 RTX 4090 上完成。

上述技术框架奠定了本小节后续方法的研究基线：STDLoc 的 Feature Gaussian 表示中颜色场与特征场紧耦合，引发了 LiteLoc 的去颜色解耦探索；$\alpha$-混合光栅化的特征渲染机制，成为 ULF-Loc 分析特征偏差的理论出发点；匹配导向的地标选取策略，启发了 SplatLoc 的显著性地标选择和 SplitGS-Loc 的一对一对应优化。

LiteLoc（2026.05）~\cite{2605_17777} 完整继承 STDLoc 的稀疏到稠密管线（匹配导向采样、场景专用检测器、两阶段求解均不做修改），将改进聚焦于 STDLoc 未触及的两个效率瓶颈——场景表示的存储冗余与稠密求解的计算冗余。

\textbf{瓶颈一——颜色-特征紧耦合引发存储膨胀。}STDLoc 的 Feature Gaussian 沿用 Feature 3DGS 的紧耦合表示 $\theta_i = \{x_i, q_i, s_i, \alpha_i, c_i, f_i\}$，颜色属性 $c_i$（球谐系数）与定位特征 $f_i$ 共享同一组高斯基元。LiteLoc 指出这种耦合存在两个问题。其一，\emph{颜色场对定位功能无用}——位姿估计依赖的是 2D-3D 匹配的几何约束，而非光度外观。其二，更关键的是\emph{频率失配}：重建高频光度细节（边缘、纹理）需要极密的高斯基元分布，但定位特征场本质上是低频的——这迫使大量仅服务于颜色精度的冗余高斯各自携带高维特征嵌入，造成巨大存储开销（典型场景约 759k 高斯、930MB），且几何优化被高频颜色约束牵制，特征场的收敛质量也受影响。

\textbf{去颜色解耦特征场。}基于上述分析，LiteLoc 构建解耦特征场，将高斯基元简化为 $\theta'_i = \{x_i, q_i, s_i, \alpha_i, f_i\}$，显式丢弃球谐系数 $c_i$。训练时仅保留特征场损失 $L_f = \|F_{gt} - \hat{F}\|_1$，不再联合优化辐射场损失。为保证仅用特征监督时渲染的稳定性，LiteLoc 将 STDLoc 的标准 $\alpha$-混合替换为双层 $L_2$ 归一化版本：先对每个高斯的特征向量做 $L_2$ 归一化后参与 $\alpha$-混合加权，再对输出做 $L_2$ 归一化，即 $\hat{F} = \text{norm}\bigl(\sum_{i \in \mathcal{N}} \text{norm}(f_i)\,\alpha_i \prod_{j=1}^{i-1}(1-\alpha_j)\bigr)$。这一改动使高斯密度完全由特征场的低频结构驱动，而非被高频颜色细节撑大。效果立竿见影：场景高斯数从 759k 降至 57k，存储从 930MB 降至 59MB——减少约 94\%——且定位精度反而从 99.1\%/90.9\% 提升至 99.6\%/92.4\%（7-Scenes 上 5cm-5$^\circ$/2cm-2$^\circ$），说明解除颜色约束后特征场获得了更好的优化空间。

\textbf{瓶颈二——稠密 PnP 的计算冗余。}STDLoc 的稠密阶段通过粗到精匹配建立上万条稠密对应 $|\mathcal{M}_{dense}| > 10^4$，全部送入 PnP+RANSAC 求解。LiteLoc 分析发现这些匹配在空间上高度相关——多数匹配提供的是\emph{饱和的}几何约束（对同一局部区域的重复约束），精度增益已严重递减，但计算开销随匹配数线性增长，稠密 PnP 单次求解耗时约 96ms，成为推理管线的主导瓶颈。

\textbf{稠密匹配蒸馏。}为此 LiteLoc 提出匹配蒸馏策略：将稠密对应嵌入由查询视角和渲染视角 2D 坐标构成的联合 4D 几何空间 $(u, \bar{u}) \in \mathbb{R}^4$，用 K-means（FAISS 实现，CPU 端执行，最多 5 次迭代）将其划分为 $K$ 个区域，从每个区域选取距聚类中心最近的匹配作为代表代理 $\mathcal{M}_{cluster}^k = (u_k, \bar{u}_k)$。代理点通过渲染深度图反投影至 3D 点 $P_k$，最终的位姿优化仅在约 $K = 1024 \approx 0.05|\mathcal{M}|$ 个代表匹配上执行（即仅保留 5\%）。聚类额外开销约 9ms，但稠密 PnP 耗时从 96ms 降至 5ms——加速近 19 倍。单独使用蒸馏（不启用解耦），定位精度几乎无损失（99.1\%/90.9\% $\to$ 99.1\%/90.7\%）。

\textbf{两项改进的协同效果。}完整 LiteLoc 结合两项改进后，在 7-Scenes 上平均平移/旋转误差从 STDLoc 的 0.76cm/0.24$^\circ$ 降至 0.73cm/0.22$^\circ$，存储仅 59MB（减少 94\%），单次稠密迭代总耗时从 230ms 降至 175ms，四次迭代从 613ms 降至 287ms（减少 53\%）。在室外 Cambridge Landmarks 上，平均平移误差从 STDLoc 的 10.1cm 降至 8.9cm，旋转误差从 0.14$^\circ$ 变为 0.19$^\circ$（精度基本持平），表明效率导向的重新设计在室外场景同样有效。

SplatLoc（2024.09）~\cite{2409_14067} 面向 AR 应用场景，在 STDLoc 的 Feature Gaussian 框架基础上重新审视了特征学习和地标选取两个核心环节。STDLoc 通过 $\alpha$-混合光栅化将高斯特征渲染为 2D 特征图并与真值对齐，这一机制虽然有效，却引入了描述符偏差——$\alpha$-混合将单个高斯的特征与邻近高斯纠缠，优化后的特征并不严格对应该高斯本身的视觉身份。SplatLoc 从这一观察出发，提出三项针对 STDLoc 的增量改进。

\textbf{无偏描述符学习——绕过 $\alpha$-混合的特征体蒸馏。}回顾 STDLoc 为每个高斯直接存储特征属性 $f_i$ 并通过 $\alpha$-混合渲染优化。SplatLoc 完全放弃这一路线，转而将多视角 2D 特征图投影融合到 3D 特征体 $\mathcal{V}$ 中，再用场景专用的 3D 解码器（多分辨率哈希编码 + 浅层 MLP）从 3D 位置直接回归高维描述符——不经过任何 $\alpha$-混合，因此是\emph{无偏}的。蒸馏仅在场景表面流形上执行（通过 Marching Cubes 提取表面后采样），避免在空洞区域学习无效描述符。这一设计使描述符学习与高斯基元参数解耦，模型存储仅需 112MB，而同样基于点式特征存储的 PNeRFLoc 需要 788MB——减少约 7 倍。

\textbf{显著性地标选择——三维评分替代匹配评分。}STDLoc 用多视角匹配评分 $s(g_i)$ 评价高斯质量，再通过最远点采样保证空间覆盖。SplatLoc 重新定义了更全面的三维显著性评分 $\mathcal{H}(p)$，融合三项指标：\emph{显著性}（可学习的地标概率分数）、\emph{泛化性}（训练过程中被观察到的最大视角张角）和\emph{几何一致性}（高斯中心到多视角表面点的距离误差）。基于该评分，采用贪婪搜索算法迭代选取最高分且与已选地标保持最小距离的高斯，确保地标在场景中均匀分布——约 500 个地标即可维持定位精度，在 Replica 室内数据集上中位平移/旋转误差均低于 1.6cm/0.6$^\circ$。

\textbf{关键高斯正则化——抑制光度优化引起的几何漂移。}STDLoc 未约束高斯基元在光度优化过程中的几何变化。SplatLoc 发现，由 SuperPoint 关键点初始化的\emph{关键高斯}在光度损失驱动下会发生中心漂移和各向异性膨胀，导致其 3D 位置偏离真实表面，直接影响基于这些高斯的定位精度。为此引入位置冻结（阻止关键高斯中心更新）和尺度正则化（惩罚缩放参数偏离 isotropic 阈值），两项正则化在四个 Replica 场景上平均降低平移误差 1.57cm、旋转误差 0.32$^\circ$。GSplatLoc（2024.09）~\cite{2409_16502} 将 XFeat 轻量特征提取器的稠密关键点描述符集成到 3DGS 中。

SplitGS-Loc（2026.05）~\cite{2605_07351} 沿用 STDLoc 的"特征高斯场 $\to$ 2D-3D 匹配 $\to$ PnP"定位范式，但从一个被忽视的结构性矛盾出发重新审视了这条路线：光度优化产生的高斯具有空间体积，多个像素可映射到同一个高斯中心，而 PnP 求解器要求严格的点对点对应——这种多对一映射直接导致重投影误差累积、RANSAC 难以达成共识。为解决该矛盾，SplitGS-Loc 对 STDLoc 的特征获取、高斯采样和场景表示三个环节均做了针对性改进。

\textbf{直接特征提升——替换渲染优化。}STDLoc 以渲染损失 $\|F_t - \hat{F}_s\|_1$ 迭代优化每个高斯的特征属性 $f_i$，梯度沿射线分散到数百个高斯，削弱了像素-高斯间的对应精度。SplitGS-Loc 放弃渲染优化，转而利用光栅化的组合权重（即 $\alpha$-混合中每个高斯对每个像素的贡献份额）直接聚合像素特征：对每个高斯 $g_n$，收集其在所有可见训练图像中贡献权重超过阈值 $\tau$（例如 0.1，即贡献超过 10\%）的像素，以 softmax 归一化的权重对这些像素的 SuperPoint 特征加权求和，得到聚合特征 $\bar{z}_n$。组合权重天然编码了遮挡与可见性，因此无需额外训练即可在不到一分钟内完成整个特征场构建。

\textbf{权重评分采样——替换余弦相似度。}在高斯采样环节，SplitGS-Loc 不沿用 STDLoc 的余弦相似度匹配评分，而是以每个高斯在多视角中的平均组合权重 $\psi(g_n) = \frac{1}{|\mathcal{W}_n|}\sum_{\tilde{w}\in\mathcal{W}_n}\tilde{w}$ 作为重要性评分。该评分偏好那些跨视角持续且显著地参与光栅化的高斯，自然筛选出多视角一致的场景地标。采样仍沿用 STDLoc 的锚点-近邻分区结构（随机锚点 + $k$-近邻分组），但评分标准的替换使筛选更直接地反映高斯在新视角下的可靠性。

\textbf{混合高斯分裂——消除多对一映射。}上述两个改进解决了特征获取与地标筛选问题，但高斯体积引起的多对一映射仍未消除。SplitGS-Loc 的核心设计是将每个高斯沿其主轴（最长轴，即匹配歧义最大的方向）分解为三个子高斯：中心子高斯位于原中心、两个侧子高斯沿主轴对称偏移 $\pm\beta s$。分裂参数由矩匹配严格推导——要求子高斯混合分布的前四阶矩与原始高斯分布一致——由此确定中心子高斯权重 $\lambda_\text{o}=\frac{2}{3}$、侧子高斯权重 $\lambda=\frac{1}{6}$、子高斯方差 $\sigma^2 = (1-\frac{\beta^2}{3})s^2$（$\beta \in (0, \sqrt{3})$ 为可调偏移因子，实验取 $\beta=1.4$）。一个关键的设计细节是采样阶段的特殊处理：若直接对子高斯评分采样，同一父高斯的多个子高斯会在同一局部区域内互相淘汰，破坏分裂带来的一对一对应；因此采样时先将子高斯分数聚合回父高斯 $\bar{\psi}(g_n) = \frac{1}{3}(\psi(g_{n_-})+\psi(g_{n_o})+\psi(g_{n_+}))$，按父高斯中心做锚点采样，选中的父高斯的所有子高斯再全部保留，在特征注册阶段各自独立配对不同像素特征。如此，原本由单个高斯覆盖的显著区域由多个细粒度高斯分别描述，多对一映射转化为精确的一对一对应。

\textbf{实验表现。}在 Cambridge Landmarks 上，SplitGS-Loc 在最大场景 GreatCourt 上表现尤为突出，无需逐场景训练或位姿精化即达到最优。消融实验清晰展示了每个改进的独立贡献：以 STDLoc 的 13.1cm/0.22$^\circ$ 为基线，仅替换特征场构建（直接提升）即降至 9.6cm/0.17$^\circ$，加上权重采样进一步降至 8.7cm/0.15$^\circ$。PnP 收敛分析定量验证了分裂策略的效果：多对一匹配输入从 313 降至 199，RANSAC 内点数增加，位姿更精确，收敛时间从 314ms 缩短至 215ms。完整管线在预训练 3DGS 上的构建时间不到一分钟，分裂步骤处理一百万高斯不到一秒。

ULF-Loc（2026.05）~\cite{2605_04730} 沿用 STDLoc 的粗到精定位管线框架（稀疏匹配求初始位姿、稠密匹配精化），但从理论层面揭示了 STDLoc 场景表示的一个根本缺陷，并基于此重新设计了特征构建和地标选取两个环节。

\textbf{核心发现——$\alpha$-混合优化引入固有特征偏差。}回顾 STDLoc 的 Feature Gaussian 训练：通过 $\alpha$-混合光栅化渲染特征图，再以 L1 损失反传优化每个高斯的特征向量 $f_i$。ULF-Loc 对这一过程做了严格的偏差分析。对目标高斯 $g_t$ 在视角 $k$ 中像素 $u_k$ 处的渲染特征进行分解：$F_s(u_k) = w_k f_t + (1 - w_k) B_k$，其中 $w_k = \alpha_t T_t$ 是目标高斯的累积权重，$B_k$ 是其余高斯的归一化背景特征。假设 2D 观测特征 $f_k^{2D} = \mu + \epsilon_k$（$\mu$ 为真实 3D 特征，$\epsilon_k$ 为视角相关噪声），ULF-Loc 推导出优化后特征的期望偏差为 $bias = \mathbb{E}[(1-w_k)/w_k \cdot (\mu - B_k)]$。该公式表明，无偏估计需要满足两个严格条件——目标高斯完全主导渲染（$w_k = 1$）或背景特征恰好等于目标特征（$B_k = \mu$）——但实际中部分遮挡和视角变化使两者均不成立。因此偏差是 Feature Gaussian 表示的固有属性，而非可通过更多训练迭代消除的收敛问题。

\textbf{几何加权特征融合——绕过偏差的替代方案。}基于上述分析，ULF-Loc 放弃 STDLoc 的特征场优化范式（即不再为每个高斯训练可微特征属性），改为直接从多视角 2D 特征聚合构造无偏 3D 特征。具体做法：对每个地标高斯 $\tilde{g}_i$，从其可见训练视角中提取 2D 特征 $f_k^{2D}$，以归一化几何权重加权求和得到融合特征 $f^{fus} = \sum_{k=1}^{K} w_{i,k} f_k^{2D}$（$\sum_k w_{i,k} = 1$）。权重定义为 $w_{i,k} = \bm{n}_i \cdot \bm{d}_{i,k}$，即地标表面法向量 $\bm{n}_i$（取最小尺度轴方向，消歧至 $\bm{n}_i \cdot \bm{d}_{i,k} > 0$）与视线方向 $\bm{d}_{i,k} = (C_k - \bm{\mu}_i)/\|C_k - \bm{\mu}_i\|$ 的点积——正对视角权重高，掠射视角权重低，因为正对视角捕获的 2D 特征更稳定、更具代表性。这一线性融合保证了 $\mathbb{E}[f^{fus}] = \mu$（理论无偏），从根源规避了 $\alpha$-混合偏差。更重要的是，它完全不需要为高斯训练特征属性，因此也无需特征场蒸馏的额外训练开销——这是 ULF-Loc 效率大幅提升的直接原因。

\textbf{关键点共识地标采样——STDLoc 匹配评分的替代。}STDLoc 通过余弦相似度评分加最远点采样选取地标，依赖已经训练好的 Feature Gaussian 特征。ULF-Loc 不再训练特征场，因此设计了不依赖 3D 特征的地标选取策略：对每个高斯 $g_i$，将其中心投影到所有训练视角，检查投影位置是否与 2D 关键点检测器的输出对齐（距离小于阈值 $\tau_D$），统计对齐次数作为共识评分 $\mathcal{S}^i$——在多个视角中都投影到关键点附近的高斯，说明其几何位置稳定且视觉上具有区分性。再用随机 $k$-近邻采样确保空间均匀性，最终选取约 20,000 个地标。与 STDLoc 的评分策略相比，共识采样无需预先训练特征场，直接从 2D 关键点几何一致性出发，因此可以在 3DGS 训练完成后立即执行，不引入额外训练阶段。

\textbf{局部几何一致性验证——拒绝渲染伪影导致的误匹配。}在稠密精化阶段，3DGS 渲染图像的模糊和浮游伪影会导致稠密匹配产生大量误匹配。ULF-Loc 提出基于拓扑结构的误匹配剔除算法 LGCV：对每个候选匹配 $(x_i, y_i)$，在其 8-近邻匹配集中构建三角形对 $T_i = \{(\triangle_{x_i x_j x_k}, \triangle_{y_i y_j y_k}) \mid j, k \in \mathcal{N}_8(i)\}$，基于局部刚性假设检查两个约束——角度一致性（对应角的余弦差 $|\cos\theta_x - \cos\theta_y| < 1 - \tau_a$）和尺度一致性（对应边长比差异 $|s_a - s_b| < \tau_s$，其中 $s = l_x / l_y$）。统计每个匹配周围满足两个约束的三角形数量作为支持分数，支持不足的匹配被剔除。这一机制有效缓解了渲染伪影对位姿精化的干扰。

\textbf{实验表现。}在 7-Scenes 上，ULF-Loc 取得平均中位误差 0.7cm/0.2$^\circ$，相比顶级 SCR 方法 ACE 平移误差降低约 36\%；1cm/1$^\circ$ 召回率达 75\%，分别超越 ACE+GS-CPR 和 STDLoc 约 9\% 和 2.2\%。在 12Scenes 上取得平均中位误差 0.3cm/0.15$^\circ$，1cm/1$^\circ$ 召回率 94.1\%。在 Cambridge Landmarks 上，平均中位平移误差相比 STDLoc 降低 17\%，在 Court 等大尺度外观变化场景中仍保持鲁棒。效率方面，由于无需训练特征场，训练时间仅为 STDLoc 的 1/10，GPU 内存为 1/6。

\subsection{大规模与特殊场景}
\label{ss:15_7}

本节介绍两条技术路线：（1）大规模与远程感知定位；（2）语义引导与特殊场景定位。

\subsubsection{大规模与远程感知定位}

本路线介绍将 3DGS 定位扩展到无人机航拍等大规模远程感知场景的方法，解决尺度估计和渲染伪影等新挑战。

将 3DGS 定位扩展到大规模和远程感知场景需要解决尺度估计、渲染伪影和计算效率等新挑战。GSLoc 在室内和小规模场景中证明了光度优化范式的有效性，但其 NetVLAD 检索仅提供粗初始位姿、6-DoF 联合优化在大基线场景下容易陷入局部最优，且缺少对 3DGS 重建伪影的防护机制——这些问题在无人机航拍的大尺度、高空俯视场景中被显著放大。LSGS-Loc（2026.04）~\cite{2604_05402} 沿用 GSLoc 的光度优化框架（$L_1$ 光度损失 + Adam 迭代精化），但在初始化和精化两个阶段都针对大规模场景做了关键改进。

\textbf{尺度感知的位姿初始化。}GSLoc 用 NetVLAD 检索最近参考图像获取粗位姿，无法解决相对位姿的尺度模糊问题。LSGS-Loc 的初始化分为两步。首先，检索阶段用基础模型 AnyLoc（替代 NetVLAD）提取全局描述子，并在参考数据库中插入 3DGS 渲染的合成视角以加密视点覆盖。然后，对于检索到的参考图像，用 ReLoc3r 估计查询与参考之间的相对旋转 $\Delta R$ 和平移方向 $\Delta \mathbf{t}$——但相对平移仅提供方向而丢失尺度信息。为恢复度量尺度，LSGS-Loc 利用 ReLoc3r 内部的跨图像注意力图定位查询图像中心块在渲染参考图像中的对应位置，再通过渲染深度图将该 2D 对应反投影为 3D 锚点 $P$，从而为平移搜索提供一个确定的几何锚定。在此基础上，以锚点 $P$ 为起点沿平移方向做一维光度搜索（在 $[0.2 d_P, 2.0 d_P]$ 区间内采样 $N_s = 30$ 个候选位置），对每个候选位置渲染图像并计算 $L_1$ 损失，选择同时最小化光度损失及其离散梯度的最优位置作为初始平移。这一设计将 GSLoc 的 6-DoF 联合搜索简化为 1D 搜索 + 旋转传播，大幅降低了大规模场景中初始化失败的风险。

\textbf{拉普拉斯可靠性掩码。}在大规模无人机场景中，3DGS 重建不可避免地产生模糊区域和浮游物（floaters）等伪影。GSLoc 的光度优化对所有像素一视同仁，伪影区域的错误梯度信号会拖偏优化方向。LSGS-Loc 引入基于拉普拉斯算子的自适应掩码机制：在首次迭代时将渲染图像划分为 $M$ 个非重叠图像块，计算每个块的拉普拉斯幅值得分 $s_j$（拉普拉斯响应高的区域对应清晰纹理，低响应区域则对应模糊或伪影），通过阈值化生成二值掩码 $m(\mathbf{u})$，后续迭代中仅对掩码内的像素计算光度损失，将优化聚焦于高保真结构区域。精化阶段采用四元数参数化旋转、$\mathfrak{se}(3)$ 李代数更新增量，Adam 优化器初始学习率 0.018 配合余弦退火，上限 200 次迭代。

在 GauUscene 五个大规模无人机场景上的实验表明，LSGS-Loc 的中位平移误差在厘米级（如 HAV 场景 2.71 cm），显著优于 ACE 等场景坐标回归方法（误差常超过 50 cm）和 GS-CPR 等 3DGS-based 方法，且即使降采样至 1600px 宽度仍保持更高精度，展现出对输入变化的鲁棒性。消融实验证实，光度引导搜索阶段将 ($2^\circ$, 20 m) 召回率提升至 97\% 以上，（$1^\circ$, 10 m）召回率平均提升超 50 个百分点，拉普拉斯掩码进一步稳定了精化收敛。Hi$^2$-GSLoc（2025.07）~\cite{2507_15683} 提出双层次重定位框架，采用稀疏到稠密和粗到精的范式，充分利用高斯基元中丰富的语义信息和几何约束。稀疏阶段包含高斯特有的一致渲染感知采样策略和地标引导检测器实现鲁棒的初始位姿估计；稠密阶段通过粗到精的稠密光栅化匹配迭代精化位姿并引入可靠性验证过滤不可靠估计。通过分区高斯训练、GPU 加速并行匹配和动态内存管理策略处理大规模遥感场景。3DGS-LSR（2025.07）~\cite{2507_05661} 面向自动驾驶大规模重定位，结合多传感器数据构建高精度户外 3DGS 地图，客户端仅需单目 RGB 输入。使用 SuperPoint 和 SuperGlue 进行特征提取和匹配，核心创新是迭代优化策略通过逐步渲染精化定位结果，适用于实时自主导航。在 KITTI 数据集上城镇道路、林荫道和密集高速路段的平均定位精度分别为 0.026 m、0.029 m 和 0.081 m，显著超越其他代表性方法。

GSFeatLoc（2025.04）~\cite{2504_20379} 与上述方法形成鲜明对比：LSGS-Loc、Hi$^2$-GSLoc 等 3DGS 定位方法都基于光度优化范式——迭代渲染图像、计算光度损失、反向传播更新位姿——通常需要数十至上百次渲染迭代，单帧推理超过 10 秒。GSFeatLoc 将这一范式替换为特征对应范式：仅从初始位姿渲染一次 RGBD 图像，用 SuperPoint 提取查询图像与渲染图像的特征点、SuperGlue 建立 2D-2D 匹配，再利用渲染深度图将渲染图像侧的 2D 匹配点反投影为 3D 坐标，将 2D-2D 对应提升为 2D-3D 对应，最终通过 PnP + RANSAC 一次性求解位姿。整个过程无需迭代渲染、无需梯度优化、无需训练场景特定的特征提取器——SuperPoint/SuperGlue 作为预训练模型直接操作。在 Synthetic NeRF、Mip-NeRF 360 和 Tanks and Temples 三个数据集（38 个场景、2700+ 测试图像）上，GSFeatLoc 的推理时间稳定在 0.09--0.15 秒/帧（比 iComMa 快两个数量级），旋转和位置误差均低于 iComMa（如 Synthetic NeRF 上 RE $1.61^\circ$ vs.\ $13.29^\circ$），且对初始位姿误差的容忍度高达 $55^\circ$ 旋转和 1.1 倍场景尺度平移。方法的局限在于严重依赖渲染深度图的精度和特征匹配的数量——当匹配点不足时，PnP 无法收敛，只能退回初始位姿作为输出。GSVisLoc（2025.08）~\cite{2508_18242} 无需修改或重训练 3DGS 模型即可实现泛化定位。

\subsubsection{语义引导与特殊场景定位}

本路线介绍利用语义信息辅助定位以及面向无人机等特殊场景的 3DGS 定位方法。SGLoc（2025.07）~\cite{2507_12027} 利用 2D 图像和 3DGS 地图之间的语义关系，通过多级位姿回归策略在无初始位姿先验的条件下实现全局定位。CityLoc（2025.01）~\cite{2501_08982} 用扩散模型生成文本描述条件下的 6-DoF 相机位姿分布，利用 3DGS 渲染进行视觉推理修正。SplatHLoc（2026.03）~\cite{2603_29185} 提出自适应视点检索方法合成与查询更对齐的虚拟候选图像。PoseCompass（2026.05）~\cite{2605_12144} 设计智能合成位姿选择管线，减少自适应时间 3 倍。

在无人机定位方向，Video2BEV（2024.11）~\cite{2411_13610} 利用 3DGS 重建场景并投影为鸟瞰图，实现跨视角地理定位。Wrivinder（2026.02）~\cite{2602_14929} 结合 SfM、3DGS 和语义基础实现零样本地面到卫星图像对齐。Satellite-Free Training（2026.04）~\cite{2604_01581} 通过 3DGS 重建将无人机图像转换为正射影像，无需卫星图像训练即可实现跨视角检索。

PoI（2025.02）~\cite{2502_04843} 框架通过渐进式像素级过滤策略筛选可靠的合成像素，提升场景坐标回归定位精度。Ress et al.（2025.02）~\cite{2502_13803} 在大型工业环境中利用 3DGS 渲染的额外视角显著提升定位性能。GSPR（2024.10）~\cite{2410_00299} 将多视角 RGB 和 LiDAR 点云融合为 3DGS 场景表示，通过 3D 图卷积和 Transformer 提取时空特征用于地点识别。Pietrantoni et al.（2025.07）~\cite{2507_23569} 提出隐私保护的 3DGS 特征场定位，通过 3D 结构引导聚类将特征转换为分割。Li et al.（2024.12）~\cite{2412_00138} 通过对抗训练利用 3DGS 生成外观多样化训练数据，显著提升绝对位姿回归的鲁棒性。

\begin{takeaway}
视觉定位利用 3DGS 的可微渲染实现分析-合成式位姿优化，语义引导和合成视角增强策略显著提升了无初始位姿先验条件下的全局定位能力。
\end{takeaway}

% ======================================================================


\section{语义SLAM与建图}
\label{sec:15_5}

语义 SLAM 在稠密重建的同时为地图赋予语义理解能力——不仅知道场景"长什么样"，还知道"是什么"。3DGS 的显式高斯基元天然适合作为语义载体：每个高斯可以关联语义特征或标签，渲染时即可产生稠密的语义分割图。本节梳理 22 篇论文，覆盖从封闭集到开放词汇语义、从 2D 提升到 3D 的完整技术栈。

\subsection{语义提升与开放词汇}
\label{ss:15_8}

本节介绍两条技术路线：（1）从 2D 语义到 3D 高斯的提升；（2）开放词汇与层次化语义。

\subsubsection{从 2D 语义到 3D 高斯的提升}

本路线介绍将 2D 图像语义信息提升到 3D 高斯空间的方法，实现无需人工标注的全景语义重建。

将 2D 图像中的语义信息可靠地提升到 3D 空间，是语义 SLAM 的核心技术难题。之前的语义 SLAM 系统存在两个根本局限：一是只做语义分割或实例分割之一，无法同时恢复场景的完整"全景"（panoptic）信息；二是依赖密集人工标注，无法适用于开放世界环境。PanoSLAM（2025.01）~\cite{2501_00352} 首次在统一 SLAM 框架内同时实现几何重建、3D 语义分割和 3D 实例分割——全景 3D 重建无需任何人工标注，仅从无标签 RGB-D 视频即可完成。

\paragraph{语义高斯扩展。}
PanoSLAM 以 SplaTAM 为基线，核心改动是将每个高斯的参数从 8 维扩展到 13 维。直觉上，如果 RGB 颜色可以通过 $\alpha$-混合渲染到 2D，那么语义信息也可以走同一条路——只需为每个高斯额外赋予"语义属性"。具体而言，新增语义嵌入 $\mathbf{s} \in \mathbb{R}^3$、语义半径 $\rho_s$ 和语义不透明度 $\alpha_s$，使语义信息通过与 RGB/深度完全相同的 $\alpha$-混合公式进行可微渲染：

\begin{equation*}
S(\mathbf{P}) = \sum_{i=1}^{n} \mathbf{s}_i \, \alpha_i^s(\mathbf{P}) \prod_{j=1}^{i-1}(1 - \alpha_j^s(\mathbf{P})),
\end{equation*}

其中 $\alpha_i^s(\mathbf{P})$ 用语义半径 $\rho_s$ 和语义不透明度 $\alpha_s$ 计算，与几何渲染中的 $\alpha$ 使用完全相同的泼溅公式。初始化时，语义半径与几何半径取相同值 $\rho_s = r = d_{\mathrm{GT}}/f$（像素深度除以焦距），语义嵌入初始化为对应像素颜色，两种不透明度均初始化为 0.5。致密化条件也做了联合设计：同时检查几何轮廓图（$\alpha$ 渲染的可见性）和语义轮廓图（$\alpha^s$ 渲染的可见性）是否低于 0.5，加上深度误差阈值判定，三者任一触发就新增高斯。

\paragraph{空间-时序提升（STL）模块。}
系统使用冻结的 2D 视觉基础模型 SEEM 对每帧生成全景伪标签（包含实例掩码和语义类别）。然而这些伪标签存在严重的跨视角不一致——同一物体在不同视角可能被分配不同实例标签或语义类别。直接用这些噪声标签优化会严重降低语义地图质量。PanoSLAM 提出在线\textbf{空间-时序提升}（Spatial-Temporal Lifting, STL）模块来解决这一问题。核心思路简洁而有效：将所有帧的 2D 像素按深度反投影到 3D 世界坐标，然后对 3D 空间做均匀体素化量化。落入同一体素的像素被视为观测到同一 3D 位置，其标签在空间维度取平均统一。形式化地，对体素 $g_n$ 内所有对应像素的实例预测取平均：

\begin{equation*}
\hat{R}_t(\mathbf{P}^*) = \frac{1}{|g_n|}\sum_{\mathbf{P}^* \in g_n} \hat{R}_t(\mathbf{P}^*),
\end{equation*}

体素内的平均操作自然地消除了单帧误检带来的标签抖动。STL 模块的计算开销极低，每帧仅需约 0.979 ms。

\paragraph{全景分割公式与优化。}
全景分割采用 MaskFormer 式的分解策略，将问题拆为两步：先将图像分割为 $N$ 个区域，再为每个区域分配 $K$ 个类别的概率分布。区域预测通过语义嵌入经 MLP 提升维度后与区域嵌入矩阵 $\mathbf{M} \in \mathbb{R}^{N \times H}$ 做矩阵乘得到：$R(\mathbf{P}) = \phi(S(\mathbf{P})) \cdot \mathbf{M}$；类别预测则通过区域嵌入矩阵与分类器矩阵 $\mathbf{C} \in \mathbb{R}^{N \times K}$ 相乘：$O(\mathbf{M}) = \mathbf{M} \cdot \mathbf{C}$。像素同时具有高区域概率和高类别概率时才被赋予对应标签，低置信度预测被丢弃。整体优化包含五项损失：RGB 渲染的 $L_1$ 损失、深度渲染的 $L_1$ 损失、语义类别的交叉熵损失、实例区域的 DICE 损失和 Sigmoid Focal 损失。STL 带来的语义一致性还反过来提升了跟踪精度——一致的空间语义约束有效减少了位姿估计中的累积漂移。

\paragraph{实验表现。}
在 Replica 和 ScanNet++ 数据集上，PanoSLAM 在全景分割质量（PQ）、语义分割精度（mIoU）、跟踪精度（ATE RMSE）、渲染质量（PSNR/SSIM/LPIPS）和重建精度（Depth L1）五个维度上均超越已有语义 SLAM 方法（SNI-SLAM、DNS-SLAM、NIDS-SLAM 等）。跟踪和建图每帧开销分别为 19.9 ms 和 46.0 ms，相比 SplaTAM（15.2 ms 和 27.0 ms）的增加主要来自语义渲染和全景处理，STL 模块本身仅占不到 1 ms。值得注意的是，PanoSLAM 是首个无需任何人工标签即可完成全景 3D 重建的 SLAM 系统。

GSFF-SLAM（2025.04）~\cite{2504_19409} 沿着 PanoSLAM"为高斯赋予语义属性"的思路继续推进，但做出了一个根本性的设计转换：将语义从低维嵌入升级为高维\textbf{特征场}（feature field）。PanoSLAM 为每个高斯附加 3 维语义嵌入 $\mathbf{s} \in \mathbb{R}^3$，本质上是将语义压缩为类别级标签；GSFF-SLAM 则为每个高斯引入可训练的 $N$ 维特征嵌入 $f \in \mathbb{R}^N$（$N=128$），将语义信息表示为通用连续特征向量，而非离散类别索引。这一变化带来的核心好处是：特征场天然兼容任意形式的 2D 先验——既可以用 ground truth 标签监督（交叉熵损失），也可以用 Grounding-DINO + SAM + CLIP 管线生成的噪声文本特征监督（$L_1$ 特征距离）。

\paragraph{语义-几何解耦优化。}
第二个关键差异在于优化策略。PanoSLAM 的语义渲染使用独立的语义半径 $\rho_s$ 和语义不透明度 $\alpha_s$，致密化条件同时检查几何轮廓图和语义轮廓图——语义信号直接影响高斯的几何属性。GSFF-SLAM 选择了完全解耦的路线：在梯度反向传播中，特征嵌入 $f$ 的梯度仅来自语义监督，不回传到高斯的位置 $\mu$ 和协方差 $\Sigma$。建图过程也分为两阶段——先迭代优化几何和颜色属性直到收敛，再优化特征场。设计动机很明确：几何先收敛后，特征场可以建立在已稳定的三维结构上，迭代次数反而更少；同时避免了语义噪声信号（尤其是来自基础模型的稀疏、不一致预测）污染几何重建质量。

\paragraph{实验表现。}
在 Replica 数据集上，使用 ground truth 先验时语义分割 mIoU 达 95.03\%，跟踪精度（ATE RMSE 0.311 cm）也优于 SplaTAM 和 MonoGS 基线。得益于特征渲染的开/关切换机制——跟踪和非语义建图阶段不渲染特征图——系统在仅损失微量语义精度的前提下实现最高 2.9$\times$ 加速。当使用噪声文本先验替代 ground truth 时，GSFF-SLAM 仍能完成合理的语义重建，验证了特征场对稀疏噪声信号的鲁棒性。

RGBDS-SLAM（2024.12）~\cite{2412_01217} 走了一条截然不同的路线：不追求语义嵌入或特征场的表达能力，而是将语义信息直接作为像素级颜色附着于高斯基元（通过球谐函数表示），核心关切是场景细节恢复和多特征耦合优化。系统基于 ORB-SLAM3 框架而非 SplaTAM，跟踪线程复用 ORB 特征匹配，建图线程在独立的 GaussianMapping 线程中完成高斯优化。

\paragraph{3D 多级金字塔高斯泼溅（MLP-GS）。}
PanoSLAM 和 GSFF-SLAM 都在固定分辨率下训练高斯，对场景中不同尺度细节的捕捉能力有限。RGBDS-SLAM 引入渐进式金字塔训练策略：为 RGB、深度和语义图像各构建 $n$ 层（实验中 $n{=}3$）图像金字塔，训练从最低分辨率层开始逐步提升至原始分辨率。直觉上，低分辨率阶段快速建立全局结构一致性，高分辨率阶段逐步精化局部细节——这类似于先勾勒轮廓再描绘细节的绘画过程。每个训练迭代随机选取一个金字塔层级和对应视角，对 RGB 渲染（$\alpha$-混合颜色 $r_i$）、深度渲染（$\alpha$-混合深度值 $d_i$）和语义渲染（$\alpha$-混合语义颜色 $s_i$）使用完全相同的泼溅核，保证三种模态在空间上的对齐。

\paragraph{紧密耦合多特征重建优化（TCMF-RO）。}
与 PanoSLAM 将语义损失和几何损失分开设计不同，RGBDS-SLAM 将三路损失直接相加为联合目标 $L_{\mathrm{reconstruction}} = L_r + L_d + L_s$，其中 RGB 损失 $L_r$ 和语义损失 $L_s$ 各自混合 $L_1$ 与 SSIM（权重 $\lambda_r{=}\lambda_s{=}0.2$），深度损失 $L_d$ 仅用 $L_1$。三路损失共享同一组高斯参数做反向传播，语义和深度的梯度信号反过来修正 RGB 渲染质量——消融实验验证了这种耦合的有效性：去掉深度或语义特征后 PSNR 分别下降 0.49 和 0.23 dB，同时保留两者时语义 mIoU 也最高（94.32\%）。

\paragraph{实验表现。}
在 Replica 上，RGBDS-SLAM 取得平均 PSNR 38.85 dB、LPIPS 0.035，相比次优方法 Photo-SLAM 分别提升 11.13\% 和 68.57\%，这是 MLP-GS 渐进细节恢复的直接收益。深度 L1 误差 0.342 cm 也为所有方法中最优。语义重建 mIoU 94.32\% 超越 SGS-SLAM 的 92.72\%。跟踪和建图分别达 29.55 FPS 和 32.22 FPS，远高于 SGS-SLAM 的 5.27/3.52 FPS，这得益于 ORB-SLAM3 的特征匹配跟踪替代了梯度优化跟踪。系统还能用 SAM2 替代 ground truth 语义输入，通过多帧观测自动修正分割噪声。ATE 精度（RMSE 0.589 cm）弱于 SplaTAM（0.360 cm），因为跟踪模块直接复用 ORB-SLAM3 而未做额外优化。

GS3LAM（2026.03）~\cite{2603_27781} 同样沿"为高斯赋予语义属性"的路线，但与 PanoSLAM 和 GSFF-SLAM 不同，它的核心关切不是语义嵌入的维度或开放性，而是增量建图场景下语义 SLAM 的两个工程瓶颈——高斯尺度失配和灾难性遗忘。GS3LAM 将场景建模为语义高斯场（SG-Field），每个高斯携带低维语义特征 $\mathbf{f}_i \in \mathbb{R}^{N_{\mathrm{sem}}}$（$N_{sem}=16$），通过共享几何不透明度的 $\alpha$-混合渲染后经 CNN 解码器 $\mathcal{F}_{cnn}$ 加 softmax 输出语义标签。与 PanoSLAM 为语义渲染引入独立的语义半径 $\rho_s$ 和语义不透明度 $\alpha_s$ 不同，GS3LAM 的语义渲染完全复用几何不透明度 $\alpha$——语义特征向量的渲染公式与颜色渲染完全相同，仅替换 $\mathbf{c}_i$ 为 $\mathbf{f}_i$。这一设计简化了参数空间，代价是语义和几何的渲染质量强耦合。

\paragraph{深度自适应尺度正则化（DSR）。}
上述耦合关系带来了一个实际问题：3DGS 的高斯尺度在不同帧间差异极大（近处帧初始化为小尺度、远处帧为大尺度），不规则的大尺度高斯会模糊物体边缘，导致几何和语义在空间上错位——边缘处的高斯同时覆盖前景和背景，$\alpha$-混合后的语义预测自然不可靠。GS3LAM 提出深度自适应尺度正则化（Depth-adaptive Scale Regularization, DSR），用 $2\sigma$ 规则约束全局高斯尺度在 $\mu_s - 2\sigma_s < s < \mu_s + 2\sigma_s$ 范围内：超过上界 $s_{big} = \mu_s + 2\sigma_s$ 的高斯施加均值惩罚 $\mathcal{L}_{big}$，低于下界 $s_{small} = \mu_s - 2\sigma_s$ 的高斯施加负对数惩罚 $\mathcal{L}_{small}$。DSR 不直接强制尺度等于深度值，而是通过软约束间接将高斯对齐到几何表面，减少物体边缘模糊。消融实验表明 DSR 使跟踪误差降低 16\%、PSNR 提升 1.17 dB、语义 mIoU 提升 3.36\%。

\paragraph{随机采样关键帧建图（RSKM）。}
3DGS 增量建图的固有问题是灾难性遗忘——优化当前帧时共视高斯会过度拟合最新观测。SplaTAM 和 MonoGS 采用局部共视性关键帧建图（LCKM）：优化当前帧时选取与之共视的关键帧一起参与优化。GS3LAM 观察到 LCKM 存在优化偏置问题：共视密集区域（相机轨迹密集处）因重复优化难以收敛、PSNR 偏低且方差大，而共视稀疏区域则欠优化。RSKM 的思路截然不同——每次建图迭代以均匀概率 $1/|\mathcal{K}|$ 从全部关键帧中随机采样一帧参与优化，当前帧仅在间隔 $t_{opt}$ 的迭代步中被选中。这既无需计算共视关系，又保证了全局关键帧的均匀覆盖，有效降低优化偏置。消融实验显示 RSKM 使 PSNR 提升 3.73 dB，PSNR 方差降低 24.42 倍，显著改善全局一致性。

\paragraph{实验表现。}
在 Replica 上，GS3LAM 的渲染质量（PSNR 36.26 dB、SSIM 0.989、LPIPS 0.052）全面超越 SplaTAM、GS-SLAM 和 Point-SLAM；语义分割 mIoU 达 96.63\%，相比次优 SNI-SLAM 提升 9.22 个百分点。跟踪精度（ATE RMSE 0.37 cm）与 SplaTAM（0.36 cm）持平。渲染速度达 109 FPS（Replica，$1200 \times 680$），比 NeRF 方法快约 37 倍。

\subsubsection{开放词汇与层次化语义}

本路线介绍利用视觉-语言基础模型实现开放词汇语义查询和层次化语义建图的方法。

传统语义 SLAM 依赖预定义类别，无法识别训练集中未出现的新概念。开放词汇方法利用视觉-语言基础模型（如 CLIP）实现任意语义查询。

\paragraph{从隐式特征到显式标签：开放词汇语义 SLAM 的认知锚点。}
OpenGS-SLAM（2025.03）~\cite{2503_01646} 的核心动机源于一个关键洞察：具身智能（embodied intelligence）需要\textbf{3D 物体级}的场景理解与交互——机器人不仅要知道"这里有什么"，还要能从地图中直接提取、操作单个物体。然而，已有语义 3DGS 方法将语义信息\textbf{隐式}编码在高维特征向量中，语义标签被"埋"在连续特征里，无法直接访问每个高斯的类别归属，因而与物体级交互的需求根本不兼容。OpenGS-SLAM 选择了一条截然不同的路线：将 2D 基础模型的\textbf{显式语义标签}（而非隐式特征）集成到 3DGS 框架，使每个高斯直接携带可读的类别标识，从而原生支持 3D 物体提取、移动、旋转等交互操作。

\paragraph{高斯投票泼溅与场景表示。}
在基线做法中，已有语义 3DGS 方法将高维 CLIP 特征嵌入到每个高斯基元，通过特征渲染生成语义图——但高维特征存储（每高斯 768 维 CLIP 向量）导致巨大内存开销，且特征渲染的 $\alpha$-混合会对语义特征产生语义失真（不同物体的特征加权平均无物理意义）。OpenGS-SLAM 的核心改进是用离散标签替代连续特征嵌入：每个高斯基元存储一个一维、不可微的\textbf{GS Label}属性（而非高维语义向量），标签图通过投票机制渲染。直觉上，标准 3DGS 的 $\alpha$-混合是"加权平均"，适合颜色等连续量；但语义标签是离散的——"椅子"和"桌子"的加权平均没有意义。投票泼溅将 $\alpha$-混合从"平均"改为"计数"：每个高斯对其 GS Label 投一票，权重为 $\alpha$-混合的不透明度贡献 $w_i$，同一标签的所有高斯投票累加为该标签的总权重 $W_j$：

\begin{equation*}
W_j = \sum_{g_i \in G_p^j} w_i, \quad w_i = \alpha_i \prod_{k=1}^{i-1}(1 - \alpha_k),
\end{equation*}

其中 $G_p^j$ 表示像素 $p$ 处标签为 $j$ 的高斯子集。像素的最终标签取得票最多（$\arg\max_j W_j$，而非加权平均）的 GS Label。这一设计从根本上避免了语义特征的混合失真问题。

\paragraph{集成语义信息生成器。}
为获取 2D 语义标签，系统设计了\textbf{集成语义信息生成器}（Ensemble Semantic Information Generator）：输入 RGB 图像同时送入三条并行管线——RAM（Recognize Anything Model）生成开放集类别列表、YOLO-World（以 RAM 输出为文本提示）生成语义检测框和置信度、以及 SAM 生成分割掩码。三个模型的输出通过 IoU 匹配关联：SAM 的每个分割区域与重叠度最高的 YOLO-World 检测框配对（IoU 阈值 $> 0.5$），获得语义类别标签；无匹配的区域被标记为 None 类（置信度 0.0）。全局维护的 Label-Class 表 $M_{LC}^G$ 从空集开始动态扩展，支持开放集类别（实验中每场景最大标签数设为 2000）。值得强调的是，该生成器是可定制的——SAM 可替换为 MobileSAMv2 等轻量模型，无需额外训练即可适配不同算力平台。

\paragraph{场景致密化与在线更新。}
当帧被识别为关键帧后，系统利用 G-ICP（Gaussian-ICP）计算协方差初始化新高斯，并将其附带语义 GS Label 属性加入地图。可微参数（位置、颜色、不透明度等）通过损失函数标准优化；而不可微的 GS Label 则以显式方式更新——渲染时系统记录每个像素贡献最大的 $K = 50$ 个高斯（$G_{\mathrm{topK}} \in \mathbb{R}^{H \times W \times K}$）。若某像素的渲染标签从 $\ell_s$ 变为 $\ell_t$，系统在 $G_{\mathrm{topK}}$ 中定位此前标签为 $\ell_s$ 的贡献高斯，将其标签更新为 $\ell_t$，实现高效精准的语义标签传播。

\paragraph{置信度驱动的 2D 标签共识。}
多视角一致性是开放词汇 SLAM 的核心难题：SAM 等分割模型在不同视角下可能对同一物体给出不同标签，或将一个物体过度分割为多个部件。OpenGS-SLAM 通过\textbf{置信度驱动的 2D 标签共识}（Confidence-based 2D Label Consensus）在线解决此问题，利用 3DGS 的高效渲染将输入帧的语义信息与已有地图实时对齐。对于输入标签图 $L_s$ 和渲染标签图 $L_r$ 中的任意标签对 $(\ell_s^i, \ell_r^j)$，系统根据区域重叠程度识别四种匹配关系：
\begin{itemize}
\item \textbf{完全匹配}（Full Match）：$\ell_s^i$ 与 $\ell_r^j$ 区域高度重合，直接采用地图标签，取两者置信度的最大值。
\item \textbf{部分匹配}（Partial Match）：$\ell_s^i$ 仅与 $\ell_r^j$ 部分重叠——若地图置信度更高，则部分标签统一为 $\ell_r^j$；否则为高置信度的部分标签分配新类别，并通过 $G_{\mathrm{topK}}$ 精确更新对应高斯的标签。
\item \textbf{整体匹配}（Whole Match）：与部分匹配对称，新部件标签来自渲染图而非输入图。
\item \textbf{新标签}（New）：$\ell_s^i$ 无法匹配任何已有标签时，分配新类别并扩展 $M_{LC}^G$。
\end{itemize}
为提升共识精度，系统引入两个辅助模块。

\textbf{输入置信度更新}（Input Confidence Update）根据物体在当前视角的可见完整性调整置信度——部分遮挡的物体获得较低置信度，避免误更新。

\textbf{部分标签衰减}（Part Label Decay）针对 SAM 的已知过度分割倾向：每当一个标签被识别为"部件标签"时，其置信度被衰减 $\delta = 0.06$，确保过度分割的部件在多次准确观测后被正确合并为整体。

\paragraph{分割计数器剪枝。}
\textbf{分割计数器剪枝}（Segmentation Counter Pruning）解决视角约束不足区域高斯过度膨胀的问题——这些过大的高斯不仅阻碍新高斯的引入，还导致不准确的分割。策略追踪完全匹配标签的对称差集中的"计数高斯"：若其任一方向尺度超过阈值 $\tau$，则被剪除（Replica 上 $\tau = 0.10$，TUM 上 $\tau = 0.25$），从而精化物体轮廓并为新高斯腾出空间。

\paragraph{实验表现。}
在 Replica 上，OpenGS-SLAM 的开放词汇分割 mIoU 相比 Feature 3DGS 提升 13\%、相比 GS-Grouping 提升 3\%，同时使用 2 倍更少的可学习参数。语义渲染速度提升 10 倍（避免了高维特征渲染），存储成本降低 2 倍（每个高斯仅需 1 个整数标签 vs 768 维向量）。跟踪精度相比其他密集语义 SLAM 方法最高提升 51\%（归因于 G-ICP 跟踪模块和分割计数器剪枝对高斯分布对齐的增强）。在 3D 物体级理解方面，系统支持基于用户提示（点击或边界框）的交互式场景操作：物体移除、移动、旋转等——这在纯隐式特征的语义 3DGS 方法中无法实现。消融实验表明，部分标签衰减贡献最大（mIoU +5.96\%，Acc +10.2\%），置信度更新次之（mIoU +3.39\%），分割计数器剪枝主要精化边界（mIoU +0.26\%）。局限性在于当前框架不适用于动态场景，未来方向是利用语义信息实现动态环境中的物体级理解。

\paragraph{从离线预处理到在线实时：语言特征的在线 SLAM 融合。}
OpenGS-SLAM 用离散标签回避了高维特征渲染的语义失真问题，但其语义能力受限于 2D 基础模型（RAM + YOLO-World + SAM）的封闭集检测能力，无法响应任意自然语言查询。Online Language Splatting（2025.03）~\cite{2503_09447} 回到 CLIP 特征嵌入路线，但瞄准的是一个被忽视的系统瓶颈：现有 Lang-GS 方法（如 LangSplat）依赖离线 SAM+CLIP 预处理生成逐像素语言特征图，每帧耗时数分钟，根本无法用于在线 SLAM。其核心设计围绕三个在线挑战展开。

\paragraph{高分辨率 CLIP 嵌入与超分辨率解码。}
离线方法用 SAM 生成分割掩码再逐区域提取 CLIP 特征，精度高但极慢。Online Language Splatting 采用 ConvNeXt 像素级 CLIP 编码器直接生成粗粒度语言特征图（$24 \times 24 \times 768$），再通过轻量化超分辨率解码器（SRD）上采样至 $192 \times 192 \times 768$，整个过程仅需 18 ms/帧（RTX-3090，SRD 仅占 2 ms）。与 OpenGS-SLAM 的集成语义信息生成器（RAM + YOLO-World + SAM 三管线并行）不同，SRD 不做分割检测，只做特征图的空间分辨率恢复——两者解决的是不同层面的问题：OpenGS-SLAM 需要将连续观测离散化为标签，而 SRD 需要在保持语义连续性的同时提升空间精度。

\paragraph{两级在线自编码器：768 $\to$ 32 $\to$ 15 维。}
即使有高分辨率特征图，768 维 CLIP 向量直接存储到每个高斯仍然不可行。离线方法（如 LangSplat）可以在目标场景上预训练自编码器，但在线方法面对的是未见场景。为此，系统设计两级压缩：第一阶段为离线预训练的通用自编码器（8 层 MLP），在 COCO 等多样化数据集上训练，将 768 维压缩至 32 维——该维度不能太低，否则跨域泛化时开放词汇能力受损。第二阶段为\textbf{在线学习自编码器}（OLAE），利用当前场景的数据方差可将特征进一步压缩至 15 维。OLAE 先进行 200 次迭代初始化（6 ms/迭代），之后每帧增量更新 1 次，每次加入 2 个随机关键帧以防止灾难性遗忘。两级设计的直觉是：通用压缩器保留广度（跨场景的开放词汇能力），OLAE 捕获当前场景的紧凑表示（利用场景内语义类别远少于全局的特性）。

\paragraph{颜色-语言解耦优化。}
此前 Lang-GS 方法（包括 LangSplat）将语言特征作为额外通道绑定到同一组高斯参数（$\alpha$、$\mathbf{R}$、$\mathbf{S}$），颜色和语言的梯度耦合在统一反向传播路径中。然而，语言特征与颜色对高斯参数的偏好根本不同——语言特征是同质的（整面墙共享同一语言编码），倾向于拉伸高斯尺度和调整旋转；而颜色需要精细的纹理细节。联合优化导致两者互相干扰，渲染质量双双下降。Online Language Splatting 为颜色和语言设计独立的 $\alpha^c$、$\mathbf{R}^c$、$\mathbf{S}^c$ 与 $\alpha^f$、$\mathbf{R}^f$、$\mathbf{S}^f$，但共享位置 $\bm{\mu}$ 以避免高斯数量翻倍。语言梯度的反向传播路径被完全隔离：$\partial\mathcal{L}/\partial\alpha_i^f = \partial\mathcal{L}/\partial F \cdot \partial F/\partial\alpha_i^f$。为防止语言模式的尺度偏离过远，引入正则化 $|\mathbf{S}_i^f - \mathbf{S}_{i\bot}^c|$，以颜色模式的尺度（stop-gradient）为锚约束语言尺度。这一设计体现了与 OpenGS-SLAM 的深层差异：OpenGS-SLAM 通过投票机制彻底绕开了 $\alpha$-混合对离散标签的语义失真；Online Language Splatting 则承认语言特征是连续量、$\alpha$-混合在数学上合理，但需要为不同模态解耦参数空间。

\paragraph{实验表现。}
在 Replica 上，Online Language Splatting 的开放词汇定位 mIoU 达 0.487、定位精度（Loc）0.826，超越离线方法 LangSplat（mIoU 0.417、Loc 0.720）和 Feature3DGS（mIoU 0.359），同时每帧处理时间仅 0.8 s，比 LangSplat（2.8 min/帧）快约 210$\times$。在 TUM RGB-D 上，定位精度（0.917、0.791）同样优于 LangSplat，但 mIoU 略低——论文归因于运动模糊和低图像质量对在线跟踪的不利影响。SLAM 渲染质量（PSNR、SSIM）与 MonoGS 基线持平，验证了语言建图模块不干扰几何重建。

\paragraph{从显式标签到层次化编码：语义 SLAM 的可扩展性瓶颈。}
OpenGS-SLAM 用一维离散标签替代高维语义特征，在物体级交互和存储效率上迈出关键一步。然而，这一路线面临一个新瓶颈：当场景语义类别数从数十增长到数百甚至上千时，每个高斯的语义嵌入维度必须随类别数线性增长（扁平 Softmax 编码），导致存储和计算成本急剧膨胀。以 ScanNet 的 550 个语义类别为例，扁平编码所需的参数量在现有硬件上已无法运行。Hier-SLAM（2024.09）~\cite{2409_12518} 的核心洞察在于：语义信息天然具有\textbf{层次结构}——``椅子''是``家具''的子类，``家具''又是``物体''的子类。利用这种层次关系，可以用一棵深度为 $L$ 的树编码 $C$ 个语义类别：一棵深度为 10 的二叉树仅需 $2 \times 10 = 20$ 维编码即可覆盖 $2^{10} = 1024$ 个类别，相比扁平编码的 1024 维实现了数量级的压缩。这一从"扁平"到"层次"的范式转换，使语义 SLAM 首次具备了处理 500+ 类别复杂场景的可扩展能力。

\paragraph{LLM 辅助的层次树生成。}
构建合理的层次树本身是一个非平凡的问题——需要同时考虑语义归属关系（``椅子''属于``家具''）和几何先验（``椅子''和``桌子''都是平面类结构）。Hier-SLAM 利用大语言模型（GPT-4o-mini）的常识推理能力自动构建层次树：将叶节点的语义类别集输入 LLM，由其逐层向上聚类为更粗粒度的父类，从叶到根反复迭代直至根节点产生少于 $\kappa = 4$ 个簇。为解决 LLM 聚类时常见的遗漏和幻觉问题（遗漏某些输入类别、虚构未出现的类别），系统设计了\textbf{循环批评机制}（loop-based critic）：每次聚类后，验证器将结果与输入对比，分为成功分组、幻觉类别（移除）和遗漏类别（重新输入 LLM 聚类），循环直至无遗漏。以类 ``Wall'' 为例，4 层编码为 $\{v_0{:}\text{Background}\} \to \{v_1{:}\text{Structure}\} \to \{v_2{:}\text{Plane}\} \to \{v_3{:}\text{Wall}\}$，每一层既编码了语义归属也蕴含了几何先验。值得注意的是，树生成在 SLAM 运行前离线完成，不增加在线计算负担。

\paragraph{层次化语义编码与渲染。}
每个 3D 高斯基元携带一个层次化语义嵌入 $\psi$，由各层编码拼接而成：$\psi = f(\psi^l)$，$\psi^l \in \mathbb{R}^{n_l}$，其中 $n_l$ 为第 $l$ 层的最大节点数，总维度 $N = \sum_{l=0}^L n_l$。语义图的渲染沿用标准 $\alpha$-混合泼溅：

\begin{equation*}
H = \sum_{i=1}^n \psi_i \, G_i(\mathbf{X}) \, T_i, \quad T_i = \prod_{j=1}^{i-1}(1 - G_j(\mathbf{X})),
\end{equation*}

与 OpenGS-SLAM 的投票泼溅不同，Hier-SLAM 对语义嵌入使用连续 $\alpha$-混合——因为层次化编码是连续的概率分布（Softmax 输出），而非离散标签，加权平均具有明确的概率含义。系统采用统一的前向/后向模块同时处理语义、颜色、深度和轮廓渲染（而非分离模块），显著提升了运行效率。

\paragraph{层次化损失：层内 + 层间联合优化。}
为充分利用层次结构的信息，Hier-SLAM 提出\textbf{层次化损失}：

\begin{equation*}
\mathcal{L}_{\text{Semantic}} = \lambda_1 \mathcal{L}_{\text{Inter}} + \lambda_2 \mathcal{L}_{\text{Cross}},
\end{equation*}

在上式中，\textbf{层内损失}（即 Inter-level loss）在各层独立施加交叉熵监督：$\mathcal{L}_{\text{Inter}} = \sum_{l=0}^L \mathcal{L}_{ce}(S(\psi^l), P^l)$，确保每一层的编码都被正确优化。

\textbf{层间损失}（Cross-level loss）则通过一个共享线性层 $F$ 将层次化嵌入解码回扁平编码，再与扁平语义真值计算交叉熵：$\mathcal{L}_{\text{Cross}} = \mathcal{L}_{ce}(S(F(\psi)), P)$。直觉上，层内损失保证"每一层都学对"，而层间损失则保证"所有层组合起来也学对"——二者互补，前者提供逐层精化的梯度信号，后者确保跨层一致性。训练策略上，前 $N_s = 15$ 次迭代仅使用层内损失初始化层次编码（$\lambda_1 = 1.0, \lambda_2 = 0.0$），之后引入层间损失精化（$\lambda_1 = 1.0, \lambda_2 = 5.0$），体现"先粗后细"的训练哲学。

\paragraph{跟踪与建图。}
跟踪采用恒速模型初始化位姿，固定全局地图优化当前帧位姿，损失为颜色和深度的 L1 加权和（仅在轮廓可见区域 $M = \mathbb{1}(S > \theta)$ 上计算）。建图阶段固定位姿优化全局地图，损失融合深度 L1、颜色 SSIM+L1 和语义层次化损失。与 OpenGS-SLAM 的 G-ICP 跟踪不同，Hier-SLAM 的跟踪模块基于可微渲染的位姿优化，语义信息通过建图损失间接提升几何质量。

\paragraph{实验表现与可扩展性。}
在 Replica（102 类）上，Hier-SLAM 的扁平版本取得 mIoU 90.35\%（所有方法最高），但存储开销巨大；层次化版本的 mIoU 竞争性持平，而存储仅需 910.50 MB，比扁平版本减少 66\%。跟踪速度比 SplaTAM 快 2.4$\times$，建图速度快 2.2$\times$；相比扁平语义编码的 SLAM，层次编码使跟踪快 3$\times$、建图快 1.2$\times$。语义渲染 FPS 达到 2000（含语义）和 3000（不含语义）。最关键的突破在于\textbf{可扩展性}：在 ScanNet 的 550 个语义类别上，扁平编码因存储限制无法运行；而层次化编码通过 LLM 辅助将 550 类压缩至 72 维编码（7.6$\times$ 压缩），首次实现了 500+ 类复杂真实场景的全局 3D 语义建图，展示了从粗到精的层次化语义理解能力。

\paragraph{单目支持与增强层次树。}
Hier-SLAM++（2025.02）~\cite{2502_14931} 在两方面做出关键扩展。其一，\textbf{支持单目输入}：成为首个语义单目高斯泼溅 SLAM 系统。单目设置缺乏直接深度信息，系统引入 3D 前馈模型 DUSt3R 提供几何先验——从视频中采样稀疏视角输入 DUSt3R 获取深度预测，再通过仿射变换（尺度-偏移校正）与 SLAM 系统的全局地图对齐。校正参数与 3D 高斯联合优化，使深度先验随在线观测持续精化。其二，\textbf{增强层次树生成}：引入 3D 生成模型 MeshGPT 提取各类别的几何嵌入（文本提示输入 MeshGPT 获取量化的 3D 形状编码），用 K-means++ 聚类并辅以 LLM 摘要生成几何层次（如"box""flat""soft"），使层次树同时编码语义归属和几何形状。在编码策略上，除 one-hot 层次编码外新增\textbf{二进制编码}：将每层的 $n$ 维 one-hot 向量压缩为 $\lceil\log_2 n\rceil$ 位二进制码，进一步降低存储需求。损失函数也从 Hier-SLAM 的 Inter/Cross 两项演化为更完整的\textbf{层内}（Intra-level）和\textbf{层间}（Inter-level）优化框架，层内损失支持标准交叉熵和二进制交叉熵两种形式。

\subsection{实时建图与特殊应用}
\label{ss:15_9}

本节介绍两条技术路线：（1）实时语义建图与高效特征管理；（2）特殊语义建图应用。

\subsubsection{实时语义建图与高效特征管理}

本路线介绍在高维语义特征与实时 SLAM 之间取得平衡的方法，包括前馈高斯预测和高效特征压缩策略。

在高维语义特征与实时 SLAM 之间取得平衡是工程上的核心挑战。现有的 GS-SLAM 方法存在两个系统性瓶颈：（1）依赖逐场景的梯度优化来估计高斯参数，计算开销大且不适于实时部署；（2）采用启发式的高斯致密化与剪枝策略，往往产生过于稠密的表示，难以扩展到大场景。GS4（2025.06）~\cite{2506_06517}（Generalizable Sparse Splatting Semantic SLAM）是首个可泛化的 GS 语义 SLAM 系统，用学习到的\textbf{前馈网络}直接预测 3D 语义高斯，彻底消除了逐场景优化的需求。

\paragraph{架构概览。}
GS4 由三个协同模块构成。（1）\textbf{高斯预测网络}：以 RGB-D 图像为输入，基于预训练的 2D 分割模型（如 Mask2Former 或 AutoFocusFormer）backbone，在多个逐步下采样的层级上生成 encoder/decoder token。这些 token 经过局部注意力 transformer 层融合编码器与解码器特征后，由 MLP 解码为每个高斯的几何参数（位置偏移 $\Delta x, \Delta y, \Delta d$、尺度 $s$、旋转 $r$、不透明度 $o$）和颜色偏移。语义分割标签则由 backbone 自带的 mask decoder head 直接预测，无需外部语义分割模块——这是将语义预测与几何预测统一在同一前馈管线中的关键设计，避免了传统 SLAM 系统中``先建图再分割''的串行范式。（2）\textbf{高斯精化网络}：处理新帧预测的高斯与已有 3D 地图中可见高斯的合并。输入包括所有可见高斯的特征向量 $f_i$ 和 3D 位置，通过在世界坐标系中构建 3D 邻域的局部注意力 transformer 层融合并更新特征，再由 MLP 预测颜色、尺度、旋转、不透明度的更新量。精化后自动剪枝不透明度低于 $0.005$ 的高斯；同时引入不确定性阈值 $U$——当高斯被精化网络更新达到 $U$ 次后（实验中 $U\!=\!8$），认为其不确定性已充分降低，不再参与后续更新。这种学习式的精化策略替代了传统启发式的致密化/剪枝规则，用数量级更少的高斯实现了同等甚至更优的表示质量。（3）\textbf{少量迭代的联合高斯--位姿优化}：回环检测和全局光束法平差（BA）会修正相机位姿，但可能导致基于旧位姿插入的高斯与新位姿不一致。为此，GS4 对位姿变化最大的 $k$ 帧重新渲染 RGB-D 图像，联合优化高斯参数和相机位姿，仅执行 1--5 次迭代。优化后的位姿反馈给跟踪模块用于后续的回环和 BA，形成闭环校正与地图一致性之间的正反馈。跟踪模块本身采用现成的 GO-SLAM 或 Mast3R-SLAM。

\paragraph{训练策略。}
高斯预测网络采用低分辨率到高分辨率的课程学习（$256\!\times\!256 \to 480\!\times\!480 \to 640\!\times\!480$），在 ScanNet 上训练，监督信号包括 RGB 渲染损失（MSE + 感知损失）、深度 L1 损失、语义交叉熵损失和分割损失 $L_{\text{seg}}$。精化网络采用渐进式训练（2 帧 $\to$ 4 帧 $\to$ 8 帧），以更好地处理多帧累积。

\paragraph{实验表现。}
在 ScanNet 上，GS4 在渲染、重建、语义分割和相机跟踪四个维度均达到最佳：语义 3D mIoU 超越第二名 GS$^3$LAM 20.42\%，2D mIoU 提升 7.29\%；渲染 PSNR 比同跟踪模块的 Splat-SLAM 高 1.76 dB。整体速度比 SplaTAM 快 12 倍、比 SGS-SLAM 快 17 倍、比 GS$^3$LAM 快 24 倍，而使用的高斯数量仅为竞争方法的约 10\%（ScanNet++ 上仅 3\%--4.5\%）。消融实验表明高斯精化网络和联合高斯--位姿优化均不可或缺，高斯剪枝在不损失精度的前提下显著减少高斯数量。在 NYUv2 和 TUM RGB-D 上的零样本迁移实验进一步验证了泛化能力：在 NYUv2 上，渲染指标相对第二名提升 15\%--20\%；在 TUM RGB-D 上 SSIM 最优，其他指标与最佳方法持平，且高斯数量远少于所有对比方法。局限性在于当前框架不适用于动态场景，且依赖 RGB-D 输入，未来方向包括提升计算速度和探索纯 RGB 方案。

LEG-SLAM（2025.06）~\cite{2506_03073} 同样追求高维语义特征的在线实时蒸馏，但与上述路线的根本区别在于：OpenGS-SLAM 以离散标签回避特征混合问题，LangGS-SLAM 以 Top-K 渲染规避语义畸变，LEG-SLAM 则选择将 DINOv2 视觉特征直接嵌入高斯基元，通过标准 $\alpha$-混合统一渲染 RGB 图像、深度图和语义特征图——语义畸变问题不通过渲染机制改造来解决，而是依靠特征的压缩质量和余弦相似度损失来约束。

\paragraph{DINOv2 + PCA 可学习特征压缩。}
特征提取管线选用 DINOv2（而非 CLIP）作为视觉骨干：输入 $518 \times 518$ 图像，输出 $37 \times 37 \times 768$ 的嵌入图，每个空间单元的 768 维向量编码局部视觉语义。核心创新是基于 PCA 的可学习压缩器，将 768 维压缩到 $K$ 维。训练数据用 ImageNet 的 1000 个文本类：每个类的文本嵌入先经 CLIP 编码、再通过 Talk2DINO 模块投影到 DINOv2 空间，并附加满足余弦相似度 $\geq 0.95$ 的噪声增强——策略的意图是让 PCA 学到通用的特征投影而非特定类别的表示，从而保持对 ImageNet 以外类别的泛化能力。压缩后的 $K$ 维嵌入经双线性插值上采样回原始分辨率 $H \times W \times K$，作为语义监督信号送入 3DGS 管线。每个高斯基元携带一个 $K$ 维语义向量，初始化为零（分裂高斯则复制父高斯的语义特征），与位置、颜色、尺度等参数联合优化。

\paragraph{统一渲染管线与 Talk2DINO 开放词汇查询。}
渲染端沿用标准 $\alpha$-混合，同时输出 RGB 图像、深度图和语义特征图——三个输出共享同一条可微泼溅管线，避免多模块串联的开销。损失函数由三项组成：颜色损失 $\mathcal{L}_{\text{color}}$（L1 + SSIM）、语义余弦相似度损失 $\mathcal{L}_{\text{cos\_sim}}$（渲染特征图与压缩后 DINOv2 特征之间的余弦距离）和深度 L1 损失。开放词汇查询通过 Talk2DINO 模块实现跨模态桥接：用户输入文本经 CLIP 编码为 $1 \times 768$ 嵌入，Talk2DINO 将其从 CLIP 空间映射到 DINOv2 空间，再经同一 PCA 压缩到 $1 \times K$，最后与渲染的语义特征图逐像素做点积，生成语义热力图。这一管线使得系统支持任意文本查询而无需预定义类别或预计算静态语义地图。

\paragraph{实验表现。}
在 Replica 上以超过 10 FPS 实现实时语义建图和渲染，ScanNet 上达 18 FPS，比同类语言增强 3DGS 方法快最高 30 倍。系统无需预先准备相机运动轨迹或预计算的静态语义地图——这是与 SGS-SLAM 等需要离线语义先验的方法的关键区别。消融实验表明 PCA 压缩维度 $K$ 的选择在语义精度和速度之间提供明确的权衡空间。

LangGS-SLAM（2026.02）~\cite{2602_06991} 提出 Top-K 渲染管线实现高吞吐量、无语义畸变的高维特征图渲染——这是语言特征渲染中的独特挑战，因为标准光栅化会对语义特征产生失真。多标准地图管理策略在保持场景完整性的同时修剪冗余或不一致的高斯，混合场优化框架通过解耦几何和语义场的优化频率在实时约束下联合精化两者。LEGO-SLAM（2025.11）~\cite{2511_16144} 是首个在 3DGS SLAM 系统中实现实时开放词汇建图的框架。核心是场景自适应编码器-解码器将高维语言嵌入蒸馏为紧凑的 16 维特征空间，在线适应未见场景（不同于静态方法的固定压缩）。基于语言的修剪策略识别语义冗余将高斯数量减少超过 60\%，基于语言的回环检测重用建图特征消除单独检测模型的需求。LatentAM（2026.02）~\cite{2602_12314} 提出在线字典学习方法——每个高斯基元关联一个紧凑查询向量，通过注意力机制和可学习字典转换为近似 VLM 嵌入。模型无关且无需预训练，实现即插即用的 VLM 集成，以接近实时速度（12--35 FPS）运行。

\subsubsection{特殊语义建图应用}

本路线介绍任务驱动的概率语义建图和多模态语义融合等特殊应用场景。Bayesian Fields（2025.03）~\cite{2503_05949} 认识到对象粒度是任务相关的——如何定义"物体"取决于具体任务。提出任务驱动的概率语义建图方法，利用贝叶斯更新（而非简单的多视角嵌入平均）聚合多视角观测。通过利用底层视觉-语言基础模型的特性实现更准确的语义融合，将 3DGS 聚类为任务相关的物体，支持简单的 3D 物体提取和降低内存使用。

OpenGS-Fusion（2025.08）~\cite{2508_01150} 将 3DGS 与截断有符号距离场结合实现无损语义特征融合，支持即时的语义特征无损融合。其 MLLM 辅助自适应阈值策略根据多模态大语言模型的反馈自适应调整相似度阈值，将 3D mIoU 提升 17\%。

从物体提取转向物体级位姿估计，DQO-MAP（2025.03）~\cite{2503_02223} 利用 3DGS 进行高保真物体重建，利用二次曲面进行精确物体位姿估计。通过为物体关联唯一 ID 实现场景中物体的快速提取。

在语言嵌入高斯方向，LEGS（2024.09）~\cite{2409_18108} 增量构建语言嵌入高斯泼溅（Language-Embedded Gaussian Splat），训练速度比 LERF 快 3.5 倍，同时实现可比较的开放词汇查询成功率。

从开放词汇查询扩展到导航应用，Go-SLAM（2024.09）~\cite{2409_16944} 为每个高斯分配唯一物体标识符，支持开放词汇查询和最优路径生成模块，计算考虑障碍物和环境不确定性的高效导航路径。

在多模态融合方面，$\Psi$-Map（2026.04）~\cite{2604_10982} 利用 LiDAR 数据构建平面约束的多模态高斯混合模型，使用 2D Gaussian Surfel 作为地图表示实现高精度表面对齐。通过查询引导的端到端学习架构——利用视锥体内的局部交叉注意力将 2D 掩码特征直接提升到 3D 空间——实现全局一致的全景理解，推理速率超过 40 FPS。

从 RGB-D 扩展到纯单目设定，OpenMonoGS-SLAM（2025.12）~\cite{2512_08625} 是首个统一 3DGS 与开放词汇语义理解的单目 SLAM 框架，利用 MASt3R、SAM 和 CLIP 等视觉基础模型，不依赖任何深度输入或 3D 语义真值。

在光照变化处理方面，"Taming the Light"（2025.11）~\cite{2511_22968} 通过内在外观归一化（IAN）模块主动解耦场景固有属性（如反照率）与瞬态光照，动态辐射平衡损失（DRB-Loss）仅在曝光极端时激活，直接作用于辐射场引导目标优化。

最后，在系统评估方面，Galagain et al.（2025.05）~\cite{2505_12384} 对比评估了几何 SLAM、NeRF 和 3DGS 三类语义 SLAM 方法在嵌入式平台（NVIDIA Jetson AGX Orin）上的精度、分割质量、内存使用和能耗，揭示 NeRF 和 3DGS 方法虽语义细节丰富但计算资源需求高。

\begin{takeaway}
语义 SLAM 将 3DGS 地图从几何描述提升为语义理解：开放词汇方法通过 CLIP 特征蒸馏实现零样本语义分割，3DGS 的显式高斯基元天然适合作为语义载体，使渲染时即可产生稠密的语义分割图。
\end{takeaway}

% ======================================================================


\section{动态场景SLAM}
\label{sec:15_6}

标准 SLAM 假设场景静态——所有物体不动。但真实世界充满动态物体：行人走过、车辆驶过、门被推开。动态物体违反了光束法平差的静态假设，导致跟踪漂移和地图污染。本节梳理 17 篇论文，覆盖从"反动态"（滤除动态物体）到"拥抱动态"（同时重建静态和动态内容）的技术演进。

\subsection{动态场景处理}
\label{ss:15_6a}

本节介绍三条技术路线：（1）动态物体检测与滤除；（2）不确定性驱动的动态处理；（3）4D 建模与动态-静态联合重建。

\subsubsection{动态物体检测与滤除}

本路线介绍检测并移除动态物体以保护静态背景重建质量的方法，聚焦无预定义类别下的动态区域识别。

处理动态场景最直接的策略是检测并移除动态物体，仅重建静态背景。核心挑战在于如何在没有预定义物体类别的情况下准确识别动态区域。

DG-SLAM（2024.11）~\cite{2411_08373} 是首个基于 3DGS 的鲁棒动态视觉 SLAM 系统。

\paragraph{Baseline 与改进动机。}
此前的 3DGS-SLAM 方法（如 SplaTAM、GS-SLAM）和 NeRF-SLAM 方法（如 Co-SLAM、ESLAM）均假设场景静态，当动态物体出现时鲁棒性显著下降甚至跟踪失败。已有动态 SLAM 方法分为两类：一类通过重采样和残差优化剔除离群点（如 ORB-SLAM3、ReFusion），但仅能处理小幅运动，面对大范围连续运动时失效；另一类依赖语义分割或目标检测先验（如 DynaSLAM、DS-SLAM），但在真实场景中存在域间隙，引入预测错误。DG-SLAM 的核心思路是：不依赖单一先验，而是将\emph{几何一致性}（深度变形）与\emph{语义先验}互补融合，再配合粗-精两级位姿优化，在动态场景中实现鲁棒的跟踪和高质量重建。

\paragraph{运动掩码生成。}
运动掩码的目标是标记哪些像素属于动态区域，以便在跟踪和建图中将其排除。DG-SLAM 利用\textbf{深度变形一致性}检测动态像素，其直觉是：如果场景是静态的，将一幅深度图沿估计位姿变形到另一帧后，变形深度应与观测深度高度吻合；不一致的像素大概率属于运动物体。具体而言，对于关键帧对 $(i, j)$，将关键帧 $i$ 中像素 $\mathbf{p}_i$ 通过估计位姿重投影到关键帧 $j$：

\begin{equation}
    \mathbf{p}_{i \to j} = f_{\mathrm{warp}}(\mathbf{p}_i, \mathbf{T}_{ji}, D_i(\mathbf{p}_i)) = \mathbf{K}_j \mathbf{T}_{ji} \left(\mathbf{K}_i^{-1} D_i(\mathbf{p}_i) \mathbf{p}_i^{\mathrm{homo}}\right)
\end{equation}

其中 $\mathbf{K}$ 为内参矩阵，$\mathbf{T}_{ji}$ 为帧间变换，$\mathbf{p}_i^{\mathrm{homo}}$ 为像素齐次坐标。计算变形深度与观测深度的残差——若残差超过阈值 $e_{\mathrm{th}}$，则该像素标记为动态，得到深度变形掩码 $M_{j,i}^{\mathrm{wd}}$。

为提高鲁棒性，DG-SLAM 在滑动窗口内融合 $N$ 个连续观测的变形掩码取\textbf{交集}：

\begin{equation}
    M_j = M_{j,i}^{\mathrm{wd}} \cap M_{j,i-1}^{\mathrm{wd}} \cap \cdots \cap M_{j,i-N}^{\mathrm{wd}} \cap M_j^{\mathrm{seg}}
\end{equation}

取交集的策略有效缓解了两个问题：单次观测中的噪声（某次变形可能因位姿误差而误判）和边缘区域的误判（物体边缘处深度不连续，单次变形易产生虚假残差）。同时，当物体运动幅度很大时，DG-SLAM 仅遮蔽深度残差为正的前景像素，避免大范围像素被错误标记为动态。最终运动掩码 $M_j$ 通过结合深度变形掩码和语义分割掩码获得——语义分割可快速识别常见动态类别（行人、车辆等），而深度变形可捕获语义模型无法识别的非常规移动物体（如被推动的椅子、飘动的气球），两者互补使检测更加完整。

\paragraph{混合相机跟踪：粗-精两级策略。}
动态场景中，恒速运动模型难以推断合理的初始位姿，而初始位姿偏差会导致优化陷入局部最优，还会影响运动掩码的生成质量。为此，DG-SLAM 设计了两阶段策略：

\textbf{粗跟踪}利用 DROID-SLAM 的视觉里程计（DROID-VO）提供初始位姿估计，并在每一步执行密集光束法平差（Dense Bundle Adjustment, DBA）优化一组关键帧的位姿和深度：

\begin{equation}
    E(\mathcal{G}, d) = \sum_{(i,j)} \left\| \mathbf{p}_{i \to j}^* - \pi_{\mathbf{C}} \left(\mathbf{G}_{ij} \pi_{\mathbf{C}}^{-1}(\mathbf{p}_i, d_i) \right) \right\|_{\Sigma_{ij}}^2
\end{equation}

其中 $\Sigma_{ij}$ 为置信度权重矩阵。关键在于：通过运动掩码 $M_j$ 将动态像素的置信度权重置零，使动态物体不影响位姿优化。

\textbf{精跟踪}基于高斯泼溅渲染，在粗跟踪得到的初始位姿基础上进一步精化。但直接利用渲染图进行优化存在两个陷阱：动态物体遮挡导致渲染伪影，以及地图中尚未充分优化的高斯点导致渲染质量差。为此，DG-SLAM 引入\textbf{累积不透明度}（accumulated opacity）$\mathcal{O}$ 作为渲染可靠性指标——将每个像素处所有高斯的加权不透明度求和，当 $\mathcal{O}$ 超过阈值 $\mathcal{O}_{\mathrm{track}}$ 时，认为该像素对应的高斯点已充分优化、渲染结果可靠。精跟踪仅在可靠像素上计算光度-深度联合损失：
\begin{multline}
    \xi_i^* = \arg\min_{\xi_i} \; \lambda_1 \frac{1}{M} \sum_{k=1}^{M} \left\| \left(\hat{C}(\mathcal{G}, \xi_i) - C\right) \odot M_i \odot \mathcal{O}_i \right\|_2^2 \\
    + \lambda_2 \frac{1}{N_d} \sum_{p \in N_d} \left\| \left(\hat{D}(\mathcal{G}, \xi_i) - D\right) \odot M_i \odot \mathcal{O}_i \right\|_2^2
\end{multline}
其中 $M_i$ 为运动掩码（排除动态像素），$\mathcal{O}_i$ 为可靠渲染掩码（排除未充分优化的区域），$N_d$ 为有效深度像素集合。这种双重过滤策略确保位姿优化既不受动态物体干扰，也不被低质量渲染误导。

\paragraph{自适应高斯点管理。}
为保持几何质量，DG-SLAM 提出自适应高斯点增删策略。增点采用两阶段方案：对新视角先均匀采样确保覆盖，再对累积不透明度低于阈值或深度残差过大的欠拟合像素添加新高斯。删点基于三个准则——不透明度过低（$\alpha_i < \alpha_\epsilon$）、尺度异常（$S_i > S_{s1}$）、椭球纵横比过大（$\max(S_i)/\min(S_i) > S_{s2}$）——删除不合理的高斯点。此外，由于运动掩码边缘可能不准确，边缘附近新添加的点可能在动态区域产生伪影；DG-SLAM 将这些点投影到滑动窗口内的关键帧上重新检查，若观测次数过低则移除。

在动态场景的新视角合成、地图重建和位姿估计上，DG-SLAM 均达到当时最佳水平。

DG-SLAM 的运动掩码依赖深度变形一致性与语义分割的互补融合，并采用粗-精两级跟踪策略。DGS-SLAM（2024.11）~\cite{2411_10722} 在此基础上做出两项结构性改进：一是将动态滤波从跟踪-建图阶段扩展到高斯初始化和关键帧选择等全管线环节，从根源上阻止动态物体污染地图；二是放弃深度变形检测，改用在线视频实例分割模块（Track Anything）获取语义掩码，并引入基于光度残差的鲁棒掩码弥补分割残缺。

\paragraph{全管线动态滤波 vs. 跟踪-建图滤波。}DG-SLAM 在跟踪和建图中通过运动掩码排除动态像素，但高斯插入和关键帧选择环节未考虑动态性，导致动态区域的高斯仍会被添加到地图中。DGS-SLAM 将动态滤波扩展到整个管线——高斯仅在语义掩码标记的静态区域插入，关键帧选择也排除动态干扰——从根源上阻止动态物体污染地图。

\paragraph{鲁棒掩码：弥补语义分割的残缺。}DG-SLAM 依赖语义分割和深度变形掩码的交集；当语义分割不完美时（如阴影、分割边界模糊），残余伪影直接进入建图。DGS-SLAM 引入\textbf{鲁棒掩码}作为第二道防线：在建图阶段，将每个关键帧的光度残差 $R_k = \|\hat{C}(\mathcal{G}, T_k) - C_k\|_1$ 构建直方图，取第 $\tau_{\mathrm{robust}}$ 百分位（$\tau_{\mathrm{robust}} = 0.9$）作为内点阈值 $\epsilon_{\mathrm{robust}}$——残差低于阈值的像素视为已正确建模的内点，超过阈值的像素则为未被语义掩码捕获的离群区域（如动态物体的阴影投射、分割边缘的误标像素）。为避免高频纹理区域（如棋盘格）被误判为离群，对二值化残差图用归一化核 $\mathbf{W}$ 做空间平滑，再取阈值 0.5 生成最终鲁棒掩码。直方图在窗口优化过程中以指数移动平均迭代更新，使掩码随地图质量改善而逐步精化。最终窗口优化掩码为语义掩码与鲁棒掩码补集的乘积：$M_{\mathrm{window}}^k = M_{\mathrm{seg}}^k \otimes (1 - M_{\mathrm{robust}}^k)$。

\paragraph{回环感知窗口选择。}DG-SLAM 仅在滑动窗口内优化，移出窗口的关键帧不再影响地图，长序列中全局一致性逐渐丧失。DGS-SLAM 为每个高斯记录其创建时的唯一关键帧 ID（$k \mapsto \mathrm{ID}(\mathcal{G}_i)$）。当当前帧的可见高斯集合中存在来自旧关键帧 $k_l$ 的高斯时，系统判定发生了回环，将该旧关键帧重新纳入优化窗口，实现当前位姿与历史地图的联合优化。消融实验表明，移除回环感知机制后，在 Bonn 数据集的大规模场景中跟踪误差从 7.3 cm 上升至 8.0 cm；移除鲁棒掩码则误差升至 8.3 cm。在 TUM-RGBD 动态序列上，DGS-SLAM 的跟踪精度（ATE 3.0 cm）优于 DG-SLAM，渲染 PSNR（20.63 dB）在动态场景的 GS-SLAM 方法中达到当时最佳。

Gassidy（2024.11）~\cite{2411_15476} 沿着 DG-SLAM 的"检测动态物体-滤除-重建静态场景"路线，但将动态检测从像素级掩码转变为\emph{物体级渲染损失流分析}，核心洞察是：在位姿优化迭代过程中，静态物体的渲染损失随迭代稳定下降（因为静态物体与地图逐步对齐），而动态物体的损失持续偏高且剧烈波动（因为它们的位置在帧间发生了不可预测的变化）。具体而言，Gassidy 利用 YOLO 实例分割将场景分解为若干物体区域和背景区域（而非 DG-SLAM 的像素级掩码），对每个物体独立计算其光度-几何联合渲染损失。在粗跟踪阶段，仅用背景损失优化位姿（40 次迭代），同时记录每个物体在每次迭代中的损失变化——这一时序损失序列即"渲染损失流"。随后用高斯混合模型（GMM）对渲染损失流分类：损失持续下降的物体被判定为静态，损失波动不收敛的物体被判定为动态。

这种基于损失演化模式而非单帧残差的检测策略使 Gassidy 能有效识别微小运动的物体（其损失波动在多次迭代中被放大），不依赖预定义的动态类别语义标签，且采用自适应权重根据深度图质量动态调节光度与几何损失的比重以应对真实场景中光照不均导致的深度噪声。被判定为动态的物体通过哈希表管理生命周期，需连续三帧被分类为静态才重新纳入优化，避免误判导致的不可逆信息丢失。在 TUM 和 BONN 数据集上，Gassidy 的相机跟踪精度比现有方法提升高达 97.9\%，地图质量提升高达 6\%。GARAD-SLAM（2025.02）~\cite{2502_03228} 直接在高斯上进行动态分割并映射回前端获取动态点标签，通过高斯金字塔网络实现从 3D 高斯到 2D 前端的动态标签传播。对动态标记的高斯施加渲染惩罚而非简单剪枝以避免不可逆的错误移除。

\subsubsection{不确定性驱动的动态处理}

本路线介绍利用不确定性估计区分静态和动态区域的方法，以渲染残差和概率建模替代显式语义检测。

显式检测动态物体需要语义先验或运动模型，而利用不确定性估计区分静态和动态区域是一种更优雅的策略——静态区域通常具有低渲染残差和高不确定性一致性，动态区域则相反。

UP-SLAM（2025.05）~\cite{2505_22335} 沿着 DG-SLAM 的"运动检测-掩蔽-建图"路线继续推进，但将每个环节从硬决策升级为概率驱动的软决策，从根本上改变了动态 SLAM 的设计范式。其核心改进可归纳为四个方面：用不确定性估计替代硬掩码、用并行架构替代串行管线、用概率八叉树替代手工阈值管理、用 DINO 特征蒸馏替代原始 RGB 表征。

\paragraph{锚点回顾与改进动机。}
DG-SLAM 通过深度变形一致性+语义分割生成二值运动掩码，在跟踪和建图中直接排除动态像素。这种"硬掩码"策略存在两个局限：其一，语义分割模型只能识别训练集中预定义的动态类别（如行人、车辆），遇到开放集动态物体（如被推动的石头、飘动的气球、机械狗）便无法泛化；其二，掩码一旦生成便固定不变，无法根据后续帧的观测信息进行修正。UP-SLAM 的设计动机正在于此——用\emph{逐像素连续不确定性}替代二值掩码，使动态区域检测从"是否运动"的二元判断转变为"多大可能运动"的概率估计，从而在保持鲁棒性的同时大幅提升泛化能力。此外，DG-SLAM 的跟踪与建图紧密耦合，建图必须在跟踪之后执行，限制了实时性。UP-SLAM 将运动掩码生成与地图优化解耦，使跟踪和建图可完全并行运行。

\paragraph{免训练多模态不确定性估计。}
UP-SLAM 的动态检测不依赖任何语义标签或预训练模型，而是融合三类残差信号估计逐像素运动不确定性：\textbf{光度残差}衡量渲染颜色与观测颜色的一致性，\textbf{几何残差}衡量渲染深度与观测深度的偏差，\textbf{语义残差}利用 DINOv2 特征衡量渲染语义特征与观测语义特征的相似度。与 DG-SLAM 的深度变形一致性不同，UP-SLAM 的多模态残差由建图线程反向传播到跟踪线程——当某个区域的渲染残差持续偏高时，说明该区域的 3D 高斯无法被当前地图很好地解释，大概率属于动态物体。这三类残差通过加权融合生成逐像素不确定性图，取代了 DG-SLAM 中的二值运动掩码。这种"训练无关"（training-free）的优化方式无需额外训练语义分割网络，且天然支持开放集动态物体——无论物体类别是否在预定义列表中，只要其多模态残差表现异常就会被标记为高不确定性。

\paragraph{概率八叉树自适应高斯管理。}
DG-SLAM 的增删策略依赖多个手工阈值（累积不透明度阈值、尺度阈值、纵横比阈值），阈值的选取对场景敏感且需要经验调整。UP-SLAM 采用\textbf{概率八叉树}（probabilistic octree）管理高斯基元，从根本上消除了对手工阈值的依赖。具体而言，3DGS 场景被压缩为结构化的锚点（anchor）表示，每个锚点编码为多个浅层 MLP 的输入，由 MLP 解码出对应高斯点的属性（颜色、不透明度、尺度等）。概率八叉树根据锚点的观测历史和优化状态自适应地决定分配新锚点或剪枝冗余锚点：当动态物体移除后在其位置产生的冗余高斯会被自动剪枝，而静态区域的新视角则触发新锚点分配。这种机制不仅无需手工调参，还将模型大小显著压缩——消融实验表明概率锚点更新模块将模型大小降至原来的约 1/3，使其更适合嵌入式平台部署。

\paragraph{DINO 特征蒸馏与时序编码。}
为增强不确定性预测的鲁棒性和地图的语义丰富度，UP-SLAM 在高斯场中蒸馏 DINOv2 特征。与直接存储高维 DINO 特征不同，UP-SLAM 为每个高斯仅维护低维特征（$N_l=16$），通过浅层 MLP $F_d$ 将其解码为高维 DINO 特征空间。这种压缩-解码策略既节省存储又保留了语义表达能力。此外，\textbf{时序编码器}利用正弦位置编码将帧间时间信息嵌入到 MLP 中，使同一锚点在不同时刻的特征表示能够动态适应——这对于动态场景尤为重要，因为动态物体的外观和运动模式随时间变化，时序信息帮助 MLP 生成更准确的属性预测。消融实验显示，时序编码器平均提升 PSNR 1.4 dB。

\paragraph{并行跟踪与建图架构。}
DG-SLAM 的跟踪（粗+精两级）和建图在串行管线中执行：跟踪完成后才能更新地图，地图更新后才能生成下一帧的残差反馈。UP-SLAM 将运动掩码生成从建图优化中解耦——建图线程独立运行 3DGS 优化并生成多模态残差，跟踪线程基于最新残差实时估计不确定性并完成位姿估计，两个线程无需互相等待。这种并行化设计使系统达到 12 Hz 的处理帧率，满足机器人实时定位需求，且定位精度比 DG-SLAM 提升 59.8\%，渲染 PSNR 提升 5.47 dB，同时产生无伪影的可复用静态地图。

WildGS-SLAM（2025.04）~\cite{2504_03886} 沿 DG-SLAM 的"检测动态-排除-重建静态"路线，但将动态检测从\emph{二值掩码}替换为\emph{逐像素连续不确定性}，用 DINOv2 特征驱动的浅层 MLP 在线预测不确定性图，无需语义分割标签。这一设计的动机在于：DG-SLAM 的语义分割仅能识别预定义类别，深度变形掩码在物体边缘易误判，两者取交集仍会漏检或误检。WildGS-SLAM 的不确定性 MLP 随流式帧在线训练——对每帧提取 DINOv2 特征（采用注入了 3D 感知的微调版本），经浅层 MLP 解码为逐像素不确定性图，以修改版 SSIM 损失、方差正则项和深度一致性损失联合训练。该不确定性图在两个环节同时发挥作用：在跟踪中，基于 DROID-SLAM 的光流跟踪框架引入不确定性加权的密集光束法平差（用马氏距离替代标准 $L_2$ 范数，使高不确定性像素对位姿优化的影响最小化），并额外加入 Metric3D 单目深度正则化以稳定在线训练早期不确定性尚未收敛时的位姿估计；在建图中，不确定性图作为渲染光度损失和深度损失的像素级权重，自动降低动态区域的优化贡献。值得注意的实现细节是梯度流分离——3DGS 地图优化与不确定性 MLP 训练使用各自独立的损失函数，互不回传梯度，避免两者互相干扰导致训练不稳定。系统仅需单目 RGB 输入，在多个动态环境数据集上实现无伪影视角合成。

DAGS-SLAM（2026.02）~\cite{2602_21644} 沿 DG-SLAM 的动态场景 SLAM 路线，但将核心表示从\emph{硬二值掩码}改为\emph{逐高斯软运动概率}（Motion Probability, MP）时序状态，同时引入不确定性感知调度器按需触发语义计算，在保持实时性的前提下降低计算开销。

\paragraph{从硬掩码到软 MP 状态。}
DG-SLAM 通过运动掩码 $M_j$ 将像素分为动态/静态两类，用二值门控排除动态区域。这种硬决策在遮挡和运动边界处容易误判——将静态结构误标为动态导致地图空洞，或将动态物体漏标导致伪影。DAGS-SLAM 的核心改进是为每个高斯 $G$ 维护一个连续 MP 状态 $M_G \in [0, 1]$，通过指数平滑在时间维度上累积更新：$M_G^{t} = (1-\eta_G^{t}) M_G^{t-1} + \eta_G^{t} \hat{M}_G^{t}$。关键在于更新率 $\eta_G^{t}$ 本身是自适应的——当观测不确定性高或几何一致性差时 $\eta_G^{t}$ 自动降低，避免对噪声观测过度反应；反之则快速响应。MP 的瞬时观测 $\hat{M}_G^{t}$ 由两条互补线索融合生成：轻量 YOLOv8 实例分割提供的语义先验（可快速识别已知动态类别如行人、车辆），以及基于历史关键帧重投影误差的几何不一致性线索（可捕获语义模型无法识别的非常规移动物体）。两者以几何可靠性为权重进行特征级融合，再通过多尺度高斯核聚合到高斯级别。最终的动态标签不是硬阈值，而是取当前帧所有可见高斯 MP 的中位数作为自适应阈值——当场景中动态物体增多时阈值自动升高，反之降低，避免了固定阈值在不同场景下的过度敏感或迟钝。

\paragraph{不确定性感知按需语义调度。}
DG-SLAM 对每一帧/关键帧都运行语义分割，语义模块的计算开销往往占系统总耗时的主导部分。DAGS-SLAM 的观察是：在连续帧之间场景的语义结构变化缓慢，无需每帧刷新分割结果。它设计了联合触发信号 $S^t = w_u U^t + w_r R^t$，其中 $U^t$ 为帧级 MP 不确定性中位数（反映当前 MP 估计是否可靠），$R^t$ 为跟踪重投影残差（反映几何一致性）。仅当 $S^t$ 超过阈值时才调用实例分割，否则复用上一次分割结果，同时以更保守的更新率传播 MP。这一策略在资源受限的移动/边缘设备上尤其重要——语义计算不再是瓶颈，系统可在保持跟踪精度的同时实现实时吞吐。

\paragraph{MP 引导的前后端协同优化。}
DG-SLAM 的运动掩码仅在前端跟踪时排除动态像素。DAGS-SLAM 将 MP 作为统一接口贯穿前后端：在前端，每个高斯的 MP 传播至其关联的图像特征点，用于动态感知的对应选择——高 MP 的特征点被降权或排除，避免动态物体上的错误匹配影响位姿估计。在后端，MP 用于构建像素级静态权重 $w^t(u) = \sum_k \alpha_k^t(u) \overline{M}_k^t$（$\overline{M}_k^t = 1 - M_k^t$ 为静态置信度），使被动态高斯主导的像素在光度损失和深度损失中自动降权。此外，DAGS-SLAM 引入 MP 加权运动一致性正则 $\mathcal{L}_{\mathrm{MP}} = \sum_k \overline{M}_k^t \|\boldsymbol{\mu}_k^t - T_{t \to t-1} \boldsymbol{\mu}_k^{t-1}\|_2^2$，约束高静态置信度的高斯在帧间保持位置稳定，而对动态高斯施加较弱约束使其自然适应运动。

\paragraph{对极几何恢复被误抑制的静态结构。}
MP 和语义线索在运动边界附近可能不可靠，误将静态结构抑制为动态。DAGS-SLAM 对被 MP 抑制的高斯施加对极几何验证：计算其关联特征点的点到极线距离，若距离小于阈值则恢复为静态，并用恢复的匹配引导增密。这种"软抑制-硬验证"的双重策略在抑制动态伪影的同时减少了地图空洞，特别适用于物体边缘等语义与几何线索均模糊的区域。

\subsubsection{4D 建模与动态-静态联合重建}

本路线介绍利用 4DGS 时间维度同时重建静态背景和动态前景的方法，使 SLAM 系统从"避动态"走向"拥抱动态"。

最前沿的工作不再简单移除动态物体，而是同时重建静态背景和动态前景，利用 4DGS 的时间维度建模动态内容，使 SLAM 系统从"避动态"走向"拥抱动态"。

DG-SLAM 开创的"检测动态物体-移除-仅重建静态背景"路线有一个根本局限：丢弃了动态物体的运动信息，当动态物体占场景比例较大时地图会变得稀疏。D$^2$GSLAM（2025.12）~\cite{2512_09411} 在此基础上实现了范式转变——从"避动态"变为"拥抱动态"，是首个在 SLAM 框架中同时利用 3DGS 重建静态场景和 4DGS 建模动态物体的系统。核心洞察是：DG-SLAM 用单一 3DGS 表示整个场景，对静态部分绰绰有余但对动态部分力不从心（缺少时间维度）；将表示拆分为静态 3DGS + 动态 4DGS，两类内容各得其所。在跟踪精度上比现有动态 SLAM 方法提升达 42\%，静态重建精度平均提升 36\%。

\paragraph{复合高斯表示。}
与 DG-SLAM 仅用单一 3DGS 表示整个场景不同，D$^2$GSLAM 将表示拆分为两支：静态部分沿用标准 3D 高斯 $G_s(x) = e^{-\frac{1}{2}(x-\boldsymbol{\mu}_s)^T \boldsymbol{\Sigma}_s^{-1}(x-\boldsymbol{\mu}_s)}$，动态部分将 3D 高斯扩展到四维空间 $\mathbf{x} = (x, y, z, t) \in \mathbb{R}^4$，得到 4D 高斯 $G_d(\mathbf{x}) = e^{-\frac{1}{2}(\mathbf{x}-\boldsymbol{\mu}_d)^T \boldsymbol{\Sigma}_d^{-1}(\mathbf{x}-\boldsymbol{\mu}_d)}$。直观地说，4D 高斯就像一根"时间腊肠"——在时间轴上也有粗细，代表动态物体在不同时刻的空间分布范围。整个复合高斯集合记为 $\mathcal{G} = \{G_s^{i_s}\}_{i_s=1}^{N_s} \cup \{G_d^{i_d}\}_{i_d=1}^{N_d}$。

4D 协方差采用分块参数化 $\boldsymbol{\Sigma}_d = \begin{pmatrix} A_{3 \times 3} & B_{3 \times 1} \\ B^T & C_{1 \times 1} \end{pmatrix}$，其中 $A$ 控制空间扩散，$C$ 控制时间跨度（$C$ 越大，高斯在时间轴上存活越久），$B$ 编码空间-时间耦合关系。给定观测时刻 $t$，4D 高斯投影为 3D 高斯：时间维度的衰减因子 $e^{-\frac{1}{2}(t-t_d)^2/C}$ 控制动态高斯的时间影响范围，空间中心的时变偏移 $\boldsymbol{\mu}_s(t) = (x_d, y_d, z_d)^T + \boldsymbol{v}(t - t_d)$ 建模物体运动速度 $\boldsymbol{v}$（即高斯中心随时间线性移动），投影后的空间协方差 $\boldsymbol{\Sigma}_s = A - BC^{-1}B^T$ 比原始空间部分更紧凑（因为部分不确定性已被时间维度吸收）。4D 旋转使用 rotors（几何代数中的旋转算子，而非四元数）表征以支持高维旋转建模。渲染时，4D 高斯先投影到当前时刻的 3D 高斯，再与静态 3D 高斯一起按深度排序进行 $\alpha$-混合，渲染像素颜色为：

\begin{equation*}
C_p(t) = \sum_{i_d=1}^{N_d} c_d^{i_d} \alpha_d^{i_d}(t) \prod_{j=1}^{i_d-1}(1 - \alpha_d^{j}(t)) + \sum_{i_s=1}^{N_s} c_s^{i_s} \alpha_s^{i_s} \prod_{j=1}^{i_s-1}(1 - \alpha_s^{j}),
\end{equation*}

深度 $D_p(t)$ 按相同混合权重用深度值替换颜色即可获得。

\paragraph{几何提示的动态分离。}
DG-SLAM 的动态检测依赖"深度变形一致性 + 语义分割"的组合——语义分割能快速识别行人、车辆等已知类别，但遇到语义模型训练集之外的动态物体（如被推动的椅子）便失效。D$^2$GSLAM 放弃语义分割，改为纯几何提示策略：利用高斯基元本身的几何一致性检测动态区域，无论物体是刚体还是非刚体，只要运动就会产生几何不一致。这一选择使检测能力不受预定义类别限制。该模块分两个阶段工作：

第一阶段是\textbf{粗检测}，利用两种残差互补定位潜在动态区域。（1）几何一致性残差 $R_g = \| \pi(T_i^{i+1}p_i) - D_{i+1}\|$：将上一帧的深度反投影为 3D 点云，经相对位姿变换后投影到当前帧，与观测深度比较——静态区域投影深度应与观测一致，动态区域则出现偏差。对 $R_g$ 提取连通域并过滤小面积噪声后，即可判断场景中是否存在动态物体。（2）预测残差 $R_p = \|D_{i+1} - \hat{D}_{i+1}(G_s, T_w^{i+1})\|$：从当前静态高斯地图渲染下一帧深度，与观测深度比较——静态区域渲染准确，动态区域出现显著差异。对 $R_p$ 二值化得到粗动态掩码。

第二阶段是\textbf{精化}，从粗掩码中提取连通域质心作为点提示（point prompts），经几何一致性残差连通域交叉验证过滤误检后，连同 YOLO-world 检测框作为框提示一并输入 MobileSAM，生成精细运动掩码。最终掩码融合粗检测和精化结果。整个流程仅需 88.2ms 即可生成准确掩码，比 DG-SLAM（163ms）快近一倍。

\paragraph{复合地图扩展。}
DG-SLAM 检测到动态物体后直接移除，仅重建静态背景——这丢弃了运动信息，动态物体占比大时地图会变得稀疏。D$^2$GSLAM 反其道而行之，将动态物体纳入地图：首帧所有物体均以 3D 高斯初始化（因为单帧无法判断动态性）；后续帧检测到动态物体后，通过 KNN 搜索判断其身份——若在世界坐标高斯地图中找到近邻，说明该物体之前是静态的，将对应 3D 高斯扩展为 4D 高斯（添加时间维度）；若未找到近邻，说明是新出现的物体，直接初始化为 4D 高斯。这样地图随动态物体的出现逐步增长，低不透明度高斯在优化中被剪枝。

\paragraph{渐进式位姿精化。}
DG-SLAM 的跟踪采用粗-精两级策略，始终只用静态分量——粗跟踪用 DROID-VO 获取初始位姿，精跟踪用高斯泼溅渲染在运动掩码过滤后的可靠像素上优化。D$^2$GSLAM 将其扩展为三级渐进策略：第一步沿用 DG-SLAM 的思路，仅利用静态场景分量估计初始位姿（因为动态物体在早期尚未充分建模，会干扰位姿估计），通过最小化渲染 RGB-D 与观测的差异在运动掩码屏蔽的静态区域上优化；第二步是 D$^2$GSLAM 独有的贡献——将动态场景信息纳入联合精化，在动态物体占比大的场景中，仅靠静态元素提供的观测约束不足以可靠估计位姿，动态分量补充了几何约束；第三步执行全局光束法调整（BA），通过维护因子图 $(V, E)$ 建立帧间共视关系，仅基于静态分量进行全局位姿优化（因为动态物体天然缺乏多视角一致性），运动掩码用于抑制 BA 中的动态区域干扰。

\paragraph{运动一致性损失与回顾帧优化。}
DG-SLAM 的优化只有光度损失和深度损失，因为它的地图全是静态的，不存在运动建模需求。D$^2$GSLAM 由于引入了动态 4D 高斯，优化目标也随之分叉：静态损失 $\mathcal{L}_\text{static} = \lambda_c \mathcal{L}_c + \lambda_d \mathcal{L}_d$ 与 DG-SLAM 类似，而新增的动态损失 $\mathcal{L}_\text{dynamic} = \lambda_c \mathcal{L}_c + \lambda_\text{sp} \mathcal{L}_\text{sp} + \lambda_\text{mc} \mathcal{L}_\text{mc}$ 包含两项 DG-SLAM 所没有的时序约束：

\textbf{空间平滑损失} $\mathcal{L}_\text{sp} = \frac{1}{N_d}\sum_{i=1}^{N_d} \| \boldsymbol{v}_i - \frac{1}{M}\sum_{j \in Q_i} \boldsymbol{v}_j \|$：空间邻域内动态高斯的运动速度应当一致。选取每个动态高斯的 $M$ 个最近邻高斯，惩罚其速度偏离邻域平均速度的程度——这防止相邻高斯朝不同方向运动，保证物体的整体运动一致性。

\textbf{运动一致性损失} $\mathcal{L}_\text{mc} = \|C(G_d(t_{i+1}), T_{i+1}, t_{i+1}) - C_{i+1}\|$：利用运动的时间连续性。由于相机帧率固定，可以基于历史动态建模信息预测下一帧的运动，将预测渲染结果与实际观测比较来优化高斯参数。这种"用过去预测未来"的约束使有限顺序视角下的动态建模更加准确。

此外，SLAM 系统的逐帧优化策略会导致 4D 高斯过拟合最新观测、遗忘历史运动——即灾难性遗忘。D$^2$GSLAM 引入动态回放缓冲区（dynamic replay buffer）维护历史时刻的运动信息，在每次建图优化时回溯历史帧持续精化动态 4D 高斯，这一\textbf{回顾帧优化}（Retrospective Frame Optimization）策略使动态建模在时间和空间上保持一致。消融实验表明，去除回顾帧优化后人脸、手臂和衣物细节会丢失，先前重建的静态背景也会变模糊。

D$^2$GSLAM 将表示拆分为静态 3DGS + 动态 4DGS 两支，D4DGS-SLAM（2025.04）~\cite{2504_04844} 则走另一条路线——用\textbf{统一 4DGS 表示整个场景}（包括静态部分），不做动静拆分。与 D$^2$GSLAM 的纯几何提示动态检测不同，D4DGS-SLAM 引入预训练的长效任意点跟踪模块 LEAP（Long-term Effective Any Point Tracking）作为动态感知前端。LEAP 对每帧提取的锚点（梯度引导的空间均匀采样点）跨序列追踪，为每个点输出三重信息：\textbf{可见性}（是否被遮挡）、\textbf{动态标签}（属于静态还是动态物体）和\textbf{可见性可靠性}（可见性预测的不确定性）。跟踪时，同时满足可见性 $V_t \geq \gamma_v$、动态标签 $m_{d,t} < \gamma_d$、可靠性 $\Phi_t \geq Q(\gamma_u)$ 且观测次数不少于 $\gamma_{\mathrm{track}}$ 的点才被保留用于位姿优化——其余全部过滤。这使得跟踪完全不依赖语义分割或几何残差，而由学习到的长期运动模式直接判别，在 BONN 数据集上 ATE RMSE 为 5.1\,cm，比 DG-SLAM（5.5\,cm）低 7.3\%。

\paragraph{动态感知的差异化各向同性正则化。}
D$^2$GSLAM 通过空间平滑损失和运动一致性损失来约束动态高斯的运动行为。D4DGS-SLAM 的策略则更为简洁——根据 LEAP 提供的动态标签对高斯施加不同尺度的各向同性正则化：强动态高斯的正则项 $E_{\mathrm{iso}}^{\mathrm{dynamic}} = \sum_{i=1}^{|\mathcal{G}^{\mathrm{dynamic}}|} \|\vec{S}_i^{4D} - \tilde{\vec{S}}_i^{4D} \cdot \mathbf{1}\|_1$ 同时约束四个时空尺度参数（$s_x, s_y, s_z, s_t$），鼓励各维度尺度趋向均值，防止动态高斯在某一维度过度拉伸；弱动态（即静态）高斯的正则项 $E_{\mathrm{iso}}^{\mathrm{static}} = \sum_{i=1}^{|\mathcal{G}^{\mathrm{static}}|} \|\vec{S}_i^{3D} - \tilde{\vec{S}}_i^{3D} \cdot \mathbf{1}\|_1$ 仅约束三个空间尺度参数，时间维度不受限制。这种差异化正则化使静态区域的 4D 高斯在时间轴上自动展开为宽而平的分布（覆盖较长时段），而动态区域的 4D 高斯保持紧凑的时间宽度以精确捕获运动轨迹。消融实验表明，去除动态感知模块后 ATE 从 3.6\,cm 暴增至 27.9\,cm，PSNR 从 27.89 降至 20.23；去除 4DGS（退回 3DGS）后 ATE 为 7.2\,cm 但 PSNR 仅 18.23——两者缺一不可。

\paragraph{实验表现。}
在 BONN 数据集上，D4DGS-SLAM 渲染质量显著超越静态 3DGS 基线：平均 PSNR 28.71（SplaTAM 17.22、MonoGS 18.97），SSIM 0.854（SplaTAM 0.685），LPIPS 0.259（SplaTAM 0.304）。在更复杂的 TartanAir-Shibuya 数据集上，ATE 平均 3.6\,cm，比 SplaTAM（73.2\,cm）降低 95.1\%；PSNR 22.75（+10.3）、SSIM 0.705（+0.321）。 MonoGS 和 SplaTAM 等静态 3DGS 方法在动态物体区域产生严重的重影伪影（半透明副本），而 D4DGS-SLAM 由于时间维度的显式建模，渲染结果在帧间保持时间一致性。运行速度为 BONN 1.5\,FPS、TartanAir-Shibuya 0.9\,FPS，尚未达到实时。RU4D-SLAM（2026.02）~\cite{2602_20807} 引入不确定性感知的模糊图像合成和动态场景重建，将原本为静态场景设计的逐像素不确定性建模扩展到模糊图像处理，语义引导的重加权机制在动态场景中显著优于已有方法。

Flow4DGS-SLAM（2026.04）~\cite{2604_22339} 利用光流引导动态区域分解——通过拟合相机自运动模型分解光流生成类别无关的运动掩码，同时提供流引导的相机位姿初始化。动态 3DGS 通过在关键帧显式建模时间中心加速训练，用 3D 场景流先验传播，并通过高斯混合模型（GMM）自适应学习复杂动态特性。Dy3DGS-SLAM（2025.06）~\cite{2506_05965} 是首个面向动态场景的单目 RGB 3DGS SLAM 方法，通过概率模型融合光流掩码和深度掩码获取融合动态掩码。其新颖的运动损失约束位姿估计网络的跟踪，动态像素的渲染损失消除瞬态干扰和遮挡。

DyPho-SLAM（2025.09）~\cite{2509_00741} 集成先验图像信息生成精化掩码，自适应特征提取策略在移除动态障碍后增强优化约束。PG-SLAM（2024.11）~\cite{2411_15800} 同时建图动态前景（包括非刚体人体和刚体物品）和静态背景，利用静态背景和动态前景增加观测进行噪声补偿。LVD-GS（2025.10）~\cite{2510_22669} 融合开放世界分割和隐式残差约束生成细粒度动态掩码，由 DINO-Depth 特征的不确定性估计引导。DynaGSLAM（2025.03）~\cite{2503_11979} 是首个在动态场景中同时实现在线 GS 渲染、跟踪和运动预测的实时系统——而非简单地检测和移除动态区域。GGD-SLAM（2026.04）~\cite{2604_12837} 利用可泛化运动模型通过序列注意力机制提取动态语义特征，通过 FIFO 队列管理传入帧，设计干扰自适应 SSIM 损失增强动态环境中的系统韧性。

\begin{takeaway}
动态场景 SLAM 从"滤除动态物体"演进为"同时重建静态与动态内容"：4D 高斯引入时间维度显式建模运动轨迹，光流引导的动态区域分解和差异化正则化使系统能在动态环境中保持稳定跟踪。
\end{takeaway}

% ======================================================================


\section{SLAM系统优化与工程}
\label{sec:15_7}

将 3DGS-SLAM 从实验室原型推向实际应用，需要解决一系列工程层面的挑战：运动模糊与图像去模糊、回环检测与全局一致性、稀疏视角和内存效率等。本节梳理 23 篇论文，聚焦系统级优化。

\subsection{系统优化与工程}
\label{ss:15_7a}

本节介绍三条技术路线：（1）运动模糊处理；（2）回环检测与全局一致性；（3）加速与稀疏处理。

\subsubsection{运动模糊处理}

本路线介绍在 3DGS-SLAM 中正确建模和处理运动模糊的方法，利用物理模糊模型消除退化影响。

运动模糊的物理本质是曝光时间内相机运动导致的图像积分效应。正确建模这一过程是处理运动模糊的关键。

手持设备或快速运动场景中的运动模糊是 3DGS-SLAM 的主要退化因素，其影响通过两个机制恶性循环：运动模糊破坏了密集 SLAM 所依赖的光度一致性假设，导致位姿估计不准；不准确的位姿又引入多视角几何不一致，进一步恶化地图重建质量。MBA-SLAM（2024.11）~\cite{2411_08279} 是首个系统性地将运动模糊的物理成像过程嵌入完整 RGB-D SLAM 管线的方法，其核心思路可用一句话概括：\textbf{与其从模糊图像中去除模糊，不如从清晰渲染中合成模糊}。

\textbf{物理模型与问题定义。}
运动模糊图像 $I(\mathbf{u})$ 本质上是曝光时间 $\tau$ 内所有虚拟清晰图像 $\hat{I}_t(\mathbf{u})$ 的积分：

\begin{equation}
    I(\mathbf{u}) = \frac{1}{\tau}\int_0^{\tau} \hat{I}_t(\mathbf{u}) \, dt
    \label{eq:ch15_blur_formation}
\end{equation}

实践中用离散近似：对 $n$ 个均匀采样的虚拟位姿分别渲染清晰图像 $\hat{I}_i(\mathbf{u})$，取平均得到合成模糊图像 $I(\mathbf{u}) \approx \frac{1}{n}\sum_{i=0}^{n-1} \hat{I}_i(\mathbf{u})$。这一离散近似是后续所有计算的基础——虚拟位姿越多，模糊建模越精确，但计算开销也越大（消融实验表明 13 个虚拟图像在精度与效率间取得最佳平衡）。

\textbf{前端跟踪：运动模糊感知的直接对齐。}
传统直接法跟踪假设参考帧与当前帧满足光度一致性，但运动模糊使得对应点外观差异巨大，这一假设完全失效。MBA-SLAM 选择重模糊（re-blurring）而非去模糊作为应对策略——重模糊只需沿轨迹渲染并取平均，而去模糊需要从严重退化图像中恢复细节，前者在数学上更良态且对强模糊更具鲁棒性。具体而言，跟踪器独立参数化曝光起始和结束时刻的相机位姿 $\mathbf{T}_{\mathrm{start}}$ 和 $\mathbf{T}_{\mathrm{end}}$（而非仅参数化末端位姿再线性插值到前一帧，这种独立参数化对非线性轨迹——如手持和头戴设备的突然转向——更具鲁棒性），通过 $SE(3)$ 李代数上的线性插值得到曝光期间任意时刻 $t$ 的虚拟位姿：

\begin{equation}
    \mathbf{T}_t = \mathbf{T}_{\mathrm{start}} \exp\!\Big(\frac{t}{\tau} \log(\mathbf{T}_{\mathrm{start}}^{-1} \mathbf{T}_{\mathrm{end}})\Big)
    \label{eq:ch15_se3_interp}
\end{equation}

在像素传递环节，MBA-SLAM 假设采样像素所在的局部区域近似为前向平行面（fronto-parallel plane），从而将当前模糊帧中的像素通过平面-射线求交反向投影至各虚拟位姿，再映射回参考帧以双线性插值获取强度值，取 $n$ 个虚拟视角的平均值作为重模糊像素：$\hat{B}_{\mathrm{cur}}(\mathbf{x}) = \frac{1}{n}\sum_{i=0}^{n-1} I_{\mathrm{ref}}(\mathbf{x}_{i\tau/(n-1)})$。跟踪器通过最小化真实模糊像素与重模糊像素之间的光度误差 $\mathbf{T}_{\mathrm{start}}^*, \mathbf{T}_{\mathrm{end}}^* = \arg\min_{\mathbf{T}_{\mathrm{start}}, \mathbf{T}_{\mathrm{end}}} \sum_{i}\|B_{\mathrm{cur}}(\mathbf{x}_i) - \hat{B}_{\mathrm{cur}}(\mathbf{x}_i)\|$ 联合优化起止位姿。优化采用从粗到精策略，辅以 Huber 鲁棒损失函数和高梯度像素稀疏采样（$9 \times 9$ 局部块），整个前向和反向传播过程以纯 CUDA 实现，支持实时跟踪速度。

\textbf{后端建图：双表示联合优化。}
后端建图器将相同的模糊形成模型应用于场景表示的联合优化，同时支持 NeRF（三平面表示）和 3DGS 两种后端：NeRF 版本通过体渲染沿光线采样 3D 点，使用 7 个虚拟图像；3DGS 版本通过光栅化渲染高斯集合，使用 13 个虚拟图像。两种表示各有取舍——NeRF 版本在采样稀疏像素后可达约 22.6 FPS 的近实时速度，但渲染保真度较低；3DGS 版本全图渲染质量更高（尤其在纹理丰富和边缘区域），但速度约 2 FPS。优化目标包括光度颜色损失 $\mathcal{L}_c = \frac{1}{|R|}\sum_{\mathbf{x} \in R}\|B(\mathbf{x}) - B^{gt}(\mathbf{x})\|$（合成模糊图像与真实模糊输入之差）和几何深度损失 $\mathcal{L}_d = \frac{1}{|R|}\sum_{\mathbf{x} \in R}\|D(\mathbf{x}) - D^{gt}(\mathbf{x})\|$（从轨迹中点渲染的深度与传感器深度之差）。NeRF 版本还额外使用自由空间损失和 SDF 损失促进几何重建收敛；3DGS 版本则加入 SSIM 损失和高斯尺度正则化 $\mathcal{L}_{\mathrm{reg}} = \frac{1}{|G|}\sum_{\mathbf{g} \in G}\max\{\max(\mathbf{S}_g)/\min(\mathbf{S}_g), r\} - r$（约束高斯核的长短轴比值不超过阈值 $r$，防止过度扁平化）。

\textbf{实验结果。}
在合成模糊数据集 ArchViz 上，MBA-SLAM 的 GS 版本取得平均 ATE RMSE 0.84 cm、PSNR 28.82 dB，NeRF 版本取得 ATE RMSE 1.69 cm、PSNR 25.08 dB，均大幅超越所有 NeRF-SLAM 和 Gaussian-SLAM 对比方法（次优方法 CoSLAM 的 ATE 为 8.04 cm，Photo-SLAM 的 ATE 为 2.41 cm）。GS 版本在渲染质量上优势尤为显著：PSNR 比次优方法高出约 4 dB，F1 分数（$<$1 cm）达 88.41\%，远超 CoSLAM 的 58.50\%。在真实自采集的低光照运动模糊数据（使用 RealSense RGB-D 相机和 Vicon 动捕系统采集）上同样取得最优跟踪性能，并能从严重模糊序列中恢复高质量清晰图像。值得注意的是，即使在标准清晰数据集（Replica、ScanNet、TUM RGB-D）上不做模糊建模，MBA-SLAM 的跟踪精度也超越大多数已有方法，表明其运动模糊感知框架在非模糊场景下不引入性能退化。MBA-SLAM 整合了作者此前三项会议工作的成果：MBA-VO（ICCV 2021 oral）的运动模糊感知视觉里程计、BAD-NeRF（CVPR 2023）的模糊感知 NeRF 优化和 BAD-Gaussians（ECCV 2024）的模糊感知 3DGS 优化，首次将三者统一为完整的 SLAM 管线，并将 BAD-NeRF 的原始 NeRF 表示替换为训练效率提升约 100 倍的三平面表示。

MBA-SLAM 奠定了运动模糊物理建模的基本范式后，后续工作从不同角度进行了拓展。其中，Deblur-SLAM（2025.03）~\cite{2503_12572} 提出了一个核心问题：MBA-SLAM 的帧到模型（frame-to-model）单阶段跟踪依赖 RGB-D 深度传感器且缺乏全局一致性维护，在纯 RGB 输入和长序列场景下面临床移累积风险。Deblur-SLAM 从三个维度对 MBA-SLAM 的范式进行了增量扩展。

\textbf{从单阶段到双阶段跟踪。}
MBA-SLAM 仅采用帧到模型的直接对齐策略，用重模糊合成与真实模糊帧比较来联合优化位姿。Deblur-SLAM 的设计动机是：纯 RGB 场景下缺乏深度传感器提供的几何约束，单靠帧到模型跟踪容易因初始化偏差陷入局部最优。因此，Deblur-SLAM 引入双阶段级联跟踪架构。第一阶段是帧到帧（frame-to-frame）跟踪，沿用 SplatSLAM 的光流递归网络与 DSPO（Disparity, Scale and Pose Optimization）目标，在因子图上交替执行稠密光束法平差（DBA）和单目深度约束的尺度-偏移优化，为相机位姿和视差提供精确初始化。第二阶段将帧到帧所得位姿作为子帧轨迹的中点初始值，执行帧到模型精修——复用 MBA-SLAM 的模糊形成物理模型，从 3DGS 渲染清晰虚拟子帧图像取平均，最小化合成模糊图像与观测模糊帧的光度误差。这种双阶段设计使得帧到帧阶段提供鲁棒的初始位姿（得益于光流网络的强先验），帧到模型阶段进一步精修子帧轨迹并同步更新场景表示，两者优势互补。

\textbf{从 RGB-D 到纯 RGB：单目深度引导与在线高斯变形。}
MBA-SLAM 依赖 RGB-D 传感器的深度图提供几何监督。Deblur-SLAM 在纯 RGB 设定下需要替代的深度来源。其策略是融合多视角视差与预训练 DPT 单目深度估计模型的互补信息：在多视角视差可靠的区域直接使用视差反演深度，在视差优化误差较高的像素子集上用 DPT 深度估计填补，两者通过最小二乘拟合的尺度和偏移参数对齐。这种代理深度图（proxy depth map）既用于高斯初始化（首帧假设静止以确保清晰），也用于建图阶段的几何损失计算。此外，由于在线回环闭合和全局光束法平差会持续修正关键帧位姿和代理深度图，Deblur-SLAM 引入在线非刚性变形（online non-rigid deformation）机制：当关键帧位姿更新后，沿光轴方向平移关联高斯的均值，并同步更新旋转和尺度参数，保持 3DGS 地图与修正后轨迹的一致性。这一变形机制直接借鉴自 SplatSLAM，但在 Deblur-SLAM 中承担了更关键的职责——因为模糊场景下的深度不确定性更高，轨迹修正的幅度也更大。

\textbf{从局部优化到全局一致性：回环闭合与曝光间隔建模。}
MBA-SLAM 没有显式的回环检测机制，长序列中不可避免地累积漂移。Deblur-SLAM 集成了两项全局一致性技术。其一是在线回环闭合：通过计算当前活跃关键帧与所有历史关键帧之间的平均光流幅值，当光流低于阈值且帧间时间间隔超过阈值时，在因子图中添加回环边。其二是在线全局光束法平差（global BA），构建包含所有历史关键帧的独立因子图进行全局优化，优化前对视差和位姿的尺度进行归一化以保证数值稳定性。此外，Deblur-SLAM 还引入了 MBA-SLAM 所不具备的曝光间隔（exposure gap）建模：在相邻帧的子帧轨迹之间学习两个可优化参数 $g_{\mathrm{prev}}$ 和 $g_{\mathrm{next}}$，表示曝光间隙中运动推进的程度，并通过正则化项约束子帧轨迹的端点与全局轨迹插值点之间的平移和旋转差异，确保局部子帧轨迹与全局轨迹保持一致。这一设计解决了实际拍摄中帧间曝光不连续的物理现实——相邻帧的曝光之间可能存在短暂的时间间隔，期间相机继续运动但不被任何帧捕获。

\textbf{实验验证。}
Deblur-SLAM 在合成模糊 Replica 数据集（36 帧平均生成模糊）上的渲染质量（PSNR/SSIM/LPIPS）和轨迹精度（ATE）均超越所有对比方法，包括 MBA-SLAM 和 $I^2$-SLAM。在真实数据集（TUM-RGBD、ScanNet、ETH3D，通过多帧平均引入模糊）上同样取得最优 ATE，并能渲染出与逐帧真实清晰图像质量接近的虚拟清晰子帧图像。系统运行速度约为 0.36 FPS（仅在线过程），对应虚拟子帧视角下约 1.82 FPS。局限方面，当前方法不建模动态物体的运动模糊，且对所有帧统一执行模糊建模的计算开销较高，未来可通过自适应调整虚拟视角数量来优化效率。

Unblur-SLAM（2026.03）~\cite{2603_26810} 从三个维度对 MBA-SLAM 的运动模糊范式进行了根本性拓展：将模糊类型从单一运动模糊扩展到运动模糊 + 散焦模糊，将固定计算量改为自适应于模糊程度，将纯物理轨迹建模改为前馈去模糊网络 + 逐像素模糊核的混合管线。系统基于 SplatSLAM 的纯 RGB 框架构建，含回环闭合和全局光束法平差。

\paragraph{自适应帧分类。}
MBA-SLAM 和 Deblur-SLAM 假设所有输入帧均为模糊帧，对每一帧都执行模糊建模——但实际拍摄中往往只有部分帧模糊，对清晰帧的冗余处理既浪费算力又可能引入退化。Unblur-SLAM 使用 ARNIQA 图像质量评估指标作为模糊检测器，将输入帧分为三档：\textbf{清晰帧}（低于阈值 $\tau_\text{sharp}$）直接通过 Droid-SLAM 跟踪器并赋予更高优化权重 $w_\text{sharp}$；\textbf{可去模糊帧}（前馈网络成功去模糊）进入多尺度渐进优化但不更新位姿；\textbf{严重模糊帧}（前馈网络失败）跳过跟踪器，直接通过多子帧渲染建模。这种分级策略使系统在清晰帧占多数的场景中大幅减少计算开销，而在全模糊场景（如 ArchViz）中退化为逐帧处理。

\paragraph{物理约束训练的前馈去模糊网络。}
MBA-SLAM 从 3D 表示渲染清晰图像再合成模糊——这是"由清晰到模糊"的正向建模。Unblur-SLAM 的第一阶段反其道而行，采用前馈图像去模糊网络直接从模糊输入恢复清晰图像。关键挑战是：单图像去模糊网络缺乏多视角几何约束，输出的清晰图像在 3D 空间中不一致（I$^2$-SLAM 曾尝试类似思路但因此失败）。Unblur-SLAM 的解决方案是\textbf{物理约束训练}：训练数据按模糊形成的物理过程 $B = \int_0^T(I(t) \otimes K_\text{defocus}(f, d(t))) dt$ 构造——对视频帧序列在曝光时间窗口内取平均生成运动模糊，与散焦核卷积生成散焦模糊，以中间帧作为重建目标（避免运动方向的二义性）；先在半合成数据上学习中间帧提取，再在 DPDD 散焦数据集上微调。这种训练策略使网络输出的清晰图像满足曝光中点约束，天然具有多视角 3D 一致性。

\paragraph{深度感知的逐像素模糊核建模。}
MBA-SLAM 的模糊建模是全局的——对整幅图像均匀采样 $n$ 个虚拟位姿取平均。Unblur-SLAM 引入 Blur Proposal Network（BPN）估计\textbf{逐像素}卷积核 $h(x) \in \mathbb{R}^{K \times K}$ 和混合掩码 $m(x) \in [0,1]$，将清晰渲染图像 $C$ 变换为 $C_\text{out}(x) = (1-m(x))C(x) + m(x)(C(x) \otimes h(x))$。核 $h(x)$ 本身是深度感知的：$h(x, D(x)) = h_\text{deblur}(D(x)) + \alpha(D(x)) \cdot h_\text{sharpen}(D(x))$，其中去模糊分量 $h_\text{deblur}$ 消除残余模糊，锐化分量 $h_\text{sharpen}$ 增强边缘，深度依赖因子 $\alpha$ 使远距离区域获得更强增强以补偿角分辨率限制。掩码稀疏约束 $\|m^s\|_1$ 确保仅在有必要的像素处激活模糊建模，防止过拟合传感器噪声。这种逐像素深度感知机制使 Unblur-SLAM 能处理散焦模糊（不同深度的物体模糊程度不同），而 MBA-SLAM 的全局轨迹积分无法建模这种空间变化的模糊。

\paragraph{严重模糊帧的回退策略。}
当前馈网络无法有效去模糊时（模糊程度极端），Unblur-SLAM 回退到与 MBA-SLAM 类似的物理建模策略：沿模糊轨迹生成 $N_\text{sub}$ 个虚拟相机，每个虚拟相机渲染清晰子帧后经 BPN 施加模糊核，取平均得到合成模糊图像 $\hat{Q}^s = \frac{1}{N_\text{sub}} \sum_{j=1}^{N_\text{sub}} \text{BPN}(Q_j^s, h_j^s, m_j^s)$，与真实模糊观测计算光度与深度损失。与 MBA-SLAM 的区别在于：每个子帧也经过 BPN 处理而非直接取平均，因此能同时建模运动模糊和散焦模糊的叠加效应。优化时采用三级尺度（3, 5, 9）渐进精化，并引入可学习的曝光补偿参数 $a_\text{exposure}^s, b_\text{exposure}^s$ 消除帧间亮度差异，同时通过在线高斯变形保持 3DGS 地图与修正后轨迹的一致性。

\paragraph{实验表现。}
在 Deblur-NeRF 数据集上，Unblur-SLAM 作为在线方法在运动模糊和散焦模糊两个子集上均超越 BAD-Gaussians、Deblurring 3DGS、BAGS、CoMoGaussian 等离线去模糊方法。在 ArchViz 合成极端模糊数据集上，ATE RMSE 和 PSNR 均超越 MBA-SLAM。在 TUM-RGBD 真实数据上（19 条序列），ATE RMSE 一致低于 Droid-SLAM 和 Deblur-SLAM。运行速度方面，完整管线在 TUM 数据集上平均约 3 FPS（清晰帧跳过去模糊时更快），其中前馈去模糊网络单帧推理约 600ms 是主要瓶颈。系统的主要局限是内存消耗较高，且对全模糊序列的实时性不足。

I$^2$-SLAM（2024.07）~\cite{2407_11347} 提出逆图像形成模块，可增强已有视觉 SLAM 管线对随意拍摄场景的鲁棒性。利用线性 HDR 辐射图收集测量，通过专用变量显式建模白平衡、曝光时间和相机响应函数。此外，沿相机轨迹聚合多个位姿的图像解释运动模糊。TAMBRIDGE（2024.05）~\cite{2405_19614} 集成以跟踪为中心的 ORB 视觉里程计和以建图为中心的在线 3DGS，核心是融合桥模块——通过联合优化重投影和渲染损失以及策略性视点选择实现精确位姿初始化，增强大规模场景中的渲染收敛。

\subsubsection{回环检测与全局一致性}

本路线介绍利用 3DGS 子地图表示实现回环检测和全局一致性维护的方法。

缺乏回环检测的 SLAM 系统在长序列中必然累积漂移。3DGS 的子地图表示为回环检测和全局一致性维护提供了新的技术路径。

LoopSplat（2024.08）~\cite{2408_10154} 是首个将 3DGS 表示同时用于跟踪、建图和回环闭合的耦合式 RGB-D SLAM 系统。其出发点是一个根本性问题：能否直接利用地图表示（即 3DGS）本身完成回环闭合，而无需回退到传统点云配准？此前的方法各有局限——GO-SLAM 需要代价高昂的哈希网格特征重训练来变形地图，Photo-SLAM 需要额外的 3DGS 参数优化来消化 ORB-SLAM 追踪器的位姿更新，而最相近的 Loopy-SLAM 虽采用子地图避免存储所有帧，但其回环约束仍依赖传统的 FPFH + RANSAC + ICP 全局到局部点云配准流水线，既缓慢又无法利用场景表示本身的丰富信息。LoopSplat 的回答是肯定的——它提出直接在 3DGS 子地图间进行配准以提取回环约束，从而将 3DGS 作为跟踪、建图和全局一致性维护的统一场景表示。

在 LoopSplat 中，场景以子地图（submap）为单位组织。每个子地图以 3DGS 点云 $P^s = \{(\mu_i, \Sigma_i, o_i, C_i)\}$ 建模若干关键帧所观测的局部区域。子地图的划分采用动态触发策略：当当前帧相对子地图首帧的位移或旋转超过阈值 $d_{\max}$ 或 $\theta_{\max}$ 时初始化新子地图，避免了固定关键帧数策略可能导致的子地图过大或过小。前端采用帧到模型跟踪——先按匀速运动假设初始化位姿，再最小化渲染颜色/深度与输入的光度-深度联合损失优化位姿，并用 alpha 掩码和内点掩码剔除重建不完整区域的粗差。

LoopSplat 的核心创新是\textbf{3DGS 子地图间的直接配准}，将回环约束的计算分解为三个阶段。

\textbf{重叠估计。}两个子地图 $P$ 和 $Q$ 之间的近似重叠是鲁棒配准的前提。作者发现直接在高斯均值点云上匹配局部特征效果很差——高斯中心分布不均匀且缺乏判别性描述子。因此 LoopSplat 转而从图像层面提取重叠信息：将所有关键帧通过 NetVLAD 提取全局描述符 $d \in \mathbb{R}^{1024}$，计算两组关键帧间的余弦相似度，保留 top-$k$ 对用于配准。

\textbf{配准即关键帧定位。}这是 LoopSplat 最核心的洞察——既然 3DGS 子地图与其关联的视角可视为刚体，配准问题可转化为关键帧定位问题。对于选定的源视角 $v_i^p$，确定其在目标子地图 $Q$ 中的相机位姿 $T_i^q$ 等价于找到刚体变换 $T_{P \to Q} = T_i^q \, T_i^{p,-1}$。具体做法是固定 $Q$ 的参数，从 $Q$ 渲染 RGB-D 图像与源视角的观测计算 $\mathcal{L}_1$ 损失，优化刚体变换。这一"渲染即配准"的策略利用了 3DGS 的快速光栅化能力，同时融合了几何和外观信息，比纯几何的点云 ICP 更鲁棒。两个方向的视角（$P$ 中视角在 $Q$ 中定位，反之亦然）并行估计变换。

\textbf{多视角姿态精化。}从 $2k$ 个独立估计中获取全局一致性变换。以渲染残差的倒数作为各估计的权重——残差越小意味着配准质量越高——通过加权旋转平均求解全局旋转：

\begin{equation}
    \hat{R} = \arg\min_{R \in SO(3)} \sum_{i=1}^{k} \frac{1}{\epsilon_i}\|R - R_i\|_F^2 + \sum_{i=k+1}^{2k} \frac{1}{\epsilon_i}\|R - R_i^{-1}\|_F^2
    \label{eq:ch15_loopsplat_rotation_averaging}
\end{equation}

其中 $\|\cdot\|_F^2$ 为 Frobenius 范数，$\epsilon_i$ 为第 $i$ 个估计的渲染残差。全局位移取各估计的加权均值。

在回环检测方面，LoopSplat 同样使用 NetVLAD 全局描述符，但采用比纯图像相似度更严格的策略：先计算子地图内部关键帧间的自相似度 $s_{\mathrm{self}}^i$（取第 $p$ 百分位），再计算子地图间的交叉相似度 $s_{\mathrm{cross}}^{i,j}$，仅当 $s_{\mathrm{cross}}^{i,j} > \min(s_{\mathrm{self}}^i, s_{\mathrm{self}}^j)$ 时初步判定为回环。为抑制仅依赖视觉相似度产生的误检，额外评估两个子地图高斯之间的初始几何重叠率 $r$，仅保留 $r > 0.2$ 的回环。检测到回环后，以 3DGS 配准提供的相对位姿作为回环边约束，构建位姿图并使用基于线过程的鲁棒公式执行位姿图优化（PGO）。优化后对每个子地图应用刚体校正：更新相机位姿 $T_j \leftarrow T_c^i T_j$、高斯中心 $\mu \leftarrow R_c^i \mu + t_c^i$、协方差矩阵 $\Sigma \leftarrow R_c^i \Sigma R_c^{i,T}$。此外，为减小地图规模并提升位姿估计精度，LoopSplat 省略球谐函数，仅使用 RGB 颜色。

与传统点云配准相比，LoopSplat 的 3DGS 配准有三个关键优势：（i）直接在 3DGS 上操作，避免了将子地图网格化再配准的额外步骤；（ii）渲染即配准的方式同时利用几何和外观信息，比纯几何配准更鲁棒；（iii）无需训练或预处理即可端到端集成到 SLAM 管线中。在 TUM RGB-D、Replica、ScanNet 和 MIT 多房间场景上的全面评估表明，LoopSplat 在跟踪、建图和新视角合成方面均达到竞争或最优性能。

GLC-SLAM（2024.09）~\cite{2409_10982} 在 Splat-SLAM 的可变形高斯地图 + 光流回环闭合范式基础上，将全局一致性维护从单一的因子图 BA 重新设计为\textbf{全局--局部层次化回环闭合}架构。其核心动机是：Splat-SLAM 将所有历史关键帧纳入统一的因子图执行全局 BA，在大规模场景中计算成本随轨迹长度线性增长；GLC-SLAM 则将场景渐进分割为独立的 3DGS 子地图，每个子地图锚定于一个全局关键帧，仅保持活跃子地图在线优化，从而将单次优化代价与总轨迹长度解耦。

具体而言，GLC-SLAM 在三个层面做出了与 Splat-SLAM 不同的设计选择。

\textbf{回环检测}方面，Splat-SLAM 依赖帧间光流幅值判定视角重叠，GLC-SLAM 改用 NetVLAD 全局描述符进行场所识别，并维护全局与局部两套关键帧数据库：全局回环检测在新子地图创建时触发，从全局关键帧数据库中检索最佳匹配（相似度需超过动态阈值 $s_g$，即全局关键帧与当前活跃子地图内关键帧的最低相似度）；局部回环检测在建图过程中持续运行，从局部数据库中接受最相似的关键帧（相似度超过固定阈值 $s_l$）。两层检测均额外施加帧重叠率 $r_o$ 的几何校验以抑制室内重复物体（如椅子、桌子）导致的误检。

\textbf{全局优化}方面，Splat-SLAM 构建稠密因子图执行全局光束法平差，GLC-SLAM 则构建稀疏位姿图（节点为关键帧位姿，边为序贯相对位姿和回环边约束），通过 Levenberg-Marquardt 算法求解位姿图优化（PGO），以边不确定性 $\Omega_i$ 作为鲁棒加权。

\textbf{地图调整}方面，Splat-SLAM 通过可变形高斯主动调整均值、缩放和旋转三个属性；GLC-SLAM 采用关键帧中心的直接调整策略——每个高斯 $g_i$ 记录其创建时关联的关键帧，PGO 完成后将高斯均值 $\mu_i$ 和旋转 $r_i$ 按关联关键帧的位姿更新量做刚体变换（$\mu_i^* = T^* T^{-1} \mu_i$，$r_i^* = R^* R^{-1} r_i$），随后在固定迭代次数内禁用剪枝和增密，仅用光度-深度-正则化联合损失精化子地图。

此外，GLC-SLAM 引入了 Splat-SLAM 未涉及的不确定性建模与关键帧选择策略。每个高斯维护一个不确定性分数 $u_g$，关键帧的信息量分数定义为 $s_{\mathrm{infor}} = \frac{1}{|G|}\sum_{g \in G} u_g$，即所观测高斯的不确定性总和归一化值。系统迭代选取覆盖最高不确定性总和的 $k$ 个关键帧参与子地图优化，优先选择观测更多尚未充分优化的有价值高斯的视角。消融实验表明，移除回环闭合显著降低跟踪鲁棒性，而随机关键帧选择替代不确定性策略则降低建图精度。在 Replica 上跟踪精度相比次优方法提升 26\%，在 ScanNet 上取得所有 3DGS 方法中最高的位姿精度，同时 GPU 内存占用为对比方法中最低。

Zhu et al.（2024.09）~\cite{2409_20111} 从一个不同的场景设定出发——旋转式多 RGB-D 相机（上/中/下三个视野重叠的相机固定在旋转平台上），而非手持单相机。这一设定带来两个特殊性：一是三个相机提供帧内多视角约束，跟踪和建图损失均由三相机联合计算（$L = L_{\mathrm{up}} + L_{\mathrm{mid}} + L_{\mathrm{down}}$），增强了位姿估计的稳定性；二是旋转平台的受控运动使得系统无需 Splat-SLAM 或 GLC-SLAM 中的可变形高斯机制来补偿大范围手持抖动。

在回环检测方面，Zhu et al. 不依赖 NetVLAD、BoW 等外部全局描述符，而是提出一种\textbf{纯 3DGS 原生}的回环检测策略。核心机制是：每个高斯在创建时记录其锚定帧的时间戳 $t_i$，根据与当前帧时间戳 $t_c$ 的差异是否超过阈值 $\tau_t$，虚拟地（不额外分配内存）将所有高斯分为\textbf{历史集合} $\mathcal{S}_h^c$（时间差超过阈值，即由旧帧创建的高斯）和\textbf{新出现集合} $\mathcal{S}_n^c$（时间差在阈值以内）。回环判定依赖两个判据的级联：首先，在当前位姿下将两组高斯分别重投影到图像平面，统计各自覆盖像素数 $N_h^c$ 和 $N_n^c$，当比值 $\gamma = N_h^c / N_n^c$ 超过阈值 $\tau_r$ 时（即当前视角下的高斯中历史高斯占比异常高，说明回到了之前观测过的区域），系统收集回环候选帧；其次，在同一位姿分别用历史集合和新出现集合渲染两幅图像，计算 SSIM 相似度，选择相似度最高的候选作为最终回环帧。这一策略直接利用了 3DGS 的快速渲染能力和高斯的时间属性，无需训练任何描述符网络，但前提是历史高斯和新出现高斯的渲染差异能够可靠区分回环与非回环场景。

回环检测成功后，Zhu et al. 的优化管线与 GLC-SLAM 在结构上相似但细节不同。第一阶段是轻量级位姿图优化（PGO）：节点为各帧位姿，边包括相邻帧的相对位姿约束和回环边约束，其中回环边的相对位姿 $\mathbf{T}_{c_i c_j}$ 通过在当前帧位姿处渲染历史高斯来估计。这与 GLC-SLAM 用 NetVLAD 检测回环 + PGO 的流程类似，但 Zhu et al. 的回环边是直接从高斯渲染中获取，而非场所识别描述符。第二阶段是高斯位置更新：与 GLC-SLAM 一样，根据锚定帧位姿的变化量对高斯均值做刚体变换 $\bm{\mu}^* = \mathbf{T}_{\mathrm{opt}} \mathbf{T}_{\mathrm{origin}}^{-1} \bm{\mu}$；但由于 Zhu et al. 将高斯缩放简化为各向同性（$s_x = s_y = s_z$），协方差矩阵 $\bm{\Sigma}$ 在位姿图优化后不需要更新（旋转变化量被各向同性缩放吸收），进一步降低了回环闭合的计算开销。第三阶段是局部光束法平差（BA）：将所有关键帧和 rand-list 帧混合后按时间戳划分为若干 bin，在每个 bin 内固定高斯参数仅优化相机位姿（使用与跟踪相同的光度-几何联合损失），优化完成后沿位姿图结构将位姿校正传播到所有非关键帧。这种分 bin 策略将全局 BA 分解为一系列小规模局部问题，避免了 GLC-SLAM 中全局 PGO 一次性处理所有帧的计算压力。在合成和真实旋转相机数据集上的实验表明，该方法在相机位姿估计和新视角渲染质量上均优于当时的最优方法。

NGM-SLAM（2024.05）~\cite{2405_19614} 与上述方法最根本的区别在于采用 NeRF-3DGS \textbf{双重表示}：每个子地图同时维护一个多分辨率哈希编码的神经辐射场和一个 3DGS 地图，NeRF 子地图并非仅用于跟踪或初始化，而是作为高斯渲染的\textbf{持续监督先验}——GLC-SLAM 和 DenseSplat 的子地图仅有单一 3DGS 表示，缺乏跨视角的插值推理能力，导致遮挡区域和稀疏观测处容易出现几何空洞；NGM-SLAM 的 NeRF 先验天然具有跨视角插值能力，可推断未直接观测到的结构，从而实现精确的空洞填补。

具体而言，NGM-SLAM 采用\textbf{渐进式渲染}策略协调两种表示的互补优势：先仅用稀疏关键帧执行 NeRF 体渲染获取基础监督（NeRF 渲染速度较慢但泛化能力强），待子地图初步调整后再用全帧位姿执行高斯光栅化渲染精化细节（高斯渲染速度快且保留高频细节）。子地图划分基于 ORB 特征点跟踪和关键帧数量阈值触发，并通过 DBoW 词袋模型评估关键帧间共视性——若连续子地图的连接帧缺乏共视关系，则在当前子地图的连接帧处强制分割新子地图，以避免融合时产生过大误差。

在回环闭合方面，NGM-SLAM 同样使用 BoW 词袋模型（与 DenseSplat 一致）在全局关键帧列表中检测回环。与 GLC-SLAM 的两层回环检测 + PGO 或 DenseSplat 的分层 BA 不同，NGM-SLAM 提出\textbf{粗到细全局调整}策略：第一步优化所有锚帧（每个子地图的首帧）的位姿并执行子地图融合，防止高斯子地图边界漂移；第二步固定锚帧位姿，对全局关键帧列表执行完整 BA，并执行第二次高斯子地图融合以完成回环校正。全局 BA 损失的设计借鉴了 DenseSplat 的渲染诱导映射思路：将当前位姿处渲染的 RGB 和深度通过相对位姿变换变形到共视关键帧处，与该处直接观测计算一致性损失，从而隐式地将位姿校正传播到稠密地图。此外，BA 过程中对缩放超过阈值的高斯基元施加 $L_2$ 正则化约束以防止尺度爆炸。

在地图质量维护方面，NGM-SLAM 引入\textbf{多尺度高斯渲染}（$1\times$、$4\times$、$16\times$、$64\times$ 四级降采样分辨率）：训练过程中将像素覆盖范围过小的细层级高斯合并为粗层级高斯，以缩放因子的倒数 $f_{\max} = 1/S_k$ 确定合并阈值，抑制 NeRF 采样引入的混叠伪影；同时采用\textbf{射线引导剪枝}——以 ORB 稀疏点云为引导，统计高斯对所有输入图像采样射线的贡献 $E(\mathbf{p}_i) = \max(\alpha_i^r T_i^r)$，移除贡献低于阈值 $t_{\mathrm{prune}}$ 的不活跃基元。系统支持单目、双目和 RGB-D 输入，在大规模场景上以约 5 FPS 运行。

DenseSplat（2025.02）~\cite{2502_09111} 在 LoopSplat 的 3DGS 子地图 + 回环闭合架构之上，解决了一个此前被忽视的根本问题：LoopSplat 及其前身的高斯基元均通过像素反投影初始化，在 SLAM 典型的稀疏视角条件下，这种 2D-to-3D 的初始化产生严重的点分布不均——前景密集、背景稀疏、遮挡区域完全空白，迫使系统依赖密集关键帧列表（如每 4 帧取 1 帧）来弥补覆盖不足，代价是高昂的内存和实时性损失。DenseSplat 的核心洞察是将 NeRF 的连续体积密度场引入作为\textbf{几何先验}：NeRF 天然具有跨视角的插值能力，可推断未直接观测到的几何结构，从而仅用极稀疏的关键帧（每 20 帧取 1 帧）即可为高斯基元提供均匀、几何感知的初始化。

DenseSplat 的管线分为两个阶段，分别对应 NeRF 先验的获取和消费。

\textbf{第一阶段}是 NeRF 跟踪与几何采样：采用多分辨率哈希辐射场（Instant-NGP 架构），先对输入帧执行帧到模型跟踪（光度损失优化相机位姿，这与 LoopSplat 通过渲染配准的方式不同——DenseSplat 直接在 NeRF 体积密度场上优化位姿，而非利用 3DGS 子地图渲染），再从训练好的 NeRF 中按密度阈值 $\sigma_{\mathrm{grid}}$ 检测网格边缘的表面交叉点，通过线性插值精确定位密度跃变位置 $x = x_1 + \frac{\sigma_{\mathrm{grid}} - \sigma(x_1)}{\sigma(x_2) - \sigma(x_1)}(x_2 - x_1)$，形成密集点云。这一几何感知采样确保高斯基元沿场景表面均匀分布，而非仅聚集在已观测像素的反投影位置——这是 DenseSplat 对 LoopSplat 最根本的改进：从"仅在已观测处放置基元"变为"在整个几何表面上均匀放置基元"。

\textbf{第二阶段}是细粒度高斯建图：以 NeRF 采样点初始化各向异性 3D 高斯基元 $\{G_j\}$，每个基元包含均值 $\mu_j$、协方差 $\Sigma_j$、不透明度 $o_j$、三阶球谐系数和缩放因子。为抑制 NeRF 采样引入的混叠伪影（尤其在物体边缘处产生细小浮游高斯），DenseSplat 引入\textbf{多尺度高斯渲染}策略（借鉴 Mip-Splatting 的思路）：在 $1\times$、$4\times$、$16\times$、$64\times$ 四个降采样分辨率层级上组织高斯，训练过程中将像素覆盖范围过小的细层级高斯合并为粗层级高斯，以缩放因子 $s_j$ 的倒数作为频率指标 $f_{\max} = 1/s_j$ 确定合并阈值。此外，采用\textbf{射线引导剪枝}（借鉴 RadSplat 的贡献度评估）评估每个高斯对所有输入图像中采样射线的贡献 $E(G_i) = \sum_{r \in I_i}(\alpha_i^r T_{G_i}^r)$，移除贡献低于阈值 $\tau_{\mathrm{prune}}$ 的不活跃基元。值得注意的是，NeRF 初始化的高斯被豁免剪枝——这些高斯可能看起来对当前观测帧贡献不高，但正是它们填补了遮挡区域的几何空白，保护它们是保持地图完整性的关键。

在回环闭合与全局一致性方面，DenseSplat 对 LoopSplat 的三个核心组件逐一做了替换和增强。（i）\textbf{回环检测}：LoopSplat 使用 NetVLAD 全局描述符 + 几何重叠率双重校验（计算代价较高且需要预训练）；DenseSplat 改用轻量的词袋模型（BoW）检索关键帧相似度，以 BoW 匹配分数作为回环候选判定标准，省去了 NetVLAD 的前向推理和几何重叠率的 3D 计算。（ii）\textbf{全局优化}：LoopSplat 执行位姿图优化（PGO）后对子地图做刚体校正；DenseSplat 将 PGO 替换为分层光束法平差（BA）——先对回环涉及的关键帧做仅位姿优化，再联合优化局部共视帧，最终对所有关键帧位姿做全局调整——BA 同时优化位姿和深度，比 PGO 的纯位姿优化提供更强的约束。（iii）\textbf{地图一致性维护}：这是 DenseSplat 最独特的贡献。LoopSplat 的刚体校正（旋转平移高斯中心和协方差）是纯几何操作，无法处理高斯颜色或遮挡关系的更新；DenseSplat 利用 3DGS 的快速渲染能力构建\textbf{BA 诱导映射损失}：将当前位姿 $T_i$ 处渲染的颜色 $C_G(T_i)$ 和深度 $D_G(T_i)$ 通过相对位姿变换 $T_{i \to k}$ 变形到共视关键帧位姿 $T_k$ 处，与该处直接渲染结果计算一致性损失 $\mathcal{L}_{\mathrm{BA}}^C = \sum_{k=1}^{N-1}\sum_{i=k+1}^{N}\|C_G T_{i \to k} - C_G T_k\|$。这一策略避免了显式操作高斯参数，通过渲染-变形-对比的方式隐式地将 BA 位姿校正传播到稠密地图，同时天然处理了颜色一致性。

子地图划分方面，DenseSplat 采用固定间隔策略（每 400 帧划分一个子地图），每个子地图独立维护 NeRF 模型和高斯地图。LoopSplat 按运动距离/旋转动态触发子地图划分，需要额外的运动阈值超参；DenseSplat 的固定间隔简化了系统设计，且每个子地图同时承载 NeRF 和高斯两种表示，在子地图粒度上实现了 NeRF 的跟踪能力与 3DGS 的渲染效率的互补。子地图融合时借鉴 Mipsfusion 的锚定帧 BA 策略：每个子地图以其首帧位姿为锚点，BA 后重新锚定以实现子地图间的精确对齐和无缝融合，同时并行构建各子地图以缓解全局地图的持续内存膨胀。此外，针对 BA 过程中高斯缩放因子可能爆炸的问题，DenseSplat 引入 L2 正则化损失 $\mathcal{L}_{\mathrm{reg}}$ 对缩放超过阈值 $\tau_{\mathrm{scale}} = 1$ 的高斯基元进行约束。

在 TUM RGB-D、Replica、ScanNet 等基准以及大规模公寓场景上的实验表明，DenseSplat 在渲染质量、轨迹精度和建图完整性上均达到当时最优水平，尤其在仅需稀疏关键帧（每 20 帧取 1 帧）的条件下仍能维持高质量重建——这一关键帧密度仅为 LoopSplat 等方法的 1/5，直接验证了 NeRF 几何先验对弥补稀疏观测的有效性。

与 DenseSplat 和 LoopSplat 采用运动立体或单目深度估计不同，MVS-GS（2024.12）~\cite{2412_19130} 将预训练的多视角立体网络 MVSFormer 引入在线 SLAM 管线，从局部时间窗口内的相邻关键帧直接估计高分辨率深度图和置信图。然而，在线 MVS 因源视角有限而产生大量离群点（尤其出现在背景和边界区域），为此 MVS-GS 设计了两级深度精化管线：先用 V-Fuse 深度融合网络将参考帧与邻域帧的深度-置信图融合为时序一致的精化深度图，再通过多视角几何一致性检查——将参考帧深度投影到邻域帧并计算相对深度差异，剔除未通过至少三视图一致性检验的像素——进一步滤除离群点。精化后的稠密深度图直接用于高斯点初始化，避免了 Photo-SLAM 等方法从稀疏特征扩展高斯的低质量问题。系统架构上，前端使用 DROID-SLAM 的稠密光束法平差进行跟踪，后端以独立线程并行运行，采用广义指数 splatting（GES）替代标准 3DGS 表示以提升大场景下的效率，且仅对深度图中未探索区域生成新高斯，避免重复初始化。值得注意的是，MVS-GS 不划分子地图也不包含显式回环检测，属于全量地图方案；其优势在于 MVS 深度为高斯提供了比运动立体更精确的几何先验，使得系统在挑战性户外场景（如 Mill 19 数据集）上超越当时最优稠密 SLAM 方法。

\subsubsection{加速与稀疏处理}

本路线介绍从像素采样、帧间冗余消除到频域优化等多维度的 3DGS-SLAM 加速策略。

在资源受限平台上实现实时运行是 3DGS-SLAM 走向实用的关键。本小节讨论从像素采样、帧间冗余消除到频域优化等多维度的加速策略。

Splatonic（2025.11）~\cite{2511_18755} 从算法-硬件协同设计的角度系统解决了 3DGS-SLAM 在移动平台上的计算瓶颈。以代表性算法 SplaTAM 为例，其在移动端 NVIDIA Orin SoC 上的帧率仅约 0.1 Hz，远低于实时要求（10--30 Hz）。通过性能剖析发现，3DGS-SLAM 的 94.7\% 执行时间集中在正向和反向光栅化两个阶段，其瓶颈可分为三类。

\textbf{线程发散}（正向光栅化的主要瓶颈）：当前瓦片级渲染管线中，每个线程负责一个像素，但同一 GPU warp 内不同像素需要积分的高斯集合不同，导致线程掩码和空闲——平均线程利用率仅 28.3\%。

\textbf{聚合}（反向光栅化的主要瓶颈）：高斯需从不同像素收集梯度，当前依赖 atomicAdd 操作导致频繁流水线停顿，聚合占反向光栅化 63.5\% 的时间。

\textbf{$\alpha$-checking}（两个阶段的共同拖累）：每次 $\alpha$-checking 需要调用特殊功能单元（SFU）计算指数函数，该操作在正向光栅化中占 43.4\%、在反向光栅化中占 33.6\%。

Splatonic 的核心思路是：既然光栅化的计算量与渲染像素数成正比，通过稀疏采样即可将计算量降低一至两个数量级。但关键挑战在于——直接将稀疏采样应用于现有瓦片级管线，在移动 GPU 上仅获得 4.2\% 的光栅化加速，远低于理论上的 256 倍。根本原因是稀疏像素导致 GPU warp 利用率极低：一个 16$\times$16 瓦片中只有一个像素活跃，其余 255 个线程全部空闲。因此，Splatonic 提出\emph{三层}协同方案：软件层自适应稀疏像素采样、算法层像素级渲染管线、硬件层专用加速器架构。

\textbf{第一层：自适应稀疏像素采样。}核心洞察来自经典 SLAM——传统 SLAM 仅用稀疏特征点即可定位，3DGS-SLAM 跟踪仅需优化一个 $4 \times 4$ 位姿矩阵，不应需要渲染所有像素。Splatonic 针对跟踪和建图的不同需求设计了差异化的采样策略。

\emph{跟踪采样}：在每个 $w_t \times w_t$（默认 $w_t=16$）瓦片内随机采样一个像素，将渲染像素数减少最高 256 倍。值得注意的是，实验表明简单的随机采样与基于 Harris 角点的特征采样精度相当，且显著优于低分辨率下采样和 GauSPU 的瓦片级采样——因为逐像素采样保证全局覆盖，而瓦片级采样会在同一瓦片内引入冗余计算。稀疏采样甚至\emph{略微提升}了跟踪精度：平均 ATE 在 Replica 上比原始算法低 0.01 cm，在 TUM RGB-D 上低 0.03 cm，原因是稀疏像素可以帮助消除重复纹理区域中的虚假匹配，类似于传统 SLAM 中关键点描述子的作用。

\emph{建图采样}：与跟踪不同，建图的目标是重建尚未观察到的场景区域。Splatonic 优先采样\textbf{未见像素}——通过首次正向渲染计算的累积透射率 $T_{\mathrm{final}}(p) > 0.5$ 判断像素是否被充分覆盖：高透射率意味着该像素几乎没有高斯贡献，属于尚未重建的新区域，这些像素集中在先前被遮挡的物体边界和未探索区域。然而，仅采样未见像素会导致像素过于稀疏，影响跟踪精度。因此 Splatonic 还引入\textbf{纹理加权随机采样}作为补充：使用 Sobel 滤波器计算每个像素的梯度幅值 $w_R(p) = \sqrt{G_x^2 + G_y^2}$ 作为纹理丰富度的近似度量，赋予纹理丰富像素更高的采样概率，在每个 $w_m \times w_m$（默认 $w_m=4$）瓦片内选一个像素。两种策略的叠加使建图阶段在降低计算量的同时将 PSNR 平均提升 0.8 dB。

\textbf{第二层：像素级渲染管线。}为解决稀疏采样与瓦片级管线的根本不兼容问题，Splatonic 将渲染粒度从瓦片级改为像素级，并引入\emph{抢占式 $\alpha$-checking}，具体包含三项关键改造。

\emph{像素级投影与排序}：原始管线在瓦片粒度上投影，即识别与整个瓦片相交的高斯——对于稀疏像素，这会引入大量不与采样像素相交的冗余高斯。Splatonic 改为在像素粒度上投影，确保只有真正与采样像素相交的高斯进入后续排序和渲染阶段，消除不必要的计算。

\emph{高斯并行渲染}（Gaussian-parallel rasterization）：原始管线中每个处理单元（PE）渲染一个像素，不同像素独立执行。Splatonic 将并行粒度翻转为\emph{同一像素内的多个高斯}由一个 warp 的线程协作处理——一个 warp 内所有线程共享目标像素的高斯列表，均匀分配给各线程，每个线程计算一部分高斯的偏色，最后通过规约操作累加得到最终像素颜色。这种设计从根本上消除了线程发散：所有线程始终处理同一像素的同一批高斯。

\emph{反向渲染的片上缓存}：反向渲染需要为每个高斯计算累积透射率 $T_i$ 和前缀颜色 $C_i$（从第 1 到第 $i$ 个高斯的偏色累加），这在瓦片级渲染下缓存代价过高（每个像素都要为所有贡献高斯存储这两个值）。像素级渲染下只需存储单个像素的数据量（约 8 KB），完全可放于片上存储器。因此，正向渲染的中间结果 $T_i$ 和 $C_i$ 可直接缓存供反向渲染复用，省去了瓦片级渲染下反向渲染中额外的一轮跨线程规约操作，消除了处理单元间的依赖，使每个高斯的偏梯度可完全并行计算。

\textbf{第三层：专用加速器架构。}像素级渲染管线将瓶颈从光栅化转移到了投影阶段（$\alpha$-checking 移至投影使投影占正向时间从 2.1\% 激增至 63.8\%）和聚合阶段（atomicAdd 导致的流水线停顿依然严重）。Splatonic 设计了流水线化的专用加速器来应对这两项新瓶颈。

\emph{投影单元}：两项优化加速抢占式 $\alpha$-checking。其一，\textbf{直接索引}：利用每个高斯包围盒的四角坐标直接计算采样像素列表中的索引范围，避免遍历整个像素列表进行不必要的 $\alpha$-checking。其二，\textbf{查找表近似}：用 64 项的查找表（LUT）近似指数函数计算，经验证该精度足以保持跟踪和重建精度不变。

\emph{聚合单元}：设计四通道并行聚合架构以隐藏片外访存延迟。每个通道维护一个\textbf{记分板}（scoreboard）记录当前高斯梯度的累积状态——哪些高斯的已累积梯度已从片外加载到片上缓存（Gaussian cache）。聚合时，合并单元先对同一批次中具有相同高斯 ID 的偏梯度进行片内规约，然后仅对记分板标记为"就绪"的高斯执行片外读写更新。这种设计通过\emph{交错处理不同高斯的梯度更新}来隐藏片外访存延迟：当等待某个高斯的片外数据时，其他高斯的梯度可继续在片上累积。

\textbf{整体架构}包含 8 个投影单元（各配 4 个 $\alpha$-filter）、4 个层次化排序单元、4 个光栅化引擎（各含 $2\times 2$ 正向/反向渲染单元和 1 个颜色规约单元）以及 1 个聚合单元（配 32 KB 高斯缓存和 8 KB 记分板），总面积仅 1.07 mm$^2$（16nm 节点），显著小于 GSCore（1.77 mm$^2$）和 GSArch（3.42 mm$^2$）。

\textbf{实验评估。}在四个 3DGS-SLAM 算法（SplaTAM、MonoGS、GS-SLAM、FlashSLAM）上评估。软件版本（Splatonic-sw，仅使用采样算法+像素级管线，在移动 Ampere GPU 上运行）端到端加速 14.6 倍，能耗降低 86.1\%，且跟踪精度和重建质量与原始算法持平或略有提升。硬件版本（Splatonic-hw，TSMC 16nm RTL 实现，缩放至 8nm 节点与 Orin GPU 对比）跟踪阶段加速 274.9 倍，能耗节省 4738.5 倍；即使与先前专用加速器（GSArch、GauSPU）对比，加速 25.2 倍，能耗节省 241.1 倍。光栅化阶段的瓶颈加速尤为显著：正向光栅化加速 103.1 倍、反向光栅化加速 95.0 倍。

Splatonic 通过稀疏像素采样减少光栅化的计算量，而 AGS（2025.09）~\cite{2509_00433} 则从一个完全不同的角度切入加速：利用\textbf{相邻帧之间的高度相似性}——SLAM 中 $63.8\%$ 的相邻帧具有高共视性（frame covisibility, FC），这意味着大量计算是冗余的。AGS 的核心洞察是：高共视性帧不仅可跳过不必要的跟踪迭代，其高斯的贡献模式也高度相似，因此可从前一帧预测当前帧中哪些高斯对渲染无贡献并直接跳过。为获取帧共视性这一关键信号，AGS 没有选择额外计算图像相似度（如 Splatonic 依赖 Sobel 梯度等图像特征），而是从\textbf{视频 CODEC 的运动估计中间数据}中零成本提取——CODEC 在编码时已为每个宏块计算了 SAD（绝对差值和）用于运动补偿，AGS 直接复用这些值判断帧间相似度，无需额外计算开销。

基于帧共视性，AGS 在跟踪和建图两个阶段分别设计了与 Splatonic 策略本质不同的加速方案。

\textbf{运动自适应跟踪}：Splatonic 通过稀疏像素采样减少每帧渲染的计算量，但每帧仍执行完整的跟踪流程；AGS 则根据 FC 值动态决定跟踪精度——高共视性帧（运动微小）仅需轻量级粗粒度位姿估计（基于 Droid-SLAM 的 CNN 特征提取 + GRU 位姿更新，无需 3DGS 训练），低共视性帧（运动较大）才额外执行少量（$Iter_T = 20$ 次）3DGS 训练迭代做细粒度精化。这一设计将跟踪从"每帧都做大量 3DGS 迭代"变为"大部分帧仅做轻量 CNN 推理"，从根本上改变了跟踪的计算模式。

\textbf{高斯贡献感知建图}：Splatonic 通过优先采样未见像素引导建图；AGS 则利用帧间高斯贡献的相似性——关键帧执行完整建图并记录每个高斯对像素的贡献度（以 $\alpha$ 值低于阈值 $Thresh_\alpha$ 的像素数作为非贡献指标），后续高共视性非关键帧直接复用前一关键帧的贡献记录，跳过被标记为非贡献的高斯（平均 $85.1\%$ 的高斯对渲染无贡献）。这种跨帧信息共享的思路是 Splatonic 所没有的，且假阳性率仅 $5.7\%$，几乎不影响渲染质量。

在硬件架构方面，AGS 与 Splatonic 虽都设计专用加速器，但针对的瓶颈和架构组织方式迥异。Splatonic 的三层架构围绕像素级渲染管线展开（投影单元、排序单元、光栅化引擎、聚合单元），核心是消除线程发散和聚合停顿；AGS 则按 SLAM 任务流组织为三个独立引擎：\textbf{FC 检测引擎}从 CODEC 的 DRAM 中读取 SAD 值并累积计算帧共视性，与 CODEC 共享硬件资源，面积开销极小；\textbf{位姿跟踪引擎}包含脉动阵列（用于 CNN 粗跟踪）和轻量 GS 数组（用于 3DGS 细跟踪），由 FC 检测结果动态选择执行路径；\textbf{建图引擎}配备完整 GS 数组、GS 记录表和 GS 跳过表——记录表在关键帧建图时捕获高斯贡献信息，跳过表在非关键帧建图时根据贡献信息过滤非贡献高斯。三个引擎拥有独立硬件资源，使跟踪和建图可流水线并行执行，而传统 GPU 上的 3DGS-SLAM 因跟踪依赖建图更新的高斯而必须串行。此外，AGS 针对建图引擎中"热高斯"（跨瓦片频繁出现的高斯）的 DRAM 访问瓶颈，设计了推测式片上缓存策略，将热高斯的贡献信息保留在片上缓冲器中直至所有瓦片处理完毕，大幅减少重复 DRAM 访问。同时设计了 GPE 调度器，将渲染的 $\alpha$ 计算和颜色累积两阶段解耦，空闲 GPE 可为繁忙 GPE 预计算 $\alpha$ 值并存入 alpha 缓冲区，缓解因跳过非贡献高斯加剧的工作负载不均衡问题。

在 TUM RGB-D、Replica 和 ScanNet++ 上评估，AGS 相比 SplaTAM 基线跟踪精度提升 1.97 cm ATE RMSE，映射质量损失仅 $2.36\%$ PSNR。AGS-Edge（28nm, 7.25 mm$^2$, 500MHz）相对 NVIDIA Jetson AGX Xavier 加速 17.12 倍、能效提升 42.28 倍；AGS-Server（14.38 mm$^2$）相对 A100 GPU 加速 6.71 倍、能效提升 22.58 倍。值得注意的是，纯算法层面在 GPU 上实施仅获得 1.12 倍加速——FC 检测与 SLAM 管线在 GPU 上串行执行，且高斯贡献信息的额外内存操作效率低下——这印证了专用硬件协同设计的必要性。

MG-SLAM（2024.05）~\cite{2405_20031} 利用曼哈顿世界假设增强室内环境几何精度和完整性。从结构化场景提取融合线段确保弱纹理室内区域的鲁棒跟踪，利用提取的线和平面表面假设在缺失几何区域策略性插值新高斯。SpectralSplats（2026.03）~\cite{2603_24036} 将优化目标从空间域转移到频域，通过频谱矩监督构建全局吸引域解决"梯度消失"问题——当严重位姿错位导致渲染对象落在目标局部足迹之外时，标准空间损失梯度严格为零，而频谱矩在全局范围内保持非零梯度信息。从一阶原理推导频率退火策略，从低频分量的全局凸性开始优化，逐步退火到高频分量实现精确空间对齐。RA-NeRF（2025.06）~\cite{2506_15242} 利用流驱动位姿调节和隐式位姿滤波处理复杂相机轨迹，采用增量管线结合光度一致性和隐式位姿滤波器消除噪声。

在基准和框架方面，SLAM\&Render（2025.04）~\cite{2504_13713} 提供专门面向 SLAM 和神经渲染交叉领域的基准数据集。Noisy-Replica（2024.06）~\cite{2406_16850} 系统评估多模态扰动下 SLAM 系统的鲁棒性。XRDSLAM（2024.10）~\cite{2410_23690} 提供灵活的模块化 SLAM 框架，集成多种 NeRF 和 3DGS 算法。Wang et al.（2026.02）~\cite{2602_04251} 综述 3DGS-SLAM 的性能优化和鲁棒性增强方法。ReefMapGS（2026.04）~\cite{2604_11992} 在水下珊瑚礁场景中实现多模态 SLAM 和 3DGS 的回环闭合。

\begin{takeaway}
系统级工程优化将 3DGS-SLAM 推向实用：专用硬件加速器实现相比 GPU 17 倍加速和 42 倍能效提升，频域优化解决大位姿偏差下的梯度消失问题，模块化框架则为方法对比和快速迭代提供基础设施。
\end{takeaway}

% ======================================================================

\subsection{近期方法脉络：Memory-aware Gaussian map maintenance}

这一组近期工作延续 3DGS-SLAM，主要把改动集中在优化等环节。下面只保留相对既有路线的变化，便于和前文方法直接对照。

\paragraph{Pocket-SLAM}
Pocket-SLAM（2026.06）~\cite{2606_24796} 可作为“Memory-aware Gaussian map maintenance”中的新锚点候选。相较于 3DGS-SLAM，主要改变是：基于渲染有效区域贡献的 Gaussian pruning 策略。3DGS 的显式高斯表示和可微渲染能力保持为基础。
其设计动机是：将剪枝依据从单个高斯的不透明度等局部启发式转向对有效渲染区域的贡献，从而更适合大场景 SLAM。
具体而言，在 3DGS-SLAM 流程中监测高斯对当前和潜在渲染区域的贡献，选择性移除低贡献高斯，减少地图膨胀并保持视图合成质量。
这一锚点判断来自当前分类中的机制关系：候选锚点：优化，且相关性 0.94；后续方法可围绕 针对 Gaussian map 的持续增长加入渲染区域感知剪枝。 继续比较和扩展。
适用边界是：剪枝策略可能依赖轨迹覆盖和可见性估计，对未来未观测视角的保持能力需要谨慎评估。
\section{大规模与多智能体SLAM}
\label{sec:15_8}

将 3DGS-SLAM 从单房间场景扩展到城市级大规模环境，从单机器人扩展到多智能体协同，需要解决 GPU 内存瓶颈、通信效率、跨智能体一致性等根本性挑战。本节梳理 13 篇论文，涵盖大规模建图、分布式协作和系统级解决方案。

\subsection{大规模建图}
\label{ss:15_8a}

本节介绍三条技术路线：（1）大规模场景建图；（2）多智能体协同 SLAM；（3）数据集与评测。

\subsubsection{大规模场景建图}

本路线介绍突破 GPU 内存限制、处理公里级大规模场景的方法，涵盖核外计算、子地图策略和二值化高斯表示。

标准 3DGS-SLAM 受限于 GPU 内存，难以处理公里级场景。核外计算、子地图策略和二值化高斯表示是突破这一瓶颈的三条主要技术路径。

上述 Gaussian-LIC 系列和 GS-LIVO 虽然通过 CPU-GPU 双层存储缓解了显存压力，但全局地图仍须完整驻留在系统 RAM 中——当场景从单房间扩展到公里级城市街道时，高斯数量可达数百万乃至数千万，系统 RAM 容量同样成为硬上限。DiskChunGS（2025.11）~\cite{2511_23030} 沿大规模场景建图方向做出根本性改进：借鉴操作系统虚拟内存的思想，将场景按空间划分为若干"块"（chunk），仅将当前相机视野涉及的活跃区域保留在 GPU 中，其余区域以序列化形式存储在\textbf{磁盘}上按需加载——与 GS-LIVO 将冷高斯留在 CPU RAM 不同，DiskChunGS 将冷数据下沉到磁盘，利用 NVMe SSD 的大容量和可扩展性彻底将可重建场景规模与任何内存上限解耦。同时，DiskChunGS 保留完整的三阶球谐各向异性高斯表示（不做简化），且前端采用 ORB-SLAM3 提供位姿估计而非 LiDAR-惯性融合，使其兼容纯视觉输入（单目、立体和 RGB-D）。

技术架构围绕空间分块与 GPU-磁盘三级调度展开，包含以下核心组件：

(1) \textbf{空间分块}：将 3D 世界划分为大小为 $s$（实验中取 10m）的立方体块。每个高斯根据其位置 $\mathbf{p}$ 被分配到对应块：

\[
\phi(\mathbf{p}) = \left\lfloor \frac{p_i + s/2}{s} \right\rfloor, \quad i \in \{x,y,z\}
\]

块坐标编码为 64 位整数——每轴分配 21 位，并加入 $2^{20}$ 的偏移量处理负坐标——从而将三维块定位转化为单次哈希查询，实现 $O(1)$ 的块查找。

(2) \textbf{高斯采样策略}：与 Gaussian-LIC 系列依赖 LiDAR 点云直接初始化高斯不同，DiskChunGS 面向纯视觉输入，采用直接基元采样（direct primitive sampling）替代 3DGS 原始的迭代致密化。具体而言，对于每帧关键帧，系统从两个来源产生高斯：ORB-SLAM3 提供的稀疏点云贡献一组高斯；同时，基于图像的 Laplacian of Gaussian（LoG）算子范数估计每个像素应生成高斯的概率，并通过渲染惩罚图避免在已良好重建的区域放置冗余高斯。采样概率为 $P_s(x,y) = \max(P_L(x,y) - \hat{P}(x,y),\, 0)$，其中 $P_L$ 是输入图像的 LoG 范数，$\hat{P}$ 是当前渲染场景的 LoG 范数——这确保新增高斯仅出现在细节不足的区域。深度估计采用 Depth-Anything-2（单目模式）或 Fast-ACVNet（立体模式），并通过三角化匹配对齐和引导式立体匹配修正单目深度误差，弥补了缺乏 LiDAR 度量深度的不足。

(3) \textbf{层次化视锥剔除}：系统从视图-投影矩阵提取视锥平面，递归测试块区域——在粗粒度层级早期拒绝整个视锥外区域，对相交区域逐级细分直至单个块。该过程使用 OpenMP 并行化，并引入距离过滤和位姿感知缓存以减少冗余计算。

(4) \textbf{LRU 淘汰机制}：与 GS-LIVO 的滑动窗口增量更新协议（帧间重叠复用保持内存连续性）不同，DiskChunGS 采用更简洁的页面置换策略：系统在可配置的 VRAM 预算下运行（实验设定为 150 万高斯），当加载新块将超出预算时，按最近最少使用（LRU）原则将旧块序列化写入磁盘。淘汰过程以贪心策略运作：将可淘汰块按访问时间排序（最旧优先），迭代选取最少数量块以容纳新数据，确保最近访问的区域得以保留并维持空间连贯性。这一设计虽然不如 GS-LIVO 的五步协议精细，但在核外场景下实现更简单——因为磁盘 I/O 延迟远高于 CPU-GPU 数据搬运，简洁的批量淘汰策略反而减少了碎片化的磁盘读写。系统另设独立 LRU 队列管理关键帧内存（预算 400 帧），超限时将最久未用关键帧写回磁盘。

(5) \textbf{空间网格关键帧选择}：大规模场景中的关键帧选择面临一个根本张力——空间局部性（最小化 I/O 开销）与梯度多样性（确保稳定优化）之间的矛盾。选择视野与已加载块重叠率低的关键帧会触发大量块换入换出，导致训练近乎停滞；而选择视野过于集中的关键帧则可能因观测约束不足而引发灾难性遗忘。为平衡这一张力，DiskChunGS 将关键帧组织在分辨率为 $g = 200$m 的空间网格中（注意：此关键帧选择网格独立于且远粗于用于高斯划分的空间块网格）。在最近关键帧所在网格单元内选择优化关键帧，最大化块重叠率：

\[
\omega(k) = \frac{|C_v(k) \cap C_{\mathrm{active}}|}{|C_v(k)|}
\]

其中 $C_v(k)$ 是关键帧 $k$ 的可见块集合，$C_{\mathrm{active}}$ 是当前已加载块集合。高重叠值（接近 1.0）意味着候选关键帧所需的大部分块已在内存中，加载开销极小。在空间候选集内部，系统采用与 CaRtGS 相同的基于使用量和损失加权的选取策略，优先为高损失关键帧分配额外迭代。这一设计的核心优势在于：当系统重访先前已建图区域时，该空间单元内的所有历史关键帧——无论时间距离多远——均可贡献其观测约束，从而有效防止灾难性遗忘。

(6) \textbf{回环闭合处理}：传统基于体素的系统在回环闭合时通常需要昂贵的全局优化或受影响区域的完整重建。DiskChunGS 利用块化架构实现更精准的局部修正。当回环闭合更新关键帧位姿后，系统仅对受影响关键帧执行视锥剔除以识别相关块，并对其中高斯施加刚体变换（同时调整位置和旋转），保持场景表示与校正后相机轨迹的一致性。为减少 I/O 开销，系统汇总所有受影响块的唯一高斯集合，若在内存预算内则批量加载并使用全局变换掩码防止同一高斯被重复处理（可见于多个关键帧时）。变换后，跨块边界的高斯通过重新计算块分配迁移到正确块中，确保空间层次的一致性。

(7) \textbf{后修正细化}：在几何变换和块重分配之后，系统执行针对性优化以细化受影响区域。此时暂停新帧摄入，对回环闭合检测点处——即先前独立重建的子地图汇合处——重置高斯不透明度和优化器状态，允许交汇区域从头重新优化。该细化阶段运行 1000 次迭代（约 10 秒），确保先前断开的地图段之间实现平滑融合。

在损失函数方面，系统采用图像损失与深度损失的组合：$L = L_{\mathrm{image}} + \lambda_{\mathrm{depth}} L_{\mathrm{depth}}$，其中图像损失为 L1 与 SSIM 的加权组合 $L_{\mathrm{image}} = (1-\lambda_s)\|\hat{I} - I_{\mathrm{gt}}\|_1 + \lambda_s(1 - \mathrm{SSIM}(\hat{I}, I_{\mathrm{gt}}))$，深度损失为渲染深度与估计深度之间的绝对差 $L_{\mathrm{depth}} = \|\hat{D} - D_{\mathrm{gt}}\|_1$。

系统与 ORB-SLAM3 无缝集成进行位姿估计，支持独立模式（使用 ORB-SLAM3 的视觉里程计）和 ROS 集成模式（接受外部位姿输入），兼容单目、立体和 RGB-D 相机输入。与 Gaussian-LIC 系列依赖 LiDAR-惯性紧密耦合获得高精度位姿不同，DiskChunGS 的纯视觉前端在 KITTI 等大规模户外场景中牺牲了部分定位精度（ATE 在 0.17\,m 至 4.50\,m 之间），但换来了无需专用传感器的广泛适用性。在室内（Replica、TUM-RGBD）、城市驾驶（KITTI）和资源受限的 NVIDIA Jetson Orin 平台上验证。DiskChunGS 是唯一成功完成所有 11 个 KITTI 序列而无内存溢出或跟踪丢失的方法——对比之下，CaRtGS 在 3 个序列上因 VRAM 不足崩溃，GigaSLAM 在 4 个序列上崩溃，On-the-Fly 在 6 个序列上丢失跟踪。在 KITTI 最长序列 02（5.07 km）上，块 I/O 操作仅占总处理时间的 5.5\%（5 km/h 摄入速率）至 19.6\%（20 km/h 摄入速率），证明磁盘调度开销可控。在 Replica 上取得最优渲染质量（PSNR 34.59、SSIM 0.94、LPIPS 0.05），并在 KITTI 上实现速度-质量 Pareto 前沿的全面领先。

VPGS-SLAM（2025.05）~\cite{2505_18992} 与 DiskChunGS 的核心差异在于\textbf{场景表示的参数化方式}和\textbf{跟踪-回环的 2D-3D 融合策略}。DiskChunGS 直接存储高斯基元属性（位置、颜色、协方差等），而 VPGS-SLAM 采用 Scaffold-GS 风格的\textbf{体素锚点 + 神经高斯}表示：场景被体素化后，每个体素中心初始化为一个锚点 $x_i^a$，锚点携带 32 维特征 $f_i^a$、缩放因子 $l_i$ 和 $k$ 个可学习偏移量 $\mathcal{O}_i \in \mathbb{R}^{k \times 3}$；每个锚点生成 $k$ 个神经高斯，其位置为 $\mu_i^m = x_i^a + \mathcal{O}_i^m \cdot l_i$，其余属性（不透明度、颜色、旋转、缩放）均由独立 MLP 从锚点特征 $f_i^a$、视距 $\delta_i$ 和视角方向 $d_i$ 解码得到。这种隐式-显式混合表示使参数量仅为 $O(N)$（$N$ 为锚点数），相比 DiskChunGS 的直接存储方式，在 Replica 上地图内存仅 70.81 MB（SplaTAM 为 273 MB）。体素采用多分辨率设计——距相机近处使用细粒度体素、远处使用粗粒度体素——并通过顺序化的渐进增长与剪枝策略动态管理：新关键帧到来时在未覆盖区域体素化并初始化锚点；当相机移动到子地图边界时，计算每个体素内神经高斯的平均梯度 $\nabla_g$，在梯度超过阈值 $\tau_g$ 的体素中生长新锚点；同时累积锚点关联神经高斯的不透明度值，移除不活跃锚点。

在子地图管理方面，VPGS-SLAM 同样采用多子地图架构，但策略与 DiskChunGS 的磁盘核外块管理不同。VPGS-SLAM 以距离阈值 $d$（室内 0.5\,m、户外 10\,m）和旋转阈值 $\omega = 50^\circ$ 触发新子地图初始化，每个子地图以其首帧为锚帧。已优化完毕的非激活子地图参数可从 GPU 内存卸载，仅保留激活子地图及其重叠区域的参数在线，峰值 GPU 占用仅 7.8 GB——相比 DiskChunGS 的 LRU 磁盘换入换出机制，VPGS-SLAM 的策略更轻量，不依赖磁盘 I/O 基础设施。

在相机跟踪方面，与 DiskChunGS 依赖 ORB-SLAM3 外部前端不同，VPGS-SLAM 自研了\textbf{粗到细的 2D-3D 融合跟踪}。粗级阶段利用 2D 光度损失（RGB L1 + SSIM + 深度 L1）优化位姿，提供良好初始估计；细级阶段引入体素化 3D 高斯 ICP——将输入点云变换到全局坐标后，在体素化高斯地图上执行最近邻搜索建立对应关系，以 Geman-McClure 鲁棒核最小化点对点距离残差。关键设计是\textbf{自适应 2D 信息评估}：当渲染损失超过阈值 $\zeta$ 时（如户外运动模糊和曝光变化场景），系统跳过粗级 2D 优化、直接使用 3D ICP 定位。消融实验证实，纯 2D 跟踪在 KITTI 上 ATE 高达 5.96\,m，纯 3D 跟踪为 1.15\,m，而 2D-3D 融合降至 1.09\,m。

在回环闭合方面，VPGS-SLAM 使用 NetVLAD 提取视觉描述符并通过余弦相似度检测回环（DiskChunGS 未集成回环检测模块）。回环校正同样采用 2D-3D 融合约束——同时利用渲染损失和体素 ICP 构建位姿图优化——校正后的位姿变换同时应用于相机位姿和高斯锚点的四元数。VPGS-SLAM 独特的贡献是\textbf{在线蒸馏子地图融合}：回环检测后，在两个子地图的重叠区域识别高度相似的图像对，以各自子地图的位姿分别渲染 RGB 和深度图像，构建蒸馏损失 $\mathcal{L}_{\alpha-\beta} = \frac{1}{m}\sum_{\alpha,\beta}((\hat{c}_\alpha - \hat{c}_\beta)^2 + (\hat{d}_\alpha - \hat{d}_\beta)^2)$，将不同子地图的 3D 高斯信息隐式整合。相比 DiskChunGS 的刚体变换 + 后修正细化策略，在线蒸馏通过渲染-对比的方式自然处理跨子地图的颜色和几何不一致性，无需显式操作高斯参数。

实验在 Replica、ScanNet、KITTI、VKITTI2 和 Newer College 五个基准上验证。在 Replica 上跟踪精度 (ATE RMSE 0.21 cm) 和渲染质量 (PSNR 35.82) 均优于 LoopSplat 等现有 3DGS 方法；在 KITTI 上 PSNR 达 21.37（LoopSplat 为 16.77），ATE RMSE 1.40\,m，是唯一成功运行所有 11 个序列的 3DGS 方法。渲染速度 322 FPS，峰值 GPU 内存 7.8 GB。

GigaSLAM（2025.03）~\cite{2503_08071} 是首个面向公里级户外环境的 RGB NeRF/3DGS SLAM 框架，在 KITTI、KITTI 360、4 Seasons 和 A2D2 数据集上验证。采用层次化稀疏体素地图表示，高斯由神经网络在多级细节上解码，支持高效可扩展建图和高保真视点渲染。前端利用度量深度模型结合对极几何和 PnP 算法估计位姿，并引入基于词袋的回环闭合机制。LSG-SLAM（2025.05）~\cite{2505_09915} 面向大规模户外场景，引入连续 GS 子地图和特征对齐变形约束处理无界场景和有限内存。通过地点识别检测子地图间回环，渲染和特征变形损失优化回环关键帧间的相对位姿。在 EuRoC 和 KITTI 上超越传统和神经方法。

VINGS-Mono（2025.01）~\cite{2501_08286} 在 DiskChunGS 的核外建图路线上做出了系统性重新设计，从表示、前端、内存管理到回环闭合全面替换，使其成为首个可仅凭单目相机和低频 IMU 在户外公里级场景运行的 Gaussian SLAM 系统。以下按模块逐一对比 DiskChunGS 的设计选择，阐明 VINGS-Mono 的改进动机与具体差异。

\textbf{（1）2D Surfel 替代 3D 椭球}。DiskChunGS 采用 3D 高斯椭球作为场景基元，每个高斯需要位置、旋转四元数、缩放、颜色、不透明度等大量参数。VINGS-Mono 改用 2D Gaussian Surfel——一种扁平化的面元表示，参数量显著少于 3D 椭球。这一替换的动机在于：户外大规模场景中高斯数量可达数千万级，每个基元节省少量参数在大规模下累积为可观的内存收益；同时 2D Surfel 的面元性质天然适合表面重建，可渲染法向量图以增强几何约束。系统在单张 RTX 4090 上成功处理了高达 5170 万个高斯椭球的公里级城市场景。

\textbf{（2）VIO 前端替代纯视觉前端}。DiskChunGS 依赖 ORB-SLAM3 提供位姿估计，但在户外快速运动、纹理匮乏的场景（如大面积地面）中 ORB 特征点容易丢失。VINGS-Mono 重新设计了基于 DROID-SLAM 的稠密束调整（Dense BA）前端：通过相关体积（correlation volume）提取相邻帧特征，GRU 迭代更新光流残差与权重，构建关于逆深度和位姿的联合优化问题，利用 Schur 补消除深度状态以高效求解位姿更新。在此基础上融合 IMU 预积分因子（基于 GTSAM 实现），构建视觉-惯性因子图，使系统在单目尺度歧义下仍能获得度量级一致的轨迹。此外，前端从 DBA 的 Hessian 矩阵中直接推导逐像素逆深度的不确定性估计（marginal covariance），为下游地图模块提供可靠的几何置信度，用于过滤浮点和天空区域。

\textbf{（3）分数管理器替代 LRU 淘汰}。DiskChunGS 的 LRU 机制以"最近是否访问"作为淘汰判据——简单高效但粗糙，无法区分"高斯质量差应删除"与"高斯质量好只是暂时不在视野中"。VINGS-Mono 为每个高斯计算两个分数：\textit{贡献分} $S_C(g)$ 衡量高斯在关键帧列表中的渲染权重之和（采用求和以保留对所有关键帧的贡献），\textit{误差分} $S_E(g)$ 取高斯在各帧上造成的最大 RGB 损失（采用最大值以避免高斯因"被多帧看到但每帧损失很小"而被误判）。在此基础上引入双状态机制——高斯处于 \textit{stable} 或 \textit{unstable} 状态：贡献分低的不稳定高斯被标记为稳定（保护历史高斯不被误删），误差分高的稳定高斯被重置为不稳定（允许系统重访时重新优化）。每 200 次迭代执行存储控制，删除不稳定且贡献分低于阈值的高斯。实验表明该机制可将高斯数量减半而不降低渲染质量。内存管理方面，VINGS-Mono 采用 CPU-GPU 双端调度（而非 DiskChunGS 的 GPU-磁盘三级调度），利用 $ID(g)$ 变量（记录高斯贡献最大的关键帧 ID）计算高斯与当前位姿的距离，按距离阈值将高斯在 GPU 和 CPU 之间迁移。

\textbf{（4）采样光栅化器}。原始 3DGS 的反向传播中，每个 GPU 块负责一个 tile（$16\times16$ 像素），256 个线程各处理一个像素——线程迭代次数取决于该像素关联的高斯数量，最慢线程成为瓶颈。VINGS-Mono 受 Taming-3DGS 启发，在前向传播中每 32 个高斯存储一次中间变量，反向传播时以 32 线程的 warp 为单位处理高斯，使线程迭代次数降至 tile 内关联像素数（而非像素关联高斯数），并在每个 tile 内选择损失率最高的 50\% 像素进行反向传播，实现 273\% 的反向加速（11.55 ms $\to$ 4.23 ms），PSNR 仅下降 0.56 dB。

\textbf{（5）单帧到多帧位姿精化}。现有 GS-SLAM（如 MonoGS、SplaTAM）通过将梯度传递到高斯属性、再由高斯属性传递到当前帧位姿来优化定位——单次渲染仅优化一帧。VINGS-Mono 利用分数管理器中已建立的高斯-关键帧关联（$ID(g)$），将不同高斯的梯度成对传递到各自所属的关键帧位姿，实现单次渲染同时优化多帧位姿。位姿精化结果直接替换前端位姿缓冲区，形成"前端粗估计 $\to$ 渲染精化 $\to$ 反馈前端"的闭环。消融实验显示该策略在定位精度上全面超越单帧优化方式。

\textbf{（6）NVS 回环闭合}。DiskChunGS 的回环闭合依赖外部回环检测器提供约束，在检测到回环后对受影响块内的高斯施加刚体变换。VINGS-Mono 则创造性地利用高斯渲染本身作为回环检测工具：首先通过 LightGlue 匹配当前帧与历史帧的特征点，用渲染深度图和 PnP 求解相对位姿，然后从该新视角渲染高斯地图，计算渲染图像与当前帧的 L1 损失——若损失低于阈值或低于其他候选帧中位数的十分之一，则确认回环。这一设计避免了词袋方法的词汇表依赖，直接检验"当前帧是否处于已建图区域的某个新视角"。回环校正阶段，系统为每个高斯匹配其贡献分最高的关键帧位姿（约 1000 帧耗时 2 秒），构建位姿图优化全局关键帧位姿，然后对高斯施加对应位姿变换并微调 100 次迭代，无需从头重训地图。

\textbf{（7）动态擦除器}。户外场景中车辆和行人是不可避免的动态干扰物——离线 4DGS 方法通过建模时间维度来处理动态，但 SLAM 要求在线增量处理。VINGS-Mono 设计了一种启发式分割方法：用 FastSAM 生成语义掩码，对新关键帧渲染颜色图，计算 SSIM 与 L1 的乘积作为重渲染损失（SSIM 敏感于纹理结构，L1 敏感于颜色值差异），并融合深度不确定性以增强边缘区域的动态检测能力。当某个语义掩码区域内超过 20\% 像素的重渲染损失超过阈值时，该掩码被标记为动态并在后续高斯添加和渲染中过滤。

在实验验证方面，VINGS-Mono 在 ScanNet、BundleFusion（室内）、KITTI、KITTI-360、Waymo、MegaNeRF、Hierarchical（户外）等全部真实数据集上进行了评测。定位方面，室内 ATE 与传统 SLAM（ORB-SLAM3、DROID-SLAM）持平，显著优于现有单目 GS-SLAM（MonoGS、PhotoSLAM）；户外 KITTI/KITTI-360 上以 VIO 模式运行，超越 VINS-Mono、ORB-SLAM3 及学习方法 iSLAM、Selective-VIO。渲染方面，ScanNet 上取得 PSNR 22.43 dB、SSIM 0.79、LPIPS 0.22 的最优综合表现；在公里级户外场景（KITTI、KITTI-360、Waymo）中，无其他 NeRF/3DGS SLAM 方法能成功完成完整轨迹，VINGS-Mono 是唯一可行方案。系统还开发了基于 Flutter 的移动端 App，验证仅凭智能手机（480$\times$720 @ 30Hz RGB + 内置 IMU）即可生成高质量高斯地图，并在约 1 km$\times$0.4 km 校园场景（自行车采集，仅 20Hz RGB + 1Hz GPS）中展示了近乎零尺度漂移的鲁棒性。

RTG-SLAM（2024.04）~\cite{2404_19706} 强制每个高斯要么不透明要么近乎透明——不透明高斯拟合表面和主色，透明高斯拟合残差颜色。这一二值化策略的核心优势在于：通过不同的深度渲染方式使单个不透明高斯即可良好拟合局部表面区域，大幅减少所需高斯数量（避免了多个半透明高斯堆叠拟合一个表面的冗余），配合高效在线优化方案实现大规模实时重建。

\subsubsection{多智能体协同 SLAM}

本路线介绍多机器人协同探索和建图的方法，解决单目尺度歧义、跨智能体一致性和通信效率等核心挑战。

多智能体 SLAM 允许机器人团队协同探索和建图，显著提高覆盖效率。核心挑战包括单目尺度歧义、跨智能体一致性和通信效率。

MAC-Ego3D（2024.12）~\cite{2412_09723} 提出\textbf{多智能体高斯共识}（Multi-Agent Gaussian Consensus）机制，将协作建图问题建模为联合优化：各智能体最小化局部目标 $f^{a_i}(T_{1:t}^{a_i}, M^{a_i})$（自身建图与位姿估计误差），同时通过全局目标 $g^{a_i,a_j}$ 对齐重叠区域的局部地图以形成一致的全局地图。与 DiskChunGS 将全局地图管理集中在单实体不同，MAC-Ego3D 将共识分解为智能体内（Intra-Agent）和智能体间（Inter-Agent）两个层次，统一遵循四步循环：选择（Selection）$\to$ 关联（Association）$\to$ 对齐（Alignment）$\to$ 优化（Optimization），两层通过异步多进程并行执行。

\textbf{智能体内高斯共识}实现单智能体的增量建图与跟踪。\textit{选择}步从局部地图 $M^{a_i}$ 中提取当前视野内且位于估计位姿 $T_k^{a_i}$ 邻近范围内的候选高斯，利用 SLAM 的时序连续性将搜索空间聚焦于最可能匹配当前观测的地图区域。\textit{关联}步将新观测的深度点建模为高斯分布 $G_k^{a_i} \sim \mathcal{N}(x_k^{a_i}, \Sigma_k^{a_i})$，与候选地图高斯按空间邻近性建立对应关系 $\{(G_k^{a_i}, G_j^{M^{a_i}})\}$。\textit{对齐}步通过最小化变换观测点与地图高斯之间的马氏距离来精化位姿 $T_k^{a_i}$，实质是最大化观测似然 $p(X_k^{a_i} \mid M^{a_i}, T_k^{a_i})$。\textit{优化}步将新高斯融入局部地图，并通过可微渲染最小化光度损失以精化高斯参数，确保地图与 RGB-D 观测持续对齐。

\textbf{智能体间高斯共识}解决跨智能体的地图对齐问题。\textit{选择}步中各智能体按固定帧间隔 $T_{\mathrm{comm}}$ 将观测图像压缩为一维向量表示 $v^{a_i}$，通过向量相似度 $S(v^{a_i}, v^{a_j})$ 检测重叠区域——超过阈值 $\tau$ 时触发回环。\textit{关联}步建立不同智能体局部地图中高斯的跨智能体对应关系。\textit{对齐}步求解使重叠区域高斯差异最小的相对变换 $T^{a_i \to a_j}$。\textit{优化}步通过跨智能体的可微渲染联合精化全局地图，损失函数同时约束视觉一致性（L1 + SSIM）、几何一致性（深度 L1）和位姿对齐：$J = \mathcal{L}(M) + \sum_{(i,j)} \|T^{a_i} T_{i,j} - T^{a_j}\|_{\Omega_{i,j}}^2$。这一异步设计使智能体间共识不会阻塞各自的局部建图流程。

在 Multi-agent Replica 数据集上，MAC-Ego3D 的平均 ATE 为 0.14 cm，较次优方法（CP-SLAM 的 1.23 cm）提升近一个数量级；PSNR 达 40.04 dB，较 CP-SLAM（29.33 dB）提升超过 10 dB。在真实 7-Scenes 数据集上，PSNR 提升 4 dB，SSIM 提升 31\%。运行效率方面，跟踪和建图分别仅需 0.02 s 和 0.05 s 每帧，渲染速度 0.005 s 每帧，整体推理速度较 CP-SLAM 快 15 倍，且在 RGB/深度噪声干扰下仍保持鲁棒（RGB 污染下 ATE 仅 0.14 cm vs.\ CP-SLAM 的 1.08 cm）。

GRAND-SLAM（2025.06）~\cite{2506_18885} 与 MAC-Ego3D 共同关注多智能体协作建图，但采取了截然不同的技术路线：MAC-Ego3D 以全局高斯地图为表示，通过共识机制直接在重叠区域的高斯层级进行对齐；GRAND-SLAM 则采用\textbf{子地图局部坐标系优化}策略——每个子地图定义在局部坐标系中（首关键帧为原点），跟踪和优化均在局部坐标系内执行，仅在建图完成后通过位姿图优化将子地图变换到全局坐标系。这一设计的关键优势在于：优化梯度和曲率在以当前位姿为中心的局部坐标系中更加均衡，避免了在世界坐标系原点远离当前位姿时导致的梯度失衡和收敛退化。GRAND-SLAM 的跟踪模块也采用两阶段策略——先通过混合颜色-深度里程计提供粗初始化 $T_i = T_{i-1} T_{i-1,i}$，再在局部坐标系中分别优化旋转和平移分量以精化位姿。

在子地图管理方面，GRAND-SLAM 的子地图划分由平移阈值 $d_{\max}$ 或旋转阈值 $\theta_{\max}$ 触发，新子地图首帧重定义为原点以确保全局一致性。子地图优化使用渲染损失，并引入两个掩码提高鲁棒性：\textit{软 alpha 掩码} $M_\alpha$（alpha 图的多项式，抑制稀疏重建区域）和\textit{误差内点掩码} $M_{\mathrm{in}}$（二值掩码，丢弃颜色或深度误差超过帧内阈值的像素）。新高斯仅添加到渲染不透明度低于阈值 $\alpha_{\min}$ 的稀疏区域。

在回环闭合方面，GRAND-SLAM 实现了\textbf{帧内和帧间两类回环闭合}，这是其区别于仅做单智能体建图方法的核心能力。回环检测基于 NetVLAD 全局描述符：对查询关键帧检索 top-$k$ 最近邻，通过余弦相似度阈值 $\tau_{\mathrm{sim}}$ 和最小空间基线 $b_{\min}$ 筛选候选。对通过筛选的候选，采用\textbf{粗到精}的配准流程：帧内回环直接用已跟踪位姿计算初始变换 $T_{\mathrm{init}} = T_j^{-1} T_q$；帧间回环或长距离帧内匹配则先通过多尺度 RGB-D 配准（联合光度-几何损失）获取粗变换，再用稠密彩色点云上的点到面 ICP 精化至收敛——以对齐质量指标（适应度分数 $f$ 和内点 RMSE $\epsilon$）过滤低质量匹配。

回环闭合约束被纳入\textbf{全局位姿图优化}（GTSAM 实现 Levenberg-Marquardt 求解）：节点为子地图首帧位姿，边包括跟踪约束、帧内回环约束和帧间回环约束。优化后，每个子地图内所有关键帧位姿和高斯参数（均值、协方差）统一变换到全局坐标系，颜色值保持不变。在 Multi-agent Replica 数据集上，GRAND-SLAM 在跟踪精度和渲染质量上均优于或持平 MAGiC-SLAM 和 CP-SLAM；在 Kimera-Multi 大规模户外数据集（6 个智能体，总计 1.85 km 轨迹）上，GRAND-SLAM 是唯一成功完成全部序列的方法——MAGiC-SLAM 和 Gaussian-SLAM 在部分序列上中途崩溃，多智能体跟踪误差降低 91\%，渲染质量提升 28\% PSNR。

Compact Keyframe-Optimized Multi-Agent GS-SLAM（2026.04）~\cite{2604_00804} 通过紧凑化步骤移除冗余 3DGS 而不降低渲染质量，在纯渲染深度模式（无需 3DGS 之外的数据）和相机深度模式（包含轻量深度图像提升配准精度和额外高斯剪枝）下均实现 85--95\% 的数据传输减少，使 3DGS 多智能体 SLAM 更接近实际部署。MAGS-SLAM（2026.05）~\cite{2605_10760} 是首个 RGB-only 多智能体 3DGS SLAM 框架，每个智能体独立构建局部单目高斯子地图。为解决单目尺度歧义问题，集成紧凑子地图通信、几何-外观感知的回环验证和占据感知的高斯融合，实现一致性全局重建。同时发布 ReplicaMultiagent Plus 基准用于评估协同高斯 SLAM。HAMMER（2025.01）~\cite{2501_14147} 基于 ROS 通信架构，无需初始位姿先验即可从异步多机器人数据流生成 3D 语义-度量地图。其帧对齐模块无需先验相对位姿知识即可将局部 SLAM 位姿和图像数据转换到全局帧，在线模块从流数据训练语义 3DGS 地图，支持混合感知模式并蒸馏 CLIP 语义码用于开放词汇语言查询。

\subsubsection{数据集与评测}

本路线介绍面向 3DGS SLAM 评测的专用数据集和基准工作。WildfireX-SLAM（2025.10）~\cite{2510_27133} 构建大规模森林火灾 SLAM 数据集——WildfireX-SLAM 包含 5.5k 低空 RGB-D 航拍图像，总覆盖面积 16 km$^2$，基于虚幻引擎 5 Electric Dreams 环境构建。数据集提供灵活的灯光、天气和火灾类型/状态控制，支持森林建图、火灾应急响应等多种任务。在该数据集上的基准评测不仅揭示了 3DGS-SLAM 在森林环境中的独特挑战，还凸显了未来改进的方向。

Ress et al.（2024.04）~\cite{2404_17215} 在大型工厂车间中部署配备四个双目相机和 3D 激光扫描仪的机器人系统，应用 LiDAR 和视觉 SLAM 进行轨迹估计和稠密地图生成，并利用 3DGS 生成稠密深度图支持建筑施工监测。Nguyen et al.（2025.10）~\cite{2510_23988} 综述多机器人协同 3DGS SLAM 的架构分类（集中式 vs. 分布式），系统分析多智能体一致性、通信效率、语义蒸馏和实时可扩展性等核心组件，并指出终身建图、语义关联和多模态鲁棒性等关键开放挑战。

\begin{takeaway}
大规模与多智能体 SLAM 将 3DGS 从单房间扩展到城市级环境：紧凑化步骤实现 85--95\% 的数据传输减少，多智能体框架通过子地图通信和回环验证实现一致性全局重建，正向实际部署迈进。
\end{takeaway}

% ======================================================================

\section{近期研究（2026 年 8--11 月）}
\paragraph{Stipple：重新利用双目惯导前端的关键帧做增量泼溅}
2026.08~\cite{2608_00931}。Stipple 以 Basalt 跟踪并产出位姿、关键帧和闭环。当前地图先渲染关键帧，系统据尚未覆盖的像素决定何时添加原语，并把昂贵的双目深度计算限于部分锚点帧。高斯训练与双目惯导跟踪并行。

作者在 EuRoC 与 Monado 的公开序列和相应双目基线上报告渲染质量与吞吐量。高帧率的比较应区分前端跟踪速度与增量建图频率。此外，这条路线依赖双目标定和惯导输入。

\paragraph{EvTrajGS：将无位姿事件流的相机估计连续化}
2026.08~\cite{2608_08585}。离散事件姿态独立优化容易在时间上出现不一致。EvTrajGS 从通用管线得到粗位姿，用连续时间函数表示相机轨迹；渲染时汇总邻近时刻的轨迹证据，联合更新相机与高斯。重加权采样让训练关注监督不足的区间。

在由 Replica 五个场景合成的事件上，论文相对各指标第二名报告平均 PSNR 增益 3.84 dB。真实 TUM-VIE 部分给出轨迹误差与渲染定性图，不能把 Replica 的渲染数字转用于真实事件。方法仍使用粗轨迹初始化，不能理解为完全没有相机估计输入。

\paragraph{EndoMD-SLAM：退化观测下暂停静态内镜地图的更新}
2026.08~\cite{2608_08949}。水滴和污物可能遮住内镜镜头，暂态假象也可能被回填进长期地图。EndoMD-SLAM 用观测可靠性控制分支选择。退化帧触发历史干净关键帧检索和重定位，并暂停静态地图写入。自监督的静态与暂态分离有助于把持续解剖结构与短时光学伪影分开。

论文从 C3VDv2 选出十段含镜头水或碎屑标签的序列比较跟踪和重建。结果说明了这类镜头遮挡下的表现，重定位仍需可匹配的历史关键帧，不能直接推广到长期没有干净视线的操作。

\paragraph{在线高斯建图中的透射率保持、尺度初值和残差增密}
2026.08~\cite{2608_14902}。在线增密的每个阶段迭代有限，直接复制原语可能累积不透明度偏差。作者在 Photo-SLAM 的建图端加入保持合成透射率的增密，按深度和相机内参估计单像素的米制足迹来初始化尺度，并将新原语优先放入可从深度回投的高残差位置。

实验同时测量地图渲染和相机跟踪，并逐项添加模块比较。它沿用分离的跟踪架构，主要改动在建图侧。论文按 Photo-SLAM 工具的全轨迹评估图像，不能把数值与只评估训练视图的表格直接并列。

\paragraph{MotionGS-SLAM：事件流调节曝光积分中的高斯投影}
2026.08~\cite{2608_15024}。MotionGS-SLAM 以事件流估计局部运动，让投影高斯沿运动方向拉伸，同时随整体运动增加曝光内位姿的采样密度。两类调节进入渲染侧，用来解释观测帧的运动模糊；这个处理无需把底层三维高斯全部拉长来拟合拖影。

在合成 EventReplica 上，论文比较可用清晰真值的渲染质量和轨迹误差；真实测试包括三段 Color-DAVIS346 暗光采集。复用该流程需要与帧同步的事件数据，合成图像指标不能代替更多真实场景的定量验证。

\paragraph{AquaFlow：在水下单目流中同步估计位姿、地图与介质}
2026.08~\cite{2608_22906}。AquaFlow 从水下单目流进入在线高斯建图。用水下数据适配的几何模型提供位姿跟踪与深度初始化线索，同时将可见颜色随距离变化拆给衰减和后向散射模型，使场景本征外观和介质现象不必都储存为视角相关颜色。

论文组织了 62 段视频，但定量重建只是其中可评测的子集，且作者报告 SeaThru-NeRF 上低于离线 UW-GS。应分开比较实时与离线约束，也不能把整个测试套件看作均有米制轨迹和参考几何。

\paragraph{CGS-SLAM：让单目惯导局部高斯子图在服务器对齐}
2026.08~\cite{2608_26868}。各客户端从单目与 IMU 生成局部视觉惯导轨迹，用单目米制深度辅助高斯子图尺度。协作阶段共享关键帧编码以发现重叠，再在服务器用 VGGT 帮助子图配准，合成全局表示。

作者报告笔记本跟踪建图约每秒十帧，原实验将较重的 Depth Pro 放在独立服务器，关键帧通信量取决于采样节奏，合并地图也需要传输。端侧运行的结果与完全离线或不依赖服务器的系统应分别理解。

\paragraph{SCOUT-SLAM：在动态场景中分别估计跟踪和建图的不确定性}
2026.09~\cite{2609_14634}。动态场景中，错误位姿会损坏地图，损坏的地图又会误导下一次跟踪。SCOUT-SLAM 从各关键帧提取 DINOv2 特征，由共享网络估计建图不确定性，再用叠加在共享网络上的低秩适配器估计跟踪不确定性，按像素重新加权位姿优化残差。两项估计通过不同监督更新，共用图像特征，同时避免直接把当前地图误差当成跟踪可靠度。

论文在 TUM RGB-D、Bonn Dynamic 和 Wild-SLAM MoCap 上比较轨迹或建图效果。TUM 的静态渲染指标排除了用 SAM3 标记的人和椅子区域，不能解读为对整幅动态场景的重建。附录另在 DROID-W 评估，其中 Downtown 1 和 4 两条超过 100 米的轨迹因显存不足未报告。

\paragraph{PanoGS-SLAM：在球面域优化全景光度与相机轨迹}
2026.09~\cite{2609_17387}。PanoGS-SLAM 将一张全景与统一三维高斯地图匹配来估计相机位姿，并在球面条件下增量写入新观测。等距矩形投影在高纬有采样面积失衡，光度损失据此调整贡献，地图则由深度先验引导初始化。

作者在合成全景 SynPano 及实采宽视场 PALVIO 对比定位，并用受控视场实验分析宽角覆盖。结果需结合输入的球面相机几何；同名全景设备未必覆盖相同的竖直角度。

\paragraph{VGGT-GS SLAM：联合细化未标定单目子图的几何与内参}
2026.09~\cite{2609_19628}。VGGT-GS SLAM 先用前馈几何给每个时间窗初始化高斯子图，再用可微束调整共同优化高斯、相机外参、子图共享内参和畸变。其后在全局位姿图中使用高斯子图间的相似变换约束，降低各窗独立初始化导致的尺度与漂移不一致。

论文在 TUM-RGBD、7-Scenes、ScanNet 等室内视频上评价定位，输入仅用 RGB。作者指出性能受几何先验质量与标定可观测性影响，窗口也带来一定延迟。

\paragraph{EliGSiR：按地图当前需要动态分配 RGB-D 在线预算}
2026.09~\cite{2609_20348}。EliGSiR 对持续 RGB-D 建图中的视图、分辨率和原语增长分别分配预算。地图变化后可重新排序历史视图，输入负载升高时可降低监督分辨率再恢复，在重复观测证实几何缺失的位置才追加容量。

它在 Replica、TUM RGB-D 等序列评估留出视图，并把输入流内、有限清尾和额外后优化的状态明确区分。多数受控比较供应同一参考位姿，因此绘图误差不能解释为实时定位精度，跨不同试验组也不宜直接比较绝对 PSNR。

\paragraph{RawSLAM：用线性高动态范围观测进行高斯跟踪与建图}
2026.09~\cite{2609_20589}。RawSLAM 将相机 RAW 转成线性 HDR 输入，用对数方式参数化高斯颜色，并调整光度梯度的亮度权重，使跟踪和高斯训练能保留高光与阴影的信息。这与将八位色调映射图像直接作为观测的做法不同。

作者在十段室内序列中比较原生 HDR 与常规图像处理，报告改善随场景照明而变化。PSNR、SSIM 和 LPIPS 在色调映射后的图像上计算，因此这些指标不能当作直接的线性辐亮度误差。

\paragraph{Cube-Splat：将全景图像分面后共用一个位姿变量}
2026.09~\cite{2609_21347}。直接将全景图像送入针孔高斯渲染器会受到投影模型限制。Cube-Splat 将全景帧分解成固定朝向、共享光心的虚拟透视视图，并复用透视高斯渲染。每个分面产生的梯度用伴随变换映射到同一个位姿变量，因而不会把各分面误当成可以独立运动的相机。

作者在 SynPano、OmniBlender 和 PALVIO 共 31 个序列中报告轨迹误差和视图质量，并给出逐序列结果。该表示依赖准确的全景标定和固定分面朝向，拼接误差会影响接缝附近的光度一致性。系统针对大体静态的环境，未融合惯性或深度观测，缺纹理与重复结构下仍可能漂移。

\paragraph{Elevator-VIGS：把电梯轿厢运动与轿厢内相机运动分开}
2026.09~\cite{2609_23491}。相机在轿厢内看到的是相对于轿厢的运动，而惯导响应包含轿厢整体运动。Elevator-VIGS 在束调整中加入沿重力方向的运输状态，通过视觉语言和深度线索识别电梯事件，使轨迹与高斯地图能在移动坐标系下更新。

作者测试十六段真实电梯记录，并用模拟电梯序列补充过程真值。真实轨迹主要在上下电梯前后的控制点评价竖向误差，中间全程轨迹的证据来自仿真。设计针对沿重力方向运动的电梯。

\paragraph{BayesianGS-SLAM：把预测不确定性用于跟踪、建图和选帧}
2026.09~\cite{2609_24140}。BayesianGS-SLAM 保留 LoopSplat 的 RGB-D 高斯地图和在线更新流程，将颜色及深度的传感器噪声与不透明度引起的地图表示不确定性合并为预测量。跟踪按预测不确定性调整残差，选帧则看当前观测相对地图预测有多意外，从而减少对已解释视角的重复更新。

原文报告 TUM 的平均绝对轨迹误差为 3.12 厘米，ScanNet 为 6.7 厘米；三条 TUM 序列上的插入关键帧和建图调用均约减少七成。渲染指标在训练关键帧上计算，只反映这些视图上的拟合效果。机器人采集部分为定性演示，不能代替大规模独立轨迹验证。

\paragraph{双协方差高斯 SLAM：让渲染与配准共享中心而各自建模}
2026.09~\cite{2609_25746}。光度优化得到的高斯协方差表示适合渲染的形状，不一定反映深度测量的可信度。该方法为共享中心的高斯保存两套协方差，一套按光度优化，一套表示 RGB-D 测量不确定性。后者用于深度的广义 ICP 残差，也给锚定到附近高斯的图像角点提供置信度，补充几何较弱时的跟踪线索。

在 TUM RGB-D、ScanNet 和 Replica 中，作者分别统计绝对轨迹误差、相对姿态误差与渲染指标，并报告困难手持或室外例子。系统缺少大场景长期一致性所需的全局束调整，且深度传感误差随距离增大，因此不能将近距离的跟踪收益外推到远景主导的环境。

\paragraph{ArborSplat：以几何与类别预算维护果园语义地图}
2026.09~\cite{2609_26315}。ArborSplat 用 LiDAR 里程计提供相机位姿，在线维护带类别分数的高斯。图像掩码结合双目点和行间地面几何形成语义监督，并在预算受限时按类别分配部分新原语。高斯地图不反馈到跟踪端。

论文在苹果与梨的三个阶段比较在线地图。薄树干在训练视图上可识别，但留出视图的树干 IoU 在完整路线中仍低于 0.07，这一差距限制了新视角的细结构语义稳定性。

\paragraph{HyperLoop：用负结果界定高光谱闭环检测的实际收益}
2026.10~\cite{2610_19742}。HyperLoop 在基于高斯的稠密 SLAM 回环模块中并列使用 RGB 描述子与光谱描述子，以严格同意或加权机制筛选闭环。它改变的是地点识别和回环决策，高斯场的建图观测没有直接改用高光谱。

在 MMOTS-WS 上，最佳融合的对齐 ATE RMSE 是 5.04 厘米，而单独使用 SIFT-VLAD 为 4.95 厘米。作者把相对 NetVLAD 的轨迹收益归因于 RGB 描述子替换。这个负结果针对 600 至 860 纳米与当前场地，不能据此断言其他波段无效。

\paragraph{GOFusion：在占据子图中锚定高斯并积累信息量}
2026.10~\cite{2610_25261}。GOFusion 用视觉惯导估计位姿和占据子图，再以体素与表面法向引导活动子图的高斯初始化。占据层支持规划空间查询，高斯层保留外观。光栅化过程中累计的对角 Fisher 近似用作原语裁剪和选视的信息指标。

论文在 TUM RGB-D、uHumans2 和真实飞行记录中分别评价。飞行演示是给定平台和传感器的实例，Fisher 近似用于相对信息量排序，并不是导航安全保证。

\paragraph{SG²SLAM：用流式点图和三维关联连接户外高斯子图}
2026.10~\cite{2610_36312}。SG²SLAM 让冻结的前馈模型按连续帧窗口预测局部位姿、点图与置信度，用三维先验同时服务跟踪和高斯建图。高斯核受视线条件调节，窗口间通过全局相似变换图纠正单目尺度和累计漂移。

论文在 KITTI 和 Waymo 等户外序列用 Sim(3) 对齐后的轨迹误差评价，不能把对齐误差解释为原生绝对尺度验证。作者也把雨夜等非结构化场景的稳定性留作后续研究。

\paragraph{LingGS：把流式前馈模型当作虚拟 RGB-D 传感器}
2026.10~\cite{2610_36853}。LingGS 让冻结的流式模型预测内参、深度和图像特征，然后以显式对应、重投影和深度误差估计相机运动。高斯后端消费连续图像、深度和这些位姿，不直接采用前馈模型相机头输出的轨迹。

作者在 7-Scenes、Oxford Spires 等室内外数据上比较定位和渲染，也检验不同流式主干下的显式运动求解。该设计把位姿误差与几何预测解耦，但仍依赖深度、匹配特征及相机内参的可用性。

\paragraph{DPGS-SLAM：区分真实观测与暂存的生成背景}
2026.11~\cite{2611_12038}。DPGS-SLAM 用光流差异和 SAM2 找动态区域，由真实静态区域支持相机跟踪。异步扩散修补可为遮挡背景生成暂存高斯，但这些原语带来源记录，不参与位姿与观测地图优化；有效真实静态观测到达后可替换它们。

作者用 BONN、TUM 评价轨迹，并用 StaDy4D 检验静态背景。不可见区域仍可能长期保留错误外观或几何，异步补全也不保证实时。保留来源的意义在于让真实测量和生成内容承担不同角色。

\paragraph{CoSphere-SLAM：球面闭环与多机器人全景子图融合}
2026.11~\cite{2611_22235}。CoSphere-SLAM 把每台机器人的全景转到立方体视面用于局部估计，并在服务器按球面方向匹配不同朝向的关键帧。全局相似变换一致性用于过滤对称外观导致的伪闭环，同时协调各局部地图的尺度与位姿，再融合高斯。

论文比较多个室内和跨采集会话全景数据，在表中分别列出有有效轨迹的场景统计和全部场景统计。比较时应读清失败处理与是否多机器人联合运行，不能只拿成功子集的轨迹均值反映整体稳定性。


\section*{本章小结}
\addcontentsline{toc}{section}{本章小结}

\begin{futurebox}{本章总结与展望}{ch15_summary}
3DGS 为 SLAM 与稠密建图领域带来了范式转换：地图不再只是稀疏点云或体素网格，而是可实时渲染的照片级三维场景。本章系统梳理了 183 篇论文，覆盖八个核心方向。

在\textbf{RGB-D SLAM}（第~\ref{sec:15_1}~节）中，3DGS-SLAM 已在跟踪精度和重建质量上展现出明显优势，从 SplaTAM 的奠基性工作到 2D Surfel 表示、不确定性建模和实时加速，技术路线日趋成熟。

\textbf{单目与 RGB SLAM}（第~\ref{sec:15_2}~节）更具挑战性，前馈范式和多视角基础模型的融合正在突破传统优化方法的瓶颈。

\textbf{多模态 SLAM}（第~\ref{sec:15_3}~节）通过融合 IMU、LiDAR、雷达等互补传感器大幅提升鲁棒性，LiDAR-视觉融合已成为户外场景的事实标准。

\textbf{视觉定位与重定位}（第~\ref{sec:15_4}~节）利用 3DGS 的可微渲染和特征场实现精确的 6-DoF 位姿估计，从稀疏到稠密的定位范式正在取代传统的粗到精策略。

\textbf{语义 SLAM}（第~\ref{sec:15_5}~节）通过开放词汇理解赋予地图语义能力，层次化表示和高效特征压缩解决了参数膨胀问题。

\textbf{动态场景 SLAM}（第~\ref{sec:15_6}~节）正在从简单的动态滤除演化为动态-静态联合重建，4DGS 的引入为动态内容建模提供了新的技术路径。

\textbf{系统优化与工程}（第~\ref{sec:15_7}~节）解决了运动模糊、回环检测和实时加速等实际部署中的关键问题。

\textbf{大规模与多智能体 SLAM}（第~\ref{sec:15_8}~节）通过核外计算、子地图策略和分布式共识将 3DGS-SLAM 扩展到城市级规模。

展望未来，3DGS-SLAM 仍面临几个开放挑战：（1）\textbf{长期运行与终身建图}——如何在数小时甚至数天的运行中维护和更新高斯地图；（2）\textbf{语义-几何联合优化}——如何让语义理解反哺几何重建精度；（3）\textbf{边缘部署}——如何进一步压缩模型以适应微无人机和 AR 眼镜等严格受限平台；（4）\textbf{Sim2Real 迁移}——如何弥合仿真训练和真实部署之间的差距。这些问题的解决将推动 3DGS-SLAM 从学术研究走向广泛工业应用。
\end{futurebox}
